
\Data\ID{1003261901}
\Updated{2020-07-15 03:07:19}
\Level{A}
\SubArea{Applications of derivatives}
\LANGen{
\Question{
Find the second derivative of the function 
\[ f(x)=\frac{x^2}{1-x} \]
at the point \( x_0=2 \).}
{\( -2 \)}{\( 2 \)}{\( -\frac14 \)}{\( \frac14 \)}{\( -4 \)}{\( 4 \)}}
\LANGpl{
\Question{
Wyznacz drugą pochodną funkcji 
\[ f(x)=\frac{x^2}{1-x} \]
w punkcie \( x_0=2 \).}
{\( -2 \)}{\( 2 \)}{\( -\frac14 \)}{\( \frac14 \)}{\( -4 \)}{\( 4 \)}}
\LANGcs{
\Question{
Určete druhou derivaci funkce
\[ f(x)=\frac{x^2}{1-x} \]
v bodě \( x_0=2 \).}
{\( -2 \)}{\( 2 \)}{\( -\frac14 \)}{\( \frac14 \)}{\( -4 \)}{\( 4 \)}}
\LANGsk{
\Question{
Určte druhú deriváciu funkcie
\[ f(x)=\frac{x^2}{1-x} \]
v bode \( x_0=2 \).}
{\( -2 \)}{\( 2 \)}{\( -\frac14 \)}{\( \frac14 \)}{\( -4 \)}{\( 4 \)}}
\LANGes{
\Question{
Halla la segunda derivada de la función
\[ f(x)=\frac{x^2}{1-x} \]
en el punto \( x_0=2 \).}
{\( -2 \)}{\( 2 \)}{\( -\frac14 \)}{\( \frac14 \)}{\( -4 \)}{\( 4 \)}}
\End

\Data\ID{1003261902}
\Updated{2020-07-15 03:07:19}
\Level{A}
\SubArea{Applications of derivatives}
\LANGen{
\Question{
Find the second derivative of the function
\[ f(x)=\sin^2 x \]
at the point \( x_0=-\frac{\pi}6 \).}
{\( 1 \)}{\( \frac12 \)}{\( -\frac12 \)}{\( -1 \)}{\( \sqrt3 \)}{\( -\frac{\sqrt3}2 \)}}
\LANGpl{
\Question{
Wyznacz drugą pochodną funkcji 
\[ f(x)=\sin^2 x \]
w punkcie \( x_0=-\frac{\pi}6 \).}
{\( 1 \)}{\( \frac12 \)}{\( -\frac12 \)}{\( -1 \)}{\( \sqrt3 \)}{\( -\frac{\sqrt3}2 \)}}
\LANGcs{
\Question{
Určete druhou derivaci funkce
\[ f(x)=\sin^2 x \]
v bodě \( x_0=-\frac{\pi}6 \).}
{\( 1 \)}{\( \frac12 \)}{\( -\frac12 \)}{\( -1 \)}{\( \sqrt3 \)}{\( -\frac{\sqrt3}2 \)}}
\LANGsk{
\Question{
Určte druhú deriváciu funkcie
\[ f(x)=\sin^2 x \]
v bode \( x_0=-\frac{\pi}6 \).}
{\( 1 \)}{\( \frac12 \)}{\( -\frac12 \)}{\( -1 \)}{\( \sqrt3 \)}{\( -\frac{\sqrt3}2 \)}}
\LANGes{
\Question{
Halla la segunda derivada de la función
\[ f(x)=\sin^2 x \]
en el punto \( x_0=-\frac{\pi}6 \).}
{\( 1 \)}{\( \frac12 \)}{\( -\frac12 \)}{\( -1 \)}{\( \sqrt3 \)}{\( -\frac{\sqrt3}2 \)}}
\End

\Data\ID{1003261903}
\Updated{2021-08-12 10:08:21}
\Level{A}
\SubArea{Applications of derivatives}
\LANGen{
\Question{
Given the function 
\[ f(x)=x^3-3x^2+2\text{ ,} \]
determine the set of all \( x \), \( x\in\mathbb{R} \), such that \( f''(x)-f'(x)=3 \).}
{\( \{1;3\} \)}{\( \{-1;-3\} \)}{\( \{-\sqrt3;\sqrt3\} \)}{\( \{\sqrt3\} \)}{\( \emptyset \)}{}}
\LANGpl{
\Question{
Dana jest funkcja
\[ f(x)=x^3-3x^2+2\text{ ,} \]
wskaż zbiór wszystkich \( x \), \( x\in\mathbb{R} \) tak, aby\( f''(x)-f'(x)=3 \).}
{\( \{1;3\} \)}{\( \{-1;-3\} \)}{\( \{-\sqrt3;\sqrt3\} \)}{\( \{\sqrt3\} \)}{\( \emptyset \)}{}}
\LANGcs{
\Question{
Je dána funkce
\[ f(x)=x^3-3x^2+2\text{ ,} \]
najděte množinu všech \( x \), \( x\in\mathbb{R} \), pro která platí \( f''(x)-f'(x)=3 \).}
{\( \{1;3\} \)}{\( \{-1;-3\} \)}{\( \{-\sqrt3;\sqrt3\} \)}{\( \{\sqrt3\} \)}{\( \emptyset \)}{}}
\LANGsk{
\Question{
Je daná funkcia
\[ f(x)=x^3-3x^2+2\text{ ,} \]
nájdite množinu všetkých \( x \), \( x\in\mathbb{R} \), pre ktoré platí \( f''(x)-f'(x)=3 \).}
{\( \{1;3\} \)}{\( \{-1;-3\} \)}{\( \{-\sqrt3;\sqrt3\} \)}{\( \{\sqrt3\} \)}{\( \emptyset \)}{}}
\LANGes{
\Question{
Dada la función 
\[ f(x)=x^3-3x^2+2\text{ ,} \]
halla el conjunto de todos los \( x \), \( x\in\mathbb{R} \), para los cuales \( f''(x)-f'(x)=3 \).}
{\( \{1;3\} \)}{\( \{-1;-3\} \)}{\( \{-\sqrt3;\sqrt3\} \)}{\( \{\sqrt3\} \)}{\( \emptyset \)}{}}
\End

\Data\ID{1003261904}
\Updated{2021-08-12 10:08:56}
\Level{A}
\SubArea{Applications of derivatives}
\LANGen{
\Question{
Given the function 
\[ f(x)=\sin x-3\cos x\text{ ,} \]
determine the set of all \( x \), \( x\in\mathbb{R} \), such  that \( f''(x)+f(x)=0 \).}
{\( \mathbb{R} \)}{\( \emptyset \)}{\( \{k\pi;\ k\in\mathbb{Z}\} \)}{\( \left\{(2k+1)\frac{\pi}2;\ k\in\mathbb{Z} \right\} \)}{}{}}
\LANGpl{
\Question{
Dana jest funkcja
\[ f(x)=\sin x-3\cos x\text{ ,} \]
wskaż zbiór wszystkich \( x \), \( x\in\mathbb{R} \) tak, aby \( f''(x)+f(x)=0 \).}
{\( \mathbb{R} \)}{\( \emptyset \)}{\( \{k\pi;\ k\in\mathbb{Z}\} \)}{\( \left\{(2k+1)\frac{\pi}2;\ k\in\mathbb{Z} \right\} \)}{}{}}
\LANGcs{
\Question{
Je dána funkce
\[ f(x)=\sin x-3\cos x\text{ ,} \]
najděte množinu všech \( x \), \( x\in\mathbb{R} \), pro která platí \( f''(x)+f(x)=0 \).}
{\( \mathbb{R} \)}{\( \emptyset \)}{\( \{k\pi;\ k\in\mathbb{Z}\} \)}{\( \left\{(2k+1)\frac{\pi}2;\ k\in\mathbb{Z} \right\} \)}{}{}}
\LANGsk{
\Question{
Je daná funkcia
\[ f(x)=\sin x-3\cos x\text{ ,} \]
nájdite množinu všetkých \( x \), \( x\in\mathbb{R} \), pre ktoré platí \( f''(x)+f(x)=0 \).}
{\( \mathbb{R} \)}{\( \emptyset \)}{\( \{k\pi;\ k\in\mathbb{Z}\} \)}{\( \left\{(2k+1)\frac{\pi}2;\ k\in\mathbb{Z} \right\} \)}{}{}}
\LANGes{
\Question{
Dada la función
\[ f(x)=\sin x-3\cos x\text{ ,} \]
halla el conjunto de todos los \( x \), \( x\in\mathbb{R} \), para los cuales \( f''(x)+f(x)=0 \).}
{\( \mathbb{R} \)}{\( \emptyset \)}{\( \{k\pi;\ k\in\mathbb{Z}\} \)}{\( \left\{(2k+1)\frac{\pi}2;\ k\in\mathbb{Z} \right\} \)}{}{}}
\End

\Data\ID{1003261905}
\Updated{2021-08-12 10:08:37}
\Level{A}
\SubArea{Applications of derivatives}
\LANGen{
\Question{
Find the local extrema of the function 
\[ f(x)=x-\ln(1+x)\text{ .} \]}
{the local minimum at \( x=0 \)}{the local minimum at \( x=0 \), the local maximum at \( x=-1 \)}{the local maximum at \( x=0 \)}{the local maximum at \( x=0 \), the local minimum at \( x=-1 \)}{do not exist}{}}
\LANGpl{
\Question{
Wskaż lokalne ekstrema funkcji
\[ f(x)=x-\ln(1+x)\text{ .} \]}
{lokalne minimum w  \( x=0 \)}{lokalne minimum w \( x=0 \), lokalne maksimum w   \( x=-1 \)}{lokalne maksimum w   \( x=0 \)}{lokalne maksimum w  \( x=0 \), lokalne minimum w \( x=-1 \)}{nie istnieje}{}}
\LANGcs{
\Question{
Najděte lokální extrémy funkce  
\[ f(x)=x-\ln(1+x)\text{ .} \]}
{lokální minimum v bodě \( x=0 \)}{lokální minimum v bodě \( x=0 \), lokální maximum v bodě \( x=-1 \)}{lokální maximum v bodě \( x=0 \)}{lokální maximum v bodě \( x=0 \), lokální minimum v bodě \( x=-1 \)}{neexistují}{}}
\LANGsk{
\Question{
Nájdite lokálne extrémy funkcie  
\[ f(x)=x-\ln(1+x)\text{ .} \]}
{lokálne minimum v bode \( x=0 \)}{lokálne minimum v bode \( x=0 \), lokálne maximum v bode \( x=-1 \)}{lokálne maximum v bode \( x=0 \)}{lokálne maximum v bode \( x=0 \), lokálne minimum v bode \( x=-1 \)}{neexistujú}{}}
\LANGes{
\Question{
Halla los extremos locales de la función:
\[ f(x)=x-\ln(1+x)\text{ .} \]}
{Mínimo local en \( x=0 \)}{Mínimo local en \( x=0 \),  máximo local en \( x=-1 \)}{Máximo local en \( x=0 \)}{Máximo local en \( x=0 \),  mínimo local en \( x=-1 \)}{No existen extremos}{}}
\End

\Data\ID{1103163602}
\Updated{2021-08-12 10:08:52}
\Level{A}
\SubArea{Applications of derivatives}
\IMG{1103163602.pdf}
\LANGen{
\Question{
The graph of \( f' \) is given in the figure. Find the interval where \( f \) is an increasing function. (The function \( f' \) is the derivative of the function \( f \).)}
{\( (-1;1) \)}{\( (-3;-1) \)}{\( (2;4) \)}{\( (0;2) \)}{}{}}
\LANGpl{
\Question{
Dany jest wykres funkcji \( f' \). Wskaż przedział w którym funkcja \( f \) jest funkcją rosnącą. (Funkcja  \( f' \) jest pochodną funkcji \( f \).)}
{\( (-1;1) \)}{\( (-3;-1) \)}{\( (2;4) \)}{\( (0;2) \)}{}{}}
\LANGcs{
\Question{
Na obrázku je dán graf funkce \( f' \). Určete, na kterém z nabídnutých intervalů je funkce \( f \) rostoucí. (Funkce  \( f' \) je derivace funkce \( f \).)}
{\( (-1;1) \)}{\( (-3;-1) \)}{\( (2;4) \)}{\( (0;2) \)}{}{}}
\LANGsk{
\Question{
Na obrázku je daný graf funkcie \( f' \). Určte, na ktorom intervale je funkcia \( f \) rastúca. (Funkcia  \( f' \) je derivácia funkcie \( f \).)}
{\( (-1;1) \)}{\( (-3;-1) \)}{\( (2;4) \)}{\( (0;2) \)}{}{}}
\LANGes{
\Question{
Dado el gráfico de la función \( f' \). Halla el intervalo donde \( f \) es creciente. (La función \( f' \) es la derivada de la función \( f \).)}
{\( (-1;1) \)}{\( (-3;-1) \)}{\( (2;4) \)}{\( (0;2) \)}{}{}}
\End

\Data\ID{1103163603}
\Updated{2021-08-12 10:08:51}
\Level{A}
\SubArea{Applications of derivatives}
\IMG{1103163603.pdf}
\LANGen{
\Question{
The graph of \( f' \) is given in the figure. Find the interval where \( f \) is an increasing function. (The function \( f' \) is the derivative of the function \( f \).)}
{\( (1;2) \)}{\( (-1;1) \)}{\( (1;3) \)}{\( (-1;0) \)}{}{}}
\LANGpl{
\Question{
Dany jest wykres funkcji\( f' \) Wskaż przedział w którym funkcja  \( f \) jest funkcją rosnącą. (Funkcja \( f' \) jest pochodną funkcji \( f \).)}
{\( (1;2) \)}{\( (-1;1) \)}{\( (1;3) \)}{\( (-1;0) \)}{}{}}
\LANGcs{
\Question{
Na obrázku je dán graf funkce \( f' \). Určete, na kterém z nabídnutých intervalů je funkce \( f \) rostoucí. (Funkce  \( f' \) je derivace funkce \( f \) .)}
{\( (1;2) \)}{\( (-1;1) \)}{\( (1;3) \)}{\( (-1;0) \)}{}{}}
\LANGsk{
\Question{
Na obrázku je daný graf funkcie \( f' \). Určte, na ktorom intervale je funkcia \( f \) rastúca. (Funkcia  \( f' \) je derivácia funkcie \( f \) .)}
{\( (1;2) \)}{\( (-1;1) \)}{\( (1;3) \)}{\( (-1;0) \)}{}{}}
\LANGes{
\Question{
Dado el gráfico de la función \( f' \). Halla el intervalo donde \( f \) es creciente. (La función \( f' \) es la derivada de la función \( f \).)}
{\( (1;2) \)}{\( (-1;1) \)}{\( (1;3) \)}{\( (-1;0) \)}{}{}}
\End

\Data\ID{1103163604}
\Updated{2021-08-12 10:08:17}
\Level{A}
\SubArea{Applications of derivatives}
\IMG{1103163604.pdf}
\LANGen{
\Question{
The graph of \( f' \) is given in the figure. Find the interval where \( f \) is a decreasing function. (The function \( f' \) is the derivative of the function \( f \).)}
{\( (-3;-2) \)}{\( (-1;1) \)}{\( (0;2) \)}{\( (-1;2) \)}{}{}}
\LANGpl{
\Question{
Dany jest wykres funkcji \( f' \) Wskaż przedział, w którym funkcja  \( f \) jest funkcją malejącą  (Funkcja \( f' \) jest pochodną funkcji \( f \).)}
{\( (-3;-2) \)}{\( (-1;1) \)}{\( (0;2) \)}{\( (-1;2) \)}{}{}}
\LANGcs{
\Question{
Na obrázku je dán graf funkce \( f' \). Určete, na kterém z nabídnutých intervalů je funkce \( f \) klesající. (Funkce  \( f' \) je derivace funkce \( f \).)}
{\( (-3;-2) \)}{\( (-1;1) \)}{\( (0;2) \)}{\( (-1;2) \)}{}{}}
\LANGsk{
\Question{
Na obrázku je daný graf funkcie \( f' \). Určte, na ktorom intervale je funkcia \( f \) klesajúca. (Funkcia  \( f' \) je derivácia funkcie \( f \).)}
{\( (-3;-2) \)}{\( (-1;1) \)}{\( (0;2) \)}{\( (-1;2) \)}{}{}}
\LANGes{
\Question{
Dado el gráfico de la función \( f' \). Halla el intervalo donde \( f \) es decreciente. (La función \( f' \) es la derivada de la función \( f \).)}
{\( (-3;-2) \)}{\( (-1;1) \)}{\( (0;2) \)}{\( (-1;2) \)}{}{}}
\End

\Data\ID{1103163605}
\Updated{2021-08-12 10:08:45}
\Level{A}
\SubArea{Applications of derivatives}
\IMG{1103163605.pdf}
\LANGen{
\Question{
The graph of \( f' \) is given in the figure. Find the interval where \( f \) is a decreasing function. (The function \( f' \) is the derivative of the function \( f \).)}
{\( (2;4) \)}{\( (-1;1) \)}{\( (1;3) \)}{\( (-4;-2) \)}{}{}}
\LANGpl{
\Question{
Dany jest wykres funkcji\( f' \) Wskaż przedział, w którym funkcja  \( f \) jest funkcją malejącą  (Funkcja \( f' \) jest pochodną funkcji \( f \).)}
{\( (2;4) \)}{\( (-1;1) \)}{\( (1;3) \)}{\( (-4;-2) \)}{}{}}
\LANGcs{
\Question{
Na obrázku je dán graf funkce \( f' \). Určete, na kterém z nabídnutých intervalů je funkce \( f \) klesající. (Funkce  \( f' \) je derivace funkce \( f \).)}
{\( (2;4) \)}{\( (-1;1) \)}{\( (1;3) \)}{\( (-4;-2) \)}{}{}}
\LANGsk{
\Question{
Na obrázku je daný graf funkcie \( f' \). Určte, na ktorom intervale je funkcia \( f \) klesajúca. (Funkcia  \( f' \) je derivácia funkcie \( f \).)}
{\( (2;4) \)}{\( (-1;1) \)}{\( (1;3) \)}{\( (-4;-2) \)}{}{}}
\LANGes{
\Question{
Dado el gráfico de la función \( f' \). Halla el intervalo donde \( f \) es decreciente. (La función \( f' \) es la derivada de la función \( f \).)}
{\( (2;4) \)}{\( (-1;1) \)}{\( (1;3) \)}{\( (-4;-2) \)}{}{}}
\End

\Data\ID{1103163606}
\Updated{2021-08-12 10:08:50}
\Level{A}
\SubArea{Applications of derivatives}
\IMG{1103163606.pdf}
\LANGen{
\Question{
The graph of \( f' \) is given in the figure. Find the local extrema of \( f \). (The function \( f' \) is the derivative of the function \( f \).)}
{local minimum at \( x=0 \), local maxima at \( x_1=-2 \) and \( x_2=3 \)}{local minimum at \( x=-1 \), local maximum at \( x=2 \)}{local minima at \( x_1=-2 \) and \( x_2=3 \), local maximum at \( x=0 \)}{local minima at \( x_1=-2 \) and \( x_2=0 \), local maximum at \( x=3 \)}{local minimum at \( x=-2 \), local maxima at \( x_1=0 \) and \( x_2=2 \)}{}}
\LANGpl{
\Question{
Dany jest wykres funkcji \( f' \). Wskaż lokalne ekstrema \( f \). (Funkcja \( f' \) jest pochodną funkcji \( f \).)}
{lokalne minimum w  \( x=0 \), lokalne maksima w  \( x_1=-2 \) i \( x_2=3 \)}{lokalne minimum w \( x=-1 \), lokalne maksimum w \( x=2 \)}{lokalne minima w \( x_1=-2 \) and \( x_2=3 \), lokalne maksimum w \( x=0 \)}{lokalne minima w  \( x_1=-2 \) i \( x_2=0 \), lokalne maksimum w  \( x=3 \)}{lokalne minimum w \( x=-2 \), lokalne maksima w  \( x_1=0 \) i \( x_2=2 \)}{}}
\LANGcs{
\Question{
Na obrázku je dán graf funkce \( f' \). Nalezněte lokální extrémy funkce \( f \). (Funkce  \( f' \) je derivace funkce \( f \).)}
{lokální minimum v bodě \( x=0 \), lokální maxima v bodech \( x_1=-2 \) a \( x_2=3 \)}{lokální minimum v bodě \( x=-1 \), lokální maximum v bodě \( x=2 \)}{lokální minima v bodech \( x_1=-2 \) a \( x_2=3 \), lokální maximum v bodě \( x=0 \)}{lokální minima v bodech \( x_1=-2 \) a \( x_2=0 \), lokální maximum v bodě \( x=3 \)}{lokální minimum v bodě \( x=-2 \), lokální maxima v bodech \( x_1=0 \) a \( x_2=2 \)}{}}
\LANGsk{
\Question{
Na obrázku je daný graf funkcie \( f' \). Nájdite lokálne extrémy funkcie \( f \). (Funkcia  \( f' \) je derivácia funkcie \( f \).)}
{lokálne minimum v bode \( x=0 \), lokálne maximá v bodoch \( x_1=-2 \) a \( x_2=3 \)}{lokálne minimum v bode \( x=-1 \), lokálne maximum v bode \( x=2 \)}{lokálne minimá v bodoch \( x_1=-2 \) a \( x_2=3 \), lokálne maximum v bode \( x=0 \)}{lokálne minimá v bodoch \( x_1=-2 \) a \( x_2=0 \), lokálne maximum v bode \( x=3 \)}{lokálne minimum v bode \( x=-2 \), lokálne maximá v bodoch \( x_1=0 \) a \( x_2=2 \)}{}}
\LANGes{
\Question{
Dado el gráfico de la función \( f' \). Halla los extremos locales de \( f \). (La función \( f' \) es la derivada de la función \( f \).)}
{mínimo local en \( x=0 \), máximos locales en \( x_1=-2 \) y \( x_2=3 \)}{mínimo local en \( x=-1 \),  máximo local en \( x=2 \)}{mínimos locales en \( x_1=-2 \) y \( x_2=3 \),  máximo local en \( x=0 \)}{mínimos locales en \( x_1=-2 \) y \( x_2=0 \),  máximo local en \( x=3 \)}{mínimo local en \( x=-2 \), máximos locales en \( x_1=0 \) y \( x_2=2 \)}{}}
\End

\Data\ID{1103163607}
\Updated{2021-08-12 10:08:59}
\Level{A}
\SubArea{Applications of derivatives}
\IMG{1103163607.pdf}
\LANGen{
\Question{
The graph of \( f' \) is given in the figure. Find the local extrema of \( f \). (The function \( f' \) is the derivative of the function \( f \).)}
{local minima at \( x_1=-1 \) and \( x_2=4 \), local maximum at \( x=1 \)}{local minimum at \( x=3 \), local maximum at \( x=0 \)}{local minimum at \( x=-1 \), local maximum at \( x=4 \)}{local minima at \( x_1=-1 \) and \( x_2=1 \), local maximum at \( x=4 \)}{local minimum at \( x=1 \), local maxima at \( x_1=-1 \) and \( x_2=4 \)}{}}
\LANGpl{
\Question{
Dany jest wykres funkcji \( f' \). Wskaż lokalne ekstrema funkcji \( f \). (Funkcja  \( f' \) jest pochodną funkcji \( f \).)}
{lokalne minima w  \( x_1=-1 \) i \( x_2=4 \), lokalne maksimum w  \( x=1 \)}{lokalne minimum w \( x=3 \),  lokalne maksimum w  \( x=0 \)}{lokalne minimum w \( x=-1 \),  lokalne maksimum w  \( x=4 \)}{lokalne minima w  \( x_1=-1 \) i \( x_2=1 \),  lokalne maksimum w  \( x=4 \)}{lokalne minimum w  \( x=1 \), lokalne maksima w  \( x_1=-1 \) i \( x_2=4 \)}{}}
\LANGcs{
\Question{
Na obrázku je dán graf funkce \( f' \). Nalezněte lokální extrémy funkce \( f \). (Funkce  \( f' \) je derivace funkce \( f \).)}
{lokální minima v bodech \( x_1=-1 \) a \( x_2=4 \), lokální maximum v bodě \( x=1 \)}{lokální minimum v bodě \( x=3 \), lokální maximum v bodě \( x=0 \)}{lokální minimum v bodě \( x=-1 \), lokální maximum v bodě \( x=4 \)}{lokální minima v bodech \( x_1=-1 \) a \( x_2=1 \), lokální maximum v bodě \( x=4 \)}{lokální minimum v bodě \( x=1 \), lokální maxima v bodech \( x_1=-1 \) a \( x_2=4 \)}{}}
\LANGsk{
\Question{
Na obrázku je daný graf funkcie \( f' \). Nájdite lokálne extrémy funkcie \( f \). (Funkcia  \( f' \) je derivácia funkcie \( f \).)}
{lokálne minimá v bodoch \( x_1=-1 \) a \( x_2=4 \), lokálne maximum v bode \( x=1 \)}{lokálne minimum v bode \( x=3 \), lokálne maximum v bode \( x=0 \)}{lokálne minimum v bode \( x=-1 \), lokálne maximum v bode \( x=4 \)}{lokálne minimá v bodoch \( x_1=-1 \) a \( x_2=1 \), lokálne maximum v bode \( x=4 \)}{lokálne minimum v bode \( x=1 \), lokálne maximá v bodoch \( x_1=-1 \) a \( x_2=4 \)}{}}
\LANGes{
\Question{
Dado el gráfico de la función \( f' \). Halla los extremos locales de \( f \). (La función \( f' \) es la derivada de la función \( f \).)}
{mínimos locales en \( x_1=-1 \) y \( x_2=4 \), máximo local en \( x=1 \)}{mínimo local en \( x=3 \), máximo local en \( x=0 \)}{mínimo local en \( x=-1 \), máximo local en \( x=4 \)}{mínimos locales en \( x_1=-1 \) y \( x_2=1 \), máximo local en \( x=4 \)}{mínimo local en \( x=1 \),  máximos locales en \( x_1=-1 \) y \( x_2=4 \)}{}}
\End

\Data\ID{1103163608}
\Updated{2021-08-12 10:08:02}
\Level{A}
\SubArea{Applications of derivatives}
\IMG{1103163608.pdf}
\LANGen{
\Question{
The graph of \( f' \) is given in the figure. Find the local extrema of \( f \). (The function \( f' \) is the derivative of the function \( f \).)}
{local minimum at \( x=3 \)}{local minimum at \( x=2 \), local maximum at \( x=0 \)}{local minimum at \( x=3 \), local maximum at \( x=0 \)}{local minimum at \( x=0 \), local maximum at \( x=3 \)}{local maximum at \( x=3 \)}{}}
\LANGpl{
\Question{
Dany jest wykres funkcji \( f' \). Wskaż lokalne ekstrema funkcji \( f \). (Funkcja  \( f' \) jest pochodną funkcji \( f \).)}
{lokalne minimum w \( x=3 \)}{lokalne minimum w  \( x=2 \), lokalne maksimum w \( x=0 \)}{lokalne minimum w  \( x=3 \), lokalne maksimum w \( x=0 \)}{lokalne minimum w  \( x=0 \), lokalne maksimum w \( x=3 \)}{lokalne maksimum w  \( x=3 \)}{}}
\LANGcs{
\Question{
Na obrázku je dán graf funkce \( f' \). Nalezněte lokální extrémy funkce \( f \). (Funkce  \( f' \) je derivace funkce \( f \).)}
{lokální minimum v bodě \( x=3 \)}{lokální minimum v bodě \( x=2 \), lokální maximum v bodě \( x=0 \)}{lokální minimum v bodě \( x=3 \), lokální maximum v bodě \( x=0 \)}{lokální minimum v bodě \( x=0 \), lokální maximum v bodě \( x=3 \)}{lokální maximum v bodě \( x=3 \)}{}}
\LANGsk{
\Question{
Na obrázku je daný graf funkcie \( f' \). Nájdite lokálne extrémy funkcie \( f \). (Funkcia  \( f' \) je derivácia funkcie \( f \).)}
{lokálne minimum v bode \( x=3 \)}{lokálne minimum v bode \( x=2 \), lokálne maximum v bode \( x=0 \)}{lokálne minimum v bode \( x=3 \), lokálne maximum v bode \( x=0 \)}{lokálne minimum v bode \( x=0 \), lokálne maximum v bode \( x=3 \)}{lokálne maximum v bode \( x=3 \)}{}}
\LANGes{
\Question{
Dado el gráfico de la función \( f' \). Halla los extremos locales de \( f \). (La función \( f' \) es la derivada de la función \( f \).)}
{mínimo local en \( x=3 \)}{mínimo local en \( x=2 \), máximo local en \( x=0 \)}{mínimo local en \( x=3 \), máximo local en \( x=0 \)}{mínimo local en \( x=0 \), máximo local en \( x=3 \)}{máximo local en \( x=3 \)}{}}
\End

\Data\ID{1103163609}
\Updated{2021-08-12 10:08:44}
\Level{A}
\SubArea{Applications of derivatives}
\IMG{1103163609.pdf}
\LANGen{
\Question{
The graph of \( f' \) is given in the figure. Find the local extrema of \( f \). (The function \( f' \) is the derivative of the function \( f \).)}
{local maximum at \( x=0 \)}{local minimum at \( x=3 \), local maximum at \( x=0 \)}{local minimum at \( x=1 \), local maximum at \( x=3 \)}{local minimum at \( x=0 \), local maximum at \( x=3 \)}{local minimum at \( x=0 \)}{}}
\LANGpl{
\Question{
Dany jest wykres funkcji \( f' \). Wskaż lokalne ekstrema funkcji \( f \). (Funkcja  \( f' \) jest pochodną funkcji \( f \).)}
{lokalne maksimum w \( x=0 \)}{lokalne minimum w  \( x=3 \), lokalne maksimum w \( x=0 \)}{lokalne minimum w  \( x=1 \), lokalne maksimum w \( x=3 \)}{lokalne minimum w  \( x=0 \), lokalne maksimum w \( x=3 \)}{lokalne minimum w  \( x=0 \)}{}}
\LANGcs{
\Question{
Na obrázku je dán graf funkce \( f' \). Nalezněte lokální extrémy funkce \( f \). (Funkce  \( f' \) je derivace funkce \( f \).)}
{lokální maximum v bodě \( x=0 \)}{lokální minimum v bodě \( x=3 \), lokální maximum v bodě \( x=0 \)}{lokální minimum v bodě \( x=1 \), lokální maximum v bodě \( x=3 \)}{lokální minimum v bodě \( x=0 \), lokální maximum v bodě \( x=3 \)}{lokální minimum v bodě \( x=0 \)}{}}
\LANGsk{
\Question{
Na obrázku je daný graf funkcie \( f' \). Nájdite lokálne extrémy funkcie \( f \). (Funkcia  \( f' \) je derivácia funkcie \( f \).)}
{lokálne maximum v bode \( x=0 \)}{lokálne minimum v bode \( x=3 \), lokálne maximum v bode \( x=0 \)}{lokálne minimum v bode \( x=1 \), lokálne maximum v bode \( x=3 \)}{lokálne minimum v bode \( x=0 \), lokálne maximum v bode \( x=3 \)}{lokálne minimum v bode \( x=0 \)}{}}
\LANGes{
\Question{
Dado el gráfico de la función \( f' \). Halla los extremos locales de \( f \). (La función \( f' \) es la derivada de la función \( f \).)}
{máximo local en \( x=0 \)}{mínimo local en \( x=3 \), máximo local en \( x=0 \)}{mínimo local en \( x=1 \),  máximo local en \( x=3 \)}{mínimo local en \( x=0 \), máximo local en \( x=3 \)}{mínimo local en \( x=0 \)}{}}
\End

\Data\ID{2010001701}
\Updated{2022-08-25 10:08:23}
\Level{A}
\SubArea{Applications of derivatives}
\LANGen{
\Question{
Find all the $x$ at which the function $f$ has local extrema.
\[ 
 f(x)= -2\left (x^2 - 4\right )^{5}
\]}
{\(x=0\)}{\(x_1=0\),
\(x_2=2\)}{\(x_1=-2\),
\(x_2=2\)}{\(x_1=- 2\),
\(x_2=0\),
\(x_3=2\)}{}{}}
\LANGpl{
\Question{
Znajdź wszystkie $x$, w których funkcja $f$ ma ekstrema lokalne.
\[ 
 f(x)= -2\left (x^2 - 4\right )^{5}
\]}
{\(x=0\)}{\(x_1=0\),
\(x_2=2\)}{\(x_1=-2\),
\(x_2=2\)}{\(x_1=- 2\),
\(x_2=0\),
\(x_3=2\)}{}{}}
\LANGcs{
\Question{
Určete všechna $x$, ve kterých má funkce $f$ lokální extrém.
\[ 
 f(x)= -2\left (x^2 - 4\right )^{5}
\]}
{\(x=0\)}{\(x_1=0\),
\(x_2=2\)}{\(x_1=-2\),
\(x_2=2\)}{\(x_1=- 2\),
\(x_2=0\),
\(x_3=2\)}{}{}}
\LANGsk{
\Question{
Určte všetky $x$, v ktorých má funkcia $f$ lokálny extrém. 
\[ 
 f(x)= -2\left (x^2 - 4\right )^{5}
\]}
{\(x=0\)}{\(x_1=0\),
\(x_2=2\)}{\(x_1=-2\),
\(x_2=2\)}{\(x_1=- 2\),
\(x_2=0\),
\(x_3=2\)}{}{}}
\LANGes{
\Question{
Halla los extremos locales de la siguiente función:
\[ 
 f(x)= -2\left (x^2 - 4\right )^{5}
\]}
{\(x=0\)}{\(x_1=0\),
\(x_2=2\)}{\(x_1=-2\),
\(x_2=2\)}{\(x_1=- 2\),
\(x_2=0\),
\(x_3=2\)}{}{}}
\End

\Data\ID{2010001702}
\Updated{2022-08-25 10:08:26}
\Level{A}
\SubArea{Applications of derivatives}
\LANGen{
\Question{
What is the function value of the function $f$ at its local maximum?
\[ 
 f(x) = \frac
{\sqrt{6x-x^2}}{2}
\]}
{\(\frac{3}{2}\)}{\(0\)}{\(2\)}{The local maximum does not exist.}{}{}}
\LANGpl{
\Question{
Jaka jest wartość funkcji $f$ przy jej lokalnym maksimum?
\[ 
 f(x) = \frac
{\sqrt{6x-x^2}}{2}
\]}
{\(\frac{3}{2}\)}{\(0\)}{\(2\)}{Lokalne maksimum nie istnieje.}{}{}}
\LANGcs{
\Question{
Jaká je hodnota funkce $f$ v jejím lokálním maximu?
\[ 
 f(x) = \frac
{\sqrt{6x-x^2}}{2}
\]}
{\(\frac{3}{2}\)}{\(0\)}{\(2\)}{Lokální maximum funkce neexistuje.}{}{}}
\LANGsk{
\Question{
Akú funkčnú hodnotu, má funkcia $f$ vo svojom lokálnom maxime?
\[ 
 f(x) = \frac
{\sqrt{6x-x^2}}{2}
\]}
{\(\frac{3}{2}\)}{\(0\)}{\(2\)}{Funkcia $f$ nemá lokálne maximum.}{}{}}
\LANGes{
\Question{
Halla el máximo local de la función:
\[ 
 f(x) = \frac
{\sqrt{6x-x^2}}{2}
\]}
{\(\frac{3}{2}\)}{\(0\)}{\(2\)}{El máximo local no existe.}{}{}}
\End

\Data\ID{2010001703}
\Updated{2022-08-25 10:08:26}
\Level{A}
\SubArea{Applications of derivatives}
\LANGen{
\Question{
Find the intervals where the following function is an increasing function.
\[ 
 f(x) = \frac{4+x^{2}} 
 {-4x}
\]}
{\([ - 2;0)\) and
\((0;2] \)}{\([ - 2;2] \)}{\((-\infty ;-2] \) and
\([2;\infty) \)}{\([ 2;\infty) \)}{}{}}
\LANGpl{
\Question{
Znajdź przedziały, w których następująca funkcja jest funkcją rosnącą.
\[ 
 f(x) = \frac{4+x^{2}} 
 {-4x}
\]}
{\(\langle - 2;0)\) i
\((0;2\rangle \)}{\(\langle - 2;2\rangle \)}{\((-\infty ;-2\rangle \) i
\(\langle2;\infty) \)}{\(\langle 2;\infty) \)}{}{}}
\LANGcs{
\Question{
Určete všechny intervaly, na kterých je následující funkce rostoucí.
\[ 
 f(x) = \frac{4+x^{2}} 
 {-4x}
\]}
{\(\langle - 2;0)\) a
\((0;2\rangle \)}{\(\langle - 2;2\rangle \)}{\((-\infty ;-2\rangle \) a
\(\langle2;\infty) \)}{\(\langle 2;\infty) \)}{}{}}
\LANGsk{
\Question{
Je daná funkcia
\[ 
 f(x) = \frac{4+x^{2}} 
 {-4x}.
\] 
V ktorom z nasledujúcich intervalov je táto funkcia rastúca?}
{\(\langle - 2;0)\) a
\((0;2\rangle \)}{\(\langle - 2;2\rangle \)}{\((-\infty ;-2\rangle \) a
\(\langle2;\infty) \)}{\(\langle 2;\infty) \)}{}{}}
\LANGes{
\Question{
Halla todos los intervalos donde la siguiente función es creciente:
\[ 
 f(x) = \frac{4+x^{2}} 
 {-4x}
\]}
{\([ - 2;0)\) y
\((0;2] \)}{\([ - 2;2] \)}{\((-\infty ;-2] \) y
\([2;\infty) \)}{\([ 2;\infty) \)}{}{}}
\End

\Data\ID{2010013909}
\Updated{2023-01-23 10:01:40}
\Level{A}
\SubArea{Applications of derivatives}
\LANGen{
\Question{
Suppose a function is decreasing on the interval \([-5;-2]\). Which one of the offered functions has that property?}
{\(h(x)=\frac{x^4}{2}-x^2-1\)}{\(f(x)=2-\sqrt{x^2+x}\)}{\(g(x)=\frac{x^2+3}{x}\)}{\(i(x)=(x+2)^3-x-1\)}{}{}}
\LANGpl{
\Question{
Załóżmy, że funkcja maleje w przedziale \(\langle-5;-2\rangle\). Która z przedstawionych funkcji ma tę właściwość?}
{\(h(x)=\frac{x^4}{2}-x^2-1\)}{\(f(x)=2-\sqrt{x^2+x}\)}{\(g(x)=\frac{x^2+3}{x}\)}{\(i(x)=(x+2)^3-x-1\)}{}{}}
\LANGcs{
\Question{
O funkci víme, že je v intervalu \(\langle-5;-2\rangle\) klesajicí. Která z nabídnutých funkcí má tuto vlastnost?}
{\(h(x)=\frac{x^4}{2}-x^2-1\)}{\(f(x)=2-\sqrt{x^2+x}\)}{\(g(x)=\frac{x^2+3}{x}\)}{\(i(x)=(x+2)^3-x-1\)}{}{}}
\LANGsk{
\Question{
Vieme, že daná funkcia je na intervale \(\langle-5;-2\rangle\) klesajúca. Ktorá z ponúknutých funkcií má túto vlastnosť?}
{\(h(x)=\frac{x^4}{2}-x^2-1\)}{\(f(x)=2-\sqrt{x^2+x}\)}{\(g(x)=\frac{x^2+3}{x}\)}{\(i(x)=(x+2)^3-x-1\)}{}{}}
\LANGes{
\Question{
Supongamos que una función es decreciente en el intervalo \([-5;-2]\). ¿Cuál de las siguientes funciones cumple esta propiedad?}
{\(h(x)=\frac{x^4}{2}-x^2-1\)}{\(f(x)=2-\sqrt{x^2+x}\)}{\(g(x)=\frac{x^2+3}{x}\)}{\(i(x)=(x+2)^3-x-1\)}{}{}}
\End

\Data\ID{2010013910}
\Updated{2023-01-23 10:01:40}
\Level{A}
\SubArea{Applications of derivatives}
\LANGen{
\Question{
Suppose a function is increasing on the interval \([-6;-3]\). Which one of the offered functions has that property?}
{\(i(x)=-x^4+2x^2-3\)}{\(h(x)=\sqrt{x^2+x}+1\)}{\(f(x)=\frac{1-x^2}{x}\)}{\(g(x)=(x+3)^3-3x-6\)}{}{}}
\LANGpl{
\Question{
Załóżmy, że funkcja rośnie w przedziale \(\langle-6;-3\rangle\). Która z przedstawionych funkcji ma tę właściwość?}
{\(i(x)=-x^4+2x^2-3\)}{\(h(x)=\sqrt{x^2+x}+1\)}{\(f(x)=\frac{1-x^2}{x}\)}{\(g(x)=(x+3)^3-3x-6\)}{}{}}
\LANGcs{
\Question{
O funkci víme, že je v intervalu \(\langle-6;-3\rangle\) rostoucí. Která z nabídnutých funkcí má tuto vlastnost?}
{\(i(x)=-x^4+2x^2-3\)}{\(h(x)=\sqrt{x^2+x}+1\)}{\(f(x)=\frac{1-x^2}{x}\)}{\(g(x)=(x+3)^3-3x-6\)}{}{}}
\LANGsk{
\Question{
Vieme, že daná funkcia je na intervale \(\langle-6;-3\rangle\) rastúca. Ktorá z ponúknutých funkcií má túto vlastnosť?}
{\(i(x)=-x^4+2x^2-3\)}{\(h(x)=\sqrt{x^2+x}+1\)}{\(f(x)=\frac{1-x^2}{x}\)}{\(g(x)=(x+3)^3-3x-6\)}{}{}}
\LANGes{
\Question{
Supongamos que una función es creciente en el intervalo \([-6;-3]\). ¿Cuál de las siguientes funciones cumple esta propiedad?}
{\(i(x)=-x^4+2x^2-3\)}{\(h(x)=\sqrt{x^2+x}+1\)}{\(f(x)=\frac{1-x^2}{x}\)}{\(g(x)=(x+3)^3-3x-6\)}{}{}}
\End

\Data\ID{2010020001}
\Updated{2023-01-23 10:01:09}
\Level{A}
\SubArea{Applications of derivatives}
\LANGen{
\Question{
A function \(f\) is given by the formula \(f(x)=1+\sqrt{\frac{x^4}{4}-2x^2+4}\). Choose the true statement about local extrema of this function.}
{The function has two local minima and one local maximum.}{The function has two local maxima and one local minimum.}{The function has two local minima and no local maximum.}{The function has two local maxima and no local minimum.}{}{}}
\LANGpl{
\Question{
Funkcja \(f\) dana jest wzorem \(f(x)=1+\sqrt{\frac{x^4}{4}-2x^2+4}\). Wybierz prawdziwe stwierdzenie dotyczące ekstremów lokalnych tej funkcji.}
{Funkcja ma dwa lokalne minima i jedno lokalne maksimum.}{Funkcja ma dwa lokalne maksima i jedno lokalne minimum.}{Funkcja ma dwa lokalne minima i nie ma lokalnego maksimum.}{Funkcja ma dwa lokalne maksima i nie ma lokalnego minimum.}{}{}}
\LANGcs{
\Question{
Funkce \(f\) je určena předpisem \(f(x)=1+\sqrt{\frac{x^4}{4}-2x^2+4}\). Vyberte pravdivé tvrzení o extrémech této funkce.}
{Funkce má dvě lokální minima a jedno lokální maximum.}{Funkce má dvě lokální maxima a jedno lokální minimum.}{Funkce má dvě lokální minima a žádné lokální maximum.}{Funkce má dvě lokální maxima a žádné lokální minimum.}{}{}}
\LANGsk{
\Question{
Funkcia \(f\) je daná predpisom \(f(x)=1+\sqrt{\frac{x^4}{4}-2x^2+4}\). Vyberte pravdivé tvrdenie o extrémoch tejto funkcie.}
{Funkcia má dve lokálne minimá a jedno lokálne maximum.}{Funkcia má dve lokálne maximá a jedno lokálne minimum.}{Funkcia má dve lokálne minimá a žiadne lokálne maximum.}{Funkcia má dve lokálne maximá a žiadne lokálne minimum.}{}{}}
\LANGes{
\Question{
La función \(f\) viene dada por la fórmula \(f(x)=1+\sqrt{\frac{x^4}{4}-2x^2+4}\). Elige la proposición lógica verdadera sobre los extremos locales de esta función.}
{La función tiene dos mínimos locales y un máximo local.}{La función tiene dos máximos locales y un mínimo local.}{La función tiene dos mínimos locales y ningún máximo local.}{La función tiene dos máximos locales y ningún mínimo local.}{}{}}
\End

\Data\ID{2010020002}
\Updated{2023-01-23 10:01:09}
\Level{A}
\SubArea{Applications of derivatives}
\LANGen{
\Question{
A function \(f\) is given by the formula \(f(x)=1-\sqrt{\frac{x^4}{4}-4x^2+16}\). Choose the true statement about local extrema of this function.}
{The function has two local maxima and one local minimum.}{The function has two local minima and one local maximum.}{The function has two local minima and no local maximum.}{The function has two local maxima and no local minimum.}{}{}}
\LANGpl{
\Question{
Funkcja \(f\) wyraża się wzorem \(f(x)=1-\sqrt{\frac{x^4}{4}-4x^2+16}\). Wybierz prawdziwe stwierdzenie dotyczące ekstremów lokalnych tej funkcji.}
{Funkcja ma dwa lokalne maksima i jedno lokalne minimum.}{Funkcja ma dwa lokalne minima i jedno lokalne maksimum.}{Funkcja ma dwa lokalne minima i nie ma lokalnego maksimum.}{Funkcja ma dwa lokalne maksima i nie ma lokalnego minimum.}{}{}}
\LANGcs{
\Question{
Funkce \(f\) je určena předpisem \(f(x)=1-\sqrt{\frac{x^4}{4}-4x^2+16}\). Vyberte pravdivé tvrzení o extrémech této funkce.}
{Funkce má dvě lokální maxima a jedno lokální minimum.}{Funkce má dvě lokální minima a jedno lokální maximum.}{Funkce má dvě lokální minima a žádné lokální maximum.}{Funkce má dvě lokální maxima a žádné lokální minimum.}{}{}}
\LANGsk{
\Question{
Funkcia \(f\) je daná predpisom \(f(x)=1-\sqrt{\frac{x^4}{4}-4x^2+16}\). Vyberte pravdivé tvrdenie o extrémoch tejto funkcie.}
{Funkcia má dve lokálne maximá a jedno lokálne minimum.}{Funkcia má dve lokálne minimá a jedno lokálne maximum.}{Funkcia má dve lokálne minimá a žiadne lokálne maximum.}{Funkcia má dve lokálne maximá a žiadne lokálne minimum.}{}{}}
\LANGes{
\Question{
La función \(f\) viene dada por la fórmula \(f(x)=1-\sqrt{\frac{x^4}{4}-4x^2+16}\). Elige la proposición lógica verdadera sobre los extremos locales de esta función.}
{La función tiene dos máximos locales y un mínimo local.}{La función tiene dos mínimos locales y un máximo local.}{La función tiene dos mínimos locales y ningún máximo local.}{La función tiene dos máximos locales y ningún mínimo local.}{}{}}
\End

\Data\ID{2110013901}
\Updated{2023-01-23 10:01:12}
\Level{A}
\SubArea{Applications of derivatives}
\LANGen{
\Question{
The pictures show parts of graphs of functions that are increasing on the interval \([1;5]\). Choose the picture showing the part of the graph of the function \(f(x)=\frac{5x-1}{x+1}\).}
{}{}{}{}{}{}}
\LANGpl{
\Question{
Rysunki przedstawiają fragmenty wykresów funkcji, które rosną na przedziale \(\langle1;5\rangle\). Wybierz rysunek przedstawiający część wykresu funkcji \(f(x)=\frac{5x-1}{x+1}\).}
{}{}{}{}{}{}}
\LANGcs{
\Question{
Na obrázcích jsou části grafů funkcí, které jsou v intervalu \(\langle1;5\rangle\) rostoucí. Vyberte obrázek, na kterém je část grafu funkce \(f(x)=\frac{5x-1}{x+1}\).}
{}{}{}{}{}{}}
\LANGsk{
\Question{
Na obrázkoch sú časti grafou funkcií, ktoré sú na intervale \(\langle1;5\rangle\) rastúce. Vyberte obrázok, na ktorom je čásť grafu funkcie \(f(x)=\frac{5x-1}{x+1}\).}
{}{}{}{}{}{}}
\LANGes{
\Question{
Las imágenes muestran una parte de gráficas de funciones que son crecientes en el intervalo \([1;5]\). Elige la imagen que muestra esta parte de la gráfica de la función \(f(x)=\frac{5x-1}{x+1}\).}
{}{}{}{}{}{}}

\IMGanswer{obr1.pdf}
\IMGanswer{obr2_1.pdf}
\IMGanswer{obr3_3.pdf}
\IMGanswer{obr4_2.pdf}\End

\Data\ID{2110013902}
\Updated{2023-01-23 10:01:12}
\Level{A}
\SubArea{Applications of derivatives}
\LANGen{
\Question{
The pictures show parts of graphs of functions that are decreasing on the interval \([1;5]\). Choose the picture showing the part of the graph of the function \(f(x)=\frac{x+7}{x+1}\).}
{}{}{}{}{}{}}
\LANGpl{
\Question{
Rysunki przedstawiają fragmenty wykresów funkcji malejących na przedziale  \(\langle1;5\rangle\). Wybierz rysunek przedstawiający część wykresu funkcj \(f(x)=\frac{x+7}{x+1}\).}
{}{}{}{}{}{}}
\LANGcs{
\Question{
Na obrázcích jsou části grafů funkcí, které jsou v intervalu \(\langle1;5\rangle\) klesajicí. Vyberte obrázek, na kterém je část grafu funkce \(f(x)=\frac{x+7}{x+1}\).}
{}{}{}{}{}{}}
\LANGsk{
\Question{
Na obrázkoch sú časti grafou funkcií, ktoré sú na intervale \(\langle1;5\rangle\) klesajúce. Vyberte obrázok, na ktorom je čásť grafu funkcie \(f(x)=\frac{x+7}{x+1}\).}
{}{}{}{}{}{}}
\LANGes{
\Question{
Las imágenes muestran una parte de gráficas de funciones que son decrecientes en el intervalo \([1;5]\). Elige la imagen que muestra esta parte de la gráfica de la función \(f(x)=\frac{x+7}{x+1}\).}
{}{}{}{}{}{}}

\IMGanswer{obr1_4.pdf}
\IMGanswer{obr2_3.pdf}
\IMGanswer{obr3_5.pdf}
\IMGanswer{obr4_3.pdf}\End

\Data\ID{2110020003}
\Updated{2023-01-23 10:01:16}
\Level{A}
\SubArea{Applications of derivatives}
\LANGen{
\Question{
A function \(f\) is given by the formula \(f(x)=(x+3)^3-3x-6\). Choose the picture that shows a part of the graph of this function.}
{}{}{}{}{}{}}
\LANGpl{
\Question{
Funkcja \(f\) wyrażona jest wzorem \(f(x)=(x+3)^3-3x-6\). Wybierz obrazek przedstawiający część wykresu tej funkcji.}
{}{}{}{}{}{}}
\LANGcs{
\Question{
Funkce \(f\) je dána předpisem \(f(x)=(x+3)^3-3x-6\). Vyberte obrázek, na kterém by mohla být část grafu této funkce.}
{}{}{}{}{}{}}
\LANGsk{
\Question{
Funkcia \(f\) je daná predpisom \(f(x)=(x+3)^3-3x-6\). Vyberte obrázok, na ktorom je časť grafu tejto funkcie.}
{}{}{}{}{}{}}
\LANGes{
\Question{
La función \(f\) viene dada por la fórmula \(f(x)=(x+3)^3-3x-6\). Elige la imagen que representa una parte de la gráfica de esta función.}
{}{}{}{}{}{}}

\IMGanswer{obr1_0.pdf}
\IMGanswer{obr2.pdf}
\IMGanswer{obr3.pdf}
\IMGanswer{obr4.pdf}\End

\Data\ID{2110020004}
\Updated{2023-01-23 10:01:16}
\Level{A}
\SubArea{Applications of derivatives}
\LANGen{
\Question{
A function \(f\) is given by the formula \(f(x)=-(x-3)^3+3x-12\). Choose the picture that shows a part of the graph of this function.}
{}{}{}{}{}{}}
\LANGpl{
\Question{
Funkcja \(f\) przedstawiona jest wzorem  \(f(x)=-(x-3)^3+3x-12\). Wybierz obrazek przedstawiający część wykresu tej funkcji.}
{}{}{}{}{}{}}
\LANGcs{
\Question{
Funkce \(f\) je dána předpisem \(f(x)=-(x+3)^3+3x-12\). Vyberte obrázek, na kterém by mohla být část grafu této funkce.}
{}{}{}{}{}{}}
\LANGsk{
\Question{
Funkcia \(f\) je daná predpisom \(f(x)=-(x-3)^3+3x-12\). Vyberte obrázok, na ktorom je časť grafu tejto funkcie.}
{}{}{}{}{}{}}
\LANGes{
\Question{
La función \(f\) viene dada por la fórmula \(f(x)=-(x-3)^3+3x-12\). Elige la imagen que representa una parte de la gráfica de esta función.}
{}{}{}{}{}{}}

\IMGanswer{obr1_1.pdf}
\IMGanswer{obr2_0.pdf}
\IMGanswer{obr3_0.pdf}
\IMGanswer{obr4_0.pdf}\End

\Data\ID{9000062409}
\Updated{2021-08-15 08:08:10}
\Level{A}
\SubArea{Applications of derivatives}
\LANGen{
\Question{
Find the first derivative of the function
\(f(x) = x^{2} + x - 6\) at each of the
\(x\)-intercepts.}
{\(- 5;\ 5\)}{\(- 3;\ 7\)}{\(- 7;\ 6\)}{\(- 9;\ 6\)}{}{}}
\LANGpl{
\Question{
Wyznacz pochodną funkcji
\(f\colon y = x^{2} + x - 6\) w każdym punkcie przecięcia z osią 
\(x\).}
{\(- 5;\ 5\)}{\(- 3;\ 7\)}{\(- 7;\ 6\)}{\(- 9;\ 6\)}{}{}}
\LANGcs{
\Question{
Vypočtěte hodnoty první derivace funkce
\(f\colon y = x^{2} + x - 6\) v průsečících
jejího grafu s osou \(x\).}
{\(- 5;\ 5\)}{\(- 3;\ 7\)}{\(- 7;\ 6\)}{\(- 9;\ 6\)}{}{}}
\LANGsk{
\Question{
Vypočítajte hodnoty prvej derivácie funkcie
\(f\colon y = x^{2} + x - 6\) v priesečníkoch
jej grafu s osou \(x\).}
{\(- 5;\ 5\)}{\(- 3;\ 7\)}{\(- 7;\ 6\)}{\(- 9;\ 6\)}{}{}}
\LANGes{
\Question{
Halla la primera derivada de la función
\(f(x) = x^{2} + x - 6\) en cada uno de los puntos de corte con el eje 
\(x\).}
{\(- 5;\ 5\)}{\(- 3;\ 7\)}{\(- 7;\ 6\)}{\(- 9;\ 6\)}{}{}}
\End

\Data\ID{9000070401}
\Updated{2022-04-23 05:04:57}
\Level{A}
\SubArea{Applications of derivatives}
\LANGen{
\Question{
Given the function \(f(x) = x^{2} + x - 2\), find
the intervals where \(f\) 
is an increasing function.}
{\(\left (-\frac{1}
{2};\infty \right )\)}{\(\left (-3;\infty \right )\)}{\(\left (-2;\infty \right )\)}{\(\left (-1;\infty \right )\)}{}{}}
\LANGpl{
\Question{
Dana jest funkcja  \(f\colon y = x^{2} + x - 2\), wskaż przedziały, gdzie funkcja 
 \(f\) 
jest funkcją rosnącą.}
{\(\left (-\frac{1}
{2};\infty \right )\)}{\(\left (-3;\infty \right )\)}{\(\left (-2;\infty \right )\)}{\(\left (-1;\infty \right )\)}{}{}}
\LANGcs{
\Question{
Je dána funkce \(f\colon y = x^{2} + x - 2\).
Ve kterém z následujících intervalů je tato funkce rostoucí?}
{\(\left (-\frac{1}
{2};\infty \right )\)}{\(\left (-3;\infty \right )\)}{\(\left (-2;\infty \right )\)}{\(\left (-1;\infty \right )\)}{}{}}
\LANGsk{
\Question{
Je daná funkcia \(f\colon y = x^{2} + x - 2\).
V ktorom z následujúcich intervalov je táto funkcia rastúca?}
{\(\left (-\frac{1}
{2};\infty \right )\)}{\(\left (-3;\infty \right )\)}{\(\left (-2;\infty \right )\)}{\(\left (-1;\infty \right )\)}{}{}}
\LANGes{
\Question{
Dada la función \(f(x) = x^{2} + x - 2\), halla el intervalo donde \(f\) 
es creciente.}
{\(\left (-\frac{1}
{2};\infty \right )\)}{\(\left (-3;\infty \right )\)}{\(\left (-2;\infty \right )\)}{\(\left (-1;\infty \right )\)}{}{}}
\End

\Data\ID{9000070402}
\Updated{2021-08-12 11:08:28}
\Level{A}
\SubArea{Applications of derivatives}
\LANGen{
\Question{
Given the function \(f(x)= x^{2} - 2x - 8\), find
the intervals where \(f\) 
is a decreasing function.}
{\(\left (-\infty ;1\right )\)}{\(\left (-\infty ;8\right )\)}{\(\left (-\infty ;2\right )\)}{\(\left (-\infty ;4\right )\)}{}{}}
\LANGpl{
\Question{
Dana jest funkcja  \(f\colon y = x^{2} - 2x - 8\), wskaż przedziały, gdzie funkcja 
 \(f\) 
jest funkcją malejącą.}
{\(\left (-\infty ;1\right )\)}{\(\left (-\infty ;8\right )\)}{\(\left (-\infty ;2\right )\)}{\(\left (-\infty ;4\right )\)}{}{}}
\LANGcs{
\Question{
Je dána funkce \(f\colon y = x^{2} - 2x - 8\).
Ve kterém z následujících intervalů je tato funkce klesající?}
{\(\left (-\infty ;1\right )\)}{\(\left (-\infty ;8\right )\)}{\(\left (-\infty ;2\right )\)}{\(\left (-\infty ;4\right )\)}{}{}}
\LANGsk{
\Question{
Je daná funkcia \(f\colon y = x^{2} - 2x - 8\).
V ktorom z následujúcich intervalov je táto funkcia klesajúca?}
{\(\left (-\infty ;1\right )\)}{\(\left (-\infty ;8\right )\)}{\(\left (-\infty ;2\right )\)}{\(\left (-\infty ;4\right )\)}{}{}}
\LANGes{
\Question{
Dada la función \(f(x)= x^{2} - 2x - 8\), halla el intervalo donde \(f\) 
es decreciente.}
{\(\left (-\infty ;1\right )\)}{\(\left (-\infty ;8\right )\)}{\(\left (-\infty ;2\right )\)}{\(\left (-\infty ;4\right )\)}{}{}}
\End

\Data\ID{9000070403}
\Updated{2021-08-12 11:08:01}
\Level{A}
\SubArea{Applications of derivatives}
\LANGen{
\Question{
Given function \(f(x) = -x^{2} + 2x + 3\), find
the intervals where \(f\) 
is an increasing function.}
{\(\left (-\infty ;1\right )\)}{\(\left (-\infty ;2\right )\)}{\(\left (-\infty ;3\right )\)}{\(\left (-\infty ;6\right )\)}{}{}}
\LANGpl{
\Question{
Dana jest funkcja \(f\colon y = -x^{2} + 2x + 3\), wskaż przedziały, gdzie funkcja 
 \(f\) 
jest funkcją rosnącą.}
{\(\left (-\infty ;1\right )\)}{\(\left (-\infty ;2\right )\)}{\(\left (-\infty ;3\right )\)}{\(\left (-\infty ;6\right )\)}{}{}}
\LANGcs{
\Question{
Je dána funkce \(f\colon y = -x^{2} + 2x + 3\).
Ve kterém z následujících intervalů je tato funkce rostoucí?}
{\(\left (-\infty ;1\right )\)}{\(\left (-\infty ;2\right )\)}{\(\left (-\infty ;3\right )\)}{\(\left (-\infty ;6\right )\)}{}{}}
\LANGsk{
\Question{
Je daná funkcia \(f\colon y = -x^{2} + 2x + 3\).
V ktorom z následujúcich intervalov je táto funkcia rastúca?}
{\(\left (-\infty ;1\right )\)}{\(\left (-\infty ;2\right )\)}{\(\left (-\infty ;3\right )\)}{\(\left (-\infty ;6\right )\)}{}{}}
\LANGes{
\Question{
Dada la función \(f(x) = -x^{2} + 2x + 3\), halla el intervalo donde \(f\) 
es creciente.}
{\(\left (-\infty ;1\right )\)}{\(\left (-\infty ;2\right )\)}{\(\left (-\infty ;3\right )\)}{\(\left (-\infty ;6\right )\)}{}{}}
\End

\Data\ID{9000070404}
\Updated{2021-08-12 11:08:29}
\Level{A}
\SubArea{Applications of derivatives}
\LANGen{
\Question{
Given function \(f(x)= -x^{2} + 4x + 12\), find
the intervals where \(f\) 
is a decreasing function.}
{\(\left (2;\infty \right )\)}{\(\left (1;\infty \right )\)}{\(\left (-4;\infty \right )\)}{\(\left (-6;\infty \right )\)}{}{}}
\LANGpl{
\Question{
Dana jest funkcja  \(f\colon y = -x^{2} + 4x + 12\), wskaż przedziały, gdzie funkcja 
 \(f\) 
jest funkcją malejącą.}
{\(\left (2;\infty \right )\)}{\(\left (1;\infty \right )\)}{\(\left (-4;\infty \right )\)}{\(\left (-6;\infty \right )\)}{}{}}
\LANGcs{
\Question{
Je dána funkce \(f\colon y = -x^{2} + 4x + 12\).
Ve kterém z následujících intervalů je tato funkce klesající?}
{\(\left (2;\infty \right )\)}{\(\left (1;\infty \right )\)}{\(\left (-4;\infty \right )\)}{\(\left (-6;\infty \right )\)}{}{}}
\LANGsk{
\Question{
Je daná funkcia \(f\colon y = -x^{2} + 4x + 12\).
V ktorom z následujúcich intervalov je táto funkcia klesajúca?}
{\(\left (2;\infty \right )\)}{\(\left (1;\infty \right )\)}{\(\left (-4;\infty \right )\)}{\(\left (-6;\infty \right )\)}{}{}}
\LANGes{
\Question{
Dada la función \(f(x)= -x^{2} + 4x + 12\), halla el intervalo donde \(f\) 
es decreciente.}
{\(\left (2;\infty \right )\)}{\(\left (1;\infty \right )\)}{\(\left (-4;\infty \right )\)}{\(\left (-6;\infty \right )\)}{}{}}
\End

\Data\ID{9000070405}
\Updated{2021-08-12 11:08:27}
\Level{A}
\SubArea{Applications of derivatives}
\LANGen{
\Question{
Given function \(f(x)= x^{3} + 3x^{2} - 24x + 5\), find
the intervals where \(f\) 
is a decreasing function.}
{\(\left (-4;2\right )\)}{\(\left (-3;5\right )\)}{\(\left (-3;3\right )\)}{\(\left (-5;1\right )\)}{}{}}
\LANGpl{
\Question{
Dana jest funkcja  \(f\colon y = x^{3} + 3x^{2} - 24x + 5\), wskaż przedziały, gdzie funkcja 
 \(f\) 
jest funkcją malejącą.}
{\(\left (-4;2\right )\)}{\(\left (-3;5\right )\)}{\(\left (-3;3\right )\)}{\(\left (-5;1\right )\)}{}{}}
\LANGcs{
\Question{
Je dána funkce \(f\colon y = x^{3} + 3x^{2} - 24x + 5\).
Ve kterém z následujících intervalů je tato funkce klesající?}
{\(\left (-4;2\right )\)}{\(\left (-3;5\right )\)}{\(\left (-3;3\right )\)}{\(\left (-5;1\right )\)}{}{}}
\LANGsk{
\Question{
Je daná funkcia \(f\colon y = x^{3} + 3x^{2} - 24x + 5\).
V ktorom z následujúcich intervalov je táto funkcia klesajúca?}
{\(\left (-4;2\right )\)}{\(\left (-3;5\right )\)}{\(\left (-3;3\right )\)}{\(\left (-5;1\right )\)}{}{}}
\LANGes{
\Question{
Dada la función \(f(x)= x^{3} + 3x^{2} - 24x + 5\), halla el intervalo donde \(f\) 
es decreciente.}
{\(\left (-4;2\right )\)}{\(\left (-3;5\right )\)}{\(\left (-3;3\right )\)}{\(\left (-5;1\right )\)}{}{}}
\End

\Data\ID{9000070406}
\Updated{2021-08-12 11:08:50}
\Level{A}
\SubArea{Applications of derivatives}
\LANGen{
\Question{
Given function \(f(x)= x^{3} + 6x^{2} - 15x + 7\), find
the intervals where \(f\) 
is an increasing function.}
{\(\left (-\infty ;-5\right )\)}{\(\left (-\infty ;-3\right )\)}{\(\left (-1;\infty \right )\)}{\(\left (-3;\infty \right )\)}{}{}}
\LANGpl{
\Question{
Dana jest funkcja  \(f\colon y = x^{3} + 6x^{2} - 15x + 7\), wskaż przedziały, gdzie funkcja 
 \(f\) 
jest funkcją rosnącą.}
{\(\left (-\infty ;-5\right )\)}{\(\left (-\infty ;-3\right )\)}{\(\left (-1;\infty \right )\)}{\(\left (-3;\infty \right )\)}{}{}}
\LANGcs{
\Question{
Je dána funkce \(f\colon y = x^{3} + 6x^{2} - 15x + 7\).
Ve kterém z následujících intervalů je tato funkce rostoucí?}
{\(\left (-\infty ;-5\right )\)}{\(\left (-\infty ;-3\right )\)}{\(\left (-1;\infty \right )\)}{\(\left (-3;\infty \right )\)}{}{}}
\LANGsk{
\Question{
Je daná funkcia \(f\colon y = x^{3} + 6x^{2} - 15x + 7\).
V ktorom z následujúcich intervalov je táto funkcia rastúca?}
{\(\left (-\infty ;-5\right )\)}{\(\left (-\infty ;-3\right )\)}{\(\left (-1;\infty \right )\)}{\(\left (-3;\infty \right )\)}{}{}}
\LANGes{
\Question{
Dada la función \(f(x)= x^{3} + 6x^{2} - 15x + 7\), halla el intervalo donde \(f\) 
es creciente.}
{\(\left (-\infty ;-5\right )\)}{\(\left (-\infty ;-3\right )\)}{\(\left (-1;\infty \right )\)}{\(\left (-3;\infty \right )\)}{}{}}
\End

\Data\ID{9000070407}
\Updated{2021-08-12 11:08:18}
\Level{A}
\SubArea{Applications of derivatives}
\LANGen{
\Question{
Given the function \(f(x) = -x^{3} + 3x^{2} + 9x - 1\), find
the intervals where \(f\) 
is a decreasing function.}
{\(\left (3;\infty \right )\)}{\(\left (-\infty ;1\right )\)}{\(\left (-1;3\right )\)}{\(\left (1;\infty \right )\)}{}{}}
\LANGpl{
\Question{
Dana jest funkcja  \(f\colon y = -x^{3} + 3x^{2} + 9x - 1\), wskaż przedziały, gdzie funkcja 
 \(f\) 
jest funkcją malejącą.}
{\(\left (3;\infty \right )\)}{\(\left (-\infty ;1\right )\)}{\(\left (-1;3\right )\)}{\(\left (1;\infty \right )\)}{}{}}
\LANGcs{
\Question{
Je dána funkce \(f\colon y = -x^{3} + 3x^{2} + 9x - 1\).
Ve kterém z následujících intervalů je tato funkce klesající?}
{\(\left (3;\infty \right )\)}{\(\left (-\infty ;1\right )\)}{\(\left (-1;3\right )\)}{\(\left (1;\infty \right )\)}{}{}}
\LANGsk{
\Question{
Je daná funkcia \(f\colon y = -x^{3} + 3x^{2} + 9x - 1\).
V ktorom z následujúcich intervalov je táto funkcia klesajúca?}
{\(\left (3;\infty \right )\)}{\(\left (-\infty ;1\right )\)}{\(\left (-1;3\right )\)}{\(\left (1;\infty \right )\)}{}{}}
\LANGes{
\Question{
Dada la función \(f(x) = -x^{3} + 3x^{2} + 9x - 1\), halla el intervalo donde \(f\) 
es decreciente.}
{\(\left (3;\infty \right )\)}{\(\left (-\infty ;1\right )\)}{\(\left (-1;3\right )\)}{\(\left (1;\infty \right )\)}{}{}}
\End

\Data\ID{9000070408}
\Updated{2022-04-23 05:04:57}
\Level{A}
\SubArea{Applications of derivatives}
\LANGen{
\Question{
Given the function \(f(x) = -x^{3} + 3x^{2} + 45x - 12\), find
the intervals where \(f\) 
is an increasing function.}
{\(\left (-3;5\right )\)}{\(\left (-\infty ;-3\right )\)}{\(\left (5;\infty \right )\)}{\(\left (-12;45\right )\)}{}{}}
\LANGpl{
\Question{
Dana jest funkcja  \(f\colon y = -x^{3} + 3x^{2} + 45x - 12\), wskaż przedziały, gdzie funkcja 
 \(f\) 
jest funkcją rosnącą.}
{\(\left (-3;5\right )\)}{\(\left (-\infty ;-3\right )\)}{\(\left (5;\infty \right )\)}{\(\left (-12;45\right )\)}{}{}}
\LANGcs{
\Question{
Je dána funkce \(f\colon y = -x^{3} + 3x^{2} + 45x - 12\).
Ve kterém z následujících intervalů je tato funkce rostoucí?}
{\(\left (-3;5\right )\)}{\(\left (-\infty ;-3\right )\)}{\(\left (5;\infty \right )\)}{\(\left (-12;45\right )\)}{}{}}
\LANGsk{
\Question{
Je daná funkcia \(f\colon y = -x^{3} + 3x^{2} + 45x - 12\).
V ktorom z následujúcich intervalov je táto funkcia rastúca?}
{\(\left (-3;5\right )\)}{\(\left (-\infty ;-3\right )\)}{\(\left (5;\infty \right )\)}{\(\left (-12;45\right )\)}{}{}}
\LANGes{
\Question{
Dada la función \(f(x) = -x^{3} + 3x^{2} + 45x - 12\), halla el intervalo donde \(f\) 
es creciente.}
{\(\left (-3;5\right )\)}{\(\left (-\infty ;-3\right )\)}{\(\left (5;\infty \right )\)}{\(\left (-12;45\right )\)}{}{}}
\End

\Data\ID{9000070409}
\Updated{2022-04-23 05:04:57}
\Level{A}
\SubArea{Applications of derivatives}
\LANGen{
\Question{
Given the function \(f(x) = \frac{x^{2}} 
{x-2}\), find
the intervals where \(f\) 
is a decreasing function.}
{\(\left (0;2\right )\)}{\(\left (-\infty ;-1\right )\)}{\(\left (5;\infty \right )\)}{\(\left (2;5\right )\)}{}{}}
\LANGpl{
\Question{
Dana jest funkcja  \(f\colon y = \frac{x^{2}} 
{x-2}\), wskaż przedziały, gdzie funkcja 
 \(f\) 
jest funkcją malejącą.}
{\(\left (0;2\right )\)}{\(\left (-\infty ;-1\right )\)}{\(\left (5;\infty \right )\)}{\(\left (2;5\right )\)}{}{}}
\LANGcs{
\Question{
Je dána funkce \(f\colon y = \frac{x^{2}} 
{x-2}\).
Ve kterém z následujících intervalů je tato funkce klesající?}
{\(\left (0;2\right )\)}{\(\left (-\infty ;-1\right )\)}{\(\left (5;\infty \right )\)}{\(\left (2;5\right )\)}{}{}}
\LANGsk{
\Question{
Je daná funkcia \(f\colon y = \frac{x^{2}} 
{x-2}\).
V ktorom z nasledujúcich intervalov je táto funkcia klesajúca?}
{\(\left (0;2\right )\)}{\(\left (-\infty ;-1\right )\)}{\(\left (5;\infty \right )\)}{\(\left (2;5\right )\)}{}{}}
\LANGes{
\Question{
Dada la función \(f(x) = \frac{x^{2}} 
{x-2}\), halla el intervalo donde \(f\) 
es decreciente.}
{\(\left (0;2\right )\)}{\(\left (-\infty ;-1\right )\)}{\(\left (5;\infty \right )\)}{\(\left (2;5\right )\)}{}{}}
\End

\Data\ID{9000070410}
\Updated{2021-08-12 11:08:40}
\Level{A}
\SubArea{Applications of derivatives}
\LANGen{
\Question{
Given the function \(f(x) = - \frac{x^{2}} 
{x+3}\), find
the intervals where \(f\) 
is an increasing function.}
{\(\left (-3;0\right )\)}{\(\left (-\infty ;-6\right )\)}{\(\left (0;\infty \right )\)}{\(\left (-3;4\right )\)}{}{}}
\LANGpl{
\Question{
Dana jest funkcja  \(f\colon y = - \frac{x^{2}} 
{x+3}\), wskaż przedziały, gdzie funkcja 
 \(f\) 
jest funkcją rosnącą.}
{\(\left (-3;0\right )\)}{\(\left (-\infty ;-6\right )\)}{\(\left (0;\infty \right )\)}{\(\left (-3;4\right )\)}{}{}}
\LANGcs{
\Question{
Je dána funkce \(f\colon y = - \frac{x^{2}} 
{x+3}\).
Ve kterém z následujících intervalů je tato funkce rostoucí?}
{\(\left (-3;0\right )\)}{\(\left (-\infty ;-6\right )\)}{\(\left (0;\infty \right )\)}{\(\left (-3;4\right )\)}{}{}}
\LANGsk{
\Question{
Je daná funkcia \(f\colon y = - \frac{x^{2}} 
{x+3}\).
V ktorom z následujúcich intervalov je táto funkcia rastúca?}
{\(\left (-3;0\right )\)}{\(\left (-\infty ;-6\right )\)}{\(\left (0;\infty \right )\)}{\(\left (-3;4\right )\)}{}{}}
\LANGes{
\Question{
Dada la función \(f(x) = - \frac{x^{2}} 
{x+3}\), halla el intervalo donde \(f\) 
es creciente.}
{\(\left (-3;0\right )\)}{\(\left (-\infty ;-6\right )\)}{\(\left (0;\infty \right )\)}{\(\left (-3;4\right )\)}{}{}}
\End

\Data\ID{9000079101}
\Updated{2020-07-15 03:07:37}
\Level{A}
\SubArea{Applications of derivatives}
\LANGen{
\Question{
Find the intervals of monotonicity for the following function. 
\[ 
 f(x)= \frac{3x + 1} 
{2x - 5}
\]}
{Decreasing on \(\left (-\infty ; \frac{5} 
{2}\right )\) 
and \(\left (\frac{5}
{2};\infty \right )\).}{Decreasing on \(\left (-\infty ; \frac{5} 
{2}\right )\cup \left (\frac{5}
{2};\infty \right )\).}{Decreasing on \(\left (-\infty ; \frac{5} 
{2}\right )\),
increasing on \(\left (\frac{5}
{2};\infty \right )\).}{Increasing on \(\left (-\infty ; \frac{5} 
{2}\right )\),
decreasing on \(\left (\frac{5}
{2};\infty \right )\).}{}{}}
\LANGpl{
\Question{
Wyznacz przedziały monotoniczności podanej  funkcji.
\[ 
 f\colon y = \frac{3x + 1} 
{2x - 5}
\]}
{Funkcja jest malejąca w przedziale \(\left (-\infty ; \frac{5} 
{2}\right )\) 
i \(\left (\frac{5}
{2};\infty \right )\).}{Funkcja jest malejąca w przedziale  \(\left (-\infty ; \frac{5} 
{2}\right )\cup \left (\frac{5}
{2};\infty \right )\).}{Funkcja jest malejąca w przedziale  \(\left (-\infty ; \frac{5} 
{2}\right )\),
rosnąca w przedziale  \(\left (\frac{5}
{2};\infty \right )\).}{Funkcja jest rosnąca  w przedziale  \(\left (-\infty ; \frac{5} 
{2}\right )\),
malejąca w przedziale \(\left (\frac{5}
{2};\infty \right )\).}{}{}}
\LANGcs{
\Question{
Určete intervaly monotonie funkce \(f\colon y = \frac{3x+1} 
{2x-5}\):}
{Funkce \(f\) je klesající
na intervalech \(\left (-\infty ; \frac{5} 
{2}\right )\) 
a \(\left (\frac{5}
{2};\infty \right )\).}{Funkce \(f\) je klesající
na množině \(\left (-\infty ; \frac{5} 
{2}\right )\cup \left (\frac{5}
{2};\infty \right )\).}{Funkce \(f\) je klesající
na intervalu \(\left (-\infty ; \frac{5} 
{2}\right )\),
rostoucí na intervalu \(\left (\frac{5}
{2};\infty \right )\).}{Funkce \(f\) je rostoucí na
intervalu \(\left (-\infty ; \frac{5} 
{2}\right )\), klesající
na intervalu \(\left (\frac{5}
{2};\infty \right )\).}{}{}}
\LANGsk{
\Question{
Určte intervaly monotónnosti funkcie \(f\colon y = \frac{3x+1} 
{2x-5}\):}
{Funkcia \(f\) je klesajúca
na intervaloch \(\left (-\infty ; \frac{5} 
{2}\right )\) 
a \(\left (\frac{5}
{2};\infty \right )\).}{Funkcia \(f\) je klesajúca
na množine \(\left (-\infty ; \frac{5} 
{2}\right )\cup \left (\frac{5}
{2};\infty \right )\).}{Funkcia \(f\) je klesajúca
na intervale \(\left (-\infty ; \frac{5} 
{2}\right )\),
rastúca na intervale \(\left (\frac{5}
{2};\infty \right )\).}{Funkcia \(f\) je rastúca na
intervale \(\left (-\infty ; \frac{5} 
{2}\right )\), klesajúca
na intervale \(\left (\frac{5}
{2};\infty \right )\).}{}{}}
\LANGes{
\Question{
Halla los intervalos de monotonía de la siguiente función:
\[ 
 f(x)= \frac{3x + 1} 
{2x - 5}
\]}
{Decreciente en \(\left (-\infty ; \frac{5} 
{2}\right )\) 
y \(\left (\frac{5}
{2};\infty \right )\).}{Decreciente en \(\left (-\infty ; \frac{5} 
{2}\right )\cup \left (\frac{5}
{2};\infty \right )\).}{Decreciente en \(\left (-\infty ; \frac{5} 
{2}\right )\),
creciente en \(\left (\frac{5}
{2};\infty \right )\).}{Creciente en \(\left (-\infty ; \frac{5} 
{2}\right )\),
decreciente en \(\left (\frac{5}
{2};\infty \right )\).}{}{}}
\End

\Data\ID{9000079102}
\Updated{2022-04-23 05:04:08}
\Level{A}
\SubArea{Applications of derivatives}
\LANGen{
\Question{
Find the intervals where the following function is a decreasing function.
\[ 
 f(x) = \frac{x^{2} + 1} 
 {x}
\]}
{\([ - 1;0)\) and
\((0;1] \)}{\([ - 1;1] \)}{\((-\infty ;-1] \) and
\([1;\infty) \)}{\([1;\infty) \)}{}{}}
\LANGpl{
\Question{
Wskaż przedziały, gdzie funkcja
\[ 
 f\colon y = \frac{x^{2} + 1} 
 {x}
\] jest funkcją malejącą.}
{\([ - 1;0)\) i
\((0;1] \)}{\([ - 1;1] \)}{\((-\infty ;-1] \) i
\([1;\infty) \)}{\([1;\infty) \)}{}{}}
\LANGcs{
\Question{
Určete všechny intervaly, na kterých je funkce
\(f\colon y = \frac{x^{2}+1} 
 {x} \) 
klesající:}
{\(\langle - 1;0)\) a
\((0;1\rangle \)}{\(\langle - 1;1\rangle \)}{\((-\infty ;-1\rangle \) a
\(\langle1;\infty) \)}{\(\langle1;\infty) \)}{}{}}
\LANGsk{
\Question{
Určte všetky intervaly, na ktorých je funkcia
\(f\colon y = \frac{x^{2}+1} 
 {x} \) 
klesajúca:}
{\(\langle - 1;0)\) a
\((0;1\rangle \)}{\(\langle - 1;1\rangle \)}{\((-\infty ;-1\rangle \) a
\(\langle1;\infty) \)}{\(\langle1;\infty) \)}{}{}}
\LANGes{
\Question{
Halla todos los intervalos donde la siguiente función es decreciente.
\[ 
 f(x) = \frac{x^{2} + 1} 
 {x}
\]}
{\([ - 1;0)\) y
\((0;1] \)}{\([ - 1;1] \)}{\((-\infty ;-1] \) y
\([1;\infty) \)}{\([1;\infty) \)}{}{}}
\End

\Data\ID{9000079103}
\Updated{2020-07-15 03:07:37}
\Level{A}
\SubArea{Applications of derivatives}
\LANGen{
\Question{
Find the $x$ at which has the function $f$ a local maximum. 
\[ 
 f(x) = x^{3} - 3x^{2} - 9x + 2
\]}
{\(x=- 1\)}{\(x=- 3\)}{\(x=1\)}{\(x=3\)}{}{}}
\LANGpl{
\Question{
Wyznacz maksimum lokalne funkcji.
\[ 
 f\colon y = x^{3} - 3x^{2} - 9x + 2
\]}
{\(x=- 1\)}{\(x=- 3\)}{\(x=1\)}{\(x=3\)}{}{}}
\LANGcs{
\Question{
Funkce \(f\colon y = x^{3} - 3x^{2} - 9x + 2\) 
má lokální maximum v bodě:}
{\(x=- 1\)}{\(x=- 3\)}{\(x=1\)}{\(x=3\)}{}{}}
\LANGsk{
\Question{
Funkcia \(f\colon y = x^{3} - 3x^{2} - 9x + 2\) 
má lokálne maximum v bode:}
{\(x=- 1\)}{\(x=- 3\)}{\(x=1\)}{\(x=3\)}{}{}}
\LANGes{
\Question{
Halla el máximo local de la siguiente función:
\[ 
 f(x) = x^{3} - 3x^{2} - 9x + 2
\]}
{\(x=- 1\)}{\(x=- 3\)}{\(x=1\)}{\(x=3\)}{}{}}
\End

\Data\ID{9000079104}
\Updated{2020-07-15 03:07:37}
\Level{A}
\SubArea{Applications of derivatives}
\LANGen{
\Question{
Find the $x$ at which has the function $f$ a local minimum.
\[ 
 f(x) = \frac{\ln x} 
{x}
\]}
{does not exist}{\(x = 0\)}{\(x = 1\)}{\(x =\mathrm{e}\)}{}{}}
\LANGpl{
\Question{
Wyznacz minimum  lokalne funkcji.
\[ 
 f\colon y = \frac{\ln x} 
{x}
\]}
{nie istnieje}{\(x = 0\)}{\(x = 1\)}{\(x =\mathrm{e}\)}{}{}}
\LANGcs{
\Question{
Doplňte správné tvrzení: „Lokální minimum funkce
\(f\colon y = \frac{\ln x} 
{x}\)...”}
{neexistuje.}{nastává v bodě \(0\).}{nastává v bodě \(1\).}{nastává v bodě \(\mathrm{e}\).}{}{}}
\LANGsk{
\Question{
Doplňte správne tvrdenie: „Lokálne minimum funkcie
\(f\colon y = \frac{\ln x} 
{x}\)...”}
{neexistuje.}{nastáva v bode \(0\).}{nastáva v bode \(1\).}{nastáva v bode \(\mathrm{e}\).}{}{}}
\LANGes{
\Question{
Halla el mínimo local de la siguiente función:
\[ 
 f(x) = \frac{\ln x} 
{x}
\]}
{no existe}{\(x = 0\)}{\(x = 1\)}{\(x =\mathrm{e}\)}{}{}}
\End

\Data\ID{9000079105}
\Updated{2022-04-23 05:04:08}
\Level{A}
\SubArea{Applications of derivatives}
\LANGen{
\Question{
Find all the $x$ at which the function $f$ has local extrema.
\[ 
 f(x)= \left (1 - x^{2}\right )^{3}
\]}
{\(x=0\)}{\(x_1=0\),
\(x_2=1\)}{\(x_1=- 1\),
\(x_2=1\)}{\(x_1=- 1\),
\(x_2=0\),
\(x_3=1\)}{}{}}
\LANGpl{
\Question{
Wyznacz ekstrema lokalne funkcji.
\[ 
 f\colon y = \left (1 - x^{2}\right )^{3}
\]}
{\(x=0\)}{\(x_1=0\),
\(x_2=1\)}{\(x_1=- 1\),
\(x_2=1\)}{\(x_1=- 1\),
\(x_2=0\),
\(x_3=1\)}{}{}}
\LANGcs{
\Question{
Funkce \(f\colon y = \left (1 - x^{2}\right )^{3}\) 
má lokální extrémy v bodě (bodech):}
{\(x=0\)}{\(x_1=0\),
\(x_2=1\)}{\(x_1=- 1\),
\(x_2=1\)}{\(x_1=- 1\),
\(x_2=0\),
\(x_3=1\)}{}{}}
\LANGsk{
\Question{
Funkcia \(f\colon y = \left (1 - x^{2}\right )^{3}\) 
má lokálne extrémy v bode (bodoch):}
{\(x=0\)}{\(x_1=0\),
\(x_2=1\)}{\(x_1=- 1\),
\(x_2=1\)}{\(x_1=- 1\),
\(x_2=0\),
\(x_3=1\)}{}{}}
\LANGes{
\Question{
Halla los extremos locales de la siguiente función:
\[ 
 f(x)= \left (1 - x^{2}\right )^{3}
\]}
{\(x=0\)}{\(x_1=0\),
\(x_2=1\)}{\(x_1=- 1\),
\(x_2=1\)}{\(x_1=- 1\),
\(x_2=0\),
\(x_3=1\)}{}{}}
\End

\Data\ID{9000079106}
\Updated{2020-07-15 03:07:37}
\Level{A}
\SubArea{Applications of derivatives}
\LANGen{
\Question{
Given function \(f(x)= x\mathrm{e}^{\frac{1} 
{x} }\),
identify a true statement.}
{The local minimum of the function \(f\) 
is at the point \(x = 1\),
the function does not have a local maximum.}{The local maximum of the function \(f\) 
is at the point \(x = 0\), the
local minimum at \(x = 1\).}{The local maximum of the function \(f\) 
is at the point \(x = 1\),
the function does not have a local minimum.}{The function \(f\) 
has neither local minimum nor maximum.}{}{}}
\LANGpl{
\Question{
Dana jest funkcja  \(f\colon y = x\mathrm{e}^{\frac{1} 
{x} }\),
wskaż zdanie prawdziwe.}
{Lokalne minimum funkcji  \(f\) jest 
w punkcie \(x = 1\),
lokalne maksimum funkcji nie istnieje.}{Lokalne maksimum funkcji \(f\) jest 
w punkcie \(x = 0\), 
lokalne minimum jest  w punkcie \(x = 1\).}{Lokalne maksimum funkcji  \(f\) jest 
w punkcie \(x = 1\),
lokalne minimum funkcji nie istnieje.}{Lokalne minimum i maksimum funkcji \(f\) 
nie istnieje.}{}{}}
\LANGcs{
\Question{
Je dána funkce \(f\colon y = x\mathrm{e}^{\frac{1} 
{x} }\).
Vyberte správné tvrzení:}
{Funkce \(f\) má lokální
minimum v bodě \(x=1\),
lokální maximum neexistuje.}{Funkce \(f\) má lokální
maximum v bodě \(x=0\) a
lokální minimum v bodě \(x=1\).}{Funkce \(f\) má lokální
maximum v bodě \(x=1\),
lokální minimum neexistuje.}{Lokální extrémy funkce \(f\) 
neexistují.}{}{}}
\LANGsk{
\Question{
Je daná funkcia \(f\colon y = x\mathrm{e}^{\frac{1} 
{x} }\).
Vyberte správne tvrdenie:}
{Funkcia \(f\) má lokálne
minimum v bode \(x=1\),
lokálne maximum neexistuje.}{Funkcia \(f\) má lokálne
maximum v bode \(x=0\) a
lokálne minimum v bode \(x=1\).}{Funkcia \(f\) má lokálne
maximum v bode \(x=1\),
lokálne minimum neexistuje.}{Lokálne extrémy funkcie \(f\) 
neexistujú.}{}{}}
\LANGes{
\Question{
Dada la función \(f(x)= x\mathrm{e}^{\frac{1} 
{x} }\),
Identifica la proposición lógica.}
{El mínimo local de la función \(f\) 
está en el punto \(x = 1\),
la función no tiene ningún máximo local.}{El máximo local de la función \(f\) 
está en el punto \(x = 0\), el mínimo local en \(x = 1\).}{El máximo local de la función \(f\) 
está en el punto  \(x = 1\),
la función no tiene ningún mínimo local.}{La función \(f\) 
no tiene mínimos ni máximos locales.}{}{}}
\End

\Data\ID{9000079107}
\Updated{2022-04-23 05:04:08}
\Level{A}
\SubArea{Applications of derivatives}
\LANGen{
\Question{
What is the function value of the function $f$ at its local minimum?
\[ 
 f(x) = \frac{2} 
{\sqrt{4x - x^{2}}}
\]}
{\(1\)}{\(2\)}{\(0\)}{the local minimum does not exist}{}{}}
\LANGpl{
\Question{
Wyznacz lokalne minimum funkcji.
\[ 
 f\colon y = \frac{2} 
{\sqrt{4x - x^{2}}}
\]}
{\(1\)}{\(2\)}{\(0\)}{lokalne minimum nie istnieje}{}{}}
\LANGcs{
\Question{
Doplňte správné tvrzení: „Funkce
\(f\colon y = \frac{2} 
{\sqrt{4x-x^{2}}} \)...”}
{má v bodě lokálního minima funkční hodnotu
\(1\).}{má v bodě lokálního minima funkční hodnotu
\(2\).}{má v bodě lokálního minima funkční hodnotu
\(0\).}{nemá lokální minimum.}{}{}}
\LANGsk{
\Question{
Doplňte správne tvrdenie: „Funkcia
\(f\colon y = \frac{2} 
{\sqrt{4x-x^{2}}} \)...”}
{má v bode lokálneho minima funkčnú hodnotu
\(1\).}{má v bode lokálneho minima funkčnú hodnotu
\(2\).}{má v bode lokálneho minima funkčnú hodnotu
\(0\).}{nemá lokálne minimum.}{}{}}
\LANGes{
\Question{
Halla el mínimo local de la función:
\[ 
 f(x) = \frac{2} 
{\sqrt{4x - x^{2}}}
\]}
{\(1\)}{\(2\)}{\(0\)}{el mínimo local no existe}{}{}}
\End

\Data\ID{1003263403}
\Updated{2022-04-23 05:04:54}
\Level{B}
\SubArea{Applications of derivatives}
\LANGen{
\Question{
Find the global extrema of the following function on the interval \( [0;3] \).
\[ f(x)=2x^3-3x^2-12x \]}
{the global minimum at \( x=2 \), the global maximum at \( x=0 \)}{the global minimum at \( x=2 \), the global maximum at \( x=-1 \)}{the global minimum at \( x=0 \), the global maximum at \( x=2 \)}{the global minimum at \( x=3 \), the global maximum at \( x=0 \)}{}{}}
\LANGpl{
\Question{
Wskaż  globalne ekstrema danej funkcji w przedziale  \( [0;3] \).
\[ f(x)=2x^3-3x^2-12x \]}
{globalne minimum w  \( x=2 \),  globalne maksimum w \( x=0 \)}{globalne minimum w  \( x=2 \), globalne maksimum w \( x=-1 \)}{globalne minimum w t \( x=0 \), globalne maksimum w \( x=2 \)}{globalne minimum w  \( x=3 \), globalne maksimum w \( x=0 \)}{}{}}
\LANGcs{
\Question{
Najděte globální extrémy funkce \( f \) na intervalu \( [0;3] \).
\[ f(x)=2x^3-3x^2-12x \]}
{globální minimum v bodě \( x=2 \), globální maximum v bodě \( x=0 \)}{globální minimum v bodě \( x=2 \), globální maximum v bodě \( x=-1 \)}{globální minimum v bodě \( x=0 \), globální maximum v bodě \( x=2 \)}{globální minimum v bodě \( x=3 \), globální maximum v bodě \( x=0 \)}{}{}}
\LANGsk{
\Question{
Nájdite globálne extrémy funkcie \( f \) na intervale \( [0;3] \).
\[ f(x)=2x^3-3x^2-12x \]}
{globálne minimum v bode \( x=2 \), globálne maximum v bode \( x=0 \)}{globálne minimum v bode \( x=2 \), globálne maximum v bode \( x=-1 \)}{globálne minimum v bode \( x=0 \), globálne maximum v bode \( x=2 \)}{globálne minimum v bode \( x=3 \), globálne maximum v bode \( x=0 \)}{}{}}
\LANGes{
\Question{
Halla los extremos globales de la función en el intervalo \( [0;3] \).
\[ f(x)=2x^3-3x^2-12x \]}
{el mínimo global en \( x=2 \), el máximo global en \( x=0 \)}{el mínimo global en \( x=2 \), el máximo global en \( x=-1 \)}{el mínimo global en \( x=0 \), el máximo global en \( x=2 \)}{el mínimo global en \( x=3 \), el máximo global en \( x=0 \)}{}{}}
\End

\Data\ID{1003263404}
\Updated{2020-07-15 03:07:21}
\Level{B}
\SubArea{Applications of derivatives}
\LANGen{
\Question{
Find the global extrema of the following function on the interval \( [-1;3] \).
\[ f(x)=x^2\cdot \mathrm{e}^{-x} \]}
{the global minimum at \( x=0 \), the global maximum at \( x=-1 \)}{the global minimum at \( x=0 \), the global maximum at \( x=2 \)}{the global minimum at \( x=3 \), the global maximum at \( x=-1 \)}{the global minimum at \( x=-1 \), the global maximum at \( x=0 \)}{}{}}
\LANGpl{
\Question{
Wskaż  globalne ekstrema danej funkcji w przedziale  \( \langle-1;3\rangle \).
\[ f(x)=x^2\cdot \mathrm{e}^{-x} \]}
{globalne minimum w  \( x=0 \), globalne maksimum w \( x=-1 \)}{globalne minimum w \( x=0 \), globalne maksimum w \( x=2 \)}{globalne minimum w \( x=3 \), globalne maksimum w \( x=-1 \)}{globalne minimum w \( x=-1 \), globalne maksimum w \( x=0 \)}{}{}}
\LANGcs{
\Question{
Najděte globální extrémy funkce \( f \) na intervalu \( \langle-1;3\rangle \).
\[ f(x)=x^2\cdot \mathrm{e}^{-x} \]}
{globální minimum v bodě \( x=0 \), globální maximum v bodě \( x=-1 \)}{globální minimum v bodě \( x=0 \), globální maximum v bodě \( x=2 \)}{globální minimum v bodě \( x=3 \), globální maximum v bodě \( x=-1 \)}{globální minimum v bodě \( x=-1 \), globální maximum v bodě \( x=0 \)}{}{}}
\LANGsk{
\Question{
Nájdite globálne extrémy funkcie \( f \) na intervalu \( \langle-1;3\rangle \).
\[ f(x)=x^2\cdot \mathrm{e}^{-x} \]}
{globálne minimum v bode \( x=0 \), globálne maximum v bode \( x=-1 \)}{globálne minimum v bode \( x=0 \), globálne maximum v bode \( x=2 \)}{globálne minimum v bode \( x=3 \), globálne maximum v bode \( x=-1 \)}{globálne minimum v bode \( x=-1 \), globálne maximum v bode \( x=0 \)}{}{}}
\LANGes{
\Question{
Halla los extremos globales de la función en el intervalo \( [-1;3] \).
\[ f(x)=x^2\cdot \mathrm{e}^{-x} \]}
{el mínimo global en \( x=0 \), el máximo global en \( x=-1 \)}{el mínimo global en \( x=0 \), el máximo global en \( x=2 \)}{el mínimo global en \( x=3 \), el máximo global en \( x=-1 \)}{el mínimo global en \( x=-1 \), el máximo global en \( x=0 \)}{}{}}
\End

\Data\ID{1003263405}
\Updated{2021-08-12 10:08:14}
\Level{B}
\SubArea{Applications of derivatives}
\LANGen{
\Question{
Find a true statement about the function \( f(x)=\sin x+\frac12\cos2x \) on the interval \( [0;\pi] \).}
{The function has global minima at the points \( x=0 \), \( x=\frac{\pi}2 \) and \( x=\pi \).}{The only global minimum of \( f \) on this interval is at the point \( x=\frac{\pi}2 \).}{The only global maximum of \( f \) on this interval is at the point \( x=\frac{\pi}6 \).}{The function \( f \) has no global minimum on this interval.}{}{}}
\LANGpl{
\Question{
Wskaż zdanie prawdziwe dotyczące podanej funkcji  \( f(x)=\sin x+\frac12\cos2x \) w przedziale \( \langle0;\pi\rangle \).}
{Funkcja ma globalne minima w punktach  \( x=0 \), \( x=\frac{\pi}2 \) i \( x=\pi \).}{Jedyne globalne minimum funkcji \( f \) w tym przedziale znajduje się w punkcie  \( x=\frac{\pi}2 \).}{Jedyne globalne maksimum funkcji \( f \) w tym przedziale znajduje się w punkcie \( x=\frac{\pi}6 \).}{Funkcja \( f \) nie ma lokalnego minimum w tym przedziale.}{}{}}
\LANGcs{
\Question{
Vyberte pravdivé tvrzení o funkci \( f(x)=\sin x+\frac12\cos2x \) na intervalu \( \langle0;\pi\rangle \).}
{Funkce \( f \) má globální minima v bodech \( x=0 \), \( x=\frac{\pi}2 \) a \( x=\pi \).}{Jediné globální minimum funkce \( f \) na daném intervalu je v bodě \( x=\frac{\pi}2 \).}{Jediné globální maximum funkce \( f \) na daném intervalu je v bodě \( x=\frac{\pi}6 \).}{Funkce \( f \) nemá na daném intervalu globální minimum.}{}{}}
\LANGsk{
\Question{
Vyberte pravdivé tvrdenie o funkcii \( f(x)=\sin x+\frac12\cos2x \) na intervale \( \langle0;\pi\rangle \).}
{Funkcia \( f \) má globálne minimá v bodoch \( x=0 \), \( x=\frac{\pi}2 \) a \( x=\pi \).}{Jediné globálne minimum funkcie \( f \) na danom intervale je v bode \( x=\frac{\pi}2 \).}{Jediné globálne maximum funkcie \( f \) na danom intervale je v bode \( x=\frac{\pi}6 \).}{Funkcia \( f \) nemá na danom intervale globálne minimum.}{}{}}
\LANGes{
\Question{
Identifica la proposición lógica sobre la función \( f(x)=\sin x+\frac12\cos2x \) en el intervalo \( [0;\pi] \).}
{La función tiene mínimos globales en los puntos \( x=0 \), \( x=\frac{\pi}2 \) y \( x=\pi \).}{El único mínimo global de la función \( f \) en este intervalo está en el punto \( x=\frac{\pi}2 \).}{El único máximo global de la función \( f \) en este intervalo está en el punto \( x=\frac{\pi}6 \).}{La función \( f \) no tiene mínimos globales en este intervalo.}{}{}}
\End

\Data\ID{1103163501}
\Updated{2020-07-15 03:07:10}
\Level{B}
\SubArea{Applications of derivatives}
\LANGen{
\Question{
Choose the graph of a function $f$ which does \textbf{not} satisfy
$$ f''(0) > 0;\ f''(2) < 0 $$
($f''$ is the second derivative of the function $f$).}
{}{}{}{}{}{}}
\LANGpl{
\Question{
Wybierz wykres funkcji $f$, kóry \textbf{nie} spełnia
$$ f''(0) > 0;\ f''(2) < 0 $$
($f''$ jest drugą pochodną funkcji $f$).}
{}{}{}{}{}{}}
\LANGcs{
\Question{
Vyberte graf funkce $f$, pro kterou \textbf{neplatí}
$$ f''(0) > 0;\ f''(2) < 0 $$
($f''$ je druhá derivace funkce $f$).}
{}{}{}{}{}{}}
\LANGsk{
\Question{
Vyberte graf funkcie $f$, pre ktorú \textbf{neplatí}
$$ f''(0) > 0;\ f''(2) < 0 $$
($f''$ je druhá derivácia funkcie $f$).}
{}{}{}{}{}{}}
\LANGes{
\Question{
Elige el gráfico de una función $f$ que \textbf{no} satisface
$$ f''(0) > 0;\ f''(2) < 0 $$
($f''$ es la segunda derivada de la función $f$).}
{}{}{}{}{}{}}

\IMGanswer{1103163501a.pdf}
\IMGanswer{1103163501b.pdf}
\IMGanswer{1103163501c.pdf}
\IMGanswer{1103163501d.pdf}\End

\Data\ID{1103163502}
\Updated{2020-07-15 03:07:10}
\Level{B}
\SubArea{Applications of derivatives}
\LANGen{
\Question{
Choose the graph of a function $f$ that satisfies
$$ f''(0) > 0;\ f''(3) < 0 $$
($f''$ is the second derivative of the function $f$).}
{}{}{}{}{}{}}
\LANGpl{
\Question{
Wybierz wykres funkcji $f$, który spełnia
$$ f''(0) > 0;\ f''(3) < 0 $$
($f''$ jest drugą pochodną funkcji $f$).}
{}{}{}{}{}{}}
\LANGcs{
\Question{
Vyberte graf funkce $f$, pro kterou platí
$$ f''(0) > 0;\ f''(3) < 0 $$
($f''$ je druhá derivace funkce $f$).}
{}{}{}{}{}{}}
\LANGsk{
\Question{
Vyberte graf funkcie $f$, pre ktorý platí
$$ f''(0) > 0;\ f''(3) < 0 $$
($f''$ je druhá derivácia funkcie $f$).}
{}{}{}{}{}{}}
\LANGes{
\Question{
Elige el gráfico de una función $f$ que satisface
$$ f''(0) > 0;\ f''(3) < 0 $$
($f''$ es la segunda derivada de la función $f$).}
{}{}{}{}{}{}}

\IMGanswer{1103163502a.pdf}
\IMGanswer{1103163502b.pdf}
\IMGanswer{1103163502c.pdf}
\IMGanswer{1103163502d.pdf}\End

\Data\ID{1103163503}
\Updated{2022-05-29 12:05:45}
\Level{B}
\SubArea{Applications of derivatives}
\LANGen{
\Question{
Choose the graph of a function $f$ that satisfies
\begin{gather*} 
f'(0)=f'(3)=0; \\
f''(0) < 0;\ f''(3) > 0
\end{gather*}
($f'$ is the derivative of the function $f$, $f''$ is the second derivative of the function $f$).}
{}{}{}{}{}{}}
\LANGpl{
\Question{
Wybierz wykres funkcji $f$ który, spełnia
\begin{gather*} 
f'(0)=f'(3)=0; \\
f''(0) < 0;\ f''(3) > 0
\end{gather*}
($f'$ jest  pochodną funkcji $f$, $f''$ jest drugą pochodną funkcji $f$).}
{}{}{}{}{}{}}
\LANGcs{
\Question{
Vyberte graf funkce $f$, pro kterou platí
\begin{gather*} 
f'(0)=f'(3)=0; \\
f''(0) < 0;\ f''(3) > 0
\end{gather*}
($f'$ je derivace funkce $f$, $f''$ je druhá derivace funkce $f$).}
{}{}{}{}{}{}}
\LANGsk{
\Question{
Vyberte graf funkcie $f$, pre ktorú platí
\begin{gather*} 
f'(0)=f'(3)=0; \\
f''(0) < 0;\ f''(3) > 0
\end{gather*}
($f'$ je derivácia funkcie $f$, $f''$ je druhá derivácia funkcie $f$).}
{}{}{}{}{}{}}
\LANGes{
\Question{
Elige el gráfico de una función $f$ que satisface
\begin{gather*} 
f'(0)=f'(3)=0; \\
f''(0) < 0;\ f''(3) > 0
\end{gather*}
($f'$ es la derivada de la función $f$, $f''$ es la segunda derivada de la función $f$).}
{}{}{}{}{}{}}

\IMGanswer{1103163503a.pdf}
\IMGanswer{1103163503b.pdf}
\IMGanswer{1103163503c.pdf}
\IMGanswer{1103163503d.pdf}\End

\Data\ID{1103163504}
\Updated{2020-07-15 03:07:10}
\Level{B}
\SubArea{Applications of derivatives}
\LANGen{
\Question{
Choose the graph of a function $f$ that satisfies
\begin{gather*}
f'(0)=f'(3)=0; \\
f''(0)=0;\ f''(3) < 0
\end{gather*}
($f'$ is the derivative of a function $f$, $f''$ is the second derivative of a function $f$).}
{}{}{}{}{}{}}
\LANGpl{
\Question{
Wybierz wykres funkcji $f$, który spełnia
\begin{gather*}
f'(0)=f'(3)=0; \\
f''(0)=0;\ f''(3) < 0
\end{gather*}
($f'$ jest  pochodną funkcji $f$, $f''$ jest drugą pochodną funkcji $f$).}
{}{}{}{}{}{}}
\LANGcs{
\Question{
Vyberte graf funkce $f$, pro kterou platí
\begin{gather*}
f'(0)=f'(3)=0; \\
f''(0)=0;\ f''(3) < 0
\end{gather*}
($f'$ je derivace funkce $f$, $f''$ je druhá derivace funkce $f$).}
{}{}{}{}{}{}}
\LANGsk{
\Question{
Vyberte graf funkcie $f$, pre ktorú platí
\begin{gather*}
f'(0)=f'(3)=0; \\
f''(0)=0;\ f''(3) < 0
\end{gather*}
($f'$ je derivácia funkcie $f$, $f''$ je druhá derivácia funkcie $f$).}
{}{}{}{}{}{}}
\LANGes{
\Question{
Elige el gráfico de una función $f$ que satisface
\begin{gather*}
f'(0)=f'(3)=0; \\
f''(0)=0;\ f''(3) < 0
\end{gather*}
($f'$ es la derivada de la función $f$, $f''$ es la segunda derivada de la función $f$).}
{}{}{}{}{}{}}

\IMGanswer{1103163504a.pdf}
\IMGanswer{1103163504b.pdf}
\IMGanswer{1103163504c.pdf}
\IMGanswer{1103163504d.pdf}\End

\Data\ID{1103163505}
\Updated{2022-07-20 09:07:18}
\Level{B}
\SubArea{Applications of derivatives}
\LANGen{
\Question{
Choose the graph of a function $f$ that satisfies 
\begin{gather*}
f'(0)  \text{ does not exist}; \\
f''(x) > 0 \text{ if } x < 0 ; \\
f''(x) > 0 \text{ if } x > 1; \\
f''(x) < 0 \text{ if } 0 < x < 1
\end{gather*}
($f'$ is the derivative of a function $f$, $f''$ is the second derivative of a function $f$).}
{}{}{}{}{}{}}
\LANGpl{
\Question{
Wybierz wykres funkcji $f$, który spełnia
\begin{gather*}
f'(0)  \text{ nie istnieje}; \\
f''(x) > 0 \text{ jeśli } x < 0 ; \\
f''(x) > 0 \text{ jeśli } x > 1; \\
f''(x) < 0 \text{ jeśli } 0 < x < 1
\end{gather*}
($f'$ jest pochodną funkcji $f$, $f''$ jest drugą pochodną funkcji $f$).}
{}{}{}{}{}{}}
\LANGcs{
\Question{
Vyberte graf funkce $f$, pro kterou platí 
\begin{gather*}
f'(0)  \text{ neexistuje}; \\
f''(x) > 0 \text{ pro } x < 0 ; \\
f''(x) > 0 \text{ pro } x > 1; \\
f''(x) < 0 \text{ pro } 0 < x < 1
\end{gather*}
($f'$ je derivace funkce $f$, $f''$ je druhá derivace funkce $f$).}
{}{}{}{}{}{}}
\LANGsk{
\Question{
Vyberte graf funkcie $f$, pre ktorú platí 
\begin{gather*}
f'(0)  \text{ neexistuje}; \\
f''(x) > 0 \text{ pro } x < 0 ; \\
f''(x) > 0 \text{ pro } x > 1; \\
f''(x) < 0 \text{ pro } 0 < x < 1
\end{gather*}
($f'$ je derivácie funkcie $f$, $f''$ je druhá derivácia funkcie $f$).}
{}{}{}{}{}{}}
\LANGes{
\Question{
Elige el gráfico de una función $f$ que satisface
\begin{gather*}
f'(0)  \text{ no existe}; \\
f''(x) > 0 \text{ si } x < 0 ; \\
f''(x) > 0 \text{ si } x > 1; \\
f''(x) < 0 \text{ si } 0 < x < 1
\end{gather*}
($f'$ es la derivada de la función $f$, $f''$ es la segunda derivada de la función $f$).}
{}{}{}{}{}{}}

\IMGanswer{1103163505a.pdf}
\IMGanswer{1103163505b.pdf}
\IMGanswer{1103163505c.pdf}
\IMGanswer{1103163505d.pdf}\End

\Data\ID{1103263401}
\Updated{2020-07-15 03:07:21}
\Level{B}
\SubArea{Applications of derivatives}
\IMG{1103263401.pdf}
\LANGen{
\Question{
The graph of \( f \) is given in the figure. Choose which of the following statements about the function \( f \) are true.
\[
\begin{array}{l}
\text{A: The global maximum of } f \text{ on the interval } [-4;4] \text{ is at } x=4. \\
\text{B: The only global minimum of } f \text{ on the interval } [-4;4] \text{ is at } x=2. \\
\text{C: On } (-2;3] \text{ there is the global minimum of } f \text{ at } x=2 \text{ and the global maximum of } f  \text{ at } x=-2. \\
\text{D: The function } f \text{ has no global maximum on the interval } [-3;4). \\
\text{E: The function } f \text{ has no global minimum on the interval } [-4;2) \text{ .}
\end{array}
\]
The only true statements are:}
{A, D}{B, C}{B, D, E}{A, D, E}{A, B, E}{C, D}}
\LANGpl{
\Question{
Dany jest wykres funkcji \( f \). Wskaż zdania prawdziwe dotyczące  funkcji \( f \).
\[
\begin{array}{l}
\text{A: Maksimum globalne funkcji  } f \text{ w przedziale } [-4;4] \text{znajduje się w punkcie } x=4. \\
\text{B: Jedyne globalne minimum funkcji } f \text{ w przedziale } [-4;4] \text{ znajduje się w punkcie } x=2. \\
\text{C: W przedziale } (-2;3] \text{ globalne minimum funkcji } f \text{ znajduje się w punkcie } x=2 \text{ i globalne maksimum funkcji } f  \text{ znajduje się w punkcie } x=-2. \\
\text{D: Funkcja } f \text{ nie ma globalnego maksimum w przedziale } [-3;4). \\
\text{E: Funkcja } f \text{ nie ma globalnego minimum w przedziale } [-4;2) \text{ .}
\end{array}
\]
Jedyne prawdziwe zdania to:}
{A, D}{B, C}{B, D, E}{A, D, E}{A, B, E}{C, D}}
\LANGcs{
\Question{
Na obrázku je dán graf funkce \( f \). Vyberte, která z následujících tvrzení o funkci \( f \) jsou pravdivá.
\[
\begin{array}{l}
\text{A: Daná funkce } f \text{ má na intervalu } \langle-4;4\rangle \text{ globální maximum v bodě } x=4. \\
\text{B: Jediné globální minimum funkce } f \text{ na intervalu } \langle-4;4\rangle \text{ je v bodě } x=2. \\
\text{C: Na intervalu} (-2;3\rangle \text{ má daná funkce } f \text{ globální minimum v bodě } x=2 \text{ a globální maximum v bodě } x=-2. \\
\text{D: Daná funkce } f \text{ nemá globální maximum na intervalu } \langle-3;4). \\
\text{E: Daná funkce } f \text{ nemá globální minimum na intervalu } \langle-4;2) \text{ .}
\end{array}
\]
Jedinými pravdivými tvrzeními jsou:}
{A, D}{B, C}{B, D, E}{A, D, E}{A, B, E}{C, D}}
\LANGsk{
\Question{
Na obrázku je daný graf funkcie \( f \). Vyberte, ktoré z nasledujúcich tvrdení o funkcii \( f \) sú pravdivé.
\[
\begin{array}{l}
\text{A: Daná funkcia } f \text{ má na intervale } \langle-4;4\rangle \text{ globálne maximum v bode } x=4. \\
\text{B: Jediné globálne minimum funkcie } f \text{ na intervale } \langle-4;4\rangle \text{ je v bode } x=2. \\
\text{C: Na intervale } (-2;3\rangle \text{ má daná funkcia } f \text{ globálne minimum v bode } x=2 \text{ a globálne maximum v bode } x=-2. \\
\text{D: Daná funkcia } f \text{ nemá globálne maximum na intervale } \langle-3;4). \\
\text{E: Daná funkcia } f \text{ nemá globálne minimum na intervale } \langle-4;2) \text{ .}
\end{array}
\]
Jedinými správnymi tvrdeniami sú:}
{A, D}{B, C}{B, D, E}{A, D, E}{A, B, E}{C, D}}
\LANGes{
\Question{
Dado el gráfico de la función \( f \). Elige las proposiciones lógicas sobre la función  \( f \).
\[
\begin{array}{l}
\text{A: El máximo global de } f \text{ en el intervalo } [-4;4] \text{ está en el punto } x=4. \\
\text{B: El único mínimo global de } f \text{ en el intervalo } [-4;4] \text{ está en el punto } x=2. \\
\text{C: En } (-2;3] \text{ el mínimo global de } f \text{ está en el punto } x=2 \text{ y el máximo global de } f  \text{ está en el punto } x=-2. \\
\text{D: La función } f \text{ no tiene máximo global en el intervalo } [-3;4). \\
\text{E: La función } f \text{ no tiene mínimo global en el intervalo } [-4;2) \text{ .}
\end{array}
\]}
{A, D}{B, C}{B, D, E}{A, D, E}{A, B, E}{C, D}}
\End

\Data\ID{1103263402}
\Updated{2020-07-15 03:07:21}
\Level{B}
\SubArea{Applications of derivatives}
\IMG{1103263402_0.pdf}
\LANGen{
\Question{
The graph of \( f \) is given in the figure. Choose which of the following statements about the function \( f \) are true.
\[
\begin{array}{l}
\text{A: The global minimum of } f \text{ on the interval } (-3;3) \text{ is at } x=0. \\
\text{B: The global maxima of } f \text{ on the interval } [-3;3] \text{ are at } x=-2 \text{ and } x=2. \\
\text{C: On } (-2;3] \text{ there is the global minimum of } f \text{ at } x=3 \text{ and the global maximum of } f \text{ at } x=2. \\
\text{D: The function } f \text{ has no global minimum on the interval } (-3;3). \\
\text{E: The function } f \text{ has no global maximum on the interval } (-3;3) .
\end{array}
\]
The only true statements are:}
{B, C, D}{C, D, E}{A, B, C}{A, B}{C, D}{A, E}}
\LANGpl{
\Question{
Dany jest wykres funkcji \( f \). Wskaż zdanie prawdziwe dotyczące funkcji \( f \).
\[
\begin{array}{l}
\text{A: Minimum globalne funkcji } f \text{ w przedziale } (-3;3) \text{ znajduje się w punkcie } x=0. \\
\text{B: Maksima globalne funkcji  } f \text{ w przedziale } [-3;3] \text{ są w punkcie } x=-2 \text{ i } x=2. \\
\text{C: W przedziale } (-2;3] \text{ globalne minimum funkcji } f \text{ znajduje się w punkcie } x=3 \text{ i globalne maksimum funkcji } f \text{ znajduje się w punkcie } x=2. \\
\text{D: Funkcja } f \text{ nie ma globalnego minimum w przedziale } (-3;3). \\
\text{E: Funkcja } f \text{ nie ma globalnego maksimum w przedziale } (-3;3) .
\end{array}
\]
Jedyne prawdziwe zdania to:}
{B, C, D}{C, D, E}{A, B, C}{A, B}{C, D}{A, E}}
\LANGcs{
\Question{
Na obrázku je dán graf funkce \( f \). Vyberte, která z následujících tvrzení o funkci \( f \) jsou pravdivá.
\[
\begin{array}{l}
\text{A: Daná funkce } f \text{ má na intervalu } (-3;3) \text{ globální minimum v bodě } x=0. \\
\text{B: Daná funkce } f \text{ má na intervalu } \langle-3;3\rangle \text{ globální maxima v bodech } x=-2 \text{ a } x=2. \\
\text{C: Na intervalu } (-2;3\rangle \text{ má daná funkce } f \text{ globální minimum v bodě } x=3 \text{ a globální maximum v bodě } x=2. \\
\text{D: Daná funkce } f \text{ nemá globální minimum na intervalu } (-3;3). \\
\text{E: Daná funkce } f \text{ nemá globální maximum na intervalu } (-3;3) .
\end{array}
\]
Jedinými pravdivými tvrzeními jsou:}
{B, C, D}{C, D, E}{A, B, C}{A, B}{C, D}{A, E}}
\LANGsk{
\Question{
Na obrázku je daný graf funkcie \( f \). Vyberte, ktoré z nasledujúcich tvrdení o funkcii \( f \) sú pravdivé.
\[
\begin{array}{l}
\text{A: Daná funkcia } f \text{ má na intervale } (-3;3) \text{ globálne minimum v bode } x=0. \\
\text{B: Daná funkcia } f \text{ má na intervale } \langle-3;3\rangle \text{ globálne maximá  v bodoch } x=-2 \text{ a } x=2. \\
\text{C: Na intervale } (-2;3\rangle \text{ má daná funkcia } f \text{ globálne minimum v bode } x=3 \text{ a globálne maximum v bode } x=2. \\
\text{D: Daná funkcia } f \text{ nemá globálne minimum na intervale } (-3;3). \\
\text{E: Daná funkcia } f \text{ nemá globálne maximum na intervale } (-3;3) .
\end{array}
\]
Jedinými správnymi tvrdeniami sú:}
{B, C, D}{C, D, E}{A, B, C}{A, B}{C, D}{A, E}}
\LANGes{
\Question{
Dado el gráfico de la función \( f \). Elige las proposiciones lógicas sobre la función \( f \).
\[
\begin{array}{l}
\text{A: El mínimo global de } f \text{ en el intervalo } (-3;3) \text{ está en el punto } x=0. \\
\text{B: Los máximos globales } f \text{ en el intervalo } [-3;3] \text{ están en los puntos } x=-2 \text{ y } x=2. \\
\text{C: En } (-2;3] \text{ el mínimo global de } f \text{ está en el punto } x=3 \text{ y el máximo global de } f \text{ está en el punto } x=2. \\
\text{D: La función } f \text{ no tiene mínimo global en el intervalo } (-3;3). \\
\text{E: La función } f \text{ no tiene máximo global en el intervalo } (-3;3) .
\end{array}
\]}
{B, C, D}{C, D, E}{A, B, C}{A, B}{C, D}{A, E}}
\End

\Data\ID{2010010901}
\Updated{2022-08-25 10:08:41}
\Level{B}
\SubArea{Applications of derivatives}
\LANGen{
\Question{
Find the minimum value of the function \(f(x)=-x^2+x-1\) on the interval \( [ 0;2 ]\).}
{\( -3\)}{\(-1\)}{\( -\frac34\)}{\(-7\)}{}{}}
\LANGpl{
\Question{
Wyznacz minimalną wartość funkcji \(f(x)=-x^2+x-1\) na przedziale \( \langle 0;2 \rangle\).}
{\( -3\)}{\(-1\)}{\( -\frac34\)}{\(-7\)}{}{}}
\LANGcs{
\Question{
Jaká je nejmenší hodnota funkce \(f(x)=-x^2+x-1\) v intervalu \( \langle 0;2 \rangle\)?}
{\( -3\)}{\(-1\)}{\( -\frac34\)}{\(-7\)}{}{}}
\LANGsk{
\Question{
Určte najmenšiu hodnotu kvadratickej funkcie \(f(x)=-x^2+x-1\) na intervale \( \langle 0;2 \rangle\).}
{\( -3\)}{\(-1\)}{\( -\frac34\)}{\(-7\)}{}{}}
\LANGes{
\Question{
Halla el valor mínimo de la función \(f(x)=-x^2+x-1\) en el intervalo \( [ 0;2 ]\).}
{\( -3\)}{\(-1\)}{\( -\frac34\)}{\(-7\)}{}{}}
\End

\Data\ID{2010010902}
\Updated{2022-08-25 10:08:41}
\Level{B}
\SubArea{Applications of derivatives}
\LANGen{
\Question{
Find the minimum value of the function \(f(x)=-x^2-x+2\) on the interval \( [ -2;0 ]\).}
{\( 0\)}{\(2\)}{\( \frac94\)}{\(-4\)}{}{}}
\LANGpl{
\Question{
Wyznacz minimalną wartość funkcji  \(f(x)=-x^2-x+2\) na przedziale  \( \langle -2;0 \rangle\).}
{\( 0\)}{\(2\)}{\( \frac94\)}{\(-4\)}{}{}}
\LANGcs{
\Question{
Jaká je nejmenší hodnota funkce \(f(x)=-x^2-x+2\) v intervalu \( \langle -2;0 \rangle\)?}
{\( 0\)}{\(2\)}{\( \frac94\)}{\(-4\)}{}{}}
\LANGsk{
\Question{
Určte najmenšiu hodnotu kvadratickej funkcie \(f(x)=-x^2-x+2\) na intervale \( \langle -2;0 \rangle\).}
{\( 0\)}{\(2\)}{\( \frac94\)}{\(-4\)}{}{}}
\LANGes{
\Question{
Halla el valor mínimo de la función \(f(x)=-x^2-x+2\) en el intervalo \( [ -2;0 ]\).}
{\( 0\)}{\(2\)}{\( \frac94\)}{\(-4\)}{}{}}
\End

\Data\ID{2010010903}
\Updated{2022-08-25 10:08:41}
\Level{B}
\SubArea{Applications of derivatives}
\LANGen{
\Question{
Find the minimum value of the function \(f(x)=2x-4\) on the interval \( (0;4)\).}
{The function \(f\) does not have the minimum on the interval \((0; 4)\).}{\(-4\)}{\(4\)}{\(-6\)}{}{}}
\LANGpl{
\Question{
Wyznacz minimalną wartość funkcji  \(f(x)=2x-4\) na przedziale \( (0;4)\).}
{Funkcja \(f\) nie posiada minimum w przedziale \((0; 4)\).}{\(-4\)}{\(4\)}{\(-6\)}{}{}}
\LANGcs{
\Question{
Jaká je nejmenší hodnota funkce  \(f(x)=2x-4\) v intervalu \( (0;4)\)?}
{Funkce \(f\) nemá v intervalu \((0; 4)\) minimum.}{\(-4\)}{\(4\)}{\(-6\)}{}{}}
\LANGsk{
\Question{
Určte najmenšiu hodnotu funkcie \(f(x)=2x-4\) na intervale \( (0;4)\).}
{Funkcia \(f\) nemá na intervale \((0; 4)\) minimum.}{\(-4\)}{\(4\)}{\(-6\)}{}{}}
\LANGes{
\Question{
Halla el valor mínimo de la función \(f(x)=2x-4\) en el intervalo \( (0;4)\).}
{La función \(f\) no tiene ningún mínimo en el intervalo \((0; 4)\).}{\(-4\)}{\(4\)}{\(-6\)}{}{}}
\End

\Data\ID{2010010904}
\Updated{2022-08-25 10:08:41}
\Level{B}
\SubArea{Applications of derivatives}
\LANGen{
\Question{
Find the minimum value of the function \(f(x)=-3x+6\) on the interval \( (0;4)\).}
{The function \(f\) does not have the minimum on the interval \((0; 4)\).}{\(6\)}{\(-6\)}{\(-9\)}{}{}}
\LANGpl{
\Question{
Wyznacz minimalną wartość funkcji  \(f(x)=-3x+6\) na przedziale \( (0;4)\).}
{Funkcja  \(f\) nie posiada minimum w przedziale \((0; 4)\).}{\(6\)}{\(-6\)}{\(-9\)}{}{}}
\LANGcs{
\Question{
Jaká je nejmenší hodnota funkce \(f(x)=-3x+6\) v intervalu \( (0;4)\)?}
{Funkce \(f\) nemá v intervalu \((0; 4)\) minimum.}{\(6\)}{\(-6\)}{\(-9\)}{}{}}
\LANGsk{
\Question{
Určte najmenšiu hodnotu funkcie \(f(x)=-3x+6\) na intervale \( (0;4)\).}
{Funkcia \(f\) nemá na intervale \((0; 4)\) minimum.}{\(6\)}{\(-6\)}{\(-9\)}{}{}}
\LANGes{
\Question{
Halla el valor mínimo de la función \(f(x)=-3x+6\) en el intervalo \( (0;4)\).}
{La función \(f\) no tiene ningún mínimo en el intervalo \((0; 4)\).}{\(6\)}{\(-6\)}{\(-9\)}{}{}}
\End

\Data\ID{2010010905}
\Updated{2022-08-25 10:08:41}
\Level{B}
\SubArea{Applications of derivatives}
\LANGen{
\Question{
Find the maximum of the function \(g(x)=3x-\frac{1}{1-3x}\) when \(x\in \left[ -1;\frac19\right]\).}
{\(x=0\)}{\(x=-\frac{13}4\)}{\( x=-\frac76\)}{The function \(g\) does not have the maximum on the interval \(\left[ -1;\frac19\right]\).}{}{}}
\LANGpl{
\Question{
Wyznacz maksimum funkcji \(g(x)=3x-\frac{1}{1-3x}\) jeśli \(x\in \left\langle -1;\frac19\right\rangle\).}
{\(x=0\)}{\(x=-\frac{13}4\)}{\( x=-\frac76\)}{Funkcja  \(g\) nie posiada maksimum w przedziale \(\left\langle -1;\frac19\right\rangle\).}{}{}}
\LANGcs{
\Question{
Určete, pro které \(x\in \left\langle -1;\frac19\right\rangle\) je hodnota \(g(x)=3x-\frac{1}{1-3x}\) největší.}
{\(x=0\)}{\(x=-\frac{13}4\)}{\( x=-\frac76\)}{Funkce \(g\) nemá v intervalu \(\left\langle -1;\frac19\right\rangle\) maximum.}{}{}}
\LANGsk{
\Question{
Určte maximum funkcie \(g(x)=3x-\frac{1}{1-3x}\), pre \(x\in \left\langle -1;\frac19\right\rangle\).}
{\(x=0\)}{\(x=-\frac{13}4\)}{\( x=-\frac76\)}{Funkcia \(g\) nemá maximum na intervale \(\left\langle -1;\frac19\right\rangle\).}{}{}}
\LANGes{
\Question{
Halla el máximo de la función \(g(x)=3x-\frac{1}{1-3x}\) si \(x\in \left[ -1;\frac19\right]\).}
{\(x=0\)}{\(x=-\frac{13}4\)}{\( x=-\frac76\)}{La función \(g\) no tiene ningún máximo en el intervalo \(\left[ -1;\frac19\right]\).}{}{}}
\End

\Data\ID{2010010906}
\Updated{2022-08-25 10:08:41}
\Level{B}
\SubArea{Applications of derivatives}
\LANGen{
\Question{
Find the maximum of the function \(g(x)=\frac{-1}{1-2x}+2x\) when \(x\in \left[ -2;\frac14\right]\).}
{\(x=0\)}{\(x=-\frac{17}4\)}{\( x=-\frac32\)}{The function \(g\) does not have the maximum on the interval \(\left[ -2;\frac14\right]\).}{}{}}
\LANGpl{
\Question{
Wyznacz maksimum funkcji  \(g(x)=\frac{-1}{1-2x}+2x\) jeśli \(x\in \left\langle -2;\frac14\right\rangle\).}
{\(x=0\)}{\(x=-\frac{17}4\)}{\( x=-\frac32\)}{Funkcja  \(g\) nie posiada maksimum w przedziale \(\left\langle -2;\frac14\right\rangle\).}{}{}}
\LANGcs{
\Question{
Určete, pro které \(x\in \left\langle -2;\frac14\right\rangle\) je hodnota funkce \(g(x)=\frac{-1}{1-2x}+2x\) největší.}
{\(x=0\)}{\(x=-\frac{17}4\)}{\( x=-\frac32\)}{Funkce \(g\) nemá v intervalu \(\left\langle -2;\frac14\right\rangle\) maximum.}{}{}}
\LANGsk{
\Question{
Určte maximum funkcie \(g(x)=\frac{-1}{1-2x}+2x\), pre \(x\in \left\langle -2;\frac14\right\rangle\).}
{\(x=0\)}{\(x=-\frac{17}4\)}{\( x=-\frac32\)}{Funkcia \(g\) nemá maximum na intervale \(\left\langle -2;\frac14\right\rangle\).}{}{}}
\LANGes{
\Question{
Halla el máximo de la función \(g(x)=\frac{-1}{1-2x}+2x\) si \(x\in \left[ -2;\frac14\right]\).}
{\(x=0\)}{\(x=-\frac{17}4\)}{\( x=-\frac32\)}{La función \(g\) no tiene ningún máximo en el intervalo \(\left[ -2;\frac14\right]\).}{}{}}
\End

\Data\ID{2010010907}
\Updated{2022-08-25 10:08:46}
\Level{B}
\SubArea{Applications of derivatives}
\LANGen{
\Question{
How many of these functions do have a global maximum on the interval \( (-\infty;-1]\) at the point \(x=-1\)?
\[f(x)=x+\frac1x+2\]
\[g(x)=-x^2+6x-9\]
\[h(x)=2x-3\]}
{\(3\)}{\(1\)}{\(2\)}{\(0\)}{}{}}
\LANGpl{
\Question{
Ile z tych funkcji ma globalne maksimum w przedziale \( (-\infty;-1\rangle\) w punkcie \(x=-1\)?
\[f(x)=x+\frac1x+2\]
\[g(x)=-x^2+6x-9\]
\[h(x)=2x-3\]}
{\(3\)}{\(1\)}{\(2\)}{\(0\)}{}{}}
\LANGcs{
\Question{
Určete, kolik z těchto funkcí má v bodě \(x=-1\) globální maximum na intervalu \( (-\infty;-1\rangle\).
\[f(x)=x+\frac1x+2\]
\[g(x)=-x^2+6x-9\]
\[h(x)=2x-3\]}
{\(3\)}{\(1\)}{\(2\)}{\(0\)}{}{}}
\LANGsk{
\Question{
Určte, koľko z daných funkcií má v bode \(x=-1\) globálne maximum na intervale \( (-\infty;-1\rangle\).
\[f(x)=x+\frac1x+2\]
\[g(x)=-x^2+6x-9\]
\[h(x)=2x-3\]}
{\(3\)}{\(1\)}{\(2\)}{\(0\)}{}{}}
\LANGes{
\Question{
¿Cuántas de las siguientes funciones tienen un máximo global en el intervalo \( (-\infty;-1]\) en el punto \(x=-1\)?
\[f(x)=x+\frac1x+2\]
\[g(x)=-x^2+6x-9\]
\[h(x)=2x-3\]}
{\(3\)}{\(1\)}{\(2\)}{\(0\)}{}{}}
\End

\Data\ID{2010010908}
\Updated{2022-08-25 10:08:46}
\Level{B}
\SubArea{Applications of derivatives}
\LANGen{
\Question{
How many of these functions do have a global minimum on the interval \( [ -1;\infty) \) at the point \(x=-1\)?
\[f(x)=x+\frac1{x+2}\]
\[g(x)=x^2+4x+4\]
\[h(x)=3x+1\]}
{\(3\)}{\(1\)}{\(2\)}{\(0\)}{}{}}
\LANGpl{
\Question{
Ile z tych funkcji ma globalne minimum w przedziale \( \langle -1;\infty) \) w punkcie \(x=-1\)?
\[f(x)=x+\frac1{x+2}\]
\[g(x)=x^2+4x+4\]
\[h(x)=3x+1\]}
{\(3\)}{\(1\)}{\(2\)}{\(0\)}{}{}}
\LANGcs{
\Question{
Určete, kolik z těchto funkcí má v bodě \(x=-1\) globální minimum na intervalu \( \langle -1;\infty) \).
\[f(x)=x+\frac1{x+2}\]
\[g(x)=x^2+4x+4\]
\[h(x)=3x+1\]}
{\(3\)}{\(1\)}{\(2\)}{\(0\)}{}{}}
\LANGsk{
\Question{
Určte, koľko z daných funkcií má v bode \(x=-1\) globálne minimum na intervale \( \langle -1;\infty) \).
\[f(x)=x+\frac1{x+2}\]
\[g(x)=x^2+4x+4\]
\[h(x)=3x+1\]}
{\(3\)}{\(1\)}{\(2\)}{\(0\)}{}{}}
\LANGes{
\Question{
¿Cuántas de las siguientes funciones tienen un mínimo global en el intervalo \( [ -1;\infty) \) en el punto \(x=-1\)?
\[f(x)=x+\frac1{x+2}\]
\[g(x)=x^2+4x+4\]
\[h(x)=3x+1\]}
{\(3\)}{\(1\)}{\(2\)}{\(0\)}{}{}}
\End

\Data\ID{2010010909}
\Updated{2022-08-25 10:08:46}
\Level{B}
\SubArea{Applications of derivatives}
\LANGen{
\Question{
Find the maximum of the function \(g(x)=\frac{x^4}2-4x^2+2\) when \( x \in [ -2;3]\).}
{\(x=3\)}{\(x=-2\)}{\(x=0\)}{The function \(g\) does not have the maximum on the interval \([ -2;3]\).}{}{}}
\LANGpl{
\Question{
Znajdź maksimum funkcji \(g(x)=\frac{x^4}2-4x^2+2\) jeśli \( x \in \langle -2;3\rangle\).}
{\(x=3\)}{\(x=-2\)}{\(x=0\)}{Funkcja  \(g\) nie ma maksimum w przedziale \(\langle -2;3\rangle\).}{}{}}
\LANGcs{
\Question{
Určete, pro které \( x \in \langle -2;3\rangle\) je hodnota funkce \(g(x)=\frac{x^4}2-4x^2+2\) největší.}
{\(x=3\)}{\(x=-2\)}{\(x=0\)}{Funkce \(g\) nemá v intervalu \(\langle -2;3\rangle\) maximum.}{}{}}
\LANGsk{
\Question{
Určte maximum funkcie \(g(x)=\frac{x^4}2-4x^2+2\), pre \( x \in \langle -2;3\rangle\).}
{\(x=3\)}{\(x=-2\)}{\(x=0\)}{Funkcia \(g\) nemá maximum na intervale \(\langle -2;3\rangle\).}{}{}}
\LANGes{
\Question{
Halla el máximo de la función \(g(x)=\frac{x^4}2-4x^2+2\) si \( x \in [ -2;3]\).}
{\(x=3\)}{\(x=-2\)}{\(x=0\)}{La función \(g\) no tiene ningún máximo en el intervalo \([ -2;3]\).}{}{}}
\End

\Data\ID{2010010910}
\Updated{2022-08-25 10:08:46}
\Level{B}
\SubArea{Applications of derivatives}
\LANGen{
\Question{
Find the minimum of the function \(g(x)=-{x^4}+2x^2+1\) when \( x \in [ -1;2]\).}
{\(x=2\)}{\(x=-1\)}{\(x=0\)}{The function \(g\) does not have the minimum on the interval \([ -1;2]\).}{}{}}
\LANGpl{
\Question{
Wyznacz minimum funkcji  \(g(x)=-{x^4}+2x^2+1\) jeśli \( x \in \langle -1;2\rangle\).}
{\(x=2\)}{\(x=-1\)}{\(x=0\)}{Funkcja \(g\) nie ma minimum w przedziale \(\langle -1;2\rangle\).}{}{}}
\LANGcs{
\Question{
Určete, pro které \( x \in \langle -1;2\rangle\) je hodnota funkce \(g(x)=-{x^4}+2x^2+1\) nejmenší.}
{\(x=2\)}{\(x=-1\)}{\(x=0\)}{Funkce \(g\) nemá v intervalu \(\langle -1;2\rangle\) minimum.}{}{}}
\LANGsk{
\Question{
Určte minimum funkcie \(g(x)=-{x^4}+2x^2+1\), pre \( x \in \langle -1;2\rangle\).}
{\(x=2\)}{\(x=-1\)}{\(x=0\)}{Funkcia \(g\) nemá minimum na intervale \(\langle -1;2\rangle\).}{}{}}
\LANGes{
\Question{
Halla el mínimo de la función \(g(x)=-{x^4}+2x^2+1\) si \( x \in [ -1;2]\).}
{\(x=2\)}{\(x=-1\)}{\(x=0\)}{La función \(g\) no tiene ningún mínimo en el intervalo \([ -1;2]\).}{}{}}
\End

\Data\ID{2010012501}
\Updated{2022-08-25 10:08:03}
\Level{B}
\SubArea{Applications of derivatives}
\LANGen{
\Question{
Find the global extrema of the following function on the interval \( [ 0;2 ] \).
\[ f(x)=x^3+3x^2-9x \]}
{the global minimum at \( x=1 \), the global maximum at \( x=2 \)}{the global minimum at \( x=1 \), the global maximum at \( x=-3 \)}{the global minimum at \( x=2 \), the global maximum at \( x=1 \)}{the global minimum at \( x=0 \), the global maximum at \( x=2 \)}{}{}}
\LANGpl{
\Question{
Znajdź ekstrema globalne następującej funkcji na przedziale \( \langle 0;2 \rangle \).
\[ f(x)=x^3+3x^2-9x \]}
{globalne minimum w \( x=1 \), globalne maksimum w \( x=2 \)}{globalne minimum w  \( x=1 \), globalne maksimum w \( x=-3 \)}{globalne minimum w \( x=2 \), globalne maksimum w  \( x=1 \)}{globalne minimum w  \( x=0 \), globalne maksimum w  \( x=2 \)}{}{}}
\LANGcs{
\Question{
Najděte globální extrémy funkce \( f \) na intervalu \( \langle 0;2 \rangle \).
\[ f(x)=x^3+3x^2-9x \]}
{Globální minimum v bodě \( x=1 \), globální maximum v bodě  \( x=2 \)}{Globální minimum v bodě \( x=1 \), globální maximum v bodě \( x=-3 \)}{Globální minimum v bodě \( x=2 \), globální maximum v bodě \( x=1 \)}{Globální minimum v bodě \( x=0 \), globální maximum v bodě \( x=2 \)}{}{}}
\LANGsk{
\Question{
Najdite globálne extrémy funkcie \( f \) na intervale \( \langle 0;2 \rangle \).
\[ f(x)=x^3+3x^2-9x \]}
{Globálne minimum v bode \( x=1 \), globálne maximum v bode \( x=2 \).}{Globálne minimum v bode \( x=1 \), globálne maximum v bode \( x=-3 \).}{Globálne minimum v bode \( x=2 \), globálne maximum v bode \( x=1 \).}{Globálne minimum v bode \( x=0 \), globálne maximum v bode \( x=2 \).}{}{}}
\LANGes{
\Question{
Halla los extremos globales de la función en el intervalo \( [ 0;2 ] \).
\[ f(x)=x^3+3x^2-9x \]}
{el mínimo global en \( x=1 \), el máximo global en \( x=2 \)}{el mínimo global en \( x=1 \), el máximo global en \( x=-3 \)}{el mínimo global en \( x=2 \), el máximo global en \( x=1 \)}{el mínimo global en \( x=0 \), el máximo global en \( x=2 \)}{}{}}
\End

\Data\ID{2010012502}
\Updated{2022-08-25 10:08:03}
\Level{B}
\SubArea{Applications of derivatives}
\LANGen{
\Question{
Identify a true statement about the function
\(f(x) = x^{3} +6x^{2} + 12x -1\).}
{There is neither local minimum nor maximum of \(f\).}{The function \(f\) has a local maximum at the point \(x = -2\).}{The function \(f\) has a local minimum at the point \(x = -2\).}{The global minimum of \(f\) on \(\mathbb{R}\) is at \(x = -2\).}{}{}}
\LANGpl{
\Question{
Wybierz prawdziwe stwierdzenie dotyczące funkcji
\(f(x) = x^{3} +6x^{2} + 12x -1\).}
{Funkcja \(f\) nie ma ekstremów lokalnych.}{Funkcja  \(f\) ma lokalne maksimum w  \(x = -2\).}{Funkcja \(f\) ma lokalne minimum w punkcie \(x = -2\).}{Globalne minimum funkcji \(f\) w \(\mathbb{R}\) jest w punkcie \(x = -2\).}{}{}}
\LANGcs{
\Question{
Vyberte pravdivé tvrzení o funkci
\(f(x) = x^{3} +6x^{2} + 12x -1\).}
{Funkce \(f\) nemá žádný lokální extrém.}{Funkce \(f\) má lokální minimum v bodě \(x = -2\).}{Funkce \(f\) má lokální maximum v bodě \(x = -2\).}{Globální minimum funkce \(f\) na množině \(\mathbb{R}\) je v bodě \(x = -2\).}{}{}}
\LANGsk{
\Question{
Vyberte pravdivé tvrdenie o nasledujúcej funkcii \(f(x) = x^{3} +6x^{2} + 12x -1\).}
{Funkcia \(f\) nemá žiadny lokálny extrém.}{Funkcia \(f\) má lokálne maximum v bode \(x = -2\).}{Funkcia \(f\) má lokálne minimum v bode \(x = -2\).}{Globálne minimum funkcie \(f\) na množine \(\mathbb{R}\) je v bode \(x = -2\).}{}{}}
\LANGes{
\Question{
Identifica la proposición lógica sobre la función
\(f(x) = x^{3} +6x^{2} + 12x -1\).}
{No hay ningún mínimo ni máximo local \(f\).}{La función \(f\) tiene un máximo local en el punto \(x = -2\).}{La función \(f\) tiene un mínimo local en el punto \(x = -2\).}{El mínimo global de \(f\) en \(\mathbb{R}\) está en \(x = -2\).}{}{}}
\End

\Data\ID{2010012505}
\Updated{2022-08-25 10:08:03}
\Level{B}
\SubArea{Applications of derivatives}
\LANGen{
\Question{
Identify a true statement about the function
\(f(x) = -\frac{3} 
{4}x^{4} +2x^{3}\).}
{The function \(f\) has
a local maximum at \(x = 2\).}{The function \(f\) has
a local minimum at \(x = 0\).}{The function \(f\) 
has two local extrema. These extrema are at
\(x = 0\) and
\(x = 2\).}{The function \(f\) 
has neither local minimum nor local maximum.}{}{}}
\LANGpl{
\Question{
Wybierz prawdziwe zdanie o funkcji
\(f(x) = -\frac{3} 
{4}x^{4} +2x^{3}\).}
{Funkcja \(f\) ma
lokalne maksimum w \(x = 2\).}{Funkcja \(f\) ma lokalne minimum w 
 \(x = 0\).}{Funkcja \(f\) 
ma dwa lokalne ekstrema. Te ekstrema są w
\(x = 0\) i 
\(x = 2\).}{Funkcja  \(f\) 
nie ma lokalnego minimum ani maksimum.}{}{}}
\LANGcs{
\Question{
Vyberte pravdivé tvrzení o následující funkci
\(f(x) = -\frac{3} 
{4}x^{4} +2x^{3}\).}
{Funkce \(f\) má lokální maximum v bodě \(x = 2\).}{Funkce \(f\) má lokální minimum v bodě \(x = 0\).}{Funkce \(f\) 
má dva lokální extrémy v bodech
\(x = 0\) a
\(x = 2\).}{Funkce \(f\) 
nemá lokální extrém v žádném bodě.}{}{}}
\LANGsk{
\Question{
Vyberte pravdivé tvrdenie o nasledujúcej funkcii \(f(x) = -\frac{3} 
{4}x^{4} +2x^{3}\).}
{Funkcia \(f\) má
lokálne maximum v bode \(x = 2\).}{Funkcia \(f\) má
lokálne minimum v bode \(x = 0\).}{Funkcia \(f\) 
má dva lokálne extrémy v bodoch
\(x = 0\) a
\(x = 2\).}{Funkcia \(f\) 
nemá lokálne extrémy v žiadnom bode.}{}{}}
\LANGes{
\Question{
Identifica la proposición lógica sobre la función
\(f(x) = -\frac{3} 
{4}x^{4} +2x^{3}\).}
{La función \(f\) tiene un máximo local en \(x = 2\).}{La función \(f\) tiene un mínimo local en \(x = 0\).}{La función \(f\) 
tiene dos extremos locales. Están en
\(x = 0\) y
\(x = 2\).}{La función \(f\) 
no tiene ningún mínimo ni máximo local.}{}{}}
\End

\Data\ID{2010013905}
\Updated{2023-01-23 10:01:40}
\Level{B}
\SubArea{Applications of derivatives}
\LANGen{
\Question{
Suppose a function is concave down on the interval \([-1;2]\). Which one of the offered functions has that property?}
{\(g(x)=\sqrt{-2x+8}\)}{\(f(x)=\frac{-x+2}{(x+3)^2}\)}{\(h(x)=\frac{x^2-1}{x}\)}{\(k(x)=\frac12x^3+3x^2-x+2\)}{}{}}
\LANGpl{
\Question{
Załóżmy, że funkcja jest wklęsła w przedziale \(\langle-1;2\rangle\). Która z przedstawionych funkcji ma tę właściwość?}
{\(g(x)=\sqrt{-2x+8}\)}{\(f(x)=\frac{-x+2}{(x+3)^2}\)}{\(h(x)=\frac{x^2-1}{x}\)}{\(k(x)=\frac12x^3+3x^2-x+2\)}{}{}}
\LANGcs{
\Question{
O dané funkci víme, že je v intervalu \(\langle-1;2\rangle\) konkávní. Která z nabídnutých funkcí má tuto vlastnost?}
{\(g(x)=\sqrt{-2x+8}\)}{\(f(x)=\frac{-x+2}{(x+3)^2}\)}{\(h(x)=\frac{x^2-1}{x}\)}{\(k(x)=\frac12{x^3+3x^2-x+2}\)}{}{}}
\LANGsk{
\Question{
Vieme, že daná funkcia je na intervale \(\langle-1;2\rangle\) konkávna. Ktorá z ponúknutých funkcií má túto vlastnosť?}
{\(g(x)=\sqrt{-2x+8}\)}{\(f(x)=\frac{-x+2}{(x+3)^2}\)}{\(h(x)=\frac{x^2-1}{x}\)}{\(k(x)=\frac12x^3+3x^2-x+2\)}{}{}}
\LANGes{
\Question{
Supongamos que una función es cóncava hacia abajo en el intervalo \([-1;2]\). ¿Cuál de las siguientes funciones tiene esta propiedad?}
{\(g(x)=\sqrt{-2x+8}\)}{\(f(x)=\frac{-x+2}{(x+3)^2}\)}{\(h(x)=\frac{x^2-1}{x}\)}{\(k(x)=\frac12x^3+3x^2-x+2\)}{}{}}
\End

\Data\ID{2010013906}
\Updated{2023-01-23 10:01:40}
\Level{B}
\SubArea{Applications of derivatives}
\LANGen{
\Question{
Suppose a function is convex up on the interval \([-2;1]\). Which one of the offered functions has that property?}
{\(h(x)=-\sqrt{x+5}\)}{\(f(x)=\frac{x-2}{(x+5)^2}\)}{\(g(x)=\frac{x^2-2}{2x}\)}{\(k(x)=-\frac15x^3-2x^2+x+1\)}{}{}}
\LANGpl{
\Question{
Załóżmy, że funkcja jest wypukła na przedziale \(\langle-2;1\rangle\). Która z przedstawionych funkcji ma tę właściwość?}
{\(h(x)=-\sqrt{x+5}\)}{\(f(x)=\frac{x-2}{(x+5)^2}\)}{\(g(x)=\frac{x^2-2}{2x}\)}{\(k(x)=-\frac15x^3-2x^2+x+1\)}{}{}}
\LANGcs{
\Question{
O dané funkci víme, že je v intervalu \(\langle-2;1\rangle\) konvexní. Která z nabídnutých funkcí má tuto vlastnost?}
{\(h(x)=-\sqrt{x+5}\)}{\(f(x)=\frac{x-2}{(x+5)^2}\)}{\(g(x)=\frac{x^2-2}{2x}\)}{\(k(x)=-\frac15x^3-2x^2+x+1\)}{}{}}
\LANGsk{
\Question{
Vieme, že daná funkcia je na intervale \(\langle-2;1\rangle\) konvexná. Ktorá z ponúknutých funkcií má túto vlastnosť?}
{\(h(x)=-\sqrt{x+5}\)}{\(f(x)=\frac{x-2}{(x+5)^2}\)}{\(g(x)=\frac{x^2-2}{2x}\)}{\(k(x)=-\frac15x^3-2x^2+x+1\)}{}{}}
\LANGes{
\Question{
Supongamos que una función es cóncava hacia arriba en el intervalo \([-2;1]\). ¿Cuál de las siguientes funciones tiene esta propiedad?}
{\(h(x)=-\sqrt{x+5}\)}{\(f(x)=\frac{x-2}{(x+5)^2}\)}{\(g(x)=\frac{x^2-2}{2x}\)}{\(k(x)=-\frac15x^3-2x^2+x+1\)}{}{}}
\End

\Data\ID{2010013907}
\Updated{2023-01-23 10:01:40}
\Level{B}
\SubArea{Applications of derivatives}
\LANGen{
\Question{
How many of the given functions have exactly one inflection point?
\[f(x)=(x+2)^5+(2+x)^3-2,\quad g(x)=\frac1{6(x+4)^4},\quad h(x)=\frac{x^3+2x^2+x+2}{x}\]}
{2}{1}{3}{None of these functions has exactly one inflection point.}{}{}}
\LANGpl{
\Question{
Ile z podanych funkcji ma dokładnie jeden punkt przegięcia?
\[f(x)=(x+2)^5+(2+x)^3-2,\quad g(x)=\frac1{6(x+4)^4},\quad h(x)=\frac{x^3+2x^2+x+2}{x}\]}
{2}{1}{3}{Żadna z tych funkcji nie ma dokładnie jednego punktu przegięcia.}{}{}}
\LANGcs{
\Question{
Kolik z uvedených funkcí má právě jeden inflexní bod?
\[f(x)=(x+2)^5+(2+x)^3-2,\quad g(x)=\frac1{6(x+4)^4},\quad h(x)=\frac{x^3+2x^2+x+2}{x}\]}
{2}{1}{3}{Žádná z uvedených funkcí nemá právě jeden inflexní bod.}{}{}}
\LANGsk{
\Question{
Koľko z daných funkcií má práve jeden inflexný bod?
\[f(x)=(x+2)^5+(2+x)^3-2,\quad g(x)=\frac1{6(x+4)^4},\quad h(x)=\frac{x^3+2x^2+x+2}{x}\]}
{2}{1}{3}{Žiadna z uvedených funkcií nemá práve jeden inflexný bod.}{}{}}
\LANGes{
\Question{
¿Cuántas de las siguientes funciones tienen exactamente un punto de inflexión?
\[f(x)=(x+2)^5+(2+x)^3-2,\quad g(x)=\frac1{6(x+4)^4},\quad h(x)=\frac{x^3+2x^2+x+2}{x}\]}
{2}{1}{3}{Ninguna de estas funciones tiene exactamente un punto de inflexión.}{}{}}
\End

\Data\ID{2010013908}
\Updated{2023-01-23 10:01:40}
\Level{B}
\SubArea{Applications of derivatives}
\LANGen{
\Question{
How many of the given functions have exactly one inflection point?
\[f(x)=\frac1{-2(x+4)^4}+6,\quad g(x)=-(x-3)^5-(-3+x)^3+1,\quad h(x)=\frac{x^3-3x^2+6x+9}{-3x}\]}
{2}{1}{3}{None of these functions has exactly one inflection point.}{}{}}
\LANGpl{
\Question{
Ile z podanych funkcji ma dokładnie jeden punkt przegięcia?
\[f(x)=\frac1{-2(x+4)^4}+6,\quad g(x)=-(x-3)^5-(-3+x)^3+1,\quad h(x)=\frac{x^3-3x^2+6x+9}{-3x}\]}
{2}{1}{3}{Żadna z tych funkcji nie ma dokładnie jednego punktu przegięcia.}{}{}}
\LANGcs{
\Question{
Kolik z uvedených funkcí má právě jeden inflexní bod?
\[f(x)=\frac1{-2(x+4)^4}+6,\quad g(x)=-(x-3)^5-(-3+x)^3+1,\quad h(x)=\frac{x^3-3x^2+6x+9}{-3x}\]}
{2}{1}{3}{Žádná z uvedených funkcí nemá právě jeden inflexní bod.}{}{}}
\LANGsk{
\Question{
Koľko z daných funkcií má práve jeden inflexný bod?
\[f(x)=\frac1{-2(x+4)^4}+6,\quad g(x)=-(x-3)^5-(-3+x)^3+1,\quad h(x)=\frac{x^3-3x^2+6x+9}{-3x}\]}
{2}{1}{3}{Žiadna z uvedených funkcií nemá práve jeden inflexný bod.}{}{}}
\LANGes{
\Question{
¿Cuántas de las siguientes funciones tienen exactamente un punto de inflexión?
\[f(x)=\frac1{-2(x+4)^4}+6,\quad g(x)=-(x-3)^5-(-3+x)^3+1,\quad h(x)=\frac{x^3-3x^2+6x+9}{-3x}\]}
{2}{1}{3}{Ninguna de estas funciones tiene exactamente un punto de inflexión.}{}{}}
\End

\Data\ID{2110012503}
\Updated{2022-08-25 10:08:02}
\Level{B}
\SubArea{Applications of derivatives}
\LANGen{
\Question{
Choose the graph of a function $f$ that satisfies
\begin{gather*} 
f'(-2)=f'(0)=0; \\
f''(-2) < 0;\ f''(0) > 0
\end{gather*}
($f'$ is the derivative of the function $f$, $f''$ is the second derivative of the function $f$).}
{}{}{}{}{}{}}
\LANGpl{
\Question{
Wybierz wykres funkcji  $f$ który spełnia
\begin{gather*} 
f'(-2)=f'(0)=0; \\
f''(-2) < 0;\ f''(0) > 0
\end{gather*}
($f'$ jest pochodną funkcji  $f$, $f''$ jest drugą pochodną funkcji  $f$).}
{}{}{}{}{}{}}
\LANGcs{
\Question{
Vyberte graf funkce $f$ pro kterou platí:
\begin{gather*} 
f'(-2)=f'(0)=0; \\
f''(-2) < 0;\ f''(0) > 0
\end{gather*}
($f'$ je derivace funkce $f$, $f''$ je druhá derivace funkce $f$).}
{}{}{}{}{}{}}
\LANGsk{
\Question{
Vyberte graf funkcie $f$, pre ktorú platí: 
\begin{gather*} 
f'(-2)=f'(0)=0; \\
f''(-2) < 0;\ f''(0) > 0
\end{gather*}
($f'$ je derivácia funkcie $f$, $f''$ je druhá derivácia funkcie $f$).}
{}{}{}{}{}{}}
\LANGes{
\Question{
Elige el gráfico de una función $f$ que satisface
\begin{gather*} 
f'(-2)=f'(0)=0; \\
f''(-2) < 0;\ f''(0) > 0
\end{gather*}
($f'$ es la derivada de la función $f$, $f''$ es la segunda derivada de la función $f$).}
{}{}{}{}{}{}}

\IMGanswer{2010012501-figure0.pdf}
\IMGanswer{2010012501-figure1.pdf}
\IMGanswer{2010012501-figure2.pdf}
\IMGanswer{2010012501-figure3.pdf}\End

\Data\ID{2110012504}
\Updated{2022-08-25 10:08:02}
\Level{B}
\SubArea{Applications of derivatives}
\LANGen{
\Question{
Choose the graph of a function $f$ that satisfies 
\begin{gather*}
f'(1)  \text{ does not exist}; \\
f''(x) < 0 \text{ if } x < 1 ; \\
f''(x) < 0 \text{ if } x > 2; \\
f''(x) > 0 \text{ if } 1 < x < 2
\end{gather*}
($f'$ is the derivative of a function $f$, $f''$ is the second derivative of a function $f$).}
{}{}{}{}{}{}}
\LANGpl{
\Question{
Wybierz wykres funkcji $f$ który spełnia
\begin{gather*}
f'(1)  \text{ nie istnieje}; \\
f''(x) < 0 \text{ jeśli } x < 1 ; \\
f''(x) < 0 \text{ jeśli } x > 2; \\
f''(x) > 0 \text{ jeśli } 1 < x < 2
\end{gather*}
($f'$ jest pochodną funkcji  $f$, $f''$ jest drugą pochodną funkcji $f$).}
{}{}{}{}{}{}}
\LANGcs{
\Question{
Vyberte graf funkce $f$, pro kterou platí:
\begin{gather*}
f'(1)  \text{ neexistuje}; \\
f''(x) < 0 \text{ pro } x < 1 ; \\
f''(x) < 0 \text{ pro } x > 2; \\
f''(x) > 0 \text{ pro } 1 < x < 2
\end{gather*}
($f'$ je derivace funkce $f$, $f''$ je druhá derivace funkce $f$).}
{}{}{}{}{}{}}
\LANGsk{
\Question{
Vyberte graf funkcie $f$, pre ktorú platí: 
\begin{gather*}
f'(1)  \text{ neexistuje}; \\
f''(x) < 0 \text{ ak } x < 1 ; \\
f''(x) < 0 \text{ ak } x > 2; \\
f''(x) > 0 \text{ ak } 1 < x < 2
\end{gather*}
($f'$ je derivácia funkcie $f$, $f''$ je druhá derivácia funkcie $f$).}
{}{}{}{}{}{}}
\LANGes{
\Question{
Elige el gráfico de una función $f$ que satisface
\begin{gather*}
f'(1)  \text{ no existe }; \\
f''(x) < 0 \text{ si } x < 1 ; \\
f''(x) < 0 \text{ si } x > 2; \\
f''(x) > 0 \text{ si } 1 < x < 2
\end{gather*}
($f'$ es la derivada de la función $f$, $f''$ es la segunda derivada de la función $f$).}
{}{}{}{}{}{}}

\IMGanswer{2010012501-figure4.pdf}
\IMGanswer{2010012501-figure5.pdf}
\IMGanswer{2010012501-figure6.pdf}
\IMGanswer{2010012501-figure7.pdf}\End

\Data\ID{2110013903}
\Updated{2023-01-23 10:01:12}
\Level{B}
\SubArea{Applications of derivatives}
\LANGen{
\Question{
The pictures show parts of graphs of even functions. Choose the picture showing the part of the graph of the function \(f(x)=\frac4{1+x^2}\).}
{}{}{}{}{}{}}
\LANGpl{
\Question{
Zdjęcia przedstawiają fragmenty wykresów funkcji parzystych. Wybierz rysunek przedstawiający część wykresu funkcji \(f(x)=\frac4{1+x^2}\).}
{}{}{}{}{}{}}
\LANGcs{
\Question{
Na obrázcích jsou části grafů sudých funkcí. Vyberte obrázek, na kterém je část grafu funkce \(f(x)=\frac4{1+x^2}\).}
{}{}{}{}{}{}}
\LANGsk{
\Question{
Na obrázkoch sú časti grafou párnych funkcií.  Vyberte obrázok, na ktorom je čásť grafu funkcie \(f(x)=\frac4{1+x^2}\).}
{}{}{}{}{}{}}
\LANGes{
\Question{
Las imágenes muestran una parte de gráficas de funciones pares. Elige la imagen que muestra esta parte de la gráfica de la función \(f(x)=\frac4{1+x^2}\).}
{}{}{}{}{}{}}

\IMGanswer{obr1_8.pdf}
\IMGanswer{obr2_6.pdf}
\IMGanswer{obr3_7.pdf}
\IMGanswer{obr4_5.pdf}\End

\Data\ID{2110013904}
\Updated{2023-01-23 10:01:16}
\Level{B}
\SubArea{Applications of derivatives}
\LANGen{
\Question{
The pictures show part of graphs of even functions. Choose the picture showing the part of the graph of the function \(f(x)=-\frac4{2+x^2}\).}
{}{}{}{}{}{}}
\LANGpl{
\Question{
Zdjęcia przedstawiają część wykresów funkcji parzystych. Wybierz obrazek przedstawiający część wykresu funkcji \(f(x)=-\frac4{2+x^2}\).}
{}{}{}{}{}{}}
\LANGcs{
\Question{
Na obrázcích jsou části grafů sudých funkcí. Vyberte obrázek, na kterém je část grafu funkce \(f(x)=-\frac4{2+x^2}\).}
{}{}{}{}{}{}}
\LANGsk{
\Question{
Na obrázkoch sú časti grafou párnych funkcií.  Vyberte obrázok, na ktorom je čásť grafu funkcie \(f(x)=-\frac4{2+x^2}\).}
{}{}{}{}{}{}}
\LANGes{
\Question{
Las imágenes muestran una parte de gráficas de funciones pares. Elige la imagen que muestra esta parte de la gráfica de la función \(f(x)=-\frac4{2+x^2}\).}
{}{}{}{}{}{}}

\IMGanswer{obr1_12.pdf}
\IMGanswer{obr2_10.pdf}
\IMGanswer{obr3_11.pdf}
\IMGanswer{obr4_9.pdf}\End

\Data\ID{9000070301}
\Updated{2020-07-15 03:07:37}
\Level{B}
\SubArea{Applications of derivatives}
\LANGen{
\Question{
Given function \(f(x)= x^{3} - 9x^{2} + 12x + 6\), find
the intervals where \(f\) 
is a strictly concave down function.}
{\((-\infty ;3)\)}{\((-\infty ;4)\)}{\((-\infty ;6)\)}{\((-\infty ;12)\)}{}{}}
\LANGpl{
\Question{
Dana jest funkcja  \(f\colon y = x^{3} - 9x^{2} + 12x + 6\), wskaż przedziały, w których  funkcja 
 \(f\) 
jest funkcją ściśle wklęsłą.}
{\((-\infty ;3)\)}{\((-\infty ;4)\)}{\((-\infty ;6)\)}{\((-\infty ;12)\)}{}{}}
\LANGcs{
\Question{
Je dána funkce \(f\colon y = x^{3} - 9x^{2} + 12x + 6\).
Ve kterém z následujících intervalů je tato funkce ryze konkávní?}
{\((-\infty ;3)\)}{\((-\infty ;4)\)}{\((-\infty ;6)\)}{\((-\infty ;12)\)}{}{}}
\LANGsk{
\Question{
Je daná funkcia \(f\colon y = x^{3} - 9x^{2} + 12x + 6\).
V ktorom z následujúcich intervalov je táto funkcia rýdzo konkávna?}
{\((-\infty ;3)\)}{\((-\infty ;4)\)}{\((-\infty ;6)\)}{\((-\infty ;12)\)}{}{}}
\LANGes{
\Question{
Dada la función \(f(x)= x^{3} - 9x^{2} + 12x + 6\), halla el intervalo donde \(f\) 
es una función estrictamente cóncava hacia abajo.}
{\((-\infty ;3)\)}{\((-\infty ;4)\)}{\((-\infty ;6)\)}{\((-\infty ;12)\)}{}{}}
\End

\Data\ID{9000070302}
\Updated{2022-04-23 05:04:05}
\Level{B}
\SubArea{Applications of derivatives}
\LANGen{
\Question{
Given function \(f(x) = x^{3} + 3x^{2} + 12x + 4\), find
the intervals where \(f\) 
is a strictly concave down function.}
{\((-\infty ;-1)\)}{\((-\infty ;0)\)}{\((-\infty ;2)\)}{\((-\infty ;4)\)}{}{}}
\LANGpl{
\Question{
Dana jest funkcja  \(f\colon y = x^{3} + 3x^{2} + 12x + 4\), wskaż przedziały, w których  funkcja 
 \(f\) 
jest funkcją ściśle wklęsłą.}
{\((-\infty ;-1)\)}{\((-\infty ;0)\)}{\((-\infty ;2)\)}{\((-\infty ;4)\)}{}{}}
\LANGcs{
\Question{
Je dána funkce \(f\colon y = x^{3} + 3x^{2} + 12x + 4\).
Ve kterém z následujících intervalů je tato funkce ryze konkávní?}
{\((-\infty ;-1)\)}{\((-\infty ;0)\)}{\((-\infty ;2)\)}{\((-\infty ;4)\)}{}{}}
\LANGsk{
\Question{
Je daná funkcie \(f\colon y = x^{3} + 3x^{2} + 12x + 4\).
V ktorom z následujúcich intervalov je táto funkcia rýdzo konkávna?}
{\((-\infty ;-1)\)}{\((-\infty ;0)\)}{\((-\infty ;2)\)}{\((-\infty ;4)\)}{}{}}
\LANGes{
\Question{
Dada la función \(f(x) = x^{3} + 3x^{2} + 12x + 4\), halla el intervalo donde \(f\) 
es una función estrictamente cóncava hacia abajo.}
{\((-\infty ;-1)\)}{\((-\infty ;0)\)}{\((-\infty ;2)\)}{\((-\infty ;4)\)}{}{}}
\End

\Data\ID{9000070303}
\Updated{2022-01-05 05:01:36}
\Level{B}
\SubArea{Applications of derivatives}
\LANGen{
\Question{
Given function \(f(x) = -x^{3} + 6x^{2} + 6x + 1\), find
the intervals where \(f\) 
is a strictly concave up function.}
{\((-\infty ;2)\)}{\((-1;\infty )\)}{\((-\infty ;3)\)}{\((-\infty ;6)\)}{}{}}
\LANGpl{
\Question{
Dana jest funkcja   \(f\colon y = -x^{3} + 6x^{2} + 6x + 1\),  wskaż przedziały, w których  funkcja 
 \(f\) 
jest funkcją ściśle wypukłą.}
{\((-\infty ;2)\)}{\((-1;\infty )\)}{\((-\infty ;3)\)}{\((-\infty ;6)\)}{}{}}
\LANGcs{
\Question{
Je dána funkce \(f\colon y = -x^{3} + 6x^{2} + 6x + 1\).
Ve kterém z následujících intervalů je tato funkce ryze konvexní?}
{\((-\infty ;2)\)}{\((-1;\infty )\)}{\((-\infty ;3)\)}{\((-\infty ;6)\)}{}{}}
\LANGsk{
\Question{
Je daná funkcia \(f\colon y = -x^{3} + 6x^{2} + 6x + 1\).
V ktorom z následujúcich intervalov je tato funkcia rýdzo konvexná?}
{\((-\infty ;2)\)}{\((-1;\infty )\)}{\((-\infty ;3)\)}{\((-\infty ;6)\)}{}{}}
\LANGes{
\Question{
Dada la función \(f(x) = -x^{3} + 6x^{2} + 6x + 1\), halla el intervalo donde \(f\) 
es una función estrictamente convexa.}
{\((-\infty ;2)\)}{\((-1;\infty )\)}{\((-\infty ;3)\)}{\((-\infty ;6)\)}{}{}}
\End

\Data\ID{9000070304}
\Updated{2020-07-15 03:07:37}
\Level{B}
\SubArea{Applications of derivatives}
\LANGen{
\Question{
Given function \(f(x) = -x^{3} - 12x^{2} + 12x - 2\), find
the intervals where \(f\) 
is a strictly concave up function.}
{\((-\infty ;-4)\)}{\((-\infty ;2)\)}{\((-\infty ;6)\)}{\((-\infty ;12)\)}{}{}}
\LANGpl{
\Question{
Dana jest funkcja  \(f\colon y = -x^{3} - 12x^{2} + 12x - 2\), wskaż przedziały, w których  funkcja 
 \(f\) 
jest funkcją ściśle wypukłą.}
{\((-\infty ;-4)\)}{\((-\infty ;2)\)}{\((-\infty ;6)\)}{\((-\infty ;12)\)}{}{}}
\LANGcs{
\Question{
Je dána funkce \(f\colon y = -x^{3} - 12x^{2} + 12x - 2\).
Ve kterém z následujících intervalů je tato funkce ryze konvexní?}
{\((-\infty ;-4)\)}{\((-\infty ;2)\)}{\((-\infty ;6)\)}{\((-\infty ;12)\)}{}{}}
\LANGsk{
\Question{
Je daná funkcia \(f\colon y = -x^{3} - 12x^{2} + 12x - 2\).
V ktorom z následujúcich intervalov je tato funkcia rýdzo konvexná?}
{\((-\infty ;-4)\)}{\((-\infty ;2)\)}{\((-\infty ;6)\)}{\((-\infty ;12)\)}{}{}}
\LANGes{
\Question{
Dada la función \(f(x) = -x^{3} - 12x^{2} + 12x - 2\), halla el intervalo donde \(f\) 
es una función estrictamente convexa.}
{\((-\infty ;-4)\)}{\((-\infty ;2)\)}{\((-\infty ;6)\)}{\((-\infty ;12)\)}{}{}}
\End

\Data\ID{9000070305}
\Updated{2022-04-23 05:04:57}
\Level{B}
\SubArea{Applications of derivatives}
\LANGen{
\Question{
Given function \(f(x) = x^{4} + 2x^{3} - 36x^{2} + 36x + 2\), find
the intervals where \(f\) 
is a strictly concave down function.}
{\((-3;2)\)}{\((-3;4)\)}{\((-4;2)\)}{\((-2;3)\)}{}{}}
\LANGpl{
\Question{
Dana jest funkcja  \(f\colon y = x^{4} + 2x^{3} - 36x^{2} + 36x + 2\),  wskaż przedziały, w których  funkcja 
 \(f\) 
jest funkcją ściśle wklęsłą.}
{\((-3;2)\)}{\((-3;4)\)}{\((-4;2)\)}{\((-2;3)\)}{}{}}
\LANGcs{
\Question{
Je dána funkce \(f\colon y = x^{4} + 2x^{3} - 36x^{2} + 36x + 2\).
Ve kterém z následujících intervalů je tato funkce ryze konkávní?}
{\((-3;2)\)}{\((-3;4)\)}{\((-4;2)\)}{\((-2;3)\)}{}{}}
\LANGsk{
\Question{
Je daná funkcia \(f\colon y = x^{4} + 2x^{3} - 36x^{2} + 36x + 2\).
V ktorom z následujúcich intervalov je tato funkcia rýdzo konkávna?}
{\((-3;2)\)}{\((-3;4)\)}{\((-4;2)\)}{\((-2;3)\)}{}{}}
\LANGes{
\Question{
Dada la función \(f(x) = x^{4} + 2x^{3} - 36x^{2} + 36x + 2\), halla el intervalo donde \(f\) 
es una función estrictamente cóncava hacia abajo.}
{\((-3;2)\)}{\((-3;4)\)}{\((-4;2)\)}{\((-2;3)\)}{}{}}
\End

\Data\ID{9000070306}
\Updated{2022-01-05 06:01:17}
\Level{B}
\SubArea{Applications of derivatives}
\LANGen{
\Question{
Given function \(f(x) = x^{4} + 6x^{3} - 24x^{2} + x + 3\), find
the intervals where \(f\) 
is a strictly concave down function.}
{\((-4;1)\)}{\((-6;2)\)}{\((2;4)\)}{\((-5;4)\)}{}{}}
\LANGpl{
\Question{
Dana jest funkcja  \(f\colon y = x^{4} + 6x^{3} - 24x^{2} + x + 3\),  wskaż przedziały, w których  funkcja 
 \(f\) 
jest funkcją ściśle wklęsłą.}
{\((-4;1)\)}{\((-6;2)\)}{\((2;4)\)}{\((-5;4)\)}{}{}}
\LANGcs{
\Question{
Je dána funkce \(f\colon y = x^{4} + 6x^{3} - 24x^{2} + x + 3\).
Ve kterém z následujících intervalů je tato funkce ryze konkávní?}
{\((-4;1)\)}{\((-6;2)\)}{\((2;4)\)}{\((-5;4)\)}{}{}}
\LANGsk{
\Question{
Je daná funkcia \(f\colon y = x^{4} + 6x^{3} - 24x^{2} + x + 3\).
V ktorom z následujúcich intervalov je tato funkcia rýdzo konkávna?}
{\((-4;1)\)}{\((-6;2)\)}{\((2;4)\)}{\((-5;4)\)}{}{}}
\LANGes{
\Question{
Dada la función \(f(x) = x^{4} + 6x^{3} - 24x^{2} + x + 3\), halla el intervalo donde \(f\) 
es una función estrictamente cóncava hacia abajo.}
{\((-4;1)\)}{\((-6;2)\)}{\((2;4)\)}{\((-5;4)\)}{}{}}
\End

\Data\ID{9000079108}
\Updated{2021-08-15 08:08:26}
\Level{B}
\SubArea{Applications of derivatives}
\LANGen{
\Question{
Find the $x$ at which the function $f$ has the global minimum on the interval
\((-3;2] \).
\[ 
 f(x) = x^{3} - 3x + 4
\]}
{does not exist}{\(x = -3\)}{\(x = -2\)}{\(x = 1\)}{}{}}
\LANGpl{
\Question{
Wyznacz globalne minimum funkcji w przedziale
\((-3;2\rangle \).
\[ 
 f\colon y = x^{3} - 3x + 4
\]}
{nie istnieje}{\(x = -3\)}{\(x = -2\)}{\(x = 1\)}{}{}}
\LANGcs{
\Question{
Doplňte správné tvrzení: „Globální minimum funkce
\(f\colon y = x^{3} - 3x + 4\) na
intervalu \((-3;2\rangle \) 
...”}
{neexistuje.}{je v bodě \(x=- 3\).}{je v bodě \(x=- 2\).}{je v bodě \(x=1\).}{}{}}
\LANGsk{
\Question{
Doplňte správne tvrdenie: „Globálne minimum funkcie
\(f\colon y = x^{3} - 3x + 4\) na
intervale \((-3;2\rangle \) 
...”}
{neexistuje.}{je v bode \(x=- 3\).}{je v bode \(x=- 2\).}{je v bode \(x=1\).}{}{}}
\LANGes{
\Question{
Hallael valor de $x$ para el cual la función $f$ tiene su mínimo global en el intervalo
\((-3;2] \).
\[ 
 f(x) = x^{3} - 3x + 4
\]}
{no existe}{\(x = -3\)}{\(x = -2\)}{\(x = 1\)}{}{}}
\End

\Data\ID{9000079109}
\Updated{2021-08-15 09:08:23}
\Level{B}
\SubArea{Applications of derivatives}
\LANGen{
\Question{
Find the $x$ at which the function $f$ has the global maximum on the interval
\([ 1;\mathrm{e}] \).
\[ 
 f(x)= x - 2\ln x
\]}
{\(x=1\)}{\(x=2\)}{\(x=\mathrm{e}\)}{\(x=\mathrm{e} - 2\)}{}{}}
\LANGpl{
\Question{
Wyznacz globalne maksimum funkcji w przedziale
\(\langle 1;\mathrm{e}\rangle \).
\[ 
 f\colon y = x - 2\ln x
\]}
{\(x=1\)}{\(x=2\)}{\(x=\mathrm{e}\)}{\(x=\mathrm{e} - 2\)}{}{}}
\LANGcs{
\Question{
Funkce \(f\colon y = x - 2\ln x\) má
na intervalu \(\langle 1;\mathrm{e}\rangle \) 
globální maximum v bodě:}
{\(x=1\)}{\(x=2\)}{\(x=\mathrm{e}\)}{\(x=\mathrm{e} - 2\)}{}{}}
\LANGsk{
\Question{
Funkcia \(f\colon y = x - 2\ln x\) má
na intervale \(\langle 1;\mathrm{e}\rangle \) 
globálne maximum v bode:}
{\(x=1\)}{\(x=2\)}{\(x=\mathrm{e}\)}{\(x=\mathrm{e} - 2\)}{}{}}
\LANGes{
\Question{
Halla el valor de $x$ para el cual la función $f$ tiene su máximo global en el intervalo
\([ 1;\mathrm{e}] \).
\[ 
 f(x)= x - 2\ln x
\]}
{\(x=1\)}{\(x=2\)}{\(x=\mathrm{e}\)}{\(x=\mathrm{e} - 2\)}{}{}}
\End

\Data\ID{9000142001}
\Updated{2020-07-15 03:07:37}
\Level{B}
\SubArea{Applications of derivatives}
\IMG{001420_01-figure0.pdf}
\LANGen{
\Question{
Identify a correct statement related to the function $f$ shown in the picture.}
{concave up on \((-1;0)\) 
and \((1;\infty )\), concave
down on \((-\infty ;-1)\) 
and \((0;1)\),
inflection at \(x = 0\)}{concave up on \((-\infty ;-1)\) 
and \((0;1)\), concave
down on \((-1;0)\) 
and \((1;\infty )\),
inflection at \(x = 0\)}{concave up on \((-1;0)\) 
and \((1;\infty )\), concave
down on \((-\infty ;-1)\) 
and \((0;1)\),
no inflection}{concave up on \((-1;0)\cup (1;\infty )\),
concave down on \((-\infty ;-1)\cup (0;1)\),
inflection at \(x = 0\)}{}{}}
\LANGpl{
\Question{
Dana jest funkcja  \(f\),
wskaż zdanie prawdziwe.}
{Funkcja jest wypukła w przedziale  \((-1;0)\) 
i  \((1;\infty )\), wklęsła w przedziale 
 \((-\infty ;-1)\) 
i \((0;1)\),
punkt przegięcia funkcji to  \(x = 0\)}{Funkcja jest wypukła w przedziale  \((-\infty ;-1)\) 
i  \((0;1)\), wklęsła w przedziale 
 \((-1;0)\) 
i  \((1;\infty )\),
punkt przegięcia funkcji to \(x = 0\)}{Funkcja jest wypukła w przedziale  \((-1;0)\) 
i  \((1;\infty )\), wklęsła w przedziale 
 \((-\infty ;-1)\) 
i  \((0;1)\),
brak punktu przegięcia}{Funkcja jest wypukła w przedziale  \((-1;0)\cup (1;\infty )\),
wklęsła w przedziale  \((-\infty ;-1)\cup (0;1)\),
punkt przegięcia funkcji to \(x = 0\)}{}{}}
\LANGcs{
\Question{
Rozhodněte, které z následujících vlastností má funkce $f$ na obrázku.}
{konvexní v \((-1;0)\) 
a \((1;\infty )\),
konkávní v \((-\infty ;-1)\) 
a \((0;1)\), inflexní
bod \(x = 0\)}{konvexní v \((-\infty ;-1)\) 
a \((0;1)\),
konkávní v \((-1;0)\) 
a \((1;\infty )\), inflexní
bod \(x = 0\)}{konvexní v \((-1;0)\) 
a \((1;\infty )\),
konkávní v \((-\infty ;-1)\) 
a \((0;1)\),
inflexní bod neexistuje}{konvexní v \((-1;0)\cup (1;\infty )\),
konkávní v \((-\infty ;-1)\cup (0;1)\),
inflexní bod \(x = 0\)}{}{}}
\LANGsk{
\Question{
Rozhodnite, ktoré z nasledujúcich vlastností má funkcie $f$ na obrázku.}
{konvexná v \((-1;0)\) 
a \((1;\infty )\),
konkávna v \((-\infty ;-1)\) 
a \((0;1)\), inflexný
bod \(x = 0\)}{konvexná v \((-\infty ;-1)\) 
a \((0;1)\),
konkávna v \((-1;0)\) 
a \((1;\infty )\), inflexný
bod \(x = 0\)}{konvexná v \((-1;0)\) 
a \((1;\infty )\),
konkávna v \((-\infty ;-1)\) 
a \((0;1)\),
inflexný bod neexistuje}{konvexná v \((-1;0)\cup (1;\infty )\),
konkávna v \((-\infty ;-1)\cup (0;1)\),
inflexný bod \(x = 0\)}{}{}}
\LANGes{
\Question{
Identifica la proposición lógica sobre la función $f$ representada en la imagen.}
{convexa en \((-1;0)\) 
y \((1;\infty )\), cóncava
en \((-\infty ;-1)\) 
y \((0;1)\),
inflexión en \(x = 0\)}{convexa en \((-\infty ;-1)\) 
y \((0;1)\), cóncava
en \((-1;0)\) 
y \((1;\infty )\),
inflexión en \(x = 0\)}{convexa en \((-1;0)\) 
y \((1;\infty )\), cóncava
en \((-\infty ;-1)\) 
y \((0;1)\),
no tiene inflexión}{convexa en \((-1;0)\cup (1;\infty )\),
cóncava
en \((-\infty ;-1)\cup (0;1)\),
inflexión en \(x = 0\)}{}{}}
\End

\Data\ID{9000142002}
\Updated{2020-07-15 03:07:37}
\Level{B}
\SubArea{Applications of derivatives}
\IMG{001420_02-figure0.pdf}
\LANGen{
\Question{
Identify a correct statement related to the function $f$ shown in the picture.}
{concave up on \((-\infty ;1)\),
concave down on \((1;\infty )\),
inflection at \(x = 1\)}{concave up on \((1;\infty )\),
concave down on \((-\infty ;1)\),
inflection at \(x = 1\)}{concave up on \((-\infty ;0)\),
concave down on \((0;\infty )\),
inflection at \(x = 0\)}{concave up on \((-\infty ;1)\),
concave down on \((1;\infty )\),
inflection at \(x = \frac{2} 
{3}\)}{}{}}
\LANGpl{
\Question{
Dana jest funkcja \(f\),
 wskaż zdanie prawdziwe.}
{Funkcja jest wypukła w przedziale  \((-\infty ;1)\),
wklęsła w przedziale  \((1;\infty )\),
punkt przegięcia funkcji to \(x = 1\)}{Funkcja jest wypukła w przedziale \((1;\infty )\),
wklęsła w przedziale  \((-\infty ;1)\),
punkt przegięcia funkcji to \(x = 1\)}{Funkcja jest wypukła w przedziale  \((-\infty ;0)\),
wklęsła w przedziale \((0;\infty )\),
punkt przegięcia funkcji to \(x = 0\)}{Funkcja jest wypukła w przedziale  \((-\infty ;1)\),
wklęsła w przedziale  \((1;\infty )\),
punkt przegięcia funkcji to \(x = \frac{2} 
{3}\)}{}{}}
\LANGcs{
\Question{
Rozhodněte, které z následujících vlastností má funkce $f$ na obrázku.}
{konvexní v \((-\infty ;1)\),
konkávní \((1;\infty )\),
inflexní bod \(x = 1\)}{konvexní v \((1;\infty )\),
konkávní \((-\infty ;1)\),
inflexní bod \(x = 1\)}{konvexní v \((-\infty ;0)\),
konkávní \((0;\infty )\),
inflexní bod \(x = 0\)}{konvexní v \((-\infty ;1)\),
konkávní \((1;\infty )\),
inflexní bod \(x = \frac{2} 
{3}\)}{}{}}
\LANGsk{
\Question{
Rozhodnite, ktoré z nasledujúcich vlastností má funkcie $ f $ na obrázku.}
{konvexná v \((-\infty ;1)\),
konkávna \((1;\infty )\),
inflexný bod \(x = 1\)}{konvexná v \((1;\infty )\),
konkávna \((-\infty ;1)\),
inflexný bod \(x = 1\)}{konvexná v \((-\infty ;0)\),
konkávna \((0;\infty )\),
inflexný bod \(x = 0\)}{konvexná v \((-\infty ;1)\),
konkávna \((1;\infty )\),
inflexný bod \(x = \frac{2} 
{3}\)}{}{}}
\LANGes{
\Question{
Identifica la proposición lógica sobre la función $f$ representada en la imagen.}
{convexa en \((-\infty ;1)\),
cóncava en \((1;\infty )\),
inflexión en \(x = 1\)}{convexa en \((1;\infty )\),
cóncava en \((-\infty ;1)\),
inflexión en \(x = 1\)}{convexa en \((-\infty ;0)\),
cóncava en \((0;\infty )\),
inflexión en \(x = 0\)}{convexa en \((-\infty ;1)\),
cóncava en \((1;\infty )\),
inflexión en \(x = \frac{2} 
{3}\)}{}{}}
\End

\Data\ID{9000142003}
\Updated{2020-07-15 03:07:37}
\Level{B}
\SubArea{Applications of derivatives}
\IMG{001420_03-figure0.pdf}
\LANGen{
\Question{
Identify a correct statement related to the function $f$ shown in the picture.}
{concave up on \((-\infty ;0)\) 
and \((1;\infty )\), concave
down on \((0;1)\),
inflection at \(x_{1} = 0\) 
and \(x_{2} = 1\)}{concave up on \((-\infty ;0)\cup (1;\infty )\),
concave down on \((0;1)\),
inflection at \(x_{1} = 0\) 
and \(x_{2} = 1\)}{concave up on \((0;1)\),
concave down on \((-\infty ;0)\) 
and \((1;\infty )\),
inflection at \(x_{1} = 0\) 
and \(x_{2} = 1\)}{concave up on \((-\infty ;0)\) 
and \((1;\infty )\), concave
down on \((0;1)\), a unique
inflection at \(x = 0\)}{}{}}
\LANGpl{
\Question{
Dana jest funkcja \(f\),
 wskaż zdanie prawdziwe.}
{Funkcja jest wypukła w przedziale \((-\infty ;0)\) 
i  \((1;\infty )\), wklęsła w przedziale 
 \((0;1)\),
punkt przegięcia funkcji to \(x_{1} = 0\) 
i  \(x_{2} = 1\)}{Funkcja jest wypukła w przedziale \((-\infty ;0)\cup (1;\infty )\),
wklęsła w przedziale  \((0;1)\),
punkty przegięcia funkcji to \(x_{1} = 0\) 
i  \(x_{2} = 1\)}{Funkcja jest wypukła w przedziale \((0;1)\),
wklęsła w przedziale  \((-\infty ;0)\) 
i  \((1;\infty )\),
punkty przegięcia funkcji to \(x_{1} = 0\) 
i  \(x_{2} = 1\)}{Funkcja jest wypukła w przedziale \((-\infty ;0)\) 
i \((1;\infty )\), wklęsła w przedziale  \((0;1)\), jedyny
punkt przegięcia funkcji to \(x = 0\)}{}{}}
\LANGcs{
\Question{
Rozhodněte, které z následujících vlastností má funkce $f$ na obrázku.}
{konvexní v \((-\infty ;0)\) a
\((1;\infty )\), konkávní v
\((0;1)\), inflexní body
\(x_{1} = 0\), \(x_{2} = 1\)}{konvexní v \((-\infty ;0)\cup (1;\infty )\), konkávní
v \((0;1)\), inflexní body
\(x_{1} = 0\), \(x_{2} = 1\)}{konvexní v \((0;1)\),
konkávní v \((-\infty ;0)\) a
\((1;\infty )\), inflexní body
\(x_{1} = 0\), \(x_{2} = 1\)}{konvexní v \((-\infty ;0)\) 
a \((1;\infty )\), konkávní
v \((0;1)\), jediný
inflexní bod \(x = 0\)}{}{}}
\LANGsk{
\Question{
Rozhodnite, ktoré z nasledujúcich vlastností má funkcie $ f $ na obrázku.}
{konvexná v \((-\infty ;0)\) a
\((1;\infty )\), konkávna v
\((0;1)\), inflexné body
\(x_{1} = 0\), \(x_{2} = 1\)}{konvexná v \((-\infty ;0)\cup (1;\infty )\), konkávna
v \((0;1)\), inflexné body
\(x_{1} = 0\), \(x_{2} = 1\)}{konvexná v \((0;1)\),
konkávna v \((-\infty ;0)\) a
\((1;\infty )\), inflexné body
\(x_{1} = 0\), \(x_{2} = 1\)}{konvexná v \((-\infty ;0)\) 
a \((1;\infty )\), konkávna
v \((0;1)\), jediný
inflexný bod \(x = 0\)}{}{}}
\LANGes{
\Question{
Identifica la proposición lógica sobre la función $f$ representada en la imagen.}
{convexa en \((-\infty ;0)\) 
y \((1;\infty )\), cóncava en \((0;1)\),
inflexión en \(x_{1} = 0\) 
y \(x_{2} = 1\)}{convexa en \((-\infty ;0)\cup (1;\infty )\),
cóncava en \((0;1)\),
inflexión en \(x_{1} = 0\) 
y \(x_{2} = 1\)}{convexa en \((0;1)\),
cóncava en \((-\infty ;0)\) 
y \((1;\infty )\),
inflexión en \(x_{1} = 0\) 
y \(x_{2} = 1\)}{convexa en \((-\infty ;0)\) 
y \((1;\infty )\), cóncava en \((0;1)\), única inflexión en \(x = 0\)}{}{}}
\End

\Data\ID{9000142004}
\Updated{2020-07-15 03:07:37}
\Level{B}
\SubArea{Applications of derivatives}
\IMG{001420_04-figure0.pdf}
\LANGen{
\Question{
Identify a correct statement related to the function $f$ shown in the picture.}
{concave up on \((-\infty ;1)\),
concave down on \((1;\infty )\),
no inflection}{concave up on \((-\infty ;1)\),
concave down on \((1;\infty )\),
inflection at \(x = 1\)}{concave up on \((1;\infty )\),
concave down on \((-\infty ;1)\),
inflection at \(x = 1\)}{concave up on \((1;\infty )\),
concave down on \((-\infty ;1)\),
no inflection}{}{}}
\LANGpl{
\Question{
Dana jest funkcja \(f\),
wskaż zdanie prawdziwe.}
{Funkcja jest wypukła w przedziale \((-\infty ;1)\),
wklęsła w przedziale  \((1;\infty )\),
brak punktu przegięcia funkcji}{Funkcja jest wypukła w przedziale \((-\infty ;1)\),
wklęsła w przedziale  \((1;\infty )\),
punkt przegięcia funkcji to \(x = 1\)}{Funkcja jest wypukła w przedziale \((1;\infty )\),
wklęsła w przedziale  \((-\infty ;1)\),
punkt przegięcia funkcji to \(x = 1\)}{Funkcja jest wypukła w przedziale \((1;\infty )\),
wklęsła w przedziale \((-\infty ;1)\),
brak punktu przegięcia funkcji}{}{}}
\LANGcs{
\Question{
Rozhodněte, které z následujících vlastností má funkce $f$ na obrázku.}
{konvexní v \((-\infty ;1)\),
konkávní v \((1;\infty )\),
inflexní bod neexistuje}{konvexní v \((-\infty ;1)\),
konkávní v \((1;\infty )\),
inflexní bod \(x = 1\)}{konvexní v \((1;\infty )\),
konkávní v \((-\infty ;1)\),
inflexní bod \(x = 1\)}{konvexní v \((1;\infty )\),
konkávní v \((-\infty ;1)\),
inflexní bod neexistuje}{}{}}
\LANGsk{
\Question{
Rozhodnite, ktoré z nasledujúcich vlastností má funkcie $ f $ na obrázku.}
{konvexná v \((-\infty ;1)\),
konkávna v \((1;\infty )\),
inflexný bod neexistuje}{konvexná v \((-\infty ;1)\),
konkávna v \((1;\infty )\),
inflexný bod \(x = 1\)}{konvexná v \((1;\infty )\),
konkávna v \((-\infty ;1)\),
inflexný bod \(x = 1\)}{konvexná v \((1;\infty )\),
konkávna v \((-\infty ;1)\),
inflexný bod neexistuje}{}{}}
\LANGes{
\Question{
Identifica la proposición lógica sobre la función $f$ representada en la imagen.}
{convexa en \((-\infty ;1)\),
cóncava en  \((1;\infty )\),
no tiene inflexión}{convexa en \((-\infty ;1)\),
cóncava en  \((1;\infty )\),
inflexión en \(x = 1\)}{convexa en \((1;\infty )\),
cóncava en  \((-\infty ;1)\),
inflexión en \(x = 1\)}{convexa en \((1;\infty )\),
cóncava en  \((-\infty ;1)\),
no tiene inflexión}{}{}}
\End

\Data\ID{9000142005}
\Updated{2020-07-15 03:07:37}
\Level{B}
\SubArea{Applications of derivatives}
\IMG{001420_05-figure0.pdf}
\LANGen{
\Question{
Identify a correct statement related to the function $f$ shown in the picture.}
{concave up on \((-1;0)\) 
and \((1;\infty )\), concave
down on \((-\infty ;-1)\) 
and \((0;1)\),
inflection at \(x_{1} = -1\),
\(x_{2} = 0\) and
\(x_{3} = 1\)}{concave up on \((-1;0)\cup (1;\infty )\),
concave down on \((-\infty ;-1)\cup (0;1)\),
inflection at \(x_{1} = -1\),
\(x_{2} = 0\) and
\(x_{3} = 1\)}{concave up on \((-\infty ;-1)\) 
and \((0;1)\), concave
down on \((-1;0)\) 
and \((1;\infty )\),
inflection at \(x_{1} = -1\),
\(x_{2} = 0\) and
\(x_{3} = 1\)}{concave up on \((-\infty ;-1)\cup (0;1)\),
concave down on \((-1;0)\cup (1;\infty )\),
inflection at \(x_{1} = -1\),
\(x_{2} = 0\) and
\(x_{3} = 1\)}{}{}}
\LANGpl{
\Question{
Dana jest funkcja  \(f\),
wskaż zdanie prawdziwe.}
{Funkcja jest wypukła w przedziale \((-1;0)\) 
i \((1;\infty )\), wklęsła w przedziale
 \((-\infty ;-1)\) 
i \((0;1)\),
punkty przegięcia funkcji to \(x_{1} = -1\),
\(x_{2} = 0\) i
\(x_{3} = 1\)}{Funkcja jest wypukła w przedziale \((-1;0)\cup (1;\infty )\),
wklęsła w przedziale  \((-\infty ;-1)\cup (0;1)\),
punkty przegięcia funkcji to \(x_{1} = -1\),
\(x_{2} = 0\) i
\(x_{3} = 1\)}{Funkcja jest wypukła w przedziale \((-\infty ;-1)\) 
i \((0;1)\), wklęsła w przedziale 
 \((-1;0)\) 
i  \((1;\infty )\),
punkty przegięcia funkcji to \(x_{1} = -1\),
\(x_{2} = 0\) i
\(x_{3} = 1\)}{Funkcja jest wypukła w przedziale \((-\infty ;-1)\cup (0;1)\),
wklęsła w przedziale  \((-1;0)\cup (1;\infty )\),
punkty przegięcia funkcji to \(x_{1} = -1\),
\(x_{2} = 0\) i
\(x_{3} = 1\)}{}{}}
\LANGcs{
\Question{
Rozhodněte, které z následujících vlastností má funkce $f$ na obrázku.}
{konvexní v \((-1;0)\) 
a \((1;\infty )\),
konkávní v \((-\infty ;-1)\) 
a \((0;1)\), inflexní
body \(x_{1} = -1\),
\(\ x_{2} = 0,\ x_{3} = 1\)}{konvexní v \((-1;0)\cup (1;\infty )\),
konkávní v \((-\infty ;-1)\cup (0;1)\),
inflexní body \(x_{1} = -1\),
\(\ x_{2} = 0,\ x_{3} = 1\)}{konvexní v \((-\infty ;-1)\) 
a \((0;1)\),
konkávní v \((-1;0)\) 
a \((1;\infty )\), inflexní
body \(x_{1} = -1\),
\(\ x_{2} = 0,\ x_{3} = 1\)}{konvexní v \((-\infty ;-1)\cup (0;1)\),
konkávní v \((-1;0)\cup (1;\infty )\),
inflexní body \(x_{1} = -1\),
\(\ x_{2} = 0,\ x_{3} = 1\)}{}{}}
\LANGsk{
\Question{
Rozhodnite, ktoré z nasledujúcich vlastností má funkcie $ f $ na obrázku.}
{konvexná v \((-1;0)\) 
a \((1;\infty )\),
konkávna v \((-\infty ;-1)\) 
a \((0;1)\), inflexné
body \(x_{1} = -1\),
\(\ x_{2} = 0,\ x_{3} = 1\)}{konvexná v \((-1;0)\cup (1;\infty )\),
konkávna v \((-\infty ;-1)\cup (0;1)\),
inflexné body \(x_{1} = -1\),
\(\ x_{2} = 0,\ x_{3} = 1\)}{konvexná v \((-\infty ;-1)\) 
a \((0;1)\),
konkávna v \((-1;0)\) 
a \((1;\infty )\), inflexné
body \(x_{1} = -1\),
\(\ x_{2} = 0,\ x_{3} = 1\)}{konvexná v \((-\infty ;-1)\cup (0;1)\),
konkávna v \((-1;0)\cup (1;\infty )\),
inflexné body \(x_{1} = -1\),
\(\ x_{2} = 0,\ x_{3} = 1\)}{}{}}
\LANGes{
\Question{
Identifica la proposición lógica sobre la función $f$ representada en la imagen.}
{convexa en \((-1;0)\) 
y\((1;\infty )\), cóncava en \((-\infty ;-1)\) 
y \((0;1)\),
inflexión en \(x_{1} = -1\),
\(x_{2} = 0\) y
\(x_{3} = 1\)}{convexa en \((-1;0)\cup (1;\infty )\),
cóncava en \((-\infty ;-1)\cup (0;1)\),
inflexión en \(x_{1} = -1\),
\(x_{2} = 0\) y
\(x_{3} = 1\)}{convexa en \((-\infty ;-1)\) 
y \((0;1)\), cóncava en \((-1;0)\) 
y \((1;\infty )\),
inflexión en \(x_{1} = -1\),
\(x_{2} = 0\) y
\(x_{3} = 1\)}{convexa en \((-\infty ;-1)\cup (0;1)\),
cóncava en \((-1;0)\cup (1;\infty )\),
inflexión en \(x_{1} = -1\),
\(x_{2} = 0\) y
\(x_{3} = 1\)}{}{}}
\End

\Data\ID{9000142006}
\Updated{2020-07-15 03:07:37}
\Level{B}
\SubArea{Applications of derivatives}
\IMG{001420_06-figure0.pdf}
\LANGen{
\Question{
Identify a correct statement related to the function $f$ shown in the picture.}
{concave up on \((-\infty ;0)\) 
and \((1;\infty )\), concave
down on \((0;1)\), a unique
inflection at \(x = 0\)}{concave up on \((-\infty ;0)\) 
and \((1;\infty )\), concave
down on \((0;1)\),
inflection at \(x_{1} = 0\) 
and \(x_{2} = 1\)}{concave up on \((-\infty ;0)\cup (1;\infty )\),
concave down on \((0;1)\), a
unique inflection at \(x = 0\)}{concave up on \((0;1)\),
concave down on \((-\infty ;0)\) 
and \((1;\infty )\),
inflection at \(x_{1} = 0\) 
and \(x_{2} = 1\)}{}{}}
\LANGpl{
\Question{
Dana jest funkcja \(f\),
wskaż zdanie prawdziwe.}
{Funkcja jest wypukła w przedziale \((-\infty ;0)\) 
i \((1;\infty )\), wklęsła w przedziale 
 \((0;1)\), jedyny punkt przegięcia funkcji to
 \(x = 0\)}{Funkcja jest wypukła w przedziale \((-\infty ;0)\) 
i \((1;\infty )\), wklęsła w przedziale 
 \((0;1)\),
punkty przegięcia funkcji to \(x_{1} = 0\) 
i \(x_{2} = 1\)}{Funkcja jest wypukła w przedziale \((-\infty ;0)\cup (1;\infty )\),
wklęsła w przedziale  \((0;1)\), 
jedyny punkt przegięcia funkcji to \(x = 0\)}{Funkcja jest wypukła w przedziale\((0;1)\),
wklęsła w przedziale  \((-\infty ;0)\) 
i \((1;\infty )\),
punkty przegięcia funkcji to \(x_{1} = 0\) 
i \(x_{2} = 1\)}{}{}}
\LANGcs{
\Question{
Rozhodněte, které z následujících vlastností má funkce $f$ na obrázku.}
{konvexní v \((-\infty ;0)\) a
\((1;\infty )\), konkávní
v \((0;1)\), jediný
inflexní bod \(x = 0\)}{konvexní v \((-\infty ;0)\) 
a \((1;\infty )\), konkávní
v \((0;1)\), inflexní
body \(x_{1} = 0,\ x_{2} = 1\)}{konvexní v \((-\infty ;0)\cup (1;\infty )\),
konkávní v \((0;1)\),
jediný inflexní bod \(x = 0\)}{konvexní v \((0;1)\),
konkávní v \((-\infty ;0)\) 
a \((1;\infty )\), inflexní
body \(x_{1} = 0,\ x_{2} = 1\)}{}{}}
\LANGsk{
\Question{
Rozhodnite, ktoré z nasledujúcich vlastností má funkcie $ f $ na obrázku.}
{konvexná v \((-\infty ;0)\) a
\((1;\infty )\), konkávna
v \((0;1)\), jediný
inflexný bod \(x = 0\)}{konvexná v \((-\infty ;0)\) 
a \((1;\infty )\), konkávna
v \((0;1)\), inflexný
body \(x_{1} = 0,\ x_{2} = 1\)}{konvexná v \((-\infty ;0)\cup (1;\infty )\),
konkávna v \((0;1)\),
jediný inflexný bod \(x = 0\)}{konvexná v \((0;1)\),
konkávna v \((-\infty ;0)\) 
a \((1;\infty )\), inflexné
body \(x_{1} = 0,\ x_{2} = 1\)}{}{}}
\LANGes{
\Question{
Identifica la proposición lógica sobre la función $f$ representada en la imagen.}
{convexa en \((-\infty ;0)\) 
y \((1;\infty )\), cóncava en \((0;1)\), única inflexión en \(x = 0\)}{convexa en \((-\infty ;0)\) 
y \((1;\infty )\), cóncava en \((0;1)\),
inflexión en \(x_{1} = 0\) 
y \(x_{2} = 1\)}{convexa en \((-\infty ;0)\cup (1;\infty )\),
cóncava en \((0;1)\), única inflexión en \(x = 0\)}{convexa en \((0;1)\),
cóncava en \((-\infty ;0)\) 
y \((1;\infty )\),
inflexión en \(x_{1} = 0\) 
y \(x_{2} = 1\)}{}{}}
\End

\Data\ID{9000145401}
\Updated{2020-07-15 03:07:37}
\Level{B}
\SubArea{Applications of derivatives}
\LANGen{
\Question{
Identify a true statement on the function
\(f(x) = 2x^{3} + 3x^{2} - 12x - 12\).}
{The function \(f\) has a local
maximum at the point \(x = -2\).}{The function \(f\) has a local
minimum at the point \(x = -2\)..}{The global maximum of \(f\) 
on \(\mathbb{R}\) is
at \(x = -2\).}{The global minimum of \(f\) 
on \(\mathbb{R}\) is
at \(x = -2\).}{}{}}
\LANGpl{
\Question{
Dana jest funkcja
\(f\colon y = 2x^{3} + 3x^{2} - 12x - 12\), wskaż zdanie prawdziwe.}
{Funkcja  \(f\) ma lokalne maksimum w punkcie  \(x = -2\).}{Funkcja  \(f\) ma lokalne
minimum w punkcie \(x = -2\).}{Maksimum globalne  funkcji \(f\) 
na zbiorze  \(\mathbb{R}\) jest w punkcie
  \(x = -2\).}{Minimum globalne funkcji \(f\) 
na zbiorze  \(\mathbb{R}\) jest w punkcie
 \(x = -2\).}{}{}}
\LANGcs{
\Question{
Je dána funkce \(f\colon y = 2x^{3} + 3x^{2} - 12x - 12\).
Vyberte pravdivé tvrzení:}
{V bodě \(x = -2\) 
má funkce \(f\) 
lokální maximum.}{V bodě \(x = -2\) 
má funkce \(f\) 
lokální minimum.}{Daná funkce \(f\) má na
množině \(\mathbb{R}\) globální
maximum v bodě \(x = -2\).}{Daná funkce \(f\) má na
množině \(\mathbb{R}\) globální
minimum v bodě \(x = -2\).}{}{}}
\LANGsk{
\Question{
Je daná funkcia \(f\colon y = 2x^{3} + 3x^{2} - 12x - 12\).
Vyberte pravdivé tvrdenie:}
{V bode \(x = -2\) 
má funkcia \(f\) 
lokálne maximum.}{V bode \(x = -2\) 
má funkcia \(f\) 
lokálne minimum.}{Daná funkcia \(f\) má na
množine \(\mathbb{R}\) globálne
maximum v bode \(x = -2\).}{Daná funkcia \(f\) má na
množine \(\mathbb{R}\) globálne
minimum v bode \(x = -2\).}{}{}}
\LANGes{
\Question{
Identifica la proposición lógica sobre la función:
\(f(x) = 2x^{3} + 3x^{2} - 12x - 12\).}
{La función \(f\) tiene un máximo local en el punto \(x = -2\).}{La función \(f\) tiene un mínimo local en el punto \(x = -2\)..}{El máximo global de \(f\) 
en \(\mathbb{R}\) está en \(x = -2\).}{El mínimo global de \(f\) 
en \(\mathbb{R}\) está en \(x = -2\).}{}{}}
\End

\Data\ID{9000145402}
\Updated{2021-08-15 09:08:43}
\Level{B}
\SubArea{Applications of derivatives}
\LANGen{
\Question{
Identify a true statement about the function
\(f(x) = 2x^{2} -\frac{x^{4}} 
 {4} \).}
{The global maximum of \(f\) 
on \(\mathbb{R}\) is
at \(x = 2\) 
and \(x = -2\).}{The global minimum of \(f\) 
on \(\mathbb{R}\) is
at \(x = 2\) 
and \(x = -2\).}{The function \(f\) has a local
minimum at the point \(x = 2\).}{The function \(f\) has a local
minimum at the point \(x = -2\).}{}{}}
\LANGpl{
\Question{
Dana jest funkcja
\(f\colon y = 2x^{2} -\frac{x^{4}} 
 {4} \), wskaż zdanie prawdziwe.}
{Maksimum globalne funkcji \(f\) 
na zbiorze  \(\mathbb{R}\) jest w punkcie
 \(x = 2\) 
i \(x = -2\).}{Minimum globalne funkcji \(f\) 
na zbiorze  \(\mathbb{R}\) jest w punkcie
 \(x = 2\) 
i \(x = -2\).}{Funkcja \(f\) ma lokalne 
minimum w punkcie \(x = 2\).}{Funkcja  \(f\) ma lokalne minimum 
w punkcie \(x = -2\).}{}{}}
\LANGcs{
\Question{
Je dána funkce \(f\colon y = 2x^{2} -\frac{x^{4}} 
 {4} \).
Vyberte pravdivé tvrzení:}
{Daná funkce \(f\) má na
množině \(\mathbb{R}\) globální
maximum v bodech \(x = 2\) 
a \(x = -2\).}{Daná funkce \(f\) má na
množině \(\mathbb{R}\) globální
minimum v bodech \(x = 2\) 
a \(x = -2\).}{V bodě \(x = 2\) 
má funkce \(f\) 
lokální minimum.}{V bodě \(x = -2\) 
má funkce \(f\) 
lokální minimum.}{}{}}
\LANGsk{
\Question{
Je daná funkcia \(f\colon y = 2x^{2} -\frac{x^{4}} 
 {4} \).
Vyberte pravdivé tvrdenie:}
{Daná funkcia \(f\) má na
množine \(\mathbb{R}\) globálne
maximum v bodoch \(x = 2\) 
a \(x = -2\).}{Daná funkcia \(f\) má na
množine \(\mathbb{R}\) globálne
minimum v bodoch \(x = 2\) 
a \(x = -2\).}{V bode \(x = 2\) 
má funkcia \(f\) 
lokálne minimum.}{V bode \(x = -2\) 
má funkcia \(f\) 
lokálne minimum.}{}{}}
\LANGes{
\Question{
Identifica la proposición lógica sobre la función:
\(f(x) = 2x^{2} -\frac{x^{4}} 
 {4} \).}
{El máximo global de \(f\) 
en \(\mathbb{R}\) está en \(x = 2\) 
y \(x = -2\).}{El mínimo global de \(f\) 
en \(\mathbb{R}\) está en \(x = 2\) 
y \(x = -2\).}{La función \(f\) tiene un mínimo local en el punto \(x = 2\).}{La función \(f\) tiene un mínimo local en el punto \(x = -2\).}{}{}}
\End

\Data\ID{9000145403}
\Updated{2020-07-15 03:07:37}
\Level{B}
\SubArea{Applications of derivatives}
\LANGen{
\Question{
Identify a true statement on the function
\(f(x)= \frac{4-3x} 
{x\left (1-x\right )}\).}
{The function \(f\) has a local
minimum at the point \(x = \frac{2} 
{3}\).}{The function \(f\) has a local
maximum at the point \(x = \frac{2} 
{3}\).}{The global maximum of \(f\) 
on \(\mathbb{R}\setminus \{0.1\}\) is
at \(x = \frac{2} 
{3}\).}{The global minimum of \(f\) 
on \(\mathbb{R}\setminus \{0.1\}\) is
at \(x = \frac{2} 
{3}\).}{}{}}
\LANGpl{
\Question{
Dana jest funkcja
\(f\colon y = \frac{4-3x} 
{x\left (1-x\right )}\), wskaż zdanie prawdziwe.}
{Funkcja  \(f\) ma lokalne 
minimum w punkcie \(x = \frac{2} 
{3}\).}{Funkcja  \(f\) ma lokalne maksimum w punkcie  \(x = \frac{2} 
{3}\).}{Maksimum globalne funkcji  \(f\) 
na zbiorze  \(\mathbb{R}\setminus \{0{,}1\}\) jest w punkcie
 \(x = \frac{2} 
{3}\).}{Minimum globalne funkcji  \(f\) 
na zbiorze \(\mathbb{R}\setminus \{0{,}1\}\) jest w punkcie
 \(x = \frac{2} 
{3}\).}{}{}}
\LANGcs{
\Question{
Je dána funkce \(f\colon y = \frac{4-3x} 
{x\left (1-x\right )}\).
Vyberte pravdivé tvrzení:}
{V bodě \(x = \frac{2} 
{3}\) 
má funkce \(f\) 
lokální minimum.}{V bodě \(x = \frac{2} 
{3}\) 
má funkce \(f\) 
lokální maximum.}{Daná funkce \(f\) má na
množině \(\mathbb{R}\setminus \{0{,}1\}\) globální
maximum v bodě \(x = \frac{2} 
{3}\).}{Daná funkce \(f\) má na
množině \(\mathbb{R}\setminus \{0{,}1\}\) globální
minimum v bodě \(x = \frac{2} 
{3}\).}{}{}}
\LANGsk{
\Question{
Je daná funkcia \(f\colon y = \frac{4-3x} 
{x\left (1-x\right )}\).
Vyberte pravdivé tvrdenie:}
{V bode \(x = \frac{2} 
{3}\) 
má funkcia \(f\) 
lokálne minimum.}{V bode \(x = \frac{2} 
{3}\) 
má funkcia \(f\) 
lokálne maximum.}{Daná funkcia \(f\) má na
množine \(\mathbb{R}\setminus \{0{,}1\}\) globálne
maximum v bode \(x = \frac{2} 
{3}\).}{Daná funkcia \(f\) má na
množine \(\mathbb{R}\setminus \{0{,}1\}\) globálne
minimum v bode \(x = \frac{2} 
{3}\).}{}{}}
\LANGes{
\Question{
Identifica la proposición lógica sobre la función:
\(f(x)= \frac{4-3x} 
{x\left (1-x\right )}\).}
{La función \(f\) tiene un mínimo local en el punto \(x = \frac{2} 
{3}\).}{La función \(f\) tiene un máximo local en el punto \(x = \frac{2} 
{3}\).}{El máximo global de \(f\) 
en \(\mathbb{R}\setminus \{0.1\}\) está en \(x = \frac{2} 
{3}\).}{El mínimo global de \(f\) 
en \(\mathbb{R}\setminus \{0.1\}\) está en \(x = \frac{2} 
{3}\).}{}{}}
\End

\Data\ID{9000145404}
\Updated{2022-04-23 05:04:54}
\Level{B}
\SubArea{Applications of derivatives}
\LANGen{
\Question{
Identify a true statement about the function
\(f(x) = x^{3} - 3x^{2} + 3x + 2\).}
{There is neither local minimum nor maximum of
\(f\).}{The function \(f\) has a local
maximum at the point \(x = 1\).}{The function \(f\) has a local
minimum at the point \(x = 1\).}{The global minimum of \(f\) 
on \(\mathbb{R}\) is
at \(x = 1\).}{}{}}
\LANGpl{
\Question{
Dana jest funkcja
\(f\colon y = x^{3} - 3x^{2} + 3x + 2\), wskaż zdanie prawdziwe.}
{Brak lokalnego minimum i maksimum funkcji
\(f\).}{Funkcja  \(f\) ma lokalne maksimum w punkcie \(x = 1\).}{Funkcja \(f\) ma lokalne minimum w punkcie \(x = 1\).}{Minimum globalne funkcji \(f\) 
na zbiorze  \(\mathbb{R}\) jest w punkcie 
\(x = 1\).}{}{}}
\LANGcs{
\Question{
Je dána funkce \(f\colon y = x^{3} - 3x^{2} + 3x + 2\).
Vyberte pravdivé tvrzení:}
{Daná funkce \(f\) 
nemá žádný lokální extrém.}{V bodě \(x = 1\) 
má funkce \(f\) 
lokální maximum.}{V bodě \(x = 1\) 
má funkce \(f\) 
lokální minimum.}{Daná funkce \(f\) má na
množině \(\mathbb{R}\) globální
minimum v bodě \(x = 1\).}{}{}}
\LANGsk{
\Question{
Je daná funkcia \(f\colon y = x^{3} - 3x^{2} + 3x + 2\).
Vyberte pravdivé tvrdenie:}
{Daná funkcia \(f\) 
nemá žiadny lokálny extrém.}{V bode \(x = 1\) 
má funkcia \(f\) 
lokálne maximum.}{V bode \(x = 1\) 
má funkcia \(f\) 
lokálne minimum.}{Daná funkcia \(f\) má na
množine \(\mathbb{R}\) globálne
minimum v bode \(x = 1\).}{}{}}
\LANGes{
\Question{
Identifica la proposición lógica sobre la función: 
\(f(x) = x^{3} - 3x^{2} + 3x + 2\).}
{No hay ningún mínimo ni máximo local
\(f\).}{La función \(f\) tiene un máximo local en el punto \(x = 1\).}{La función \(f\) tiene un mínimo local en el punto \(x = 1\).}{El mínimo global de \(f\) 
en \(\mathbb{R}\) está en \(x = 1\).}{}{}}
\End

\Data\ID{9000145405}
\Updated{2020-07-15 03:07:37}
\Level{B}
\SubArea{Applications of derivatives}
\LANGen{
\Question{
Identify a true statement on the function
\(f(x) = \frac{1} 
{4}x^{4} -\frac{2} 
{3}x^{3} -\frac{3} 
{2}x^{2} + 2\text{ on }\left (-2;4\right )\).}
{The function \(f\) has a local
maximum at the point \(x = 0\).}{The function \(f\) has a local
minimum at the point \(x = 0\).}{The global maximum of \(f\) 
on this interval is at \(x = 0\).}{The global minimum of \(f\) 
on this interval is at \(x = 0\).}{}{}}
\LANGpl{
\Question{
Dana jest funkcja 
\(f\colon y = \frac{1} 
{4}x^{4} -\frac{2}
{3}x^{3} -\frac{3} 
{2}x^{2} + 2\text{ na przedziale }\left (-2;4\right )\), wskaż zdanie prawdziwe.}
{Funkcja \(f\) ma lokalne maksimum w punkcie  \(x = 0\).}{Funkcja \(f\) ma lokalne minimum  w punkcie \(x = 0\).}{Maksimum globalne funkcji  \(f\) 
na tym przedziale jest w punkcie \(x = 0\).}{Minimum globalne funkcji \(f\) 
na tym przedziale jest w punkcie \(x = 0\).}{}{}}
\LANGcs{
\Question{
Je dána funkce \(f\colon y = \frac{1} 
{4}x^{4} -\frac{2} 
{3}x^{3} -\frac{3} 
{2}x^{2} + 2\text{ na intervalu }\left (-2;4\right )\).
Vyberte pravdivé tvrzení:}
{V bodě \(x = 0\) 
má funkce \(f\) 
lokální maximum.}{V bodě \(x = 0\) 
má funkce \(f\) 
lokální minimum.}{Funkce \(f\) 
má na daném intervalu globální maximum v bodě
\(x = 0\).}{Funkce \(f\) 
má na daném intervalu globální minimum v bodě
\(x = 0\).}{}{}}
\LANGsk{
\Question{
Je daná funkcia \(f\colon y = \frac{1} 
{4}x^{4} -\frac{2} 
{3}x^{3} -\frac{3} 
{2}x^{2} + 2\text{ na intervale }\left (-2;4\right )\).
Vyberte pravdivé tvrdenie:}
{V bode \(x = 0\) 
má funkcia \(f\) 
lokálne maximum.}{V bode \(x = 0\) 
má funkcia \(f\) 
lokálne minimum.}{Funkcia \(f\) 
má na danom intervale globálne maximum v bode
\(x = 0\).}{Funkcia \(f\) 
má na danom intervale globálne minimum v bode
\(x = 0\).}{}{}}
\LANGes{
\Question{
Identifica la proposición lógica sobre la función:
\(f(x) = \frac{1} 
{4}x^{4} -\frac{2} 
{3}x^{3} -\frac{3} 
{2}x^{2} + 2\text{ en }\left (-2;4\right )\).}
{La función \(f\) tiene un máximo local en el punto \(x = 0\).}{La función \(f\) tiene un mínimo local en el punto \(x = 0\).}{El máximo global de \(f\) 
en este intervalo está en \(x = 0\).}{El mínimo global de \(f\) 
en este intervalo está en \(x = 0\).}{}{}}
\End

\Data\ID{9000145406}
\Updated{2020-07-15 03:07:37}
\Level{B}
\SubArea{Applications of derivatives}
\LANGen{
\Question{
Identify a true statement on the function
\(f(x) = x^{3} - 12x + 20\text{ on }\left (-3;4\right )\).}
{The global minimum of \(f\) 
on this interval is at \(x = 2\).}{The global maximum of \(f\) 
on this interval is at \(x = 2\).}{The function \(f\) has a local
minimum at the point \(x = -2\).}{The global minimum of \(f\) on
this interval is at the point \(x = -2\).}{}{}}
\LANGpl{
\Question{
Dana jest funkcja
\(f\colon y = x^{3} - 12x + 20\text{ na przedziale }\left (-3;4\right )\), wskaż zdanie prawdziwe.}
{Minimum globalne funkcji \(f\) na tym przedziale jest w punkcie
 \(x = 2\).}{Maksimum globalne funkcji  \(f\) 
na tym przedziale jest w punkcie \(x = 2\).}{Funkcja \(f\) ma lokalne 
minimum w punkcie \(x = -2\).}{Minimum globalne funkcji \(f\) na tym przedziale jest w punkcie
 \(x = -2\).}{}{}}
\LANGcs{
\Question{
Je dána funkce \(f\colon y = x^{3} - 12x + 20\text{ na intervalu }\left (-3;4\right )\).
Vyberte pravdivé tvrzení:}
{Funkce \(f\) 
má na daném intervalu globální minimum v bodě
\(x = 2\).}{Funkce \(f\) 
má na daném intervalu globální maximum v bodě
\(x = 2\).}{V bodě \(x = -2\) 
má funkce \(f\) 
lokální minimum.}{Funkce \(f\) 
má na daném intervalu globální minimum v bodě
\(x = -2\).}{}{}}
\LANGsk{
\Question{
Je daná funkcia \(f\colon y = x^{3} - 12x + 20\text{ na intervale }\left (-3;4\right )\).
Vyberte pravdivé tvrdenie:}
{Funkcia \(f\) 
má na danom intervale globálne minimum v bode
\(x = 2\).}{Funkcia \(f\) 
má na danom intervale globálne maximum v bode
\(x = 2\).}{V bode \(x = -2\) 
má funkcia \(f\) 
lokálne minimum.}{Funkcia \(f\) 
má na danom intervalu globálne minimum v bode
\(x = -2\).}{}{}}
\LANGes{
\Question{
Identifica la proposición lógica sobre la función:
\(f(x) = x^{3} - 12x + 20\text{ en }\left (-3;4\right )\).}
{El mínimo global de\(f\) 
en este intervalo está en \(x = 2\).}{El máximo global de \(f\) 
en este intervalo está en \(x = 2\).}{La función \(f\) tiene un mínimo local en el punto \(x = -2\).}{El mínimo global de \(f\) en este intervalo está en \(x = -2\).}{}{}}
\End

\Data\ID{9000145407}
\Updated{2020-07-15 03:07:37}
\Level{B}
\SubArea{Applications of derivatives}
\LANGen{
\Question{
Identify a true statement on the function
\(f(x) = x^{4} - 8x^{3} + 22x^{2} - 24x + 12\).}
{The global minimum of \(f\) 
on \(\mathbb{R}\) is
at \(x = 1\) 
and \(x = 3\).}{The global maximum of \(f\) 
on \(\mathbb{R}\) is
at \(x = 2\).}{The local minima of \(f\) 
are at \(x = 1\) 
and \(x = 2\).}{The local maximum of \(f\) 
is at \(x = 3\).}{}{}}
\LANGpl{
\Question{
Dana jest funkcja 
\(f\colon y = x^{4} - 8x^{3} + 22x^{2} - 24x + 12\), wskaż zdanie prawdziwe.}
{Minimum globalne funkcji  \(f\) 
na zbiorze  \(\mathbb{R}\) jest w punkcie 
\(x = 1\) 
i \(x = 3\).}{Maksimum globalne funkcji \(f\) 
na zbiorze  \(\mathbb{R}\)  jest w punkcie 
\(x = 2\).}{Lokalne minima funkcji \(f\) 
są w punkcie  \(x = 1\) 
i  \(x = 2\).}{Lokalne maksimum funkcji \(f\) 
 jest w punkcie  \(x = 3\).}{}{}}
\LANGcs{
\Question{
Je dána funkce \(f\colon y = x^{4} - 8x^{3} + 22x^{2} - 24x + 12\).
Vyberte pravdivé tvrzení:}
{Daná funkce \(f\) má na
množině \(\mathbb{R}\) globální
minimum v bodech \(x = 1\) 
a \(x = 3\).}{Daná funkce \(f\) má na
množině \(\mathbb{R}\) globální
maximum v bodě \(x = 2\).}{Daná funkce \(f\) má
lokální minima v bodech \(x = 1\) 
a \(x = 2\).}{Daná funkce \(f\) má
lokální maximum v bodě \(x = 3\).}{}{}}
\LANGsk{
\Question{
Je daná funkcia \(f\colon y = x^{4} - 8x^{3} + 22x^{2} - 24x + 12\).
Vyberte pravdivé tvrdenie:}
{Daná funkcia \(f\) má na
množine \(\mathbb{R}\) globálne
minimum v bodoch \(x = 1\) 
a \(x = 3\).}{Daná funkcia \(f\) má na
množine \(\mathbb{R}\) globálne
maximum v bode \(x = 2\).}{Daná funkcia \(f\) má
lokálne minima v bodoch \(x = 1\) 
a \(x = 2\).}{Daná funkcia \(f\) má
lokálne maximum v bode \(x = 3\).}{}{}}
\LANGes{
\Question{
Identifica la proposición lógica sobre la función:
\(f(x) = x^{4} - 8x^{3} + 22x^{2} - 24x + 12\).}
{El mínimo global de \(f\) 
en \(\mathbb{R}\) está en \(x = 1\) 
y \(x = 3\).}{El máximo global de \(f\) 
en \(\mathbb{R}\) está en \(x = 2\).}{Los mínimos locales de \(f\) 
están en \(x = 1\) 
y \(x = 2\).}{El máximo local de \(f\) 
está en \(x = 3\).}{}{}}
\End

\Data\ID{9000145408}
\Updated{2020-07-15 03:07:37}
\Level{B}
\SubArea{Applications of derivatives}
\LANGen{
\Question{
Identify a true statement on the function
\(f(x) = \left (x - 1\right )^{3}\left (x + 1\right )^{2}\).}
{The function \(f\) has neither local
minimum nor maximum at \(x = 1\).}{The global maximum of \(f\) 
on \(\mathbb{R}\) is
at \(x = -1\).}{The function \(f\) has a
local maximum at \(x = -\frac{1}
{5}\).}{The function \(f\) 
has three local extrema. These extrema are at
\(x = 1\),
\(x = -1\) and
\(x = -\frac{1}
{5}\).}{}{}}
\LANGpl{
\Question{
Dana jest funkcja 
\(f\colon y = \left (x - 1\right )^{3}\left (x + 1\right )^{2}\), wskaż zdanie prawdziwe.}
{Funkcja  \(f\) nie ma lokalnego minimum i maksimum w punkcie \(x = 1\).}{Maksimum globalne funkcji \(f\) 
na zbiorze  \(\mathbb{R}\)  jest w punkcie 
 \(x = -1\).}{Funkcja  \(f\) ma lokalne maksimum   w punkcie \(x = -\frac{1}
{5}\).}{Funkcja  \(f\) 
ma trzy lokalne ekstrema;
\(x = 1\),
\(x = -1\) i
\(x = -\frac{1}
{5}\).}{}{}}
\LANGcs{
\Question{
Je dána funkce \(f\colon y = \left (x - 1\right )^{3}\left (x + 1\right )^{2}\).
Vyberte pravdivé tvrzení:}
{Daná funkce \(f\) 
nemá v bodě \(x = 1\) 
lokální extrém.}{Daná funkce \(f\) má na
množině \(\mathbb{R}\) globální
maximum v bodě \(x = -1\).}{Daná funkce \(f\) má
lokální maximum v bodě \(x = -\frac{1}
{5}\).}{Daná funkce \(f\) 
má právě tři lokální extrémy v bodech
\(x = 1\),
\(x = -1\) a
\(x = -\frac{1}
{5}\).}{}{}}
\LANGsk{
\Question{
Je daná funkcia \(f\colon y = \left (x - 1\right )^{3}\left (x + 1\right )^{2}\).
Vyberte pravdivé tvrdenie:}
{Daná funkcia \(f\) 
nemá v bode \(x = 1\) 
lokálny extrém.}{Daná funkcia \(f\) má na
množine \(\mathbb{R}\) globálne
maximum v bode \(x = -1\).}{Daná funkcia \(f\) má
lokálne maximum v bode \(x = -\frac{1}
{5}\).}{Daná funkcia \(f\) 
má práve tri lokálne extrémy v bodoch
\(x = 1\),
\(x = -1\) a
\(x = -\frac{1}
{5}\).}{}{}}
\LANGes{
\Question{
Identifica la proposición lógica sobre la función:
\(f(x) = \left (x - 1\right )^{3}\left (x + 1\right )^{2}\).}
{La función \(f\) no tiene ningún mínimo local ni máximo local en \(x = 1\).}{El máximo global de \(f\) 
en \(\mathbb{R}\) está en \(x = -1\).}{La función \(f\) tiene un máximo local en \(x = -\frac{1}
{5}\).}{La función \(f\) 
tiene tres extremos locales. Están en
\(x = 1\),
\(x = -1\) y
\(x = -\frac{1}
{5}\).}{}{}}
\End

\Data\ID{9000145409}
\Updated{2020-07-15 03:07:37}
\Level{B}
\SubArea{Applications of derivatives}
\LANGen{
\Question{
Identify a true statement on the function
\(f(x) = 1 + 2x^{2} -\frac{1} 
{4}x^{4}\).}
{The global maximum of \(f\) 
on \(\mathbb{R}\) is
at \(x = 2\) a
\(x = -2\).}{The function \(f\) has
a global minimum on \(\mathbb{R}\).}{The function \(f\) has a
local maximum at \(x = 0\).}{The function \(f\) 
has neither local minimum nor maximum.}{}{}}
\LANGpl{
\Question{
Dana jest funkcja
\(f\colon y = 1 + 2x^{2} -\frac{1} 
{4}x^{4}\).}
{Maksimum globalne funkcji \(f\) 
na zbiorze  \(\mathbb{R}\)  jest w punkcie 
\(x = 2\) a
\(x = -2\).}{Funkcja \(f\) ma minimum globalne na zbiorze  \(\mathbb{R}\).}{Funkcja \(f\) ma lokalne maksimum w punkcie \(x = 0\).}{Funkcja  \(f\) 
nie ma lokalnego minimum i maksimum.}{}{}}
\LANGcs{
\Question{
Je dána funkce \(f\colon y = 1 + 2x^{2} -\frac{1} 
{4}x^{4}\).
Vyberte pravdivé tvrzení:}
{Daná funkce \(f\) má na
množině \(\mathbb{R}\) globální
maximum v bodech \(x = 2\) 
a \(x = -2\).}{Daná funkce \(f\) 
má na množině \(\mathbb{R}\) 
globální minimum.}{Daná funkce \(f\) má
lokální maximum v bodě \(x = 0\).}{Daná funkce \(f\) 
nemá lokální extrém v žádném bodě.}{}{}}
\LANGsk{
\Question{
Je daná funkcia \(f\colon y = 1 + 2x^{2} -\frac{1} 
{4}x^{4}\).
Vyberte pravdivé tvrdenie:}
{Daná funkcia \(f\) má na
množine \(\mathbb{R}\) globálne
maximum v bodoch \(x = 2\) 
a \(x = -2\).}{Daná funkcia \(f\) 
má na množine \(\mathbb{R}\) 
globálne minimum.}{Daná funkcia \(f\) má
lokálne maximum v bode \(x = 0\).}{Daná funkcia \(f\) 
nemá lokálny extrém v žiadnom bode.}{}{}}
\LANGes{
\Question{
Identifica la proposición lógica sobre la función:
\(f(x) = 1 + 2x^{2} -\frac{1} 
{4}x^{4}\).}
{El máximo global de \(f\) 
en \(\mathbb{R}\) está en \(x = 2\) y
\(x = -2\).}{La función \(f\) tiene un mínimo global en \(\mathbb{R}\).}{La función \(f\) tiene un máximo local en \(x = 0\).}{La función \(f\) 
no tiene ningún mínimo ni máximo local.}{}{}}
\End

\Data\ID{9000145410}
\Updated{2022-04-23 05:04:54}
\Level{B}
\SubArea{Applications of derivatives}
\LANGen{
\Question{
Identify a true statement about the function
\(f(x) = \frac{1} 
{4}x^{4} - x^{3}\).}
{The local minimum of \(f\) 
is
at \(x = 3\).}{The function \(f\) 
has neither local minimum nor local maximum.}{The function \(f\) has
a local minimum at \(x = 0\).}{The function \(f\) 
has two local extrema. These extrema are at
\(x = 3\) and
\(x = 0\).}{}{}}
\LANGpl{
\Question{
Dana jest funkcja 
\(f\colon y = \frac{1} 
{4}x^{4} - x^{3}\).}
{Minimum lokalne funkcji \(f\) 
jest w punkcie 
 \(x = 3\).}{Funkcja  \(f\) 
nie ma lokalnego minimum i maksimum.}{Funkcja  \(f\) ma lokalne minimum w punkcie \(x = 0\).}{Funkcja  \(f\) 
ma dwa lokalne ekstrema w punkcie 
\(x = 3\) i
\(x = 0\).}{}{}}
\LANGcs{
\Question{
Je dána funkce \(f\colon y = \frac{1} 
{4}x^{4} - x^{3}\).
Vyberte pravdivé tvrzení:}
{Daná funkce \(f\) má lokální
minimum v bodě \(x = 3\).}{Daná funkce \(f\) 
nemá lokální extrém v žádném bodě.}{Daná funkce \(f\) má
lokální minimum v bodě \(x = 0\).}{Daná funkce \(f\) má dva
lokální extrémy v bodech \(x = 3\) 
a \(x = 0\).}{}{}}
\LANGsk{
\Question{
Je daná funkcia \(f\colon y = \frac{1} 
{4}x^{4} - x^{3}\).
Vyberte pravdivé tvrdenie:}
{Daná funkcia \(f\) má lokálne
minimum v bode \(x = 3\).}{Daná funkcia \(f\) 
nemá lokálny extrém v žiadnom bode.}{Daná funkcia \(f\) má
lokálne minimum v bode \(x = 0\).}{Daná funkcia \(f\) má dva
lokálne extrémy v bodoch \(x = 3\) 
a \(x = 0\).}{}{}}
\LANGes{
\Question{
Identifica la proposición lógica sobre la función:
\(f(x) = \frac{1} 
{4}x^{4} - x^{3}\).}
{El mínimo local de \(f\) 
en \(\mathbb{R}\) está en \(x = 3\).}{La función \(f\) 
no tiene ningún mínimo ni máximo local.}{La función \(f\) tiene un mínimo local en \(x = 0\).}{La función \(f\) 
tiene dos extremos locales. Están en
\(x = 3\) y
\(x = 0\).}{}{}}
\End

\Data\ID{1003163901}
\Updated{2022-04-23 05:04:48}
\Level{C}
\SubArea{Applications of derivatives}
\LANGen{
\Question{
Use L'Hospital's rule to find the following limit.
\[ \lim_{x\to2}\frac{2x^3-3x^2-4}{x^2+x-6} \]}
{\( \frac{12}5 \)}{\( \frac{18}5 \)}{\( \frac{12}3 \)}{\( 0 \)}{\( \infty \)}{}}
\LANGpl{
\Question{
Korzystając z reguły de I'Hospitala oblicz granicę.
\[ \lim_{x\to2}\frac{2x^3-3x^2-4}{x^2+x-6} \]}
{\( \frac{12}5 \)}{\( \frac{18}5 \)}{\( \frac{12}3 \)}{\( 0 \)}{\( \infty \)}{}}
\LANGcs{
\Question{
Užitím L'Hospitalova pravidla vypočítejte následující limitu.
\[ \lim_{x\to2}\frac{2x^3-3x^2-4}{x^2+x-6} \]}
{\( \frac{12}5 \)}{\( \frac{18}5 \)}{\( \frac{12}3 \)}{\( 0 \)}{\( \infty \)}{}}
\LANGsk{
\Question{
Použitím L'Hospitalova pravidla vypočítajte následujúcu limitu.
\[ \lim_{x\to2}\frac{2x^3-3x^2-4}{x^2+x-6} \]}
{\( \frac{12}5 \)}{\( \frac{18}5 \)}{\( \frac{12}3 \)}{\( 0 \)}{\( \infty \)}{}}
\LANGes{
\Question{
Resuelve el siguiente límite utilizando la Regla de L'Hôpital:
\[ \lim_{x\to2}\frac{2x^3-3x^2-4}{x^2+x-6} \]}
{\( \frac{12}5 \)}{\( \frac{18}5 \)}{\( \frac{12}3 \)}{\( 0 \)}{\( \infty \)}{}}
\End

\Data\ID{1003163902}
\Updated{2020-07-15 03:07:21}
\Level{C}
\SubArea{Applications of derivatives}
\LANGen{
\Question{
Use L'Hospital's rule to find the following limit.
\[ \lim_{x\to0}\frac{\mathrm{e}^x-1}{\sin2x} \]}
{\( \frac12 \)}{\( -\frac12 \)}{\( 1 \)}{\( 0 \)}{\( -1 \)}{}}
\LANGpl{
\Question{
Korzystając z reguły de I'Hospitala oblicz granicę.
\[ \lim_{x\to0}\frac{\mathrm{e}^x-1}{\sin2x} \]}
{\( \frac12 \)}{\( -\frac12 \)}{\( 1 \)}{\( 0 \)}{\( -1 \)}{}}
\LANGcs{
\Question{
Užitím L'Hospitalova pravidla vypočítejte následující limitu.
\[ \lim_{x\to0}\frac{\mathrm{e}^x-1}{\sin2x} \]}
{\( \frac12 \)}{\( -\frac12 \)}{\( 1 \)}{\( 0 \)}{\( -1 \)}{}}
\LANGsk{
\Question{
Použitím L'Hospitalova pravidla vypočítajte následujúcu limitu.
\[ \lim_{x\to0}\frac{\mathrm{e}^x-1}{\sin2x} \]}
{\( \frac12 \)}{\( -\frac12 \)}{\( 1 \)}{\( 0 \)}{\( -1 \)}{}}
\LANGes{
\Question{
Resuelve el siguiente límite utilizando la Regla de L'Hôpital:
\[ \lim_{x\to0}\frac{\mathrm{e}^x-1}{\sin2x} \]}
{\( \frac12 \)}{\( -\frac12 \)}{\( 1 \)}{\( 0 \)}{\( -1 \)}{}}
\End

\Data\ID{1003163903}
\Updated{2022-04-23 05:04:48}
\Level{C}
\SubArea{Applications of derivatives}
\LANGen{
\Question{
Use L'Hospital's rule to find the following limit.
\[ \lim_{x\to0}\frac{\sin2x}{\mathrm{tg}\,3x} \]}
{\( \frac23 \)}{\( 1 \)}{\( 2 \)}{\( 0 \)}{\( 6 \)}{}}
\LANGpl{
\Question{
Korzystając z reguły de I'Hospitala oblicz granicę.
\[ \lim_{x\to0}\frac{\sin2x}{\mathrm{tg}\,3x} \]}
{\( \frac23 \)}{\( 1 \)}{\( 2 \)}{\( 0 \)}{\( 6 \)}{}}
\LANGcs{
\Question{
Užitím L'Hospitalova pravidla vypočítejte následující limitu.
\[ \lim_{x\to0}\frac{\sin2x}{\mathrm{tg}\,3x} \]}
{\( \frac23 \)}{\( 1 \)}{\( 2 \)}{\( 0 \)}{\( 6 \)}{}}
\LANGsk{
\Question{
Použitím L'Hospitalova pravidla vypočítajte následujúcu limitu.
\[ \lim_{x\to0}\frac{\sin2x}{\mathrm{tg}\,3x} \]}
{\( \frac23 \)}{\( 1 \)}{\( 2 \)}{\( 0 \)}{\( 6 \)}{}}
\LANGes{
\Question{
Resuelve el siguiente límite utilizando la Regla de L'Hôpital:
\[ \lim_{x\to0}\frac{\sin2x}{\mathrm{tg}\,3x} \]}
{\( \frac23 \)}{\( 1 \)}{\( 2 \)}{\( 0 \)}{\( 6 \)}{}}
\End

\Data\ID{1003163904}
\Updated{2020-07-15 03:07:21}
\Level{C}
\SubArea{Applications of derivatives}
\LANGen{
\Question{
Use L'Hospital's rule to find the following limit.
\[ \lim_{x\to\frac{\pi}3}\frac{1-2\cos x}{\pi-3x} \]}
{\( -\frac{\sqrt3}3 \)}{\( -\frac{\sqrt3}6 \)}{\( \frac{\sqrt3}3 \)}{\( \frac{\sqrt3}6 \)}{\( -\frac13 \)}{}}
\LANGpl{
\Question{
Korzystając z reguły de I'Hospitala oblicz granicę.
\[ \lim_{x\to\frac{\pi}3}\frac{1-2\cos x}{\pi-3x} \]}
{\( -\frac{\sqrt3}3 \)}{\( -\frac{\sqrt3}6 \)}{\( \frac{\sqrt3}3 \)}{\( \frac{\sqrt3}6 \)}{\( -\frac13 \)}{}}
\LANGcs{
\Question{
Užitím L'Hospitalova pravidla vypočítejte následující limitu.
\[ \lim_{x\to\frac{\pi}3}\frac{1-2\cos x}{\pi-3x} \]}
{\( -\frac{\sqrt3}3 \)}{\( -\frac{\sqrt3}6 \)}{\( \frac{\sqrt3}3 \)}{\( \frac{\sqrt3}6 \)}{\( -\frac13 \)}{}}
\LANGsk{
\Question{
Použitím L'Hospitalova pravidla vypočítajte následujúcu limitu.
\[ \lim_{x\to\frac{\pi}3}\frac{1-2\cos x}{\pi-3x} \]}
{\( -\frac{\sqrt3}3 \)}{\( -\frac{\sqrt3}6 \)}{\( \frac{\sqrt3}3 \)}{\( \frac{\sqrt3}6 \)}{\( -\frac13 \)}{}}
\LANGes{
\Question{
Resuelve el siguiente límite utilizando la Regla de L'Hôpital:
\[ \lim_{x\to\frac{\pi}3}\frac{1-2\cos x}{\pi-3x} \]}
{\( -\frac{\sqrt3}3 \)}{\( -\frac{\sqrt3}6 \)}{\( \frac{\sqrt3}3 \)}{\( \frac{\sqrt3}6 \)}{\( -\frac13 \)}{}}
\End

\Data\ID{1003163905}
\Updated{2020-07-15 03:07:21}
\Level{C}
\SubArea{Applications of derivatives}
\LANGen{
\Question{
Use L'Hospital's rule to find the following limit.
\[ \lim_{x\to1}\frac{\sqrt{x+3}-2}{\ln x} \]}
{\( \frac14 \)}{\( \frac12 \)}{\( 0 \)}{\( 1 \)}{\( \frac{\sqrt2}2 \)}{}}
\LANGpl{
\Question{
Korzystając z reguły de I'Hospitala oblicz granicę.
\[ \lim_{x\to1}\frac{\sqrt{x+3}-2}{\ln x} \]}
{\( \frac14 \)}{\( \frac12 \)}{\( 0 \)}{\( 1 \)}{\( \frac{\sqrt2}2 \)}{}}
\LANGcs{
\Question{
Užitím L'Hospitalova pravidla vypočítejte následující limitu.
\[ \lim_{x\to1}\frac{\sqrt{x+3}-2}{\ln x} \]}
{\( \frac14 \)}{\( \frac12 \)}{\( 0 \)}{\( 1 \)}{\( \frac{\sqrt2}2 \)}{}}
\LANGsk{
\Question{
Použitím L'Hospitalova pravidla vypočítajte následujúcu limitu.
\[ \lim_{x\to1}\frac{\sqrt{x+3}-2}{\ln x} \]}
{\( \frac14 \)}{\( \frac12 \)}{\( 0 \)}{\( 1 \)}{\( \frac{\sqrt2}2 \)}{}}
\LANGes{
\Question{
Resuelve el siguiente límite utilizando la Regla de L'Hôpital:
\[ \lim_{x\to1}\frac{\sqrt{x+3}-2}{\ln x} \]}
{\( \frac14 \)}{\( \frac12 \)}{\( 0 \)}{\( 1 \)}{\( \frac{\sqrt2}2 \)}{}}
\End

\Data\ID{1003163906}
\Updated{2020-07-15 03:07:21}
\Level{C}
\SubArea{Applications of derivatives}
\LANGen{
\Question{
Evaluate the following limit. (Apply L'Hospital's rule repeatedly.)
\[ \lim_{x\to\infty}\frac{x^2-1}{x^3-2x^2+x} \]}
{\( 0 \)}{\( \infty \)}{\( \frac23 \)}{\( 1 \)}{}{}}
\LANGpl{
\Question{
Wyznacz granicę. (Zastosuj regułę de I'Hospitala)
\[ \lim_{x\to\infty}\frac{x^2-1}{x^3-2x^2+x} \]}
{\( 0 \)}{\( \infty \)}{\( \frac23 \)}{\( 1 \)}{}{}}
\LANGcs{
\Question{
Vypočítejte následující limitu. (Opakovaně užijte L'Hospitalovo pravidlo.)
\[ \lim_{x\to\infty}\frac{x^2-1}{x^3-2x^2+x} \]}
{\( 0 \)}{\( \infty \)}{\( \frac23 \)}{\( 1 \)}{}{}}
\LANGsk{
\Question{
Vypočítajte následujúcu limitu. (Opakovane použite L'Hospitalovo pravidlo.)
\[ \lim_{x\to\infty}\frac{x^2-1}{x^3-2x^2+x} \]}
{\( 0 \)}{\( \infty \)}{\( \frac23 \)}{\( 1 \)}{}{}}
\LANGes{
\Question{
Resuelve el siguiente límite utilizando repetidamente la Regla de L'Hôpital:
\[ \lim_{x\to\infty}\frac{x^2-1}{x^3-2x^2+x} \]}
{\( 0 \)}{\( \infty \)}{\( \frac23 \)}{\( 1 \)}{}{}}
\End

\Data\ID{1003163907}
\Updated{2020-07-15 03:07:21}
\Level{C}
\SubArea{Applications of derivatives}
\LANGen{
\Question{
Evaluate the following limit. (Apply L'Hospital's rule repeatedly.)
\[ \lim_{x\to\infty}\frac{\mathrm{e}^x-2}{x^2} \]}
{\( \infty \)}{\( 0 \)}{\( 1 \)}{\( \frac12 \)}{}{}}
\LANGpl{
\Question{
Wyznacz granicę. (Zastosuj regułę de I'Hospitala)
\[ \lim_{x\to\infty}\frac{\mathrm{e}^x-2}{x^2} \]}
{\( \infty \)}{\( 0 \)}{\( 1 \)}{\( \frac12 \)}{}{}}
\LANGcs{
\Question{
Vypočítejte následující limitu. (Opakovaně užijte L'Hospitalovo pravidlo.)
\[ \lim_{x\to\infty}\frac{\mathrm{e}^x-2}{x^2} \]}
{\( \infty \)}{\( 0 \)}{\( 1 \)}{\( \frac12 \)}{}{}}
\LANGsk{
\Question{
Vypočítajte následujúcu limitu. (Opakovane použite L'Hospitalovo pravidlo.)
\[ \lim_{x\to\infty}\frac{\mathrm{e}^x-2}{x^2} \]}
{\( \infty \)}{\( 0 \)}{\( 1 \)}{\( \frac12 \)}{}{}}
\LANGes{
\Question{
Resuelve el siguiente límite utilizando repetidamente la Regla de L'Hôpital:
\[ \lim_{x\to\infty}\frac{\mathrm{e}^x-2}{x^2} \]}
{\( \infty \)}{\( 0 \)}{\( 1 \)}{\( \frac12 \)}{}{}}
\End

\Data\ID{1003163908}
\Updated{2020-07-15 03:07:21}
\Level{C}
\SubArea{Applications of derivatives}
\LANGen{
\Question{
Evaluate the following limit. (Apply  L'Hospital's rule repeatedly.)
\[ \lim_{x\to1}\frac{\cos(\pi x)+1}{(x-1)^2} \]}
{\( \frac{\pi^2}2 \)}{\( -\frac{\pi^2}2 \)}{\( \frac{\pi}2 \)}{\( -\frac{\pi}2 \)}{}{}}
\LANGpl{
\Question{
Wyznacz granicę. (Zastosuj regułę de I'Hospitala)
\[ \lim_{x\to1}\frac{\cos(\pi x)+1}{(x-1)^2} \]}
{\( \frac{\pi^2}2 \)}{\( -\frac{\pi^2}2 \)}{\( \frac{\pi}2 \)}{\( -\frac{\pi}2 \)}{}{}}
\LANGcs{
\Question{
Vypočítejte následující limitu. (Opakovaně užijte L'Hospitalovo pravidlo.)
\[ \lim_{x\to1}\frac{\cos(\pi x)+1}{(x-1)^2} \]}
{\( \frac{\pi^2}2 \)}{\( -\frac{\pi^2}2 \)}{\( \frac{\pi}2 \)}{\( -\frac{\pi}2 \)}{}{}}
\LANGsk{
\Question{
Vypočítajte následujúcu limitu. (Opakovane použite L'Hospitalovo pravidlo.)
\[ \lim_{x\to1}\frac{\cos(\pi x)+1}{(x-1)^2} \]}
{\( \frac{\pi^2}2 \)}{\( -\frac{\pi^2}2 \)}{\( \frac{\pi}2 \)}{\( -\frac{\pi}2 \)}{}{}}
\LANGes{
\Question{
Resuelve el siguiente límite utilizando repetidamente la Regla de L'Hôpital:
\[ \lim_{x\to1}\frac{\cos(\pi x)+1}{(x-1)^2} \]}
{\( \frac{\pi^2}2 \)}{\( -\frac{\pi^2}2 \)}{\( \frac{\pi}2 \)}{\( -\frac{\pi}2 \)}{}{}}
\End

\Data\ID{1003163909}
\Updated{2022-04-23 05:04:48}
\Level{C}
\SubArea{Applications of derivatives}
\LANGen{
\Question{
Evaluate the following limit. (Apply  L'Hospital's rule repeatedly.)
\[ \lim_{x\to0}\frac{x-\sin x}{x^3} \]}
{\( \frac16 \)}{\( -\frac16 \)}{\( \frac13 \)}{\( -\frac13 \)}{\( 0 \)}{}}
\LANGpl{
\Question{
Wyznacz granicę. (Zastosuj regułę de I'Hospitala)
\[ \lim_{x\to0}\frac{x-\sin x}{x^3} \]}
{\( \frac16 \)}{\( -\frac16 \)}{\( \frac13 \)}{\( -\frac13 \)}{\( 0 \)}{}}
\LANGcs{
\Question{
Vypočítejte následující limitu. (Opakovaně užijte L'Hospitalovo pravidlo.)
\[ \lim_{x\to0}\frac{x-\sin x}{x^3} \]}
{\( \frac16 \)}{\( -\frac16 \)}{\( \frac13 \)}{\( -\frac13 \)}{\( 0 \)}{}}
\LANGsk{
\Question{
Vypočítajte následujúcu limitu. (Opakovane použite L'Hospitalovo pravidlo.)
\[ \lim_{x\to0}\frac{x-\sin x}{x^3} \]}
{\( \frac16 \)}{\( -\frac16 \)}{\( \frac13 \)}{\( -\frac13 \)}{\( 0 \)}{}}
\LANGes{
\Question{
Resuelve el siguiente límite utilizando repetidamente la Regla de L'Hôpital:
\[ \lim_{x\to0}\frac{x-\sin x}{x^3} \]}
{\( \frac16 \)}{\( -\frac16 \)}{\( \frac13 \)}{\( -\frac13 \)}{\( 0 \)}{}}
\End

\Data\ID{1003259601}
\Updated{2020-07-15 03:07:19}
\Level{C}
\SubArea{Applications of derivatives}
\LANGen{
\Question{
Find all the asymptotes of the graph of the function \( f(x)=2x+\frac2{x-3} \).}
{\( x=3 \), \( y=2x \)}{\( y=3 \), \( y=2x \)}{\( x=3 \), \( y=-2x \)}{\( x=2 \), \( y=3x \)}{\( x=3 \), \( y=2 \)}{}}
\LANGpl{
\Question{
Wyznacz wszystkie asymptoty  funkcji  \( f(x)=2x+\frac2{x-3} \).}
{\( x=3 \), \( y=2x \)}{\( y=3 \), \( y=2x \)}{\( x=3 \), \( y=-2x \)}{\( x=2 \), \( y=3x \)}{\( x=3 \), \( y=2 \)}{}}
\LANGcs{
\Question{
Určete všechny asymptoty grafu funkce \( f(x)=2x+\frac2{x-3} \).}
{\( x=3 \), \( y=2x \)}{\( y=3 \), \( y=2x \)}{\( x=3 \), \( y=-2x \)}{\( x=2 \), \( y=3x \)}{\( x=3 \), \( y=2 \)}{}}
\LANGsk{
\Question{
Určte všetky asymptoty grafu funkcie \( f(x)=2x+\frac2{x-3} \).}
{\( x=3 \), \( y=2x \)}{\( y=3 \), \( y=2x \)}{\( x=3 \), \( y=-2x \)}{\( x=2 \), \( y=3x \)}{\( x=3 \), \( y=2 \)}{}}
\LANGes{
\Question{
Halla todas las asíntotas de la gráfica de la función \( f(x)=2x+\frac2{x-3} \).}
{\( x=3 \), \( y=2x \)}{\( y=3 \), \( y=2x \)}{\( x=3 \), \( y=-2x \)}{\( x=2 \), \( y=3x \)}{\( x=3 \), \( y=2 \)}{}}
\End

\Data\ID{1003259602}
\Updated{2020-07-15 03:07:19}
\Level{C}
\SubArea{Applications of derivatives}
\LANGen{
\Question{
Find all the asymptotes of the graph of the function \( f(x)=x+\frac{2x}{x^2-1} \).}
{\( x=-1 \), \( x=1 \), \( y=x \)}{\( x=-1 \), \( x=1 \), \( y=1 \)}{\( x=-1 \), \( x=1 \), \( y=x+1 \)}{\( x=-1 \), \( x=1 \)}{\( x=-1 \), \( x=1 \), \( y=0 \)}{}}
\LANGpl{
\Question{
Wyznacz wszystkie asymptoty funkcji \( f(x)=x+\frac{2x}{x^2-1} \).}
{\( x=-1 \), \( x=1 \), \( y=x \)}{\( x=-1 \), \( x=1 \), \( y=1 \)}{\( x=-1 \), \( x=1 \), \( y=x+1 \)}{\( x=-1 \), \( x=1 \)}{\( x=-1 \), \( x=1 \), \( y=0 \)}{}}
\LANGcs{
\Question{
Určete všechny asymptoty grafu funkce \( f(x)=x+\frac{2x}{x^2-1} \).}
{\( x=-1 \), \( x=1 \), \( y=x \)}{\( x=-1 \), \( x=1 \), \( y=1 \)}{\( x=-1 \), \( x=1 \), \( y=x+1 \)}{\( x=-1 \), \( x=1 \)}{\( x=-1 \), \( x=1 \), \( y=0 \)}{}}
\LANGsk{
\Question{
Určte všetky asymptoty grafu funkcie \( f(x)=x+\frac{2x}{x^2-1} \).}
{\( x=-1 \), \( x=1 \), \( y=x \)}{\( x=-1 \), \( x=1 \), \( y=1 \)}{\( x=-1 \), \( x=1 \), \( y=x+1 \)}{\( x=-1 \), \( x=1 \)}{\( x=-1 \), \( x=1 \), \( y=0 \)}{}}
\LANGes{
\Question{
Halla todas las asíntotas de la gráfica de la función \( f(x)=x+\frac{2x}{x^2-1} \).}
{\( x=-1 \), \( x=1 \), \( y=x \)}{\( x=-1 \), \( x=1 \), \( y=1 \)}{\( x=-1 \), \( x=1 \), \( y=x+1 \)}{\( x=-1 \), \( x=1 \)}{\( x=-1 \), \( x=1 \), \( y=0 \)}{}}
\End

\Data\ID{1003259603}
\Updated{2020-07-15 03:07:19}
\Level{C}
\SubArea{Applications of derivatives}
\LANGen{
\Question{
Find all the asymptotes of the graph of the function \( f(x)=\frac{x}{x^2+1} \).}
{\( y=0 \)}{there is no an asymptote}{\( x=0 \)}{\( x=-1 \), \( y=0 \)}{}{}}
\LANGpl{
\Question{
Wyznacz wszystkie asymptoty funkcji \( f(x)=\frac{x}{x^2+1} \).}
{\( y=0 \)}{brak asymptot}{\( x=0 \)}{\( x=-1 \), \( y=0 \)}{}{}}
\LANGcs{
\Question{
Určete všechny asymptoty grafu funkce \( f(x)=\frac{x}{x^2+1} \).}
{\( y=0 \)}{graf funkce nemá žádnou asymptotu}{\( x=0 \)}{\( x=-1 \), \( y=0 \)}{}{}}
\LANGsk{
\Question{
Určte všetky asymptoty grafu funkcie \( f(x)=\frac{x}{x^2+1} \).}
{\( y=0 \)}{graf funkcie nemá žiadnu asymptotu}{\( x=0 \)}{\( x=-1 \), \( y=0 \)}{}{}}
\LANGes{
\Question{
Halla todas las asíntotas de la gráfica de la función \( f(x)=\frac{x}{x^2+1} \).}
{\( y=0 \)}{la gráfica de la función no tiene ninguna asíntota}{\( x=0 \)}{\( x=-1 \), \( y=0 \)}{}{}}
\End

\Data\ID{1003259604}
\Updated{2020-07-15 03:07:19}
\Level{C}
\SubArea{Applications of derivatives}
\LANGen{
\Question{
Find all the asymptotes of the graph of the function \( f(x)=\frac1{1-x^2} \).}
{\( x=1 \), \( x=-1 \), \( y=0 \)}{\( x=1 \), \( x=-1 \)}{\( x=1 \), \( x=-1 \), \( y=1 \)}{\( y=1 \), \( y=-1 \), \( y=0 \)}{\( x=1 \), \( y=0 \)}{}}
\LANGpl{
\Question{
Wyznacz wszystkie asymptoty funkcji \( f(x)=\frac1{1-x^2} \).}
{\( x=1 \), \( x=-1 \), \( y=0 \)}{\( x=1 \), \( x=-1 \)}{\( x=1 \), \( x=-1 \), \( y=1 \)}{\( y=1 \), \( y=-1 \), \( y=0 \)}{\( x=1 \), \( y=0 \)}{}}
\LANGcs{
\Question{
Určete všechny asymptoty grafu funkce \( f(x)=\frac1{1-x^2} \).}
{\( x=1 \), \( x=-1 \), \( y=0 \)}{\( x=1 \), \( x=-1 \)}{\( x=1 \), \( x=-1 \), \( y=1 \)}{\( y=1 \), \( y=-1 \), \( y=0 \)}{\( x=1 \), \( y=0 \)}{}}
\LANGsk{
\Question{
Určte všetky asymptoty grafu funkcie \( f(x)=\frac1{1-x^2} \).}
{\( x=1 \), \( x=-1 \), \( y=0 \)}{\( x=1 \), \( x=-1 \)}{\( x=1 \), \( x=-1 \), \( y=1 \)}{\( y=1 \), \( y=-1 \), \( y=0 \)}{\( x=1 \), \( y=0 \)}{}}
\LANGes{
\Question{
Halla todas las asíntotas de la gráfica de la función \( f(x)=\frac1{1-x^2} \).}
{\( x=1 \), \( x=-1 \), \( y=0 \)}{\( x=1 \), \( x=-1 \)}{\( x=1 \), \( x=-1 \), \( y=1 \)}{\( y=1 \), \( y=-1 \), \( y=0 \)}{\( x=1 \), \( y=0 \)}{}}
\End

\Data\ID{1003259605}
\Updated{2020-07-15 03:07:19}
\Level{C}
\SubArea{Applications of derivatives}
\LANGen{
\Question{
Find all the asymptotes of the graph of the function \( f(x)=\frac1{x-2}+x-1 \).}
{\( x=2 \), \( y=x-1 \)}{\( x=2 \), \( y=x \)}{\( x=2 \)}{\( x=2 \), \( y=1 \)}{\( x=2 \), \( y=-1 \)}{}}
\LANGpl{
\Question{
Wyznacz wszystkie asymptoty funkcji \( f(x)=\frac1{x-2}+x-1 \).}
{\( x=2 \), \( y=x-1 \)}{\( x=2 \), \( y=x \)}{\( x=2 \)}{\( x=2 \), \( y=1 \)}{\( x=2 \), \( y=-1 \)}{}}
\LANGcs{
\Question{
Určete všechny asymptoty grafu funkce \( f(x)=\frac1{x-2}+x-1 \).}
{\( x=2 \), \( y=x-1 \)}{\( x=2 \), \( y=x \)}{\( x=2 \)}{\( x=2 \), \( y=1 \)}{\( x=2 \), \( y=-1 \)}{}}
\LANGsk{
\Question{
Určte všetky asymptoty grafu funkcie \( f(x)=\frac1{x-2}+x-1 \).}
{\( x=2 \), \( y=x-1 \)}{\( x=2 \), \( y=x \)}{\( x=2 \)}{\( x=2 \), \( y=1 \)}{\( x=2 \), \( y=-1 \)}{}}
\LANGes{
\Question{
Halla todas las asíntotas de la gráfica de la función \( f(x)=\frac1{x-2}+x-1 \).}
{\( x=2 \), \( y=x-1 \)}{\( x=2 \), \( y=x \)}{\( x=2 \)}{\( x=2 \), \( y=1 \)}{\( x=2 \), \( y=-1 \)}{}}
\End

\Data\ID{1003259606}
\Updated{2020-07-15 03:07:19}
\Level{C}
\SubArea{Applications of derivatives}
\LANGen{
\Question{
Find all the asymptotes of the graph of the function \( f(x)=\mathrm{e}^{\frac1x} \).}
{\( x=0 \), \( y=1 \)}{\( x=0 \), \( y=0 \)}{\( x=0 \)}{there is no an asymptote}{\( x=0 \), \( y=-1 \)}{}}
\LANGpl{
\Question{
Wyznacz wszystkie asymptoty funkcji \( f(x)=\mathrm{e}^{\frac1x} \).}
{\( x=0 \), \( y=1 \)}{\( x=0 \), \( y=0 \)}{\( x=0 \)}{brak asymptot}{\( x=0 \), \( y=-1 \)}{}}
\LANGcs{
\Question{
Určete všechny asymptoty grafu funkce \( f(x)=\mathrm{e}^{\frac1x} \).}
{\( x=0 \), \( y=1 \)}{\( x=0 \), \( y=0 \)}{\( x=0 \)}{graf funkce nemá žádnou asymptotu}{\( x=0 \), \( y=-1 \)}{}}
\LANGsk{
\Question{
Určte všetky asymptoty grafu funkcie \( f(x)=\mathrm{e}^{\frac1x} \).}
{\( x=0 \), \( y=1 \)}{\( x=0 \), \( y=0 \)}{\( x=0 \)}{graf funkcie nemá žiadnu asymptotu}{\( x=0 \), \( y=-1 \)}{}}
\LANGes{
\Question{
Halla todas las asíntotas de la gráfica de la función \( f(x)=\mathrm{e}^{\frac1x} \).}
{\( x=0 \), \( y=1 \)}{\( x=0 \), \( y=0 \)}{\( x=0 \)}{la gráfica de la función no tiene ninguna asíntota}{\( x=0 \), \( y=-1 \)}{}}
\End

\Data\ID{1003259607}
\Updated{2020-07-15 03:07:19}
\Level{C}
\SubArea{Applications of derivatives}
\LANGen{
\Question{
Find all the asymptotes of the graph of the function \( f(x)=x\mathrm{e}^{\frac1x} \).}
{\( x=0 \), \( y=x+1 \)}{\( x=0 \), \( y=x-1 \)}{\( x=0 \)}{there is no an asymptote}{\( x=0 \), \( y=1 \)}{}}
\LANGpl{
\Question{
Wyznacz wszystkie asymptoty funkcji \( f(x)=x\mathrm{e}^{\frac1x} \).}
{\( x=0 \), \( y=x+1 \)}{\( x=0 \), \( y=x-1 \)}{\( x=0 \)}{brak asymptot}{\( x=0 \), \( y=1 \)}{}}
\LANGcs{
\Question{
Určete všechny asymptoty grafu funkce \( f(x)=x\mathrm{e}^{\frac1x} \).}
{\( x=0 \), \( y=x+1 \)}{\( x=0 \), \( y=x-1 \)}{\( x=0 \)}{graf funkce nemá žádnou asymptotu}{\( x=0 \), \( y=1 \)}{}}
\LANGsk{
\Question{
Určte všetky asymptoty grafu funkcie \( f(x)=x\mathrm{e}^{\frac1x} \).}
{\( x=0 \), \( y=x+1 \)}{\( x=0 \), \( y=x-1 \)}{\( x=0 \)}{graf funkcie nemá žiadnu asymptotu}{\( x=0 \), \( y=1 \)}{}}
\LANGes{
\Question{
Halla todas las asíntotas de la gráfica de la función \( f(x)=x\mathrm{e}^{\frac1x} \).}
{\( x=0 \), \( y=x+1 \)}{\( x=0 \), \( y=x-1 \)}{\( x=0 \)}{la gráfica de la función no tiene ninguna asíntota}{\( x=0 \), \( y=1 \)}{}}
\End

\Data\ID{1003259608}
\Updated{2020-07-15 03:07:19}
\Level{C}
\SubArea{Applications of derivatives}
\LANGen{
\Question{
Determine the values of \( a \), \( b \) (\( a \), \( b\in\mathbb{R} \)) such that the line \( y=2x+\frac13 \)  is an asymptote of the graph of the function \( f(x)=\frac x{ax-1}+bx \).}
{\( a=3 \), \( b=2 \)}{\( a=\frac12 \), \( b=3 \)}{\( a=2 \), \( b=\frac13 \)}{\( a=\frac12 \), \( b=\frac13 \)}{\( a=\frac13 \), \( b=2 \)}{}}
\LANGpl{
\Question{
Wskaż wartość \( a \), \( b \) (\( a \), \( b\in\mathbb{R} \)) tak, aby \( y=2x+\frac13 \) była asymptotą wykresu funkcji \( f(x)=\frac x{ax-1}+bx \).}
{\( a=3 \), \( b=2 \)}{\( a=\frac12 \), \( b=3 \)}{\( a=2 \), \( b=\frac13 \)}{\( a=\frac12 \), \( b=\frac13 \)}{\( a=\frac13 \), \( b=2 \)}{}}
\LANGcs{
\Question{
Určete hodnoty parametrů \( a \), \( b \) (\( a \), \( b\in\mathbb{R} \)) tak, aby přímka \( y=2x+\frac13 \)  byla asymptotou grafu funkce \( f(x)=\frac x{ax-1}+bx \).}
{\( a=3 \), \( b=2 \)}{\( a=\frac12 \), \( b=3 \)}{\( a=2 \), \( b=\frac13 \)}{\( a=\frac12 \), \( b=\frac13 \)}{\( a=\frac13 \), \( b=2 \)}{}}
\LANGsk{
\Question{
Určte hodnoty parametrov \( a \), \( b \) (\( a \), \( b\in\mathbb{R} \)) tak, aby priamka \( y=2x+\frac13 \)  bola asymptotou grafu funkcie \( f(x)=\frac x{ax-1}+bx \).}
{\( a=3 \), \( b=2 \)}{\( a=\frac12 \), \( b=3 \)}{\( a=2 \), \( b=\frac13 \)}{\( a=\frac12 \), \( b=\frac13 \)}{\( a=\frac13 \), \( b=2 \)}{}}
\LANGes{
\Question{
Halla los valores de \( a \), \( b \) (\( a \), \( b\in\mathbb{R} \)) para que la recta \( y=2x+\frac13 \) sea una asíntota de la gráfica de la función \( f(x)=\frac x{ax-1}+bx \).}
{\( a=3 \), \( b=2 \)}{\( a=\frac12 \), \( b=3 \)}{\( a=2 \), \( b=\frac13 \)}{\( a=\frac12 \), \( b=\frac13 \)}{\( a=\frac13 \), \( b=2 \)}{}}
\End

\Data\ID{1003259609}
\Updated{2020-07-15 03:07:19}
\Level{C}
\SubArea{Applications of derivatives}
\LANGen{
\Question{
Determine the values of \( a \), \( b \) (\( a \), \( b\in\mathbb{R} \)) such that the line \( y=0 \) is an asymptote of the graph of the function \( f(x)=\frac x{ax-1}+bx \).}
{there are no such \( a \), \( b \)}{\( a\in\mathbb{R}\setminus\{1\} \), \( b=0 \)}{\( a=0 \), \( b=0 \)}{\( a\in\mathbb{R} \), \( b=0 \)}{}{}}
\LANGpl{
\Question{
Wskaż wartość  \( a \), \( b \) (\( a \), \( b\in\mathbb{R} \)) tak, aby prosta \( y=0 \) była asymptotą wykresu funkcji \( f(x)=\frac x{ax-1}+bx \).}
{brak wartośći \( a \), \( b \)}{\( a\in\mathbb{R}\setminus\{1\} \), \( b=0 \)}{\( a=0 \), \( b=0 \)}{\( a\in\mathbb{R} \), \( b=0 \)}{}{}}
\LANGcs{
\Question{
Určete hodnoty parametrů \( a \), \( b \) (\( a \), \( b\in\mathbb{R} \)) tak, aby přímka \( y=0 \) byla asymptotou grafu funkce \( f(x)=\frac x{ax-1}+bx \).}
{neexistují žádná taková \( a \), \( b \)}{\( a\in\mathbb{R}\setminus\{1\} \), \( b=0 \)}{\( a=0 \), \( b=0 \)}{\( a\in\mathbb{R} \), \( b=0 \)}{}{}}
\LANGsk{
\Question{
Určte hodnoty parametrov \( a \), \( b \) (\( a \), \( b\in\mathbb{R} \)) tak, aby priamka \( y=0 \) bola asymptotou grafu funkcie \( f(x)=\frac x{ax-1}+bx \).}
{neexistujú žiadne také \( a \), \( b \)}{\( a\in\mathbb{R}\setminus\{1\} \), \( b=0 \)}{\( a=0 \), \( b=0 \)}{\( a\in\mathbb{R} \), \( b=0 \)}{}{}}
\LANGes{
\Question{
Halla los valores de \( a \), \( b \) (\( a \), \( b\in\mathbb{R} \)) para que la recta \( y=0 \) sea una asíntota de la gráfica de la función \( f(x)=\frac x{ax-1}+bx \).}
{no hay ningunas \( a \), \( b \) así}{\( a\in\mathbb{R}\setminus\{1\} \), \( b=0 \)}{\( a=0 \), \( b=0 \)}{\( a\in\mathbb{R} \), \( b=0 \)}{}{}}
\End

\Data\ID{1003259610}
\Updated{2020-07-15 03:07:19}
\Level{C}
\SubArea{Applications of derivatives}
\LANGen{
\Question{
Let there be a function \( f(x)=\frac{ax^2}{x-b} \), \( a \), \( b\in\mathbb{R} \). Determine the values of \( a \), \( b \), such that the line \( y=3x+2 \) is an asymptote of the graph of the function \( f \).}
{\( a=3 \), \( b=\frac23 \)}{\( a=3 \), \( b=\frac43 \)}{\( a=3 \), \( b=2 \)}{\( a=2 \), \( b=\frac32 \)}{there are no such \( a \), \( b \)}{}}
\LANGpl{
\Question{
Dana jest funkcja  \( f(x)=\frac{ax^2}{x-b} \), \( a \), \( b\in\mathbb{R} \). Wskaż wartości \( a \), \( b \) tak, aby prosta \( y=3x+2 \) była asymptotą wykresu funkcji  \( f \).}
{\( a=3 \), \( b=\frac23 \)}{\( a=3 \), \( b=\frac43 \)}{\( a=3 \), \( b=2 \)}{\( a=2 \), \( b=\frac32 \)}{brak wartości \( a \), \( b \)}{}}
\LANGcs{
\Question{
Je dána funkce \( f(x)=\frac{ax^2}{x-b} \), kde \( a \), \( b\in\mathbb{R} \). Určete hodnoty parametrů \( a \), \( b \) tak, aby přímka \( y=3x+2 \) byla asymptotou grafu funkce \( f \).}
{\( a=3 \), \( b=\frac23 \)}{\( a=3 \), \( b=\frac43 \)}{\( a=3 \), \( b=2 \)}{\( a=2 \), \( b=\frac32 \)}{neexistují žádná taková \( a \), \( b \)}{}}
\LANGsk{
\Question{
Je daná funkcia \( f(x)=\frac{ax^2}{x-b} \), kde \( a \), \( b\in\mathbb{R} \). Určte hodnoty parametrov \( a \), \( b \) tak, aby priamka \( y=3x+2 \) bola asymptotou grafu funkcie \( f \).}
{\( a=3 \), \( b=\frac23 \)}{\( a=3 \), \( b=\frac43 \)}{\( a=3 \), \( b=2 \)}{\( a=2 \), \( b=\frac32 \)}{neexistujú žiadne také \( a \), \( b \)}{}}
\LANGes{
\Question{
Dada la función \( f(x)=\frac{ax^2}{x-b} \), \( a \), \( b\in\mathbb{R} \). Halla los valores de \( a \), \( b \), para que la recta \( y=3x+2 \) sea una asíntota de la gráfica de la función \( f \).}
{\( a=3 \), \( b=\frac23 \)}{\( a=3 \), \( b=\frac43 \)}{\( a=3 \), \( b=2 \)}{\( a=2 \), \( b=\frac32 \)}{no hay ningunas \( a \), \( b \) así}{}}
\End

\Data\ID{1003261906}
\Updated{2022-04-23 05:04:17}
\Level{C}
\SubArea{Applications of derivatives}
\LANGen{
\Question{
Determine the values of \( a \) and \( b \) (\( a \), \( b \in\mathbb{R} \)) such that the function
\[ f(x)=ax^3+bx^2+1 \] 
has a local extremum of \( 9 \) at \( x=-2 \).}
{\( a=2 \), \( b=6 \)}{\( a=-2 \), \( b = 6 \)}{\( a=-2 \), \( b = -6 \)}{\( a=2 \), \( b = -6 \)}{}{}}
\LANGpl{
\Question{
Wskaż wartości \( a \) i \( b \) (\( a \), \( b \in\mathbb{R} \)) tak, aby funkcja
\[ f(x)=ax^3+bx^2+1 \] 
miała lokalne ekstremum \( 9 \) w \( x=-2 \).}
{\( a=2 \), \( b=6 \)}{\( a=-2 \), \( b = 6 \)}{\( a=-2 \), \( b = -6 \)}{\( a=2 \), \( b = -6 \)}{}{}}
\LANGcs{
\Question{
Určete hodnoty \( a \), \( b \in\mathbb{R} \) tak, aby funkce
\[ f(x)=ax^3+bx^2+1 \] 
měla lokální extrém v bodě \( x=-2 \)  a jeho hodnota byla \( 9 \).}
{\( a=2 \), \( b=6 \)}{\( a=-2 \), \( b = 6 \)}{\( a=-2 \), \( b = -6 \)}{\( a=2 \), \( b = -6 \)}{}{}}
\LANGsk{
\Question{
Určte hodnoty \( a \), \( b \in\mathbb{R} \) tak, aby funkcia
\[ f(x)=ax^3+bx^2+1 \] 
mala lokálny extrém v bode \( x=-2 \)  a jeho hodnota bola \( 9 \).}
{\( a=2 \), \( b=6 \)}{\( a=-2 \), \( b = 6 \)}{\( a=-2 \), \( b = -6 \)}{\( a=2 \), \( b = -6 \)}{}{}}
\LANGes{
\Question{
Halla los valores de \( a \) y \( b \) (\( a \), \( b \in\mathbb{R} \)) para que la función
\[ f(x)=ax^3+bx^2+1 \] 
tenga un extremo local \( 9 \) en \( x=-2 \).}
{\( a=2 \), \( b=6 \)}{\( a=-2 \), \( b = 6 \)}{\( a=-2 \), \( b = -6 \)}{\( a=2 \), \( b = -6 \)}{}{}}
\End

\Data\ID{1003261907}
\Updated{2021-08-12 10:08:52}
\Level{C}
\SubArea{Applications of derivatives}
\LANGen{
\Question{
Determine the values of \( m \) and \( n \) (\( m \), \( n\in\mathbb{R} \)) such that the function
\[ f(x)=x^4+mx^3-n+2 \]
has a local minimum of \( -30 \) at \( x=3 \).}
{\( m=-4 \), \( n=5 \)}{\( m = 4 \) \( n=5 \)}{\( m=4 \), \( n=140 \)}{\( m=-4 \), \( n=-5 \)}{}{}}
\LANGpl{
\Question{
Wskaż wartości \( m \) i \( n \) (\( m \), \( n\in\mathbb{R} \)) tak, aby funkcja
\[ f(x)=x^4+mx^3-n+2 \]
miała lokalne minimum  \( -30 \) w \( x=3 \).}
{\( m=-4 \), \( n=5 \)}{\( m = 4 \) \( n=5 \)}{\( m=4 \), \( n=140 \)}{\( m=-4 \), \( n=-5 \)}{}{}}
\LANGcs{
\Question{
Určete hodnoty \( m \), \( n\in\mathbb{R} \) tak, aby funkce
\[ f(x)=x^4+mx^3-n+2 \]
měla lokální minimum v bodě \( x=3 \) a jeho hodnota byla \( -30 \).}
{\( m=-4 \), \( n=5 \)}{\( m = 4 \) \( n=5 \)}{\( m=4 \), \( n=140 \)}{\( m=-4 \), \( n=-5 \)}{}{}}
\LANGsk{
\Question{
Určte hodnoty \( m \), \( n\in\mathbb{R} \) tak, aby funkcia
\[ f(x)=x^4+mx^3-n+2 \]
mala lokálne minimum v bode \( x=3 \) a jeho hodnota bola \( -30 \).}
{\( m=-4 \), \( n=5 \)}{\( m = 4 \) \( n=5 \)}{\( m=4 \), \( n=140 \)}{\( m=-4 \), \( n=-5 \)}{}{}}
\LANGes{
\Question{
Halla los valores de \( m \) y \( n \) (\( m \), \( n\in\mathbb{R} \))  para que la función
\[ f(x)=x^4+mx^3-n+2 \]
tenga un mínimo local \( -30 \) en \( x=3 \).}
{\( m=-4 \), \( n=5 \)}{\( m = 4 \) \( n=5 \)}{\( m=4 \), \( n=140 \)}{\( m=-4 \), \( n=-5 \)}{}{}}
\End

\Data\ID{1003261908}
\Updated{2021-08-12 10:08:23}
\Level{C}
\SubArea{Applications of derivatives}
\LANGen{
\Question{
Determine all the values of \( t \), \( t\in\mathbb{R} \), such that the function
\[ f(x)=tx^3+(t+1)x^2-(t-2)x+3 \]
has local extrema.}
{\( t\in\mathbb{R}\setminus\left\{\frac12\right\} \)}{\( t\in\mathbb{R} \)}{\( t\in\left(-\frac12;\frac12\right) \)}{\( t\in\left(-\infty;-\frac12\right)\cup\left(\frac12;\infty\right) \)}{}{}}
\LANGpl{
\Question{
Wskaż wszystkie wartości \( t \), \( t\in\mathbb{R} \) tak, aby funkcja
\[ f(x)=tx^3+(t+1)x^2-(t-2)x+3 \]
miała lokalne ekstrema.}
{\( t\in\mathbb{R}\setminus\left\{\frac12\right\} \)}{\( t\in\mathbb{R} \)}{\( t\in\left(-\frac12;\frac12\right) \)}{\( t\in\left(-\infty;-\frac12\right)\cup\left(\frac12;\infty\right) \)}{}{}}
\LANGcs{
\Question{
Určete všechny hodnoty  \( t \), \( t\in\mathbb{R} \), pro něž má funkce 
\[ f(x)=tx^3+(t+1)x^2-(t-2)x+3 \]
 lokální extrémy.}
{\( t\in\mathbb{R}\setminus\left\{\frac12\right\} \)}{\( t\in\mathbb{R} \)}{\( t\in\left(-\frac12;\frac12\right) \)}{\( t\in\left(-\infty;-\frac12\right)\cup\left(\frac12;\infty\right) \)}{}{}}
\LANGsk{
\Question{
Určte všetky hodnoty  \( t \), \( t\in\mathbb{R} \), pre ktoré má funkcia 
\[ f(x)=tx^3+(t+1)x^2-(t-2)x+3 \]
 lokálne extrémy.}
{\( t\in\mathbb{R}\setminus\left\{\frac12\right\} \)}{\( t\in\mathbb{R} \)}{\( t\in\left(-\frac12;\frac12\right) \)}{\( t\in\left(-\infty;-\frac12\right)\cup\left(\frac12;\infty\right) \)}{}{}}
\LANGes{
\Question{
Halla todos los valores de \( t \), \( t\in\mathbb{R} \), para que la función
\[ f(x)=tx^3+(t+1)x^2-(t-2)x+3 \]
tenga extremos locales.}
{\( t\in\mathbb{R}\setminus\left\{\frac12\right\} \)}{\( t\in\mathbb{R} \)}{\( t\in\left(-\frac12;\frac12\right) \)}{\( t\in\left(-\infty;-\frac12\right)\cup\left(\frac12;\infty\right) \)}{}{}}
\End

\Data\ID{1003261909}
\Updated{2020-07-15 03:07:19}
\Level{C}
\SubArea{Applications of derivatives}
\LANGen{
\Question{
Determine all the values of \( a \), \( a\in\mathbb{R} \), such that the function
\[ f(x)=\frac{a^2-1}3x^3+(a-1)x^2+2x+1 \]
does not have any local extrema.}
{\( a\in(-\infty;-3]\cup[1;\infty) \)}{\( a\in(-\infty;-3)\cup(1;\infty) \)}{\( a\in(-3;1) \)}{\( a\in[-3;1] \)}{}{}}
\LANGpl{
\Question{
Wskaż wszystkie wartości  \( a \), \( a\in\mathbb{R} \) tak, aby funkcja
\[ f(x)=\frac{a^2-1}3x^3+(a-1)x^2+2x+1 \]
nie miała żadnego lokalnego ekstremum.}
{\( a\in(-\infty;-3\rangle\cup\langle1;\infty) \)}{\( a\in(-\infty;-3)\cup(1;\infty) \)}{\( a\in(-3;1) \)}{\( a\in\langle-3;1\rangle \)}{}{}}
\LANGcs{
\Question{
Určete všechny hodnoty \( a \), \( a\in\mathbb{R} \), pro něž funkce 
\[ f(x)=\frac{a^2-1}3x^3+(a-1)x^2+2x+1 \]
nemá lokální extrémy.}
{\( a\in(-\infty;-3\rangle\cup\langle1;\infty) \)}{\( a\in(-\infty;-3)\cup(1;\infty) \)}{\( a\in(-3;1) \)}{\( a\in\langle-3;1\rangle \)}{}{}}
\LANGsk{
\Question{
Určte všetky hodnoty \( a \), \( a\in\mathbb{R} \), pre ktoré funkcia 
\[ f(x)=\frac{a^2-1}3x^3+(a-1)x^2+2x+1 \]
nemá lokálne extrémy.}
{\( a\in(-\infty;-3\rangle\cup\langle1;\infty) \)}{\( a\in(-\infty;-3)\cup(1;\infty) \)}{\( a\in(-3;1) \)}{\( a\in\langle-3;1\rangle \)}{}{}}
\LANGes{
\Question{
Halla todos los valores de \( a \), \( a\in\mathbb{R} \), para que la función
\[ f(x)=\frac{a^2-1}3x^3+(a-1)x^2+2x+1 \]
no tenga ningún extremo local.}
{\( a\in(-\infty;-3]\cup[1;\infty) \)}{\( a\in(-\infty;-3)\cup(1;\infty) \)}{\( a\in(-3;1) \)}{\( a\in[-3;1] \)}{}{}}
\End

\Data\ID{1003266402}
\Updated{2021-08-15 08:08:53}
\Level{C}
\SubArea{Applications of derivatives}
\LANGen{
\Question{
The price of an Archery game experience program for groups up to $8$ participants is $12$ EUR/person. In case of a larger group (number of participants higher than $8$), each additional person reduces the price for all participants by $0.5$ $\mathrm{EUR}$/person. Find the number of participants that will bring the organizing company maximum income and calculate the total income.}
{There will be a maximum income of $128$ $\mathrm{EUR}$ for $16$ participants.}{There will be a maximum income of $128$ $\mathrm{EUR}$ for $8$ participants.}{There will be a maximum income of $192$ $\mathrm{EUR}$ for $16$ participants.}{There will be a maximum income of $192$ $\mathrm{EUR}$ for $12$ participants.}{None of the answers is correct.}{}}
\LANGpl{
\Question{
Cena kursu łucznictwa dla grup do $8$ uczestników wynosi $12$ EUR/ za osobę. W przypadku większej grupy (liczba uczestników większa niż $8$), każda dodatkowa osoba zmniejsza cenę dla wszystkich uczestników o  $0{,}5$ $\mathrm{EUR}$/ za osobę. Wskaż liczbę uczestników, która przyniesie firmie najwyższy dochód oraz oblicz całkowity dochód.}
{Maksymalny dochód wyniesie $128$ $\mathrm{EUR}$ dla $16$ uczestników.}{Maksymalny dochód wyniesie $128$ $\mathrm{EUR}$ dla $8$ uczestników.}{Maksymalny dochód wyniesie $192$ $\mathrm{EUR}$ dla $16$ uczestników.}{Maksymalny dochód wyniesie $192$ $\mathrm{EUR}$ dla $12$ uczestników.}{Brak poprawnej odpowiedzi.}{}}
\LANGcs{
\Question{
Cena zážitkového programu Archery game pro skupiny do $8$ účastníků je $12$ EUR/os. Pro větší skupiny (počet osob je větší než $8$) se s každou další osobou snižuje cena pro všechny účastníky o $0{,}5$ $\mathrm{EUR}$/os. Při jakém počtu účastníků bude mít pořádající společnost z toho programu maximální příjem a kolik tento příjem bude?}
{Maximální příjem bude $128$ $\mathrm{EUR}$ při účasti $16$ osob.}{Maximální příjem bude $128$ $\mathrm{EUR}$ při účasti $8$ osob.}{Maximální příjem bude $192$ $\mathrm{EUR}$ při účasti $16$ osob.}{Maximální příjem bude $192$ $\mathrm{EUR}$ při účasti $12$ osob.}{Žádná z odpovědí není správná.}{}}
\LANGsk{
\Question{
Cena zážitkového programu Archery game pre skupiny do $8$ účastníkov je $12$ EUR/os. Pre väčšie skupiny (počet osôb je väčší než $8$) sa s každou ďalšou osobou znižuje cena pre všetkých účastníkov o $0{,}5$ $\mathrm{EUR}$/os. Pri akom počte účastníkov bude mať usporiadateľská spoločnosť z toho programu maximálny príjem a koľko tento príjem bude?}
{Maximálny príjem bude $128$ $\mathrm{EUR}$ pri účasti $16$ osob.}{Maximálny príjem bude $128$ $\mathrm{EUR}$ pri účasti $8$ osob.}{Maximálny príjem bude $192$ $\mathrm{EUR}$ pri účasti $16$ osôb.}{Maximálny príjem bude $192$ $\mathrm{EUR}$ pri účasti $12$ osôb.}{Žiadna z odpovedí nie je správna.}{}}
\LANGes{
\Question{
El precio de un curso de tiro con arco para grupos de hasta $8$ participantes es $12$ EUR/persona. En el caso de un grupo más grande (el númeto de participantes superior a $8$), cada persona adicional reduce el precio de todos los participantes en $0.5$ $\mathrm{EUR}$/persona. Halla el número de participantes que traerá a la empresa organizadora el máximo ingreso y calcula el ingreso total.}
{El ingreso máximo sería $128$ $\mathrm{EUR}$ por $16$ participantes.}{El ingreso máximo sería $128$ $\mathrm{EUR}$ por $8$ participantes.}{El ingreso máximo sería $192$ $\mathrm{EUR}$ por $16$ participantes.}{El ingreso máximo sería $192$ $\mathrm{EUR}$ por $12$ participantes.}{Ninguna de las respuestas es correcta.}{}}
\End

\Data\ID{1003266404}
\Updated{2020-07-15 03:07:12}
\Level{C}
\SubArea{Applications of derivatives}
\LANGen{
\Question{
Find the shortest possible distance between the point $[4;0]$ and the graph of $f(x)=\sqrt x$.}
{$\frac{\sqrt{15}}2$}{$\frac72$}{$2$}{$7-\sqrt{15}$}{$\frac{7-\sqrt{15}}2$}{$\frac{1+\sqrt{65}}2$}}
\LANGpl{
\Question{
Wskaż  najkrótszą możliwą odległość pomiędzy punktem $[4;0]$ a wykresem $f(x)=\sqrt x$.}
{$\frac{\sqrt{15}}2$}{$\frac72$}{$2$}{$7-\sqrt{15}$}{$\frac{7-\sqrt{15}}2$}{$\frac{1+\sqrt{65}}2$}}
\LANGcs{
\Question{
Najděte minimální vzdálenost bodu $[4;0]$ od grafu funkce $f(x)=\sqrt x$.}
{$\frac{\sqrt{15}}2$}{$\frac72$}{$2$}{$7-\sqrt{15}$}{$\frac{7-\sqrt{15}}2$}{$\frac{1+\sqrt{65}}2$}}
\LANGsk{
\Question{
Nájdite minimálnu vzdialenosť bodu $[4;0]$ od grafu funkcie $f(x)=\sqrt x$.}
{$\frac{\sqrt{15}}2$}{$\frac72$}{$2$}{$7-\sqrt{15}$}{$\frac{7-\sqrt{15}}2$}{$\frac{1+\sqrt{65}}2$}}
\LANGes{
\Question{
Halla la distancia más corta posible entre el punto $[4;0]$ y la gráfica de $f(x)=\sqrt x$.}
{$\frac{\sqrt{15}}2$}{$\frac72$}{$2$}{$7-\sqrt{15}$}{$\frac{7-\sqrt{15}}2$}{$\frac{1+\sqrt{65}}2$}}
\End

\Data\ID{1103266401}
\Updated{2022-06-24 08:06:51}
\Level{C}
\SubArea{Applications of derivatives}
\IMG{1103266401.pdf}
\LANGen{
\Question{
A producer of sterilized canned vegetables needs to reduce the production costs of a $0.5$ liter cylindrical can. Find such radius $r$ and the height $h$ of the can (in centimeters) so that its surface (i.e. the amount of material needed) is minimal.}
{$r\doteq 4.3\,\mathrm{cm}$, $h\doteq 8.6\,\mathrm{cm}$}{$r\doteq 3.4\,\mathrm{cm}$, $h\doteq 13.8\,\mathrm{cm}$}{$r\doteq 5.4\,\mathrm{cm}$, $h\doteq 5.5\,\mathrm{cm}$}{$r\doteq 3.4\,\mathrm{cm}$, $h\doteq 8.6\,\mathrm{cm}$}{}{}}
\LANGpl{
\Question{
Producent warzyw w puszkach musi zredukować koszty produkcji  puszki w kształcie walca o pojemności $0{,}5$ litra. Wyznacz promień  $r$ i wysokość $h$ puszki (w centymetrach) tak, aby jej powierzchnia (tj. ilość potrzebnego materiału) była minimalna.}
{$r\doteq 4{,}3\,\mathrm{cm}$, $h\doteq 8{,}6\,\mathrm{cm}$}{$r\doteq 3{,}4\,\mathrm{cm}$, $h\doteq 13{,}8\,\mathrm{cm}$}{$r\doteq 5{,}4\,\mathrm{cm}$, $h\doteq 5{,}5\,\mathrm{cm}$}{$r\doteq 3{,}4\,\mathrm{cm}$, $h\doteq 8{,}6\,\mathrm{cm}$}{}{}}
\LANGcs{
\Question{
Výrobce sterilované zeleniny potřebuje snížit náklady na výrobu válcové plechovky o objemu $0{,}5$ l. Jaký by měl být poloměr $r$ a výška $h$ plechovky (v cm), aby byl její povrch (a tím i spotřeba materiálu) minimální?}
{$r\doteq 4{,}3\,\mathrm{cm}$, $h\doteq 8{,}6\,\mathrm{cm}$}{$r\doteq 3{,}4\,\mathrm{cm}$, $h\doteq 13{,}8\,\mathrm{cm}$}{$r\doteq 5{,}4\,\mathrm{cm}$, $h\doteq 5{,}5\,\mathrm{cm}$}{$r\doteq 3{,}4\,\mathrm{cm}$, $h\doteq 8{,}6\,\mathrm{cm}$}{}{}}
\LANGsk{
\Question{
Výrobca sterilizovanej zeleniny potrebuje znížiť náklady na výrobu valcovej plechovky s objemom $0{,}5$ l. Aký by mal byť polomer $r$ a výška $h$ plechovky (v cm), aby bol jej povrch (a tým i spotreba materiálu) minimálna?}
{$r\doteq 4{,}3\,\mathrm{cm}$, $h\doteq 8{,}6\,\mathrm{cm}$}{$r\doteq 3{,}4\,\mathrm{cm}$, $h\doteq 13{,}8\,\mathrm{cm}$}{$r\doteq 5{,}4\,\mathrm{cm}$, $h\doteq 5{,}5\,\mathrm{cm}$}{$r\doteq 3{,}4\,\mathrm{cm}$, $h\doteq 8{,}6\,\mathrm{cm}$}{}{}}
\LANGes{
\Question{
Un productor de vegetales enlatados necesita reducir los costos de producción de una lata cilíndrica de $0.5$ litros. Halla el radio $r$ y la altura $h$ de la lata (en centímetros) para que su área (es decir, la cantidad de material necesaria) sea mínima.}
{$r\doteq 4.3\,\mathrm{cm}$, $h\doteq 8.6\,\mathrm{cm}$}{$r\doteq 3.4\,\mathrm{cm}$, $h\doteq 13.8\,\mathrm{cm}$}{$r\doteq 5.4\,\mathrm{cm}$, $h\doteq 5.5\,\mathrm{cm}$}{$r\doteq 3.4\,\mathrm{cm}$, $h\doteq 8.6\,\mathrm{cm}$}{}{}}
\End

\Data\ID{1103266403}
\Updated{2022-04-23 05:04:17}
\Level{C}
\SubArea{Applications of derivatives}
\IMG{1103266403.pdf}
\LANGen{
\Question{
We want to create a rabbit cage in the shape of a rectangle with sides $a$ and $b$. The cage will be divided by parallel walls into four sections with the same area (see the picture). Find the dimensions $a$ and $b$ providing we have $50\,\mathrm{m}$ of fencing wire available and we want the total area to be as large as possible. (Fencing wire will also be used for the walls.)}
{$a=5\,\mathrm{m}$, $b=12.5\,\mathrm{m}$}{$a=4\,\mathrm{m}$, $b=15\,\mathrm{m}$}{$a=4.5\,\mathrm{m}$, $b=13.75\,\mathrm{m}$}{$a=6.5\,\mathrm{m}$, $b=8.75\,\mathrm{m}$}{}{}}
\LANGpl{
\Question{
Chcemy zbudować klatkę dla królika w kształcie prostokąta o bokach  $a$ i $b$. Klatka zostanie podzielona przez równoległe ściany na cztery sekcje o tej samej powierzchni (patrz zdjęcie). Znajdź wymiary $a$ i $b$ zakładając, że mamy $50\,\mathrm{m}$ drutu ogrodzeniowego i chcemy, aby całkowita powierzchnia była jak największa. (Drut ogrodzeniowy będzie również używany do ścian).}
{$a=5\,\mathrm{m}$, $b=12{,}5\,\mathrm{m}$}{$a=4\,\mathrm{m}$, $b=15\,\mathrm{m}$}{$a=4{,}5\,\mathrm{m}$, $b=13{,}75\,\mathrm{m}$}{$a=6{,}5\,\mathrm{m}$, $b=8{,}75\,\mathrm{m}$}{}{}}
\LANGcs{
\Question{
Chceme vytvořit výběh pro králíky, který bude mít tvar obdélníku se stranami $a$ a $b$. Výběh má být rozdělen pomocí rovnoběžných přepážek na čtyři části se stejným obsahem (viz obrázek). Jaké budou celkové rozměry výběhu $a$ a $b$, máme-li k dispozici $50\,\mathrm{m}$ pletiva a chceme-li celkový obsah výběhu co největší? (Pletivo bude použito i na přepážky.)}
{$a=5\,\mathrm{m}$, $b=12{,}5\,\mathrm{m}$}{$a=4\,\mathrm{m}$, $b=15\,\mathrm{m}$}{$a=4{,}5\,\mathrm{m}$, $b=13{,}75\,\mathrm{m}$}{$a=6{,}5\,\mathrm{m}$, $b=8{,}75\,\mathrm{m}$}{}{}}
\LANGsk{
\Question{
Chceme vytvoriť výbeh pre králiky, ktorý bude mať tvar obdĺžnika so stranami $a$ a $b$. Výbeh má byť rozdelený pomocou rovnobežných prepážok na štyri časti s rovnakým obsahom (viď obrázok). Aké budú celkové rozmery výbehu $a$ a $b$, ak máme k dispozícii $50\,\mathrm{m}$ pletiva a ak chceme celkový obsah výbehu čo najväčší? (Pletivo bude použité i na prepážky.)}
{$a=5\,\mathrm{m}$, $b=12{,}5\,\mathrm{m}$}{$a=4\,\mathrm{m}$, $b=15\,\mathrm{m}$}{$a=4{,}5\,\mathrm{m}$, $b=13{,}75\,\mathrm{m}$}{$a=6{,}5\,\mathrm{m}$, $b=8{,}75\,\mathrm{m}$}{}{}}
\LANGes{
\Question{
Queremos construir una jaula para conejos en forma de rectángulo de lados $a$ y $b$. Dividiremos la jaula con paredes paralelas formando cuatro secciones con la misma área (mira la imagen). Halla las dimensiones de $a$ y $b$ sabiendo que tenemos $50\,\mathrm{m}$ de alambre y queremos que el área total sea lo más grande posible. (El alambre se usará también para las paredes.)}
{$a=5\,\mathrm{m}$, $b=12.5\,\mathrm{m}$}{$a=4\,\mathrm{m}$, $b=15\,\mathrm{m}$}{$a=4.5\,\mathrm{m}$, $b=13.75\,\mathrm{m}$}{$a=6.5\,\mathrm{m}$, $b=8.75\,\mathrm{m}$}{}{}}
\End

\Data\ID{1103266405}
\Updated{2021-08-15 08:08:25}
\Level{C}
\SubArea{Applications of derivatives}
\IMG{1103266405.pdf}
\LANGen{
\Question{
Adam's House ($A$) is located at the distance of $0.9\,\mathrm{km}$ from the road. There is a bus stop ($B$) on this road at the distance of $1.5\,\mathrm{km}$ from the house (see the picture). Adam has overslept and needs to get to the bus stop as quickly as possible. At what distance $x$ from the nearest point $P$ should Adam reach the road knowing that he can move at the speed of $6\,\mathrm{km}/\mathrm{h}$ in rough terrain while being on the road he can move at the speed of $10\,\mathrm{km}/\mathrm{h}$?}
{$0.675\,\mathrm{km}$}{$0.525\,\mathrm{km}$}{$0.625\,\mathrm{km}$}{$0.575\,\mathrm{km}$}{}{}}
\LANGpl{
\Question{
Dom Adama ($A$) znajduje się $0{,}9\,\mathrm{km}$ od drogi. Przystanek autobusowy ($B$) na drodze znajduje się $1{,}5\,\mathrm{km}$ od domu (spójrz na rysunek). Adam zaspał i musi jak najszybciej dotrzeć na przystanek autobusowy. W jakiej odległośc  $x$ od najbliższego punktu $P$ Adam powinien dotrzeć do drogi wiedząc, że  może się poruszać $6\,\mathrm{km}/\mathrm{h}$  w trudnym terenie, będąc na drodze porusza się z prędkością  $10\,\mathrm{km}/\mathrm{h}$?}
{$0{,}675\,\mathrm{km}$}{$0{,}525\,\mathrm{km}$}{$0{,}625\,\mathrm{km}$}{$0{,}575\,\mathrm{km}$}{}{}}
\LANGcs{
\Question{
Adamův dům ($A$) je umístěn ve vzdálenosti $0{,}9\,\mathrm{km}$ od cesty, po níž jezdí autobusy. Zastávka autobusu ($B$) je umístěna na této cestě ve vzdálenosti $1{,}5\,\mathrm{km}$ od domu (viz obrázek). Adam zaspal a potřebuje se co nejrychleji dostat na zastávku. V jaké vzdálenosti $x$ od nejbližšího bodu $P$ se má na cestu napojit, jestliže v terénu se může pohybovat rychlostí $6\,\mathrm{km}/\mathrm{h}$ a po cestě rychlostí $10\,\mathrm{km}/\mathrm{h}$?}
{$0{,}675\,\mathrm{km}$}{$0{,}525\,\mathrm{km}$}{$0{,}625\,\mathrm{km}$}{$0{,}575\,\mathrm{km}$}{}{}}
\LANGsk{
\Question{
Adamov dom ($A$) je umiestnený vo vzdialenosti $0{,}9\,\mathrm{km}$ od cesty, po ktorej jazdia autobusy. Zastávka autobusu ($B$) je umiestená na tejto ceste vo vzdialenosti $1{,}5\,\mathrm{km}$ od domu (viď obrázok). Adam zaspal a potrebuje sa čo najrýchlejšie dostať na zastávku. V akej vzdialenosti $x$ od najbližšieho bodu $P$ sa má na cestu napojiť, ak v teréne sa môže pohybovať rýchlosťou $6\,\mathrm{km}/\mathrm{h}$ a po ceste rýchlosťou $10\,\mathrm{km}/\mathrm{h}$?}
{$0{,}675\,\mathrm{km}$}{$0{,}525\,\mathrm{km}$}{$0{,}625\,\mathrm{km}$}{$0{,}575\,\mathrm{km}$}{}{}}
\LANGes{
\Question{
La casa de Adam ($A$) se encuentra a una distancia de $0.9\,\mathrm{km}$ de la carretera. En la carretera hay una parada de autobuses ($B$) a una distancia de $1.5\,\mathrm{km}$ de la casa (mira la imagen). Adam se ha quedado dormido y necesita llegar a la parada lo más rápido posible. ¿A qué distancia $x$ desde el punto más cercano $P$  debe llegar Adam a la carretera sabiendo que se puede mover con una velocidad de $6\,\mathrm{km}/\mathrm{h}$ en el terreno accidentado y en la carretera con una velocidad de $10\,\mathrm{km}/\mathrm{h}$?}
{$0.675\,\mathrm{km}$}{$0.525\,\mathrm{km}$}{$0.625\,\mathrm{km}$}{$0.575\,\mathrm{km}$}{}{}}
\End

\Data\ID{1103266406}
\Updated{2021-08-19 10:08:52}
\Level{C}
\SubArea{Applications of derivatives}
\IMG{1103266406.pdf}
\LANGen{
\Question{
The medieval builder has a $5$-ell-long iron belt. His task is to shape the belt into a frame of the Romanesque window (that is the union of a rectangle and a semicircle, see the picture). Find the optimal width $x$ of the window to get as much light coming through the window as possible (i. e. the area of the window should be as large as possible). Express the result rounded in inches ($1\,\mathrm{ell} = 45\,\mathrm{inches}$).}
{$63$}{$140$}{$32$}{$112$}{$83$}{$20$}}
\LANGpl{
\Question{
Średniowieczny budowniczy ma żelazny pas o długości $5$ elli. Jego zadaniem jest ukształtowanie pasa w ramę romańskiego okna (to jest połączenie prostokąta i półkola, patrz zdjęcie). Znajdź optymalną szerokość $x$ okna, aby uzyskać jak najwięcej światła wpadającego przez okno (tj. Obszar okna powinien być tak duży, jak to możliwe). Wyraź wynik zaokrąglony w calach  ($1\,\mathrm{ell} = 45\,\mathrm{cali}$).}
{$63$}{$140$}{$32$}{$112$}{$83$}{$20$}}
\LANGcs{
\Question{
Středověký stavitel má železný pás o délce $5$ loktů, ze kterého potřebuje vytvarovat rám románského okna (to je sjednocením obdélníku a půlkruhu, viz obrázek). Určete optimální šířku okna $x$, aby jím procházelo co nejvíce světla (tzn. aby plocha okna byla co největší). Výsledek vyjádřete zaokrouhleně v palcích ($1$ loket = $45$ palců).}
{$63$}{$140$}{$32$}{$112$}{$83$}{$20$}}
\LANGsk{
\Question{
Stredoveký staviteľ má železný pás s dĺžkou $5$ lakťov, z ktorého potrebuje vytvarovať rám románskeho okna (to je zjednotením obdĺžnika a polkruhu, viď obrázok). Určte optimálnu šírku okna $x$, aby ním prechádzalo čo najviac svetla (tzn. aby plocha okna bola čo najväčšia). Výsledok vyjadrite zaokrúhlene v palcoch ($1$ lakeť = $45$ palcov).}
{$63$}{$140$}{$32$}{$112$}{$83$}{$20$}}
\LANGes{
\Question{
Un constructor medieval tiene una correa metálica de $5$ codos de largo. Su tarea es dar forma a la correa en un marco de ventana románica  (es decir, la unión de un rectángulo y un semicírculo, ver la imagen). Halla el ancho óptimo $x$ de la ventana para admita la cantidad más grande posible de luz (es decir, el área de la ventana debe ser lo más grande posible). Expresa el resultado redondeado en pulgadas ($1\,\mathrm{codo} = 45\,\mathrm{pulgadas}$).}
{$63$}{$140$}{$32$}{$112$}{$83$}{$20$}}
\End

\Data\ID{1103266407}
\Updated{2022-04-23 05:04:17}
\Level{C}
\SubArea{Applications of derivatives}
\IMG{1103266407.pdf}
\LANGen{
\Question{
Our task is to support a square tarp sized $4\,\mathrm{m}\times 4\,\mathrm{m}$ in the centers of its opposite sides to create a hay roof, as shown in the picture. What should be the height $v$ of the pillars to make the shelter space (volume) as large as possible?}
{$\sqrt2\,\mathrm{m}$}{$\frac{\sqrt2}2\,\mathrm{m}$}{$2\,\mathrm{m}$}{$\sqrt3\,\mathrm{m}$}{$\frac{\sqrt3}2\,\mathrm{m}$}{}}
\LANGpl{
\Question{
Naszym zadaniem jest wsparcie kwadratowej plandeki o wymiarach $4\,\mathrm{m}\times 4\,\mathrm{m}$ w środkach jej przeciwległych boków, aby utworzyć dach z siana, jak pokazano na rysunku. Jaka powinna być wysokość $v$ filarów, aby przestrzeń schronienia (objętość) była jak największa?}
{$\sqrt2\,\mathrm{m}$}{$\frac{\sqrt2}2\,\mathrm{m}$}{$2\,\mathrm{m}$}{$\sqrt3\,\mathrm{m}$}{$\frac{\sqrt3}2\,\mathrm{m}$}{}}
\LANGcs{
\Question{
Čtvercovou plachtu o rozměrech $4\,\mathrm{m}\times 4\,\mathrm{m}$ máme podepřít ve středech jejích protějších stran tak, aby vznikl přístřešek na seno, viz obrázek. Jakou výšku $v$ musejí mít podpěry, aby byl objem přístřešku co největší?}
{$\sqrt2\,\mathrm{m}$}{$\frac{\sqrt2}2\,\mathrm{m}$}{$2\,\mathrm{m}$}{$\sqrt3\,\mathrm{m}$}{$\frac{\sqrt3}2\,\mathrm{m}$}{}}
\LANGsk{
\Question{
Štvorcovú plachtu s rozmermi $4\,\mathrm{m}\times 4\,\mathrm{m}$ máme podoprieť v stredoch jej protiľahlých strán tak, aby vznikol prístrešok na seno, viď obrázok. Akú výšku $v$ musia mať podpery, aby bol objem prístreška čo najväčší?}
{$\sqrt2\,\mathrm{m}$}{$\frac{\sqrt2}2\,\mathrm{m}$}{$2\,\mathrm{m}$}{$\sqrt3\,\mathrm{m}$}{$\frac{\sqrt3}2\,\mathrm{m}$}{}}
\LANGes{
\Question{
Queremos levantar una lona cuadrada del tamaño de $4\,\mathrm{m}\times 4\,\mathrm{m}$ desde los centros de sus lados opuestos para crear un refugio (mira la imagen). ¿Cuál debería ser la altura $v$ de los pilares para que el espacio (volumen) del refugio sea lo más grande posible?}
{$\sqrt2\,\mathrm{m}$}{$\frac{\sqrt2}2\,\mathrm{m}$}{$2\,\mathrm{m}$}{$\sqrt3\,\mathrm{m}$}{$\frac{\sqrt3}2\,\mathrm{m}$}{}}
\End

\Data\ID{2000010801}
\Updated{2022-11-02 09:11:22}
\Level{C}
\SubArea{Applications of derivatives}
\LANGen{
\Question{
Consider non-uniform motion of an object whose position as a function of time is given by
\[
s=12t-\frac12 t^2,
\]
where time \(t\) is measured in seconds and position \(s\) is measured in meters. Find the instantaneous velocity of the object at \(8\) seconds. (Hint: Instantaneous velocity can be expressed as the derivative of position function with respect to time: \(v(t)=\frac{\mathrm{d}s}{\mathrm{d}t}\).)}
{\( 4 \,\mathrm{m}/\mathrm{s}\)}{\( 64\, \mathrm{m}/\mathrm{s}\)}{\( 8\,\mathrm{m}/\mathrm{s}\)}{The object will be at rest at this moment (\( v=0\, \mathrm{m}/\mathrm{s}\)).}{}{}}
\LANGpl{
\Question{
Rozważ ruch niejednostajny obiektu, którego położenie w funkcji czasu jest podane przez
\[
s=12t-\frac12 t^2,
\]
gdzie czas \(t\) mierzony jest w sekundach, a położenie \(s\) mierzony jest w metrach. Znajdź chwilową prędkość obiektu w \(8\) sekundach. (Wskazówka: Prędkość chwilową można wyrazić jako pochodną funkcji położenia względem czasu:  \(v(t)=\frac{\mathrm{d}s}{\mathrm{d}t}\).)}
{\( 4 \,\mathrm{m}/\mathrm{s}\)}{\( 64\, \mathrm{m}/\mathrm{s}\)}{\( 8\,\mathrm{m}/\mathrm{s}\)}{W tym momencie obiekt będzie w spoczynku  (\( v=0\, \mathrm{m}/\mathrm{s}\)).}{}{}}
\LANGcs{
\Question{
Pohyb tělesa, které se pohybuje nerovnoměrným pohybem je popsán rovnicí
\[
s=12t-\frac12 t^2,
\]
kde čas \(t\) je měřen v sekundách a dráha \(s\) je měřena v metrech. Určete velikost okamžité rychlosti, kterou se bude těleso pohybovat na konci osmé sekundy. (Nápověda: Okamžitou  rychlost můžeme určit pomocí derivace funkce dráhy, tj. \(v(t)=\frac{\mathrm{d}s}{\mathrm{d}t}\).)}
{\( 4 \,\mathrm{m}/\mathrm{s}\)}{\( 64\, \mathrm{m}/\mathrm{s}\)}{\( 8\,\mathrm{m}/\mathrm{s}\)}{V tom čase už bude těleso stát (\( v=0\, \mathrm{m}/\mathrm{s}\)).}{}{}}
\LANGsk{
\Question{
Pohyb telesa, ktoré sa pohybuje nerovnomerným pohybom je popísaný rovnicou
\[
s=12t-\frac12 t^2,
\]
kde čas \(t\) je nameraný v sekundách a dráha \(s\)  je nameraná v metroch. Určte veľkosť okamžitej rýchlosti, ktorou sa bude teleso pohybovať na konci \(8\) sekundy. 
(Pomôcka: Okamžitú rýchlosť môžeme určiť pomocou derivácie funkce: \(v(t)=\frac{\mathrm{d}s}{\mathrm{d}t}\).)}
{\( 4 \,\mathrm{m}/\mathrm{s}\)}{\( 64\, \mathrm{m}/\mathrm{s}\)}{\( 8\,\mathrm{m}/\mathrm{s}\)}{V tom čase už bude teleso stáť (\( v=0\, \mathrm{m}/\mathrm{s}\)).}{}{}}
\LANGes{
\Question{
Considera el movimiento no uniforme de un objeto cuya posición en función del tiempo viene dada por
\[
s=12t-\frac12 t^2,
\]
donde el tiempo \(t\) se mide en segundos y la distancia \(s\) se mide en metros. Calcula la velocidad instantánea del objeto a los \(8\) segundos. (Sugerencia: la velocidad instantánea se puede expresar como la derivada de la función de posición con respecto al tiempo: \(v(t)=\frac{\mathrm{d}s}{\mathrm{d}t}\).)}
{\( 4 \,\mathrm{m}/\mathrm{s}\)}{\( 64\, \mathrm{m}/\mathrm{s}\)}{\( 8\,\mathrm{m}/\mathrm{s}\)}{En este momento (\( v=0\, \mathrm{m}/\mathrm{s}\)) el objeto está en reposo.}{}{}}
\End

\Data\ID{2000010802}
\Updated{2022-11-02 09:11:41}
\Level{C}
\SubArea{Applications of derivatives}
\LANGen{
\Question{
Consider non-uniform motion of an object whose position as a function of time is given by
\[
s=t^3-t^2+\frac12 t,
\]
where time \(t\) is measured in seconds and position \(s\) is measured in meters. Find the instantaneous acceleration of the object at time \(t = 2\) s. (Hint: Instantaneous acceleration can be expressed as the derivative of the velocity function with respect to time and since velocity is the derivative of position function, instantaneous acceleration is its second derivative: \(a(t)=\frac{\mathrm{d}v}{\mathrm{d}t}=\frac{\mathrm{d}^2s}{\mathrm{d}t^2}\).)}
{\( 10 \,\frac{\mathrm{m}}{\mathrm{s}^2}\)}{\( 10.5 \,\frac{\mathrm{m}}{\mathrm{s}^2}\)}{\( 8.5 \,\frac{\mathrm{m}}{\mathrm{s}^2}\)}{\( 5\,\frac{\mathrm{m}}{\mathrm{s}^2}\)}{}{}}
\LANGpl{
\Question{
Rozważ ruch niejednostajny obiektu, którego położenie w funkcji czasu jest podane przez
\[
s=t^3-t^2+\frac12 t,
\]
gdzie czas \(t\) mierzony jest w sekundach, a położenie \(s\) mierzony jest w metrach. Znajdź chwilowe przyspieszenie obiektu w czasie \(t = 2\) s. (Wskazówka: Przyspieszenie chwilowe może być wyrażone jako pochodna funkcji prędkości względem czasu, a ponieważ prędkość jest pochodną funkcji położenia, przyspieszenie chwilowe jest jego drugą pochodną: \(a(t)=\frac{\mathrm{d}v}{\mathrm{d}t}=\frac{\mathrm{d}^2s}{\mathrm{d}t^2}\).)}
{\( 10 \,\frac{\mathrm{m}}{\mathrm{s}^2}\)}{\( 10{,}5 \,\frac{\mathrm{m}}{\mathrm{s}^2}\)}{\( 8{,}5 \,\frac{\mathrm{m}}{\mathrm{s}^2}\)}{\( 5\,\frac{\mathrm{m}}{\mathrm{s}^2}\)}{}{}}
\LANGcs{
\Question{
Pohyb tělesa, které se pohybuje nerovnoměrným pohybem je popsán rovnicí
\[
s=t^3-t^2+\frac12 t,
\]
kde čas \(t\) je měřen v sekundách a dráha \(s\) je měřena v metrech. Určete velikost okamžitého zrychlení tohoto tělesa na konci druhé sekundy jeho pohybu. (Nápověda: Okamžité zrychlení \(a\) můžeme určit pomocí derivace funkce rychlosti \(v(t)\). Protože rychlost můžeme určit pomocí derivace funkce \(s(t)\),  můžeme zrychlení určit pomocí její druhé derivace: \(a(t)=\frac{\mathrm{d}v}{\mathrm{d}t}=\frac{\mathrm{d}^2s}{\mathrm{d}t^2}\).)}
{\( 10 \,\frac{\mathrm{m}}{\mathrm{s}^2}\)}{\( 10{,}5 \,\frac{\mathrm{m}}{\mathrm{s}^2}\)}{\( 8{,}5 \,\frac{\mathrm{m}}{\mathrm{s}^2}\)}{\( 5\,\frac{\mathrm{m}}{\mathrm{s}^2}\)}{}{}}
\LANGsk{
\Question{
Pohyb telesa, ktoré sa pohybuje nerovnomerným pohybom je popísaný rovnicou
\[
s=t^3-t^2+\frac12 t,
\]
kde čas \(t\) je meraný v sekundách a dráha \(s\) je meraná v metroch. Určte veľkosť okamžitého zrýchlenia tohoto telesa na konci druhej sekundy jeho pohybu. (Pomôcka: Okamžité zrýchlenie \(a\) môžeme určiť pomocou derivácie funkcie rýchlosti \(v(t)\). Pretože rýchlosť môžeme určiť pomocou derivácie funkcie \(s(t)\),  môžeme zrýchlenie určiť pomocou jej druhej derivácie: \(a(t)=\frac{\mathrm{d}v}{\mathrm{d}t}=\frac{\mathrm{d}^2s}{\mathrm{d}t^2}\).)}
{\( 10 \,\frac{\mathrm{m}}{\mathrm{s}^2}\)}{\( 10{,}5 \,\frac{\mathrm{m}}{\mathrm{s}^2}\)}{\( 8{,}5 \,\frac{\mathrm{m}}{\mathrm{s}^2}\)}{\( 5\,\frac{\mathrm{m}}{\mathrm{s}^2}\)}{}{}}
\LANGes{
\Question{
Considera el movimiento no uniforme de un objeto cuya posición en función del tiempo viene dada por
\[
s=t^3-t^2+\frac12 t,
\]
donde el tiempo \(t\) se mide en segundos y la distancia \(s\) se mide en metros. Calcula la aceleración instantánea del objeto a los \(t = 2\) s. (Sugerencia: la aceleración instantánea se puede expresar como la derivada de la función de velocidad con respecto al tiempo y dado que la velocidad es la derivada de la función de posición, la aceleración instantánea es su segunda derivada: \(a(t)=\frac{\mathrm{d}v}{\mathrm{d}t}=\frac{\mathrm{d}^2s}{\mathrm{d}t^2}\).).}
{\( 10 \,\frac{\mathrm{m}}{\mathrm{s}^2}\)}{\( 10.5 \,\frac{\mathrm{m}}{\mathrm{s}^2}\)}{\( 8.5 \,\frac{\mathrm{m}}{\mathrm{s}^2}\)}{\( 5\,\frac{\mathrm{m}}{\mathrm{s}^2}\)}{}{}}
\End

\Data\ID{2000010803}
\Updated{2022-11-02 09:11:55}
\Level{C}
\SubArea{Applications of derivatives}
\IMG{2000010801-figure0.pdf}
\LANGen{
\Question{
Given the position-versus-time graph (in black) of an object in motion and the tangent line to the graph at the time point of \(10\) seconds (in red), find the instantaneous velocity of this object at \(10\) seconds. (Hint: Instantaneous velocity can be expressed as the derivative of position function with respect to time: \(v(t)=\frac{\mathrm{d}s}{\mathrm{d}t}\).)}
{\( 2 \,\frac{\mathrm{m}}{\mathrm{s}}\)}{\( 0.5 \,\frac{\mathrm{m}}{\mathrm{s}}\)}{\( 1 \,\frac{\mathrm{m}}{\mathrm{s}}\)}{\( 30\,\frac{\mathrm{m}}{\mathrm{s}}\)}{}{}}
\LANGpl{
\Question{
Mając wykres położenie w funkcji czasu (na czarno) obiektu w ruchu i linię styczną do wykresu w punkcie czasowym \(10\) sekund (na czerwono), znajdź prędkość chwilową tego obiektu w \( 10\) sekund. (Wskazówka: Prędkość chwilową można wyrazić jako pochodną funkcji położenia względem \(v(t)=\frac{\mathrm{d}s}{\mathrm{d}t}\).)}
{\( 2 \,\frac{\mathrm{m}}{\mathrm{s}}\)}{\( 0{,}5 \,\frac{\mathrm{m}}{\mathrm{s}}\)}{\( 1 \,\frac{\mathrm{m}}{\mathrm{s}}\)}{\( 30\,\frac{\mathrm{m}}{\mathrm{s}}\)}{}{}}
\LANGcs{
\Question{
Na obrázku je graf závislosti dráhy na čase pohybujícího se tělesa (černá barva). V čase  \(t=10\) sekund je sestrojena tečna ke grafu (červená barva). Pomocí obrázku určete rychlost tělesa v čase \(t=10\ \mathrm{s}\). (Nápověda: Okamžitou  rychlost můžeme určit pomocí derivace funkce dráhy, tj. \(v(t)=\frac{\mathrm{d}s}{\mathrm{d}t}\).)}
{\( 2 \,\frac{\mathrm{m}}{\mathrm{s}}\)}{\( 0{,}5 \,\frac{\mathrm{m}}{\mathrm{s}}\)}{\( 1 \,\frac{\mathrm{m}}{\mathrm{s}}\)}{\( 30\,\frac{\mathrm{m}}{\mathrm{s}}\)}{}{}}
\LANGsk{
\Question{
Na obrázku je graf závislosti dráhy na čase pohybujúceho sa telesa (čierna farba). V čase  \(t=10\) sekúnd je zostrojená dotyčnica ku grafu (červená farba). Pomocou obrázku určte rýchlosť telesa v čase \(t=10\ \mathrm{s}\). (Pomôcka: Okamžitú  rýchlosť môžeme určiť pomocou derivácie funkcie dráhy, tj. \(v(t)=\frac{\mathrm{d}s}{\mathrm{d}t}\).)}
{\( 2 \,\frac{\mathrm{m}}{\mathrm{s}}\)}{\( 0{,}5 \,\frac{\mathrm{m}}{\mathrm{s}}\)}{\( 1 \,\frac{\mathrm{m}}{\mathrm{s}}\)}{\( 30\,\frac{\mathrm{m}}{\mathrm{s}}\)}{}{}}
\LANGes{
\Question{
Dada la gráfica de distancia versus tiempo (en negro) de un objeto en movimiento y la línea tangente a la gráfica en el punto de tiempo de \(10\) segundos (en rojo), calcula la velocidad instantánea de este objeto en \(10\) segundos. (Sugerencia: la velocidad instantánea se puede expresar como la derivada de la función de posición con respecto al tiempo: \(v(t)=\frac{\mathrm{d}s}{\mathrm{d}t}\).).}
{\( 2 \,\frac{\mathrm{m}}{\mathrm{s}}\)}{\( 0.5 \,\frac{\mathrm{m}}{\mathrm{s}}\)}{\( 1 \,\frac{\mathrm{m}}{\mathrm{s}}\)}{\( 30\,\frac{\mathrm{m}}{\mathrm{s}}\)}{}{}}
\End

\Data\ID{2000010804}
\Updated{2022-11-02 09:11:17}
\Level{C}
\SubArea{Applications of derivatives}
\LANGen{
\Question{
For a given object to move with uniform acceleration, the engine must perform work that is related with time by the formula
\[
W=3t^2,
\]
where work \(W\) is measured in joules and time \(t\) is measured in seconds. Determine the instantaneous engine power at time \(t=4\,\mathrm{s}\). (Hint: Instantaneous power of a given object can be expressed as the derivative of work function with respect to time: \(P(t)=\frac{\mathrm{d}W}{\mathrm{d}t}\).)}
{\( 24 \,\mathrm{W}\)}{\( 48 \,\mathrm{W}\)}{\( 8 \,\mathrm{W}\)}{\( 12 \,\mathrm{W}\)}{}{}}
\LANGpl{
\Question{
Aby dany obiekt poruszał się z równomiernym przyspieszeniem, silnik musi wykonać pracę, która jest powiązana z czasem wzorem
\[
W=3t^2,
\]
gdzie praca \(W\) jest mierzona w dżulach, a czas \(t\) jest mierzony w sekundach. Wyznacz chwilową moc silnika w czasie \(t=4\,\mathrm{s}\). (Wskazówka: Moc chwilową danego obiektu można wyrazić jako pochodną funkcji pracy względem czasu: \(P(t)=\frac{\mathrm{d}W}{\mathrm{d}t}\).)}
{\( 24 \,\mathrm{W}\)}{\( 48 \,\mathrm{W}\)}{\( 8 \,\mathrm{W}\)}{\( 12 \,\mathrm{W}\)}{}{}}
\LANGcs{
\Question{
Aby dané těleso mohlo rovnoměrně zrychlovat, musí motor konat práci, která závisí na době pohybu vztahem
\[
W=3t^2,
\]
kde práce \(W\) je udávána v joulech a čas \(t\) v sekundách. Určete okamžitý výkon motoru v čase  \(t=4\,\mathrm{s}\). (Nápověda: Okamžitý výkon \(P\) můžeme určit pomocí derivace funkce \(W(t)\) tj. \(P(t)=\frac{\mathrm{d}W}{\mathrm{d}t}\).)}
{\( 24 \,\mathrm{W}\)}{\( 48 \,\mathrm{W}\)}{\( 8 \,\mathrm{W}\)}{\( 12 \,\mathrm{W}\)}{}{}}
\LANGsk{
\Question{
Aby dané teleso mohlo rovnomerne zrýchľovať, musí motor vykonať prácu, ktorá je závislá na čase pohybu vzťahom
\[
W=3t^2,
\]
kde práca \(W\) sa udáva v jouloch a čas \(t\) v sekundách. Určte okamžitý výkon motora v čase  \(t=4\,\mathrm{s}\). (Pomôcka: Okamžitý výkon \(P\) môžeme určiť pomocou derivácie funkcie \(W(t)\) tj. \(P(t)=\frac{\mathrm{d}W}{\mathrm{d}t}\).)}
{\( 24 \,\mathrm{W}\)}{\( 48 \,\mathrm{W}\)}{\( 8 \,\mathrm{W}\)}{\( 12 \,\mathrm{W}\)}{}{}}
\LANGes{
\Question{
Para que un objeto determinado se mueva con la aceleración uniforme, el motor debe realizar un trabajo que está relacionado con el tiempo por la fórmula
\[
W=3t^2,
\]
donde el trabajo \(W\) se mide en julios y el tiempo \(t\) se mide en segundos. Determina la potencia instantánea del motor en el momento \(t=4\,\mathrm{s}\). (Sugerencia: la potencia instantánea de un objeto dado se puede expresar como la derivada de la función de trabajo con respecto al tiempo: \(P(t)=\frac{\mathrm{d}W}{\mathrm{d}t}\).).}
{\( 24 \,\mathrm{W}\)}{\( 48 \,\mathrm{W}\)}{\( 8 \,\mathrm{W}\)}{\( 12 \,\mathrm{W}\)}{}{}}
\End

\Data\ID{2000010805}
\Updated{2022-11-02 09:11:40}
\Level{C}
\SubArea{Applications of derivatives}
\LANGen{
\Question{
A flywheel rotates such that it sweeps out an angle at the rate of
\[
\varphi = 4t^2,
\]
where an angle \(\varphi\) is measured in radians and time \(t\) is measured in seconds. At what time is instantaneous angular velocity of the flywheel equal to \(36\,\frac{\mathrm{rad}}{s}\)? (Hint: Instantaneous angular velocity can be expressed as the derivative of the function \(\varphi(t)\) with respect to time: \(\omega(t)=\frac{\mathrm{d}\varphi}{\mathrm{d}t}\).)}
{\( 4.5 \,\mathrm{s}\)}{\( 3\,\mathrm{s}\)}{\( 288 \,\mathrm{s}\)}{\( 9 \,\mathrm{s}\)}{}{}}
\LANGpl{
\Question{
Koło zamachowe obraca się tak, że wymiata kąt z szybkością
\[
\varphi = 4t^2,
\]
gdzie kąt \(\varphi\) jest mierzony w radianach, a czas \(t\) jest mierzony w sekundach. W jakim czasie chwilowa prędkość kątowa koła zamachowego jest równa \(36\,\frac{\mathrm{rad}}{s}\)? (Wskazówka: Chwilową prędkość kątową można wyrazić jako pochodną funkcji \(\varphi(t)\) względem czasu: \(\omega(t)=\frac{\mathrm{d}\varphi}{\mathrm{d}t}\).)}
{\( 4{,}5 \,\mathrm{s}\)}{\( 3\,\mathrm{s}\)}{\( 288 \,\mathrm{s}\)}{\( 9 \,\mathrm{s}\)}{}{}}
\LANGcs{
\Question{
Daný setrvačník se roztáčí tak, že úhel jeho otočení závisí na čase podle rovnice
\[
\varphi = 4t^2,
\]
kde úhel otočení \(\varphi\) udáváme v radiánech a čas \(t\) v sekundách. Za jak dlouho se bude pohybovat úhlovou rychlostí  \(36\,\frac{\mathrm{rad}}{s}\)? (Nápověda: Úhlovou rychlost můžeme určit pomocí derivace funkce \(\varphi(t)\), tj. \(\omega(t)=\frac{\mathrm{d}\varphi}{\mathrm{d}t}\).)}
{\( 4{,}5 \,\mathrm{s}\)}{\( 3\,\mathrm{s}\)}{\( 288 \,\mathrm{s}\)}{\( 9 \,\mathrm{s}\)}{}{}}
\LANGsk{
\Question{
Daný zotrvačník sa roztáča tak, že uhol jeho otočenia závisí na čase podľa rovnice
\[
\varphi = 4t^2,
\]
kde uhol otočenia \(\varphi\) udávame v radiánoch a čas \(t\) v sekundách. Za ako dlho sa bude pohybovať uhlovou rýchlosťou  \(36\,\frac{\mathrm{rad}}{s}\)? (Pomôcka: Uhlovú rýchlosť môžeme určiť pomocou derivácie funkcie \(\varphi(t)\), tj. \(\omega(t)=\frac{\mathrm{d}\varphi}{\mathrm{d}t}\).)}
{\( 4{,}5 \,\mathrm{s}\)}{\( 3\,\mathrm{s}\)}{\( 288 \,\mathrm{s}\)}{\( 9 \,\mathrm{s}\)}{}{}}
\LANGes{
\Question{
Un volante gira de tal manera que barre un ángulo a razón de
\[
\varphi = 4t^2,
\]
donde un ángulo \(\varphi\) se mide en radianes y el tiempo \(t\) se mide en segundos. ¿En qué momento la velocidad angular instantánea del volante es igual a \(36\,\frac{\mathrm{rad}}{s}\)? (Sugerencia: la velocidad angular instantánea se puede expresar como la derivada de la función \(\varphi(t)\) con respecto al tiempo: \(\omega(t)=\frac{\mathrm{d}\varphi}{\mathrm{d}t}\).).}
{\( 4.5 \,\mathrm{s}\)}{\( 3\,\mathrm{s}\)}{\( 288 \,\mathrm{s}\)}{\( 9 \,\mathrm{s}\)}{}{}}
\End

\Data\ID{2000010806}
\Updated{2022-11-02 09:11:52}
\Level{C}
\SubArea{Applications of derivatives}
\LANGen{
\Question{
Let’s have a coil of \(0.06\,\mathrm{H}\) inductance. The current flowing through the coil is given by 
\[
i=0.2\sin(100\pi t),\]
where time \(t\) is measured in seconds and current \(i\) is measured in amperes. Determine the voltage induced in the coil at time \(t=2\) seconds. (Hint: Instantaneous voltage can be expressed as the derivative of current function with respect to time: \(u(t)=-L\frac{\mathrm{d}i}{\mathrm{d}t}\). The negative sign indicates only that voltage induced opposes the change in current through the coil per unit time. It does not affect the magnitude of the voltage.)}
{\( 1.2\pi \,\mathrm{V}\)}{\( 20\pi \,\mathrm{V}\)}{\( 0 \,\mathrm{V}\)}{\( 12 \,\mathrm{V}\)}{}{}}
\LANGpl{
\Question{
Rozważmy cewkę o indukcyjności \(0{,}06\,\mathrm{H}\). Prąd płynący przez cewkę jest podawany przez
\[
i=0{,}2\sin(100\pi t),\]
gdzie czas \(t\) jest mierzony w sekundach, a prąd \(i\) jest mierzony w amperach. Określ napięcie indukowane w cewce w czasie \(t=2\) sekund. (Wskazówka: Chwilowe napięcie można wyrazić jako pochodną funkcji prądu względem czasu: \(u(t)=-L\frac{\mathrm{d}i}{\mathrm{d}t}\). Znak minus wskazuje tylko, że indukowane napięcie przeciwstawia się zmianie prądu przepływającego przez cewkę w jednostce czasu. Nie wpływa na wielkość napięcia.)}
{\( 1{,}2\pi \,\mathrm{V}\)}{\( 20\pi \,\mathrm{V}\)}{\( 0 \,\mathrm{V}\)}{\( 12 \,\mathrm{V}\)}{}{}}
\LANGcs{
\Question{
Cívkou, jejíž indukčnost je \(0{,}06\,\mathrm{H}\) protéká střídavý proud
\[
i=0{,}2\sin(100\pi t),\]
kde čas \(t\) je v sekundách a proud  \(i\) je měřen v ampérech. Určete velikost indukovaného napětí v čase  \(t=2\) s. (Nápověda: Okamžité napětí \(u\), které se indukuje v cívce o indukčnosti \(L\) průchodem střídavého proudu \(i\) můžeme určit pomocí derivace funkce \(i(t)\), tj. \(u(t)=-L\frac{\mathrm{d}i}{\mathrm{d}t}\). Pro velikost napětí není záporné znaménko podstatné.)}
{\( 1{,}2\pi \,\mathrm{V}\)}{\( 20\pi \,\mathrm{V}\)}{\( 0 \,\mathrm{V}\)}{\( 12 \,\mathrm{V}\)}{}{}}
\LANGsk{
\Question{
Cievkou, ktorej indukčnosť je \(0{,}06\,\mathrm{H}\) preteká striedavý prúd
\[
i=0{,}2\sin(100\pi t),\]
kde čas \(t\) je v sekundách a prúd  \(i\) je meraný v ampéroch. Určte veľkosť indukovaného napätia v čase  \(t=2\) s. (Pomôcka: Okamžité napätie \(u\), ktoré sa indukuje v cievke s indukčnosťou \(L\) prietokom striedavého prúdu \(i\) môžeme určiť pomocou derivácie funkcie \(i(t)\), tj. \(u(t)=-L\frac{\mathrm{d}i}{\mathrm{d}t}\). Pre veľkosť napätia nie je záporné znamienko podstatné.)}
{\( 1{,}2\pi \,\mathrm{V}\)}{\( 20\pi \,\mathrm{V}\)}{\( 0 \,\mathrm{V}\)}{\( 12 \,\mathrm{V}\)}{}{}}
\LANGes{
\Question{
Dada la bobina de \(0.06\,\mathrm{H}\) de inductancia. La corriente que fluye a través de la bobina está dada por 
\[
i=0.2\sin(100\pi t),\]
donde el tiempo \(t\) se mide en segundos y la corriente \(i\) se mide en amperios. Determina el voltaje inducido en la bobina en el tiempo \(t=2\) segundos. (Sugerencia: el voltaje instantáneo se puede expresar como la derivada de la función actual con respecto al tiempo: \(u(t)=-L\frac{\mathrm{d}i}{\mathrm{d}t}\). El signo negativo indica solo que el voltaje se opone al cambio de corriente a través de la bobina por unidad de tiempo. No afecta la magnitud del voltaje).}
{\( 1.2\pi \,\mathrm{V}\)}{\( 20\pi \,\mathrm{V}\)}{\( 0 \,\mathrm{V}\)}{\( 12 \,\mathrm{V}\)}{}{}}
\End

\Data\ID{2010011101}
\Updated{2022-08-25 10:08:46}
\Level{C}
\SubArea{Applications of derivatives}
\LANGen{
\Question{
Use L'Hospital's rule to find the following limit.
\[ \lim_{x\to2}\frac{x^3-4x^2+8}{2x^2-3x-2} \]}
{\( -\frac45\)}{\( \frac45\)}{\( -4\)}{\( 0\)}{\( \infty\)}{}}
\LANGpl{
\Question{
Użyj reguły L'Hospital, aby znaleźć następującą granicę.
\[ \lim_{x\to2}\frac{x^3-4x^2+8}{2x^2-3x-2} \]}
{\( -\frac45\)}{\( \frac45\)}{\( -4\)}{\( 0\)}{\( \infty\)}{}}
\LANGcs{
\Question{
Užitím L'Hospitalova pravidla vypočítejte následující limitu.
\[ \lim_{x\to2}\frac{x^3-4x^2+8}{2x^2-3x-2} \]}
{\( -\frac45\)}{\( \frac45\)}{\( -4\)}{\( 0\)}{\( \infty\)}{}}
\LANGsk{
\Question{
Použitím L'Hospitalovho pravidla vypočítajte nasledujúcu limitu.
\[ \lim_{x\to2}\frac{x^3-4x^2+8}{2x^2-3x-2} \]}
{\( -\frac45\)}{\( \frac45\)}{\( -4\)}{\( 0\)}{\( \infty\)}{}}
\LANGes{
\Question{
Resuelve el siguiente límite utilizando la Regla de L'Hôpital:
\[ \lim_{x\to2}\frac{x^3-4x^2+8}{2x^2-3x-2} \]}
{\( -\frac45\)}{\( \frac45\)}{\( -4\)}{\( 0\)}{\( \infty\)}{}}
\End

\Data\ID{2010011102}
\Updated{2022-08-25 10:08:46}
\Level{C}
\SubArea{Applications of derivatives}
\LANGen{
\Question{
Use L'Hospital's rule to find the following limit.
\[ \lim_{x\to0}\frac{\mathrm{tg}\,2x}{2x+\sin3x}\]}
{\( \frac25\)}{\( \frac12\)}{\( 10\)}{\( 0\)}{\(1\)}{}}
\LANGpl{
\Question{
Użyj reguły L'Hospital, aby znaleźć następującą granicę.
\[ \lim_{x\to0}\frac{\mathrm{tg}\,2x}{2x+\sin3x}\]}
{\( \frac25\)}{\( \frac12\)}{\( 10\)}{\( 0\)}{\(1\)}{}}
\LANGcs{
\Question{
Užitím L'Hospitalova pravidla vypočítejte následující limitu.
\[ \lim_{x\to0}\frac{\mathrm{tg}\,2x}{2x+\sin3x}\]}
{\( \frac25\)}{\( \frac12\)}{\( 10\)}{\( 0\)}{\(1\)}{}}
\LANGsk{
\Question{
Použitím L'Hospitalovho pravidla vypočítajte nasledujúcu limitu.
\[ \lim_{x\to0}\frac{\mathrm{tg}\,2x}{2x+\sin3x}\]}
{\( \frac25\)}{\( \frac12\)}{\( 10\)}{\( 0\)}{\(1\)}{}}
\LANGes{
\Question{
Resuelve el siguiente límite utilizando la Regla de L'Hôpital:
\[ \lim_{x\to0}\frac{\mathrm{tg}\,2x}{2x+\sin3x}\]}
{\( \frac25\)}{\( \frac12\)}{\( 10\)}{\( 0\)}{\(1\)}{}}
\End

\Data\ID{2010011103}
\Updated{2022-08-25 10:08:46}
\Level{C}
\SubArea{Applications of derivatives}
\LANGen{
\Question{
Use L'Hospital's rule to find the following limit.
\[ \lim_{x\to\infty}\frac{\mathrm{ln}\,(2x)+1}{5x+3}\]}
{\(0\)}{\( \frac15\)}{\( \frac25\)}{\( \frac1{10}\)}{\(\infty\)}{}}
\LANGpl{
\Question{
Użyj reguły L'Hospital, aby znaleźć następującą granicę.
\[ \lim_{x\to\infty}\frac{\mathrm{ln}\,(2x)+1}{5x+3}\]}
{\(0\)}{\( \frac15\)}{\( \frac25\)}{\( \frac1{10}\)}{\(\infty\)}{}}
\LANGcs{
\Question{
Užitím L'Hospitalova pravidla vypočítejte následující limitu.
\[ \lim_{x\to\infty}\frac{\mathrm{ln}\,(2x)+1}{5x+3}\]}
{\(0\)}{\( \frac15\)}{\( \frac25\)}{\( \frac1{10}\)}{\(\infty\)}{}}
\LANGsk{
\Question{
Použitím L'Hospitalovho pravidla vypočítajte nasledujúcu limitu.
\[ \lim_{x\to\infty}\frac{\mathrm{ln}\,(2x)+1}{5x+3}\]}
{\(0\)}{\( \frac15\)}{\( \frac25\)}{\( \frac1{10}\)}{\(\infty\)}{}}
\LANGes{
\Question{
Resuelve el siguiente límite utilizando la Regla de L'Hôpital:
\[ \lim_{x\to\infty}\frac{\mathrm{ln}\,(2x)+1}{5x+3}\]}
{\(0\)}{\( \frac15\)}{\( \frac25\)}{\( \frac1{10}\)}{\(\infty\)}{}}
\End

\Data\ID{2010011104}
\Updated{2022-08-25 10:08:46}
\Level{C}
\SubArea{Applications of derivatives}
\LANGen{
\Question{
Use L'Hospital's rule to find the following limit.
\[ \lim_{x\to\infty}\frac{2x+3}{\mathrm{e}^{2x}}\]}
{\(0\)}{\( \frac12\)}{\( \infty\)}{\( 1\)}{\(2\)}{}}
\LANGpl{
\Question{
Użyj reguły L'Hospital, aby znaleźć następującą granicę.
\[ \lim_{x\to\infty}\frac{2x+3}{\mathrm{e}^{2x}}\]}
{\(0\)}{\( \frac12\)}{\( \infty\)}{\( 1\)}{\(2\)}{}}
\LANGcs{
\Question{
Užitím L'Hospitalova pravidla vypočítejte následující limitu.
\[ \lim_{x\to\infty}\frac{2x+3}{\mathrm{e}^{2x}}\]}
{\(0\)}{\( \frac12\)}{\( \infty\)}{\( 1\)}{\(2\)}{}}
\LANGsk{
\Question{
Použitím L'Hospitalovho pravidla vypočítajte nasledujúcu limitu.
\[ \lim_{x\to\infty}\frac{2x+3}{\mathrm{e}^{2x}}\]}
{\(0\)}{\( \frac12\)}{\( \infty\)}{\( 1\)}{\(2\)}{}}
\LANGes{
\Question{
Resuelve el siguiente límite utilizando la Regla de L'Hôpital:
\[ \lim_{x\to\infty}\frac{2x+3}{\mathrm{e}^{2x}}\]}
{\(0\)}{\( \frac12\)}{\( \infty\)}{\( 1\)}{\(2\)}{}}
\End

\Data\ID{2010011105}
\Updated{2022-08-25 10:08:46}
\Level{C}
\SubArea{Applications of derivatives}
\LANGen{
\Question{
Evaluate the following limit. (Apply  L'Hospital's rule repeatedly.)
\[ \lim_{x\to0}\frac{x^4+2x^3}{x-\sin x} \]}
{\(12\)}{\( 0\)}{\( -12\)}{\( 6\)}{\(-6\)}{}}
\LANGpl{
\Question{
Użyj reguły L'Hospital, aby znaleźć następującą granicę.
\[ \lim_{x\to0}\frac{x^4+2x^3}{x-\sin x} \]}
{\(12\)}{\( 0\)}{\( -12\)}{\( 6\)}{\(-6\)}{}}
\LANGcs{
\Question{
Opakovaným užitím L'Hospitalova pravidla vypočítejte následující limitu.
\[ \lim_{x\to0}\frac{x^4+2x^3}{x-\sin x} \]}
{\(12\)}{\( 0\)}{\( -12\)}{\( 6\)}{\(-6\)}{}}
\LANGsk{
\Question{
Opakovaným použitím L'Hospitalovho pravidla vypočítajte nasledujúcu limitu.
\[ \lim_{x\to0}\frac{x^4+2x^3}{x-\sin x} \]}
{\(12\)}{\( 0\)}{\( -12\)}{\( 6\)}{\(-6\)}{}}
\LANGes{
\Question{
Resuelve el siguiente límite utilizando repetidamente la Regla de L'Hôpital:
\[ \lim_{x\to0}\frac{x^4+2x^3}{x-\sin x} \]}
{\(12\)}{\( 0\)}{\( -12\)}{\( 6\)}{\(-6\)}{}}
\End

\Data\ID{2010013701}
\Updated{2023-01-23 10:01:22}
\Level{C}
\SubArea{Applications of derivatives}
\LANGen{
\Question{
The motion of two bodies is given by equations
\[s_1=\frac12t^2+6t+1\mbox{,}\quad s_2=\frac13t^3+t^2+4,\]
where the positions \(s_1\) and \(s_2\) are given in meters and the time \(t\) in seconds. Determine at what time both bodies will move at the same speed.
\[\]
Hint: Instantaneous velocity can be expressed as the derivative of a position function \(s(t)\) with respect to time: \(v(t)=\frac{\mathrm{d} s}{\mathrm{d} t}\).}
{\(t=2\,\mathrm{s}\)}{\(t=\sqrt2\,\mathrm{s}\)}{\(t=3\,\mathrm{s}\)}{The speeds of these bodies will always be different.}{}{}}
\LANGpl{
\Question{
Ruch dwóch ciał jest opisany równaniami
\[s_1=\frac12t^2+6t+1\mbox{,}\quad s_2=\frac13t^3+t^2+4,\]
gdzie droga \(s\) podana jest w metrach, a czas  \(t\) w sekundach. Określ, w jakim czasie oba ciała będą poruszać się z tą samą prędkością.\[\]
Wskazówka: Możemy wyznaczyć prędkość korzystając z pochodnej funkcji \(s(t)\), tj. \(v(t)=\frac{\mathrm{d} s}{\mathrm{d} t}\).}
{\(t=2\,\mathrm{s}\)}{\(t=\sqrt2\,\mathrm{s}\)}{\(t=3\,\mathrm{s}\)}{Prędkości tych ciał zawsze będą inne.}{}{}}
\LANGcs{
\Question{
Pohyb dvou těles je popsán rovnicemi
\[s_1=\frac12t^2+6t+1\mbox{,}\quad s_2=\frac13t^3+t^2+4,\]
kde dráha \(s\) je uváděna v metrech a čas \(t\) v sekundách. Určete, v jakém čase se obě tělesa budou pohybovat stejnou rychlostí.\[\]
Nápověda: Rychlost můžeme určit pomocí derivace funkce \(s(t)\), tj. \(v(t)=\frac{\mathrm{d} s}{\mathrm{d} t}\).}
{\(t=2\,\mathrm{s}\)}{\(t=\sqrt2\,\mathrm{s}\)}{\(t=3\,\mathrm{s}\)}{Rychlosti těchto těles budou vždy rozdílné.}{}{}}
\LANGsk{
\Question{
Pohyb dvoch telies je popísaný rovnicami
\[s_1=\frac12t^2+6t+1\mbox{,}\quad s_2=\frac13t^3+t^2+4,\]
kde dráha \(s\) je uvedená v metroch a čas \(t\) v sekundách. Určte, v akom čase sa budú obidve telesá pohybovať rovnakou rychlosťou.\[\]
Pomôcka: Rýchlosť môžme určiť pomocou derivácie funkcie \(s(t)\), tj. \(v(t)=\frac{\mathrm{d} s}{\mathrm{d} t}\).}
{\(t=2\,\mathrm{s}\)}{\(t=\sqrt2\,\mathrm{s}\)}{\(t=3\,\mathrm{s}\)}{Rýchlosti daných telies budú vždy rozdielne.}{}{}}
\LANGes{
\Question{
El movimiento de dos cuerpos viene dado por las siguientes ecuaciones:
\[s_1=\frac12t^2+6t+1\mbox{,}\quad s_2=\frac13t^3+t^2+4,\]
donde las distancias \(s_1\) y \(s_2\) se dan en metros y el tiempo \(t\) en segundos. Determina en qué momento ambos cuerpos se moverán con la misma velocidad.
\[\]
Sugerencia: La velocidad instantánea se puede expresar como la derivada de una función de distancia \(s(t)\) con respecto al tiempo: \(v(t)=\frac{\mathrm{d} s}{\mathrm{d} t}\).}
{\(t=2\,\mathrm{s}\)}{\(t=\sqrt2\,\mathrm{s}\)}{\(t=3\,\mathrm{s}\)}{Las velocidades de ambos cuerpos siempre serán diferentes.}{}{}}
\End

\Data\ID{2010013702}
\Updated{2023-01-23 10:01:22}
\Level{C}
\SubArea{Applications of derivatives}
\LANGen{
\Question{
The motion of two bodies is given by equations
\[s_1=\frac32t^2+3t+2\mbox{,}\quad s_2=\frac13t^3+\frac{t^2}{2}+1,\]
where the positions \(s_1\) and \(s_2\) are given in meters and the time \(t\) in seconds. Determine at what time both bodies will move at the same speed.
\[\]
Hint: Instantaneous velocity can be expressed as the derivative of a position function \(s(t)\) with respect to time: \(v(t)=\frac{\mathrm{d} s}{\mathrm{d} t}\).}
{\(t=3\,\mathrm{s}\)}{\(t=1\,\mathrm{s}\)}{\(t=\sqrt7\,\mathrm{s}\)}{The speeds of these bodies will always be different.}{}{}}
\LANGpl{
\Question{
Ruch dwóch ciał jest określony równaniami
\[s_1=\frac32t^2+3t+2\mbox{,}\quad s_2=\frac13t^3+\frac{t^2}{2}+1,\]
gdzie drogi \(s_1\) i \(s_2\) podane są w metrach a czas \(t\) w sekundach. Określ, w jakim czasie oba ciała będą poruszać się z tą samą prędkością.
\[\]
Wskazówka:  Prędkość chwilową można wyrazić jako pochodną funkcji położenia \(s(t)\) względem czasu: \(v(t)=\frac{\mathrm{d} s}{\mathrm{d} t}\).}
{\(t=3\,\mathrm{s}\)}{\(t=1\,\mathrm{s}\)}{\(t=\sqrt7\,\mathrm{s}\)}{Prędkości tych ciał zawsze będą inne.}{}{}}
\LANGcs{
\Question{
Pohyb dvou těles je popsán rovnicemi
\[s_1=\frac32t^2+3t+2\mbox{,}\quad s_2=\frac13t^3+\frac{t^2}{2}+1,\]
kde dráha \(s\) je uváděna v metrech a čas \(t\) v sekundách. Určete, v jakém čase se obě tělesa budou pohybovat stejnou rychlostí.\[\]
Nápověda: Rychlost můžeme určit pomocí derivace funkce \(s(t)\), tj. \(v(t)=\frac{\mathrm{d} s}{\mathrm{d} t}\).}
{\(t=3\,\mathrm{s}\)}{\(t=1\,\mathrm{s}\)}{\(t=\sqrt7\,\mathrm{s}\)}{Rychlosti těchto těles budou vždy rozdílné.}{}{}}
\LANGsk{
\Question{
Pohyb dvoch telies je popísaný rovnicami
\[s_1=\frac32t^2+3t+2\mbox{,}\quad s_2=\frac13t^3+\frac{t^2}{2}+1,\]
kde dráha \(s\) je uvedená v metroch a čas \(t\) v sekundách. Určte, v akom čase sa budú obidve telesá pohybovať rovnakou rychlosťou.\[\]
Pomôcka: Rýchlosť môžme určiť pomocou derivácie funkcie \(s(t)\), tj. \(v(t)=\frac{\mathrm{d} s}{\mathrm{d} t}\).}
{\(t=3\,\mathrm{s}\)}{\(t=1\,\mathrm{s}\)}{\(t=\sqrt7\,\mathrm{s}\)}{Rýchlosti daných telies budú vždy rozdielne.}{}{}}
\LANGes{
\Question{
El movimiento de dos cuerpos viene dado por las siguientes ecuaciones:
\[s_1=\frac32t^2+3t+2\mbox{,}\quad s_2=\frac13t^3+\frac{t^2}{2}+1,\]
donde las distancias \(s_1\) y \(s_2\) se dan en metros y el tiempo \(t\) en segundos. Determina en qué momento ambos cuerpos se moverán con la misma velocidad.
\[\]
Sugerencia: La velocidad instantánea se puede expresar como la derivada de una función de distancia \(s(t)\) con respecto al tiempo: \(v(t)=\frac{\mathrm{d} s}{\mathrm{d} t}\).}
{\(t=3\,\mathrm{s}\)}{\(t=1\,\mathrm{s}\)}{\(t=\sqrt7\,\mathrm{s}\)}{Las velocidades de ambos cuerpos siempre serán diferentes.}{}{}}
\End

\Data\ID{2010013703}
\Updated{2023-01-23 10:01:26}
\Level{C}
\SubArea{Applications of derivatives}
\LANGen{
\Question{
Suppose \(A\), \(B\), \(C\) and \(D\) are bodies, which are set in motion at the same initial time \(t\). We know how the position \(s\) or speed \(v\) of these bodies changes with time:
\[\begin{aligned}
A: \, s=2t^2+12t+1,\qquad&C:\, v=10t+4,\\
B:\, s=\frac13t^3+\frac{t^2}{2}+2,\qquad&D:\, v=\frac12t^2+1.\end{aligned}\]
Position \(s\) is given in meters, time \(t\) in seconds and speed \(v\) in meters per second. Determine which body moves with the greatest acceleration at the time \(t=1\,\mathrm{s}\).
\[\]
Hint: Instantaneous velocity can be expressed as the derivative of a position function \(s(t)\) with respect to time: \(v(t)=\frac{\mathrm{d}s}{\mathrm{d}t}\), and the instantaneous acceleration can be expressed as the derivative of a function \(v(t)\) with respect to time: \(a(t)=\frac{\mathrm{d}v}{\mathrm{d}t}\). Since we can determine the velocity using the derivative of the position function \(s(t)\), we as well can determine the acceleration using the second derivative of \(s(t)\): \(\,a(t)=\frac{\mathrm{d}v}{\mathrm{d}t}=\frac{\mathrm{d}}{\mathrm{d}t} \left(\frac{\mathrm{d}s}{\mathrm{d}t}\right)=\frac{\mathrm{d}^2s}{\mathrm{d}t^2}\).}
{\(C\)}{\(B\)}{\(A\)}{\(D\)}{}{}}
\LANGpl{
\Question{
Załóżmy, że \(A\), \(B\), \(C\) i \(D\) są ciałami, które są wprawiane w ruch w tym samym momencie początkowym \(t\).  Wiemy, jak zmienia się droga \(s\) lub prędkość \(v\) tych ciał w czasie:
\[\begin{aligned}
A: \, s=2t^2+12t+1,\qquad&C:\, v=10t+4,\\
B:\, s=\frac13t^3+\frac{t^2}{2}+2,\qquad&D:\, v=\frac12t^2+1.\end{aligned}\]
Pozycja \(s\) podana jest w metrach, czas \(t\) w sekundach i prędkość \(v\) w metrach na sekundę.  Określ, które ciało porusza się z największym przyspieszeniem w danym momencie \(t=1\,\mathrm{s}\).
\[\]
Wskazówka:  Prędkość chwilową można wyrazić jako pochodną funkcji położenia \(s(t)\) względem czasu: \(v(t)=\frac{\mathrm{d}s}{\mathrm{d}t}\),  a przyspieszenie chwilowe można wyrazić jako pochodną funkcji \(v(t)\) względem czasu: \(a(t)=\frac{\mathrm{d}v}{\mathrm{d}t}\). Ponieważ możemy wyznaczyć prędkość za pomocą pochodnej funkcji położenia \(s(t)\), możemy również wyznaczyć przyspieszenie za pomocą drugiej pochodnej \(s(t)\): \(\,a(t)=\frac{\mathrm{d}v}{\mathrm{d}t}=\frac{\mathrm{d}}{\mathrm{d}t} \left(\frac{\mathrm{d}s}{\mathrm{d}t}\right)=\frac{\mathrm{d}^2s}{\mathrm{d}t^2}\).}
{\(C\)}{\(B\)}{\(A\)}{\(D\)}{}{}}
\LANGcs{
\Question{
Máme tělesa \(A\), \(B\), \(C\) a \(D\), která se dají současně do pohybu. Víme, jak se mění dráha, či rychlost těchto těles v závislosti na čase:
\[\begin{aligned}
A: \, s=2t^2+12t+1,\qquad&C:\, v=10t+4,\\
B:\, s=\frac13t^3+\frac{t^2}{2}+2,\qquad&D:\, v=\frac12t^2+1.\end{aligned}\]
Dráha \(s\) je měřena v metrech, čas \(t\) v sekundách a rychlost \(v\) v metrech za sekundu. Určete, které těleso se v čase \(t=1\,\mathrm{s}\) pohybuje s největším zrychlením.
\[\]
Nápověda: Rychlost \(v(t)\) můžeme určit pomocí derivace funkce \(s(t)\), tj. \(v(t)=\frac{\mathrm{d}s}{\mathrm{d}t}\). Zrychlení \(a(t)\) lze určit jako derivaci funkce \(v(t)\), tj. \(a(t)=\frac{\mathrm{d}v}{\mathrm{d}t}\). Protože rychlost můžeme určit pomocí derivace funkce \(s(t)\), můžeme zrychlení určit pomocí druhé derivace \(s(t)\), tj. \(\,a(t)=\frac{\mathrm{d}v}{\mathrm{d}t}=\frac{\mathrm{d}}{\mathrm{d}t} \left(\frac{\mathrm{d}s}{\mathrm{d}t}\right)=\frac{\mathrm{d}^2s}{\mathrm{d}t^2}\).}
{\(C\)}{\(B\)}{\(A\)}{\(D\)}{}{}}
\LANGsk{
\Question{
Máme telesá \(A\), \(B\), \(C\) a \(D\), ktoré sa dajú do pohybu súčasne. Vieme, ako sa mení dráha či rýchlosť daných telies v závislosti na čase:
\[\begin{aligned}
A: \, s=2t^2+12t+1,\qquad&C:\, v=10t+4,\\
B:\, s=\frac13t^3+\frac{t^2}{2}+2,\qquad&D:\, v=\frac12t^2+1.\end{aligned}\]
Dráha \(s\) je meraná v metroch, čas \(t\) v sekundách a rýchlosť \(v\) v metroch za sekundu. Určte, ktoré teleso sa v čase \(t=1\,\mathrm{s}\) pohybuje s najväčším zrychlením.
\[\]
Pomôcka: Rýchlosť \(v(t)\) môžme určiť pomocou derivácie funkcie \(s(t)\), tj. \(v(t)=\frac{\mathrm{d}s}{\mathrm{d}t}\). Zrýchlenie \(a(t)\) vieme určiť ako deriváciu funkcie \(v(t)\), tj. \(a(t)=\frac{\mathrm{d}v}{\mathrm{d}t}\). Protože rýchlosť môžme určiť pomocou derivácie funkcie \(s(t)\), môžme zrýchlenie určiť pomocou druhej derivácie \(s(t)\), tj. \(\,a(t)=\frac{\mathrm{d}v}{\mathrm{d}t}=\frac{\mathrm{d}}{\mathrm{d}t} \left(\frac{\mathrm{d}s}{\mathrm{d}t}\right)=\frac{\mathrm{d}^2s}{\mathrm{d}t^2}\).}
{\(C\)}{\(B\)}{\(A\)}{\(D\)}{}{}}
\LANGes{
\Question{
Suponemos que \(A\), \(B\), \(C\) y \(D\) son cuerpos, que se ponen en movimiento al mismo tiempo inicial \(t\). Sabemos cómo se cambia la distancia \(s\) o la velocidad \(v\) de estos cuerpos con el tiempo:
\[\begin{aligned}
A: \, s=2t^2+12t+1,\qquad&C:\, v=10t+4,\\
B:\, s=\frac13t^3+\frac{t^2}{2}+2,\qquad&D:\, v=\frac12t^2+1.\end{aligned}\]
La distancia \(s\) se da en metros, el tiempo \(t\) en segundos y la velocidad \(v\) en metros por segundo. Determina cuál de los cuerpos se mueve con la mayor aceleración en el momento \(t=1\,\mathrm{s}\).
\[\]
Sugerencia: La velocidad instantánea se puede expresar como la derivada de una función de distancia \(s(t)\) con respecto al tiempo: \(v(t)=\frac{\mathrm{d}s}{\mathrm{d}t}\), y la aceleración instantánea se puede expresar como la derivada de una función \(v(t)\) con respecto al tiempo: \(a(t)=\frac{\mathrm{d}v}{\mathrm{d}t}\). Como podemos determinar la velocidad usando la derivada de la función de distancia \(s(t)\), también podemos determinar la aceleración usando la segunda derivada de \(s(t)\): \(\,a(t)=\frac{\mathrm{d}v}{\mathrm{d}t}=\frac{\mathrm{d}}{\mathrm{d}t} \left(\frac{\mathrm{d}s}{\mathrm{d}t}\right)=\frac{\mathrm{d}^2s}{\mathrm{d}t^2}\).}
{\(C\)}{\(B\)}{\(A\)}{\(D\)}{}{}}
\End

\Data\ID{2010013704}
\Updated{2023-01-23 10:01:26}
\Level{C}
\SubArea{Applications of derivatives}
\LANGen{
\Question{
Suppose \(A\), \(B\), \(C\) and \(D\) are bodies, which are set in motion at the same initial time \(t\). We know how the position \(s\) or speed \(v\) of these bodies changes with time:
\[\begin{aligned}
A: \, s=\frac12t^2+10t+1,\qquad&C:\, v=9t+15,\\
B:\, s=\frac13t^3+t^2+4,\qquad\ \ &D:\, v=\frac52t^2+3.\end{aligned}\]
Position \(s\) is given in meters, time \(t\) in seconds and speed \(v\) in meters per second. Determine which body moves with the greatest acceleration at the time \(t=1\,\mathrm{s}\).
\[\]
Hint: Instantaneous velocity can be expressed as the derivative of a position function \(s(t)\) with respect to time: \(v(t)=\frac{\mathrm{d}s}{\mathrm{d}t}\), and the instantaneous acceleration can be expressed as the derivative of a function \(v(t)\) with respect to time: \(a(t)=\frac{\mathrm{d}v}{\mathrm{d}t}\). Since we can determine the velocity using the derivative of the position function \(s(t)\), we as well can determine the acceleration using the second derivative of \(s(t)\): \(\,a(t)=\frac{\mathrm{d}v}{\mathrm{d}t}=\frac{\mathrm{d}}{\mathrm{d}t}\cdot\frac{\mathrm{d}s}{\mathrm{d}t}=\frac{\mathrm{d}^2s}{\mathrm{d}t^2}\).}
{\(C\)}{\(B\)}{\(A\)}{\(D\)}{}{}}
\LANGpl{
\Question{
Mamy ciała \(A\), \(B\), \(C\) i \(D\), które można jednocześnie wprawić w ruch. Wiemy jak zmienia się droga \(s\) lub  prędkość danych ciał \(v\) w zależności od czasu. 
\[\begin{aligned}
A: \, s=\frac12t^2+10t+1,\qquad&C:\, v=9t+15,\\
B:\, s=\frac13t^3+t^2+4,\qquad\ \ &D:\, v=\frac52t^2+3.\end{aligned}\]
Droga \(s\) mierzona jest w metrach, czas  \(t\) w sekundach a prędkość \(v\) w metrach na sekundę. Określ, które ciało porusza się z największym przyspieszeniem w czasie \(t=1\,\mathrm{s}\).
\[\]
 Wskazówka: Możemy wyznaczyć prędkość korzystając z pochodnej funkcji \(s(t)\) względem czasu: \(v(t)=\frac{\mathrm{d}s}{\mathrm{d}t}\), a przyspieszenie chwilowe można wyrazić jako pochodną funkcji \(v(t)\) względem czasu: \(a(t)=\frac{\mathrm{d}v}{\mathrm{d}t}\). Ponieważ możemy wyznaczyć prędkość za pomocą pochodnej funkcji położenia \(s(t)\),  możemy również wyznaczyć przyspieszenie za pomocą drugiej pochodnej \(s(t)\): \(\,a(t)=\frac{\mathrm{d}v}{\mathrm{d}t}=\frac{\mathrm{d}}{\mathrm{d}t}\cdot\frac{\mathrm{d}s}{\mathrm{d}t}=\frac{\mathrm{d}^2s}{\mathrm{d}t^2}\).}
{\(C\)}{\(B\)}{\(A\)}{\(D\)}{}{}}
\LANGcs{
\Question{
Máme tělesa \(A\), \(B\), \(C\) a \(D\), která se dají současně do pohybu. Víme, jak se mění dráha, či rychlost těchto těles v závislosti na čase:
\[\begin{aligned}
A: \, s=\frac12t^2+10t+1,\qquad&C:\, v=9t+15,\\
B:\, s=\frac13t^3+t^2+4,\qquad\ \ &D:\, v=\frac52t^2+3.\end{aligned}\]
Dráha \(s\) je měřena v metrech, čas \(t\) v sekundách a rychlost \(v\) v metrech za sekundu. Určete, které těleso se v čase \(t=1\,\mathrm{s}\) pohybuje s největším zrychlením.
\[\]
Nápověda: Rychlost \(v(t)\) můžeme určit pomocí derivace funkce \(s(t)\), tj. \(v(t)=\frac{\mathrm{d}s}{\mathrm{d}t}\). Zrychlení \(a(t)\) lze určit jako derivaci funkce \(v(t)\), tj. \(a(t)=\frac{\mathrm{d}v}{\mathrm{d}t}\). Protože rychlost můžeme určit pomocí derivace funkce \(s(t)\), můžeme zrychlení určit pomocí druhé derivace \(s(t)\), tj. \(\,a(t)=\frac{\mathrm{d}v}{\mathrm{d}t}=\frac{\mathrm{d}}{\mathrm{d}t}\cdot\frac{\mathrm{d}s}{\mathrm{d}t}=\frac{\mathrm{d}^2s}{\mathrm{d}t^2}\).}
{\(C\)}{\(B\)}{\(A\)}{\(D\)}{}{}}
\LANGsk{
\Question{
Máme telesá \(A\), \(B\), \(C\) a \(D\), ktoré sa dajú do pohybu súčasne. Vieme, ako sa mení dráha či rýchlosť daných telies v závislosti na čase:
\[\begin{aligned}
A: \, s=\frac12t^2+10t+1,\qquad&C:\, v=9t+15,\\
B:\, s=\frac13t^3+t^2+4,\qquad\ \ &D:\, v=\frac52t^2+3.\end{aligned}\]
Dráha \(s\) je meraná v metroch, čas \(t\) v sekundách a rýchlosť \(v\) v metroch za sekundu. Určte, ktoré teleso sa v čase \(t=1\,\mathrm{s}\) pohybuje s najväčším zrychlením.
\[\]
Pomôcka: Rýchlosť \(v(t)\) môžme určiť pomocou derivácie funkcie \(s(t)\), tj. \(v(t)=\frac{\mathrm{d}s}{\mathrm{d}t}\). Zrýchlenie \(a(t)\) vieme určiť ako deriváciu funkcie \(v(t)\), tj. \(a(t)=\frac{\mathrm{d}v}{\mathrm{d}t}\). Protože rýchlosť môžme určiť pomocou derivácie funkcie \(s(t)\), môžme zrýchlenie určiť pomocou druhej derivácie \(s(t)\), tj. \(\,a(t)=\frac{\mathrm{d}v}{\mathrm{d}t}=\frac{\mathrm{d}}{\mathrm{d}t}\cdot\frac{\mathrm{d}s}{\mathrm{d}t}=\frac{\mathrm{d}^2s}{\mathrm{d}t^2}\).}
{\(C\)}{\(B\)}{\(A\)}{\(D\)}{}{}}
\LANGes{
\Question{
Suponemos que \(A\), \(B\), \(C\) y \(D\) son cuerpos, que se ponen en movimiento al mismo tiempo inicial \(t\). Sabemos cómo se cambia la distancia  \(s\) o la velocidad \(v\) de estos cuerpos con el tiempo:
\[\begin{aligned}
A: \, s=\frac12t^2+10t+1,\qquad&C:\, v=9t+15,\\
B:\, s=\frac13t^3+t^2+4,\qquad\ \ &D:\, v=\frac52t^2+3.\end{aligned}\]
La distancia\(s\) se da en metros, el tiempo \(t\) en segundos y la velocidad \(v\) en metros por segundo. Determina cuál de los cuerpos se mueve con la mayor aceleración en el momento \(t=1\,\mathrm{s}\).
\[\]
Sugerencia: La velocidad instantánea se puede expresar como la derivada de una función de distancia \(s(t)\) con respecto al tiempo: \(v(t)=\frac{\mathrm{d}s}{\mathrm{d}t}\), y la aceleración instantánea se puede expresar como la derivada de una función \(v(t)\) con respecto al tiempo: \(a(t)=\frac{\mathrm{d}v}{\mathrm{d}t}\). Como podemos determinar la velocidad usando la derivada de la función de distancia \(s(t)\), también podemos determinar la aceleración usando la segunda derivada de \(s(t)\): \(\,a(t)=\frac{\mathrm{d}v}{\mathrm{d}t}=\frac{\mathrm{d}}{\mathrm{d}t}\cdot\frac{\mathrm{d}s}{\mathrm{d}t}=\frac{\mathrm{d}^2s}{\mathrm{d}t^2}\).}
{\(C\)}{\(B\)}{\(A\)}{\(D\)}{}{}}
\End

\Data\ID{2010013705}
\Updated{2023-01-23 10:01:26}
\Level{C}
\SubArea{Applications of derivatives}
\LANGen{
\Question{
An electrical source is characterized by the electromotive force \(U_e=60\,\mathrm{V}\)
and the internal resistance \(R_i=2\,\Omega\). Determine the value of the electric current
for which the appliance power will be at its maximum and determine the value of this maximum power as well.
\[\]
Hint: The dependence of the power of an appliance (\(P\), unit Watt (\(\mathrm{W}\)))
on the magnitude of the folowing current (\(I\), unit Ampere (\(\mathrm{A}\))) is given by the relation
\(P=U_eI-R_iI^2\). The source properties have a role of parameters: \(U_e\) is the electromotive 
force and \(R_i\) is the internal resistance of the source.}
{\(15\,\mathrm{A},\ 450\,\mathrm{W}\)}{\(15\,\mathrm{A},\ 870\,\mathrm{W}\)}{\(30\,\mathrm{A},\ 1740\,\mathrm{W}\)}{\(10\,\mathrm{A},\ 400\,\mathrm{W}\)}{}{}}
\LANGpl{
\Question{
Źródło elektryczne charakteryzuje się napięciem elektromotorycznym  \(U_e=60\,\mathrm{V}\)
i rezystancją wewnętrzną \(R_i=2\,\Omega\). Określ wartość prądu elektrycznego, przy której będzie maksymalna moc w urządzeniu, a także odpowiednią wartość tej maksymalnej mocy.
\[\]
Wskazówka: Moc odbiornika  (\(P\), jednostka Watt (\(\mathrm{W}\)))
zależy od wielkości przepływającego prądu (\(I\), jednostka Amper (\(\mathrm{A}\)))  według wzoru 
\(P=U_eI-R_iI^2\). Własności źródła pełnią rolę parametrów: \(U_e\) jest to napięcie elektromotoryczne i  \(R_i\) wewnętrzna rezystancja źródła.}
{\(15\,\mathrm{A},\ 450\,\mathrm{W}\)}{\(15\,\mathrm{A},\ 870\,\mathrm{W}\)}{\(30\,\mathrm{A},\ 1740\,\mathrm{W}\)}{\(10\,\mathrm{A},\ 400\,\mathrm{W}\)}{}{}}
\LANGcs{
\Question{
Elektrický zdroj je charakterizován elektromotorickým napětím \(U_e=60\,\mathrm{V}\) a vnitřním odporem \(R_i=2\,\Omega\).  Určete, při jaké hodnotě elektrického proudu bude ve spotřebiči maximální výkon a příslušnou hodnotu maximálního výkonu.
\[\]
Nápověda: Výkon spotřebiče (\(P\), jednotka Watt (\(\mathrm{W}\))), závisí na velikosti protékajícího proudu (\(I\), jednotka Ampér (\(\mathrm{A}\))) vztahem \(P=U_eI-R_iI^2\). Vlastnosti zdroje mají úlohu parametrů: \(U_e\) je elektromotorické napětí a \(R_i\) vnitřní odpor zdroje.}
{\(15\,\mathrm{A},\ 450\,\mathrm{W}\)}{\(15\,\mathrm{A},\ 870\,\mathrm{W}\)}{\(30\,\mathrm{A},\ 1740\,\mathrm{W}\)}{\(10\,\mathrm{A},\ 400\,\mathrm{W}\)}{}{}}
\LANGsk{
\Question{
Elektrický zdroj je charakterizovaný elektromotorickým napätím \(U_e=60\,\mathrm{V}\) a vnútorným odporom \(R_i=2\,\Omega\).  Určte, pri akej hodnote elektrického prúdu bude v spotrebiči maximálny výkon a tiež príslušnú hodnotu maximálneho výkonu.
\[\]
Pomôcka: Výkon spotrebiča (\(P\), jednotka Watt (\(\mathrm{W}\))), závisí na veľkosti pretekajúceho prúdu (\(I\), jednotka Ampér (\(\mathrm{A}\))) vzťahom \(P=U_eI-R_iI^2\). Vlastnosti zdroja majú úlohu parametrov: \(U_e\) je elektromotorické napätie a \(R_i\) vnútorného odporu zdroja.}
{\(15\,\mathrm{A},\ 450\,\mathrm{W}\)}{\(15\,\mathrm{A},\ 870\,\mathrm{W}\)}{\(30\,\mathrm{A},\ 1740\,\mathrm{W}\)}{\(10\,\mathrm{A},\ 400\,\mathrm{W}\)}{}{}}
\LANGes{
\Question{
Una fuente eléctrica se caracteriza por la fuerza electromotriz \(U_e=60\,\mathrm{V}\)
y la resistencia interna \(R_i=2\,\Omega\). Determina el valor de la corriente eléctrica para la cual la potencia del aparato será máxima y también determina el valor de esta potencia máxima.
\[\]
Sugerencia: La dependencia de la potencia de un aparato (\(P\), unidad Watt (\(\mathrm{W}\)))
sobre la magnitud de la siguiente corriente (\(I\), unidad Amperio (\(\mathrm{A}\))) viene dada por la relación
\(P=U_eI-R_iI^2\). Las propiedades de la fuente tienen un papel de parámetros: \(U_e\) es la fuerza electromotriz, \(R_i\) es la resistencia interna de la fuente.}
{\(15\,\mathrm{A},\ 450\,\mathrm{W}\)}{\(15\,\mathrm{A},\ 870\,\mathrm{W}\)}{\(30\,\mathrm{A},\ 1740\,\mathrm{W}\)}{\(10\,\mathrm{A},\ 400\,\mathrm{W}\)}{}{}}
\End

\Data\ID{2010013706}
\Updated{2023-01-23 10:01:26}
\Level{C}
\SubArea{Applications of derivatives}
\LANGen{
\Question{
An electrical source is characterized by the electromotive force \(U_e=40\,\mathrm{V}\)
and the internal resistance \(R_i=2\,\Omega\). Determine the value of the electric current
for which the appliance power will be at its maximum and determine the value of this maximum power as well.
\[\]
Hint: The dependence of the power of an appliance (\(P\), unit Watt (\(\mathrm{W}\)))
on the magnitude of the folowing current (\(I\), unit Ampere (\(\mathrm{A}\))) is given by the relation
\(P=U_eI-R_iI^2\). The source properties have a role of parameters: \(U_e\) is the electromotive 
force and \(R_i\) is the internal resistance of the source.}
{\(10\,\mathrm{A},\ 200\,\mathrm{W}\)}{\(10\,\mathrm{A},\ 380\,\mathrm{W}\)}{\(20\,\mathrm{A},\ 760\,\mathrm{W}\)}{\(4\,\mathrm{A},\ 128\,\mathrm{W}\)}{}{}}
\LANGpl{
\Question{
Źródło elektryczne charakteryzuje napięcie elektromotoryczne \(U_e=40\,\mathrm{V}\), 
oraz opór wewnętrzny \(R_i=2\,\Omega\). Określ wartość prądu elektrycznego
dla której moc urządzenia będzie maksymalna i określ wartość tej maksymalnej mocy.
\[\]
Wskazówka: Moc urządzenia (\(P\), jednostka Wat (\(\mathrm{W}\))) zależy od wielkości przepływającego prądu (\(I\), jednostka Amper (\(\mathrm{A}\)))  według zależności 
\(P=U_eI-R_iI^2\). Własności źródła pełnią rolę parametrów:  \(U_e\) to napięcie elektromotoryczne, a \(R_i\) jest wewnętrznym oporem źródła.}
{\(10\,\mathrm{A},\ 200\,\mathrm{W}\)}{\(10\,\mathrm{A},\ 380\,\mathrm{W}\)}{\(20\,\mathrm{A},\ 760\,\mathrm{W}\)}{\(4\,\mathrm{A},\ 128\,\mathrm{W}\)}{}{}}
\LANGcs{
\Question{
Elektrický zdroj je charakterizován elektromotorickým napětím \(U_e=40\,\mathrm{V}\) a vnitřním odporem \(R_i=2\,\Omega\).  Určete, při jaké hodnotě elektrického proudu bude ve spotřebiči maximální výkon a příslušnou hodnotu maximálního výkonu.
\[\]
Nápověda: Výkon spotřebiče (\(P\), jednotka Watt (\(\mathrm{W}\))), závisí na velikosti protékajícího proudu (\(I\), jednotka Ampér (\(\mathrm{A}\))) vztahem \(P=U_eI-R_iI^2\). Vlastnosti zdroje mají úlohu parametrů: \(U_e\) je elektromotorické napětí a \(R_i\) vnitřní odpor zdroje.}
{\(10\,\mathrm{A},\ 200\,\mathrm{W}\)}{\(10\,\mathrm{A},\ 380\,\mathrm{W}\)}{\(20\,\mathrm{A},\ 760\,\mathrm{W}\)}{\(4\,\mathrm{A},\ 128\,\mathrm{W}\)}{}{}}
\LANGsk{
\Question{
Elektrický zdroj je charakterizovaný elektromotorickým napätím \(U_e=40\,\mathrm{V}\) a vnútorným odporom \(R_i=2\,\Omega\).  Určte, pri akej hodnote elektrického prúdu bude v spotrebiči maximálny výkon a tiež príslušnú hodnotu maximálneho výkonu.
\[\]
Pomôcka: Výkon spotrebiča (\(P\), jednotka Watt (\(\mathrm{W}\))), závisí na veľkosti pretekajúceho prúdu (\(I\), jednotka Ampér (\(\mathrm{A}\))) vzťahom \(P=U_eI-R_iI^2\). Vlastnosti zdroja majú úlohu parametrov: \(U_e\) je elektromotorické napätie a \(R_i\) vnútorného odporu zdroja.}
{\(10\,\mathrm{A},\ 200\,\mathrm{W}\)}{\(10\,\mathrm{A},\ 380\,\mathrm{W}\)}{\(20\,\mathrm{A},\ 760\,\mathrm{W}\)}{\(4\,\mathrm{A},\ 128\,\mathrm{W}\)}{}{}}
\LANGes{
\Question{
Una fuente eléctrica se caracteriza por la fuerza electromotriz \(U_e=40\,\mathrm{V}\)
y la resistencia interna \(R_i=2\,\Omega\). Determina el valor de la corriente eléctrica para la cual la potencia del aparato será máxima y también determina el valor de esta potencia máxima.
\[\]
Sugerencia: La dependencia de la potencia de un aparato (\(P\), unidad Watt (\(\mathrm{W}\)))
sobre la magnitud de la siguiente corriente  (\(I\), unidad Amperio (\(\mathrm{A}\))) viene dada por la relación
\(P=U_eI-R_iI^2\). Las propiedades de la fuente tienen un papel de parámetros: \(U_e\) es la fuerza electromotriz, \(R_i\) es la resistencia interna de la fuente.}
{\(10\,\mathrm{A},\ 200\,\mathrm{W}\)}{\(10\,\mathrm{A},\ 380\,\mathrm{W}\)}{\(20\,\mathrm{A},\ 760\,\mathrm{W}\)}{\(4\,\mathrm{A},\ 128\,\mathrm{W}\)}{}{}}
\End

\Data\ID{2010013707}
\Updated{2023-01-23 10:01:26}
\Level{C}
\SubArea{Applications of derivatives}
\LANGen{
\Question{
Suppose we throw an object vertically upwards at the initial speed \(v_0=60\,\mathrm{m}/\mathrm{s}\). Determine the time needed for the object to reach the maximum height and determine the corresponding maximum height as well.
\[\]
Hint: The vertical upwards motion of a body is the movement composed of uniformly rectilinear motion (vertically upwards) and free fall. The dependence of the instantaneous height of a body on the time is given by the relation \(h=v_0t-\frac12gt^2\), where \(v_0\) is the magnitude of the initial velocity and \(g\) is the gravitational acceleration. In this problem, calculate with the rounded value of \(g=10\,\frac{\mathrm{m}}{\mathrm{s}^2}\). We measure time \(t\) in second and height \(h\) in meters.}
{\(6\,\mathrm{s}\), \(180\,\mathrm{m}\)}{\(6\,\mathrm{s}\), \(330\,\mathrm{m}\)}{\(12\,\mathrm{s}\), \(660\,\mathrm{m}\)}{\(3\,\mathrm{s}\), \(135\,\mathrm{m}\)}{}{}}
\LANGpl{
\Question{
Załóżmy, że rzucamy obiektem pionowo w górę z prędkością początkową \(v_0=60\,\mathrm{m}/\mathrm{s}\). Określ czas potrzebny do osiągnięcia przez obiekt maksymalnej wysokości i określ również odpowiednią wysokość maksymalną.
\[\]
Wskazówka: Pionowy ruch ciała w górę to ruch składający się z ruchu jednostajnie prostoliniowego (pionowo w górę) i swobodnego spadania. Zależność chwilowej wysokości ciała od czasu wyraża wzór \(h=v_0t-\frac12gt^2\), gdzie \(v_0\) jest wielkością prędkości początkowej, a \(g\) to przyspieszenie grawitacyjne. W tym zadaniu przyjmij zaokrągloną wartość \(g=10\,\frac{\mathrm{m}}{\mathrm{s}^2}\). Mierzymy czas \(t\) w sekundach i wysokość \(h\) w metrach.}
{\(6\,\mathrm{s}\), \(180\,\mathrm{m}\)}{\(6\,\mathrm{s}\), \(330\,\mathrm{m}\)}{\(12\,\mathrm{s}\), \(660\,\mathrm{m}\)}{\(3\,\mathrm{s}\), \(135\,\mathrm{m}\)}{}{}}
\LANGcs{
\Question{
Těleso vystřelíme svisle vzhůru počáteční rychlostí \(v_0=60\,\mathrm{m}/\mathrm{s}\). Určete dobu, za kterou těleso vystoupá do maximální výšky a příslušnou maximální dosaženou výšku.
\[\]
Nápověda: Vrh svislý vzhůru je pohyb složený z pohybu rovnoměrně přímočarého (svisle vzhůru rychlostí \(v_0\)) a volného pádu. Okamžitá výška tělesa (\(h\)) závisí na čase (\(t\)) vztahem \(h=v_0t-\frac12gt^2\), kde \(v_0\) je velikost počáteční rychlosti a \(g\) tíhové zrychlení. V této úloze počítejte se zaokrouhlenou hodnotou \(g=10\,\frac{\mathrm{m}}{\mathrm{s}^2}\).  Čas \(t\) měříme v sekundách, výšku \(h\) měříme v metrech.}
{\(6\,\mathrm{s}\), \(180\,\mathrm{m}\)}{\(6\,\mathrm{s}\), \(330\,\mathrm{m}\)}{\(12\,\mathrm{s}\), \(660\,\mathrm{m}\)}{\(3\,\mathrm{s}\), \(135\,\mathrm{m}\)}{}{}}
\LANGsk{
\Question{
Teleso vystrelíme zvislo nahor začiatočnou rýchlosťou \(v_0=60\,\mathrm{m}/\mathrm{s}\). Určte čas, za ktorý teleso vystúpi do maximálnej výšky a tiež príslušnú maximálnu dosiahnutú výšku.
\[\]
Pomôcka: Zvislý vrh nahor je pohyb zložený z pohybu rovnomerne priamočiareho (zvislo nahor rýchlosťou \(v_0\)) a voľného pádu. Okamžitá výška telesa (\(h\)) závisí na čase (\(t\)) vzťahom \(h=v_0t-\frac12gt^2\), kde \(v_0\) je veľkosť začiatočnej rýchlosti a \(g\) tiažové zrýchlenie. V tejto úlohe počítajte so zaokrúhlenou hodnotou \(g=10\,\frac{\mathrm{m}}{\mathrm{s}^2}\).  Čas \(t\) meriame v sekundách, výšku \(h\) meriame v metroch.}
{\(6\,\mathrm{s}\), \(180\,\mathrm{m}\)}{\(6\,\mathrm{s}\), \(330\,\mathrm{m}\)}{\(12\,\mathrm{s}\), \(660\,\mathrm{m}\)}{\(3\,\mathrm{s}\), \(135\,\mathrm{m}\)}{}{}}
\LANGes{
\Question{
Suponemos que lanzamos un objeto verticalmente hacia arriba con la velocidad inicial \(v_0=60\,\mathrm{m}/\mathrm{s}\). Determina el tiempo necesario para que el objeto alcance la altura máxima y también determina la altura máxima correspondiente.
\[\]
Sugerencia: El movimiento vertical hacia arriba de un cuerpo es el movimiento compuesto por un movimiento uniformemente rectilíneo y caída libre. La dependencia de la altura instantánea de un cuerpo con el tiempo viene dada por la relación \(h=v_0t-\frac12gt^2\), donde \(v_0\) es la magnitud de la velocidad inicial y \(g\) es la aceleración de la gravedad. Calcula este problema con el valor redondeado de \(g=10\,\frac{\mathrm{m}}{\mathrm{s}^2}\). El tiempo \(t\) lo medimos en segundos y la altura \(h\) en metros.}
{\(6\,\mathrm{s}\), \(180\,\mathrm{m}\)}{\(6\,\mathrm{s}\), \(330\,\mathrm{m}\)}{\(12\,\mathrm{s}\), \(660\,\mathrm{m}\)}{\(3\,\mathrm{s}\), \(135\,\mathrm{m}\)}{}{}}
\End

\Data\ID{2010013708}
\Updated{2023-01-23 10:01:26}
\Level{C}
\SubArea{Applications of derivatives}
\LANGen{
\Question{
Suppose we throw an object vertically upwards at the initial speed \(v_0=80\,\mathrm{m}/\mathrm{s}\). Determine the time needed for the object to reach the maximum height and determine the corresponding maximum height as well.
\[\]
Hint: The vertical upwards motion of a body is the movement composed of uniformly rectilinear motion (vertically upwards) and free fall. The dependence of the instantaneous height of a body on the time is given by the relation \(h=v_0t-\frac12gt^2\), where \(v_0\) is the magnitude of the initial velocity and \(g\) is the gravitational acceleration. In this problem, calculate with the rounded value of \(g=10\,\frac{\mathrm{m}}{\mathrm{s}^2}\). We measure time \(t\) in second and height \(h\) in meters.}
{\(8\,\mathrm{s}\), \(320\,\mathrm{m}\)}{\(8\,\mathrm{s}\), \(600\,\mathrm{m}\)}{\(16\,\mathrm{s}\), \(1190\,\mathrm{m}\)}{\(4\,\mathrm{s}\), \(230\,\mathrm{m}\)}{}{}}
\LANGpl{
\Question{
Załóżmy, że rzucamy obiektem pionowo w górę z prędkością początkową \(v_0=80\,\mathrm{m}/\mathrm{s}\). Określ czas potrzebny do osiągnięcia przez obiekt maksymalnej wysokości i określ również odpowiednią wysokość maksymalną.
\[\]
Wskazówka: Pionowy ruch ciała w górę to ruch składający się z ruchu jednostajnie prostoliniowego (pionowo w górę) i swobodnego spadania. Zależność chwilowej wysokości ciała od czasu wyraża wzór \(h=v_0t-\frac12gt^2\), gdzie \(v_0\) jest wielkością prędkości początkowej, a \(g\) to przyspieszenie grawitacyjne. W tym zadaniu przyjmij zaokrągloną wartość  \(g=10\,\frac{\mathrm{m}}{\mathrm{s}^2}\). Mierzymy czas \(t\) w sekundach i wysokość \(h\) w metrach.}
{\(8\,\mathrm{s}\), \(320\,\mathrm{m}\)}{\(8\,\mathrm{s}\), \(600\,\mathrm{m}\)}{\(16\,\mathrm{s}\), \(1190\,\mathrm{m}\)}{\(4\,\mathrm{s}\), \(230\,\mathrm{m}\)}{}{}}
\LANGcs{
\Question{
Těleso vystřelíme svisle vzhůru počáteční rychlostí \(v_0=80\,\mathrm{m}/\mathrm{s}\). Určete dobu, za kterou těleso vystoupá do maximální výšky a příslušnou maximální dosaženou výšku.
\[\]
Nápověda: Vrh svislý vzhůru je pohyb složený z pohybu rovnoměrně přímočarého (svisle vzhůru rychlostí \(v_0\)) a volného pádu. Okamžitá výška tělesa (\(h\)) závisí na čase (\(t\)) vztahem \(h=v_0t-\frac12gt^2\), kde \(v_0\) je velikost počáteční rychlosti a \(g\) tíhové zrychlení. V této úloze počítejte se zaokrouhlenou hodnotou \(g=10\,\frac{\mathrm{m}}{\mathrm{s}^2}\).  Čas \(t\) měříme v sekundách, výšku \(h\) měříme v metrech.}
{\(8\,\mathrm{s}\), \(320\,\mathrm{m}\)}{\(8\,\mathrm{s}\), \(600\,\mathrm{m}\)}{\(16\,\mathrm{s}\), \(1190\,\mathrm{m}\)}{\(4\,\mathrm{s}\), \(230\,\mathrm{m}\)}{}{}}
\LANGsk{
\Question{
Teleso vystrelíme zvislo nahor začiatočnou rýchlosťou \(v_0=80\,\mathrm{m}/\mathrm{s}\). Určte čas, za ktorý teleso vystúpi do maximálnej výšky a tiež príslušnú maximálnu dosiahnutú výšku.
\[\]
Pomôcka: Zvislý vrh nahor je pohyb zložený z pohybu rovnomerne priamočiareho (zvislo nahor rýchlosťou \(v_0\)) a voľného pádu. Okamžitá výška telesa (\(h\)) závisí na čase (\(t\)) vzťahom \(h=v_0t-\frac12gt^2\), kde \(v_0\) je veľkosť začiatočnej rýchlosti a \(g\) tiažové zrýchlenie. V tejto úlohe počítajte so zaokrúhlenou hodnotou \(g=10\,\frac{\mathrm{m}}{\mathrm{s}^2}\).  Čas \(t\) meriame v sekundách, výšku \(h\) meriame v metroch.}
{\(8\,\mathrm{s}\), \(320\,\mathrm{m}\)}{\(8\,\mathrm{s}\), \(600\,\mathrm{m}\)}{\(16\,\mathrm{s}\), \(1190\,\mathrm{m}\)}{\(4\,\mathrm{s}\), \(230\,\mathrm{m}\)}{}{}}
\LANGes{
\Question{
Suponemos que lanzamos un objeto verticalmente hacia arriba con la velocidad inicial \(v_0=80\,\mathrm{m}/\mathrm{s}\). Determina el tiempo necesario para que el objeto alcance la altura máxima y también determina la altura máxima correspondiente.
\[\]
Sugerencia: El movimiento vertical hacia arriba de un cuerpo es el movimiento compuesto por un movimiento uniformemente rectilíneo y caída libre. La dependencia de la altura instantánea de un cuerpo con el tiempo viene dada por la relación \(h=v_0t-\frac12gt^2\), donde \(v_0\) es la magnitud de la velocidad inicial y \(g\) es la aceleración de la gravedad. Calcula este problema con el valor redondeado de \(g=10\,\frac{\mathrm{m}}{\mathrm{s}^2}\). El tiempo \(t\) lo medimos en segundos y la altura \(h\) en metros.}
{\(8\,\mathrm{s}\), \(320\,\mathrm{m}\)}{\(8\,\mathrm{s}\), \(600\,\mathrm{m}\)}{\(16\,\mathrm{s}\), \(1190\,\mathrm{m}\)}{\(4\,\mathrm{s}\), \(230\,\mathrm{m}\)}{}{}}
\End

\Data\ID{2010017801}
\Updated{2022-08-25 10:08:47}
\Level{C}
\SubArea{Applications of derivatives}
\LANGen{
\Question{
Let \(p\) be the tangent to the graph of the function \(f(x) = x^{2} -6x +1\) perpendicular to the line \(x - 2y + 3 = 0\). Find the point  \(A\), where \(p\) touches the graph of the function \(f\).}
{\(A = \left [2;-7\right ]\)}{\(A = \left [4;-7\right ]\)}{\(A = \left [1;-4\right ]\)}{\(A = \left [0;1\right ]\)}{}{}}
\LANGpl{
\Question{
Niech \(p\) będzie styczną do wykresu funkcji \(f(x) = x^{2} -6x +1\) prostopadłej do prostej \(x - 2y + 3 = 0\). Znajdź punkt \(A\), gdzie \(p\) dotyka wykresu funkcji \(f\).}
{\(A = \left [2;-7\right ]\)}{\(A = \left [4;-7\right ]\)}{\(A = \left [1;-4\right ]\)}{\(A = \left [0;1\right ]\)}{}{}}
\LANGcs{
\Question{
Přímka \(p\) je tečnou  grafu funkce \(f(x) = x^{2} -6x +1\) kolmá 
na přímku \(x - 2y + 3 = 0\).
Najděte bod \(A\), ve kterém se tečna
\(p\) dotýká grafu funkce \(f\).}
{\(A = \left [2;-7\right ]\)}{\(A = \left [4;-7\right ]\)}{\(A = \left [1;-4\right ]\)}{\(A = \left [0;1\right ]\)}{}{}}
\LANGsk{
\Question{
Priamka \(p\) je dotyčnica ku grafu funkcie \(f(x) = x^{2} -6x +1\) kolmá 
na priamku \(x - 2y + 3 = 0\).
Nájdite bod \(A\), v ktorom sa dotyčnica
\(p\) dotýka grafu funkcie \(f\).}
{\(A = \left [2;-7\right ]\)}{\(A = \left [4;-7\right ]\)}{\(A = \left [1;-4\right ]\)}{\(A = \left [0;1\right ]\)}{}{}}
\LANGes{
\Question{
Dada la tangente \(p\) a la gráfica de la función \(f(x) = x^{2} -6x +1\) perpendicular a la recta \(x - 2y + 3 = 0\). Halla el punto \(A\), donde \(p\) toca la gráfica de la función \(f\).}
{\(A = \left [2;-7\right ]\)}{\(A = \left [4;-7\right ]\)}{\(A = \left [1;-4\right ]\)}{\(A = \left [0;1\right ]\)}{}{}}
\End

\Data\ID{2010017802}
\Updated{2022-11-02 09:11:42}
\Level{C}
\SubArea{Applications of derivatives}
\LANGen{
\Question{
Find the tangent to the graph of the function
\(f(x) = x^{2} - 6x + 1\) perpendicular
to the line \(x -2y + 4 = 0\).}
{\(2x + y +3 = 0\)}{\(-2x - y - 1= 0\)}{\(4x +2y - 1= 0\)}{\(-4x - 2y +3 = 0\)}{}{}}
\LANGpl{
\Question{
Znajdź styczną do wykresu funkcji
\(f(x) = x^{2} - 6x + 1\) prostopadłą
do prostej \(x -2y + 4 = 0\).}
{\(2x + y +3 = 0\)}{\(-2x - y - 1= 0\)}{\(4x +2y - 1= 0\)}{\(-4x - 2y +3 = 0\)}{}{}}
\LANGcs{
\Question{
Určete tečnu grafu funkce
\(f(x) = x^{2} - 6x + 1\) kolmou k přímce \(x -2y + 4 = 0\).}
{\(2x + y +3 = 0\)}{\(-2x - y - 1= 0\)}{\(4x +2y - 1= 0\)}{\(-4x - 2y +3 = 0\)}{}{}}
\LANGsk{
\Question{
Určte dotyčnicu grafu funkcie
\(f(x) = x^{2} - 6x + 1\) kolmej
na priamku \(x -2y + 4 = 0\).}
{\(2x + y +3 = 0\)}{\(-2x - y - 1= 0\)}{\(4x +2y - 1= 0\)}{\(-4x - 2y +3 = 0\)}{}{}}
\LANGes{
\Question{
Halla la tangente a la gráfica de la función
\(f(x) = x^{2} - 6x + 1\) perpendicular a la recta \(x -2y + 4 = 0\).}
{\(2x + y +3 = 0\)}{\(-2x - y - 1= 0\)}{\(4x +2y - 1= 0\)}{\(-4x - 2y +3 = 0\)}{}{}}
\End

\Data\ID{2010017803}
\Updated{2022-08-25 10:08:47}
\Level{C}
\SubArea{Applications of derivatives}
\LANGen{
\Question{
Determine the values of \( a \) and \( b \) (\( a \), \( b \in\mathbb{R} \)) such that the function
\[ f(x)=ax^3-2bx+2 \] 
has a local extremum of \( 6 \) at \( x=-1 \).}
{\( a=2 \), \( b=3 \)}{\( a=-2 \), \( b=3 \)}{\( a=-2 \), \( b=-3 \)}{\( a=2 \), \( b=-3 \)}{}{}}
\LANGpl{
\Question{
Wyznacz wartości  \( a \) i \( b \) (\( a \), \( b \in\mathbb{R} \)) takie, że
\[ f(x)=ax^3-2bx+2 \] 
ma lokalne ekstremum \( 6 \) w \( x=-1 \).}
{\( a=2 \), \( b=3 \)}{\( a=-2 \), \( b=3 \)}{\( a=-2 \), \( b=-3 \)}{\( a=2 \), \( b=-3 \)}{}{}}
\LANGcs{
\Question{
Určete hodnoty \( a \) a \( b \) (\( a \), \( b \in\mathbb{R} \)) tak, aby funkce
\[ f(x)=ax^3-2bx+2 \] 
měla lokální extrém v bodě \( x=-1 \) a jeho hodnota byla \( 6 \).}
{\( a=2 \), \( b=3 \)}{\( a=-2 \), \( b=3 \)}{\( a=-2 \), \( b=-3 \)}{\( a=2 \), \( b=-3 \)}{}{}}
\LANGsk{
\Question{
Určte hodnoty \( a \) a \( b \) (\( a \), \( b \in\mathbb{R} \)) tak, aby funkcia
\[ f(x)=ax^3-2bx+2 \] 
mala lokálny extrém v bode \( x=-1 \) a jeho hodnota bola \( 6 \).}
{\( a=2 \), \( b=3 \)}{\( a=-2 \), \( b=3 \)}{\( a=-2 \), \( b=-3 \)}{\( a=2 \), \( b=-3 \)}{}{}}
\LANGes{
\Question{
Halla los valores de \( a \) y \( b \) (\( a \), \( b \in\mathbb{R} \)) para que la función
\[ f(x)=ax^3-2bx+2 \] 
tenga un extremo local \( 6 \) en \( x=-1 \).}
{\( a=2 \), \( b=3 \)}{\( a=-2 \), \( b=3 \)}{\( a=-2 \), \( b=-3 \)}{\( a=2 \), \( b=-3 \)}{}{}}
\End

\Data\ID{2010017804}
\Updated{2022-11-02 09:11:42}
\Level{C}
\SubArea{Applications of derivatives}
\IMG{2010017801-figure0.pdf}
\LANGen{
\Question{
Using a \(60\,\mathrm{m}\) long wire mash we shall fence a rectangular garden with two inner walls (see the picture). What will the dimensions \(a\) and \(b\) of the garden be, if there is \(2\,\mathrm{m}\) wide opening in one outside wall and the area of the garden shall be as large as possible? (The wire mesh is used to make inner walls as well.)}
{$a=7.75\,\mathrm{m}$, $b=15.5\,\mathrm{m}$}{$a=7.25\,\mathrm{m}$, $b=16.5\,\mathrm{m}$}{$a=7.5\,\mathrm{m}$, $b=16\,\mathrm{m}$}{$a=10\,\mathrm{m}$, $b=11\,\mathrm{m}$}{}{}}
\LANGpl{
\Question{
Używając siatki drucianej o długości \(60\,\mathrm{m}\) ogrodzimy prostokątny ogród z dwoma wewnętrznymi przegrodami (patrz rysunek). Jakie będą wymiary \(a\) i \(b\) ogrodu, jeśli w jednej ze ścian zewnętrznych jest  \(2\,\mathrm{m}\) szeroki otwór, a powierzchnia ogrodu powinna być jak największa? (Siatka druciana służy również do wykonywania ścian wewnętrznych.)}
{$a=7{,}75\,\mathrm{m}$, $b=15{,}5\,\mathrm{m}$}{$a=7{,}25\,\mathrm{m}$, $b=16{,}5\,\mathrm{m}$}{$a=7{,}5\,\mathrm{m}$, $b=16\,\mathrm{m}$}{$a=10\,\mathrm{m}$, $b=11\,\mathrm{m}$}{}{}}
\LANGcs{
\Question{
Pletivem o délce \(60\,\mathrm{m}\) máme ohradit zahradu tvaru obdélníku se dvěma vnitřními přepážkami (viz obrázek). Jaké rozměry \(a\) a \(b\) bude mít zahrada, jestliže v jedné stěně má být otvor o délce \(2\,\mathrm{m}\) a ohraničená plocha má být co největší? (Pletivo bude použito i na přepážky.)}
{$a=7{,}75\,\mathrm{m}$, $b=15{,}5\,\mathrm{m}$}{$a=7{,}25\,\mathrm{m}$, $b=16{,}5\,\mathrm{m}$}{$a=7{,}5\,\mathrm{m}$, $b=16\,\mathrm{m}$}{$a=10\,\mathrm{m}$, $b=11\,\mathrm{m}$}{}{}}
\LANGsk{
\Question{
Pletivom o dĺžke \(60\,\mathrm{m}\) máme oplotiť záhradu tvaru obdĺžnika s dvoma vnútornými stenami  (pozri obrázok). Aké budú rozmery \(a\) a \(b\) tejto záhrady, ak v jednej vonkajšej stene má býť otvor dĺžky \(2\,\mathrm{m}\) a ohraničená plocha má býť čo najväčšia? (Pletivo bude použité aj na vnútorné steny.)}
{$a=7{,}75\,\mathrm{m}$, $b=15{,}5\,\mathrm{m}$}{$a=7{,}25\,\mathrm{m}$, $b=16{,}5\,\mathrm{m}$}{$a=7{,}5\,\mathrm{m}$, $b=16\,\mathrm{m}$}{$a=10\,\mathrm{m}$, $b=11\,\mathrm{m}$}{}{}}
\LANGes{
\Question{
Con una malla de alambre de \(60\,\mathrm{m}\) vallaremos un jardín rectangular con dos paredes interiores (mira la imagen). Halla las dimensiones de \(a\) y \(b\) del jardín sabiendo que hay una abertura de \(2\,\mathrm{m}\) de anchura en una pared exterior y queremos que el área total sea lo más grande posible. (El alambre se usará también para las paredes interiores).}
{$a=7.75\,\mathrm{m}$, $b=15.5\,\mathrm{m}$}{$a=7.25\,\mathrm{m}$, $b=16.5\,\mathrm{m}$}{$a=7.5\,\mathrm{m}$, $b=16\,\mathrm{m}$}{$a=10\,\mathrm{m}$, $b=11\,\mathrm{m}$}{}{}}
\End

\Data\ID{2010017805}
\Updated{2023-01-23 10:01:49}
\Level{C}
\SubArea{Applications of derivatives}
\LANGen{
\Question{
What dimensions (in centimeters) must a glass aquarium in the shape of a cuboid with a square bottom have, so that its volume is \(20\) liters and the surface of the aquarium is as small as possible. (We consider the cuboid without the lid.)}
{$a\doteq 34.2\,\mathrm{cm}$, $v\doteq 17.1\,\mathrm{cm}$}{$a\doteq 27.1\,\mathrm{cm}$, $v\doteq 27.1\,\mathrm{cm}$}{$a\doteq 63.2\,\mathrm{cm}$, $v\doteq 5\,\mathrm{cm}$}{$a\doteq 13.6\,\mathrm{cm}$, $v\doteq 108.6\,\mathrm{cm}$}{}{}}
\LANGpl{
\Question{
Jakie wymiary (w centymetrach) musi mieć szklane akwarium w kształcie prostopadłościanu z kwadratowym dnem, aby jego objętość wynosiła \(20\) litrów, a powierzchnia akwarium była jak najmniejsza. (Rozważamy prostopadłościan bez pokrywy.)}
{$a\doteq 34{,}2\,\mathrm{cm}$, $v\doteq 17{,}1\,\mathrm{cm}$}{$a\doteq 27{,}1\,\mathrm{cm}$, $v\doteq 27{,}1\,\mathrm{cm}$}{$a\doteq 63{,}2\,\mathrm{cm}$, $v\doteq 5\,\mathrm{cm}$}{$a\doteq 13{,}6\,\mathrm{cm}$, $v\doteq 108{,}6\,\mathrm{cm}$}{}{}}
\LANGcs{
\Question{
Jaké rozměry (v cm) musí mít skleněné akvárium tvaru kvádru s čtvercovým dnem, aby jeho objem byl \(20\) litrů a povrch akvária byl co nejmenší?  (Uvažujeme kvádr bez horní podstavy.)}
{$a\doteq 34{,}2\,\mathrm{cm}$, $v\doteq 17{,}1\,\mathrm{cm}$}{$a\doteq 27{,}1\,\mathrm{cm}$, $v\doteq 27{,}1\,\mathrm{cm}$}{$a\doteq 63{,}2\,\mathrm{cm}$, $v\doteq 5\,\mathrm{cm}$}{$a\doteq 13{,}6\,\mathrm{cm}$, $v\doteq 108{,}6\,\mathrm{cm}$}{}{}}
\LANGsk{
\Question{
Aké rozmery (v cm) musí mať sklenené akvárium tvaru kvádra so štvorcovým dnom, aby jeho objem bol \(20\) litrov a povrch akvária aby bol čo nejmenší?  (Uvažujeme kváder bez hornej podstavy.)}
{$a\doteq 34{,}2\,\mathrm{cm}$, $v\doteq 17{,}1\,\mathrm{cm}$}{$a\doteq 27{,}1\,\mathrm{cm}$, $v\doteq 27{,}1\,\mathrm{cm}$}{$a\doteq 63{,}2\,\mathrm{cm}$, $v\doteq 5\,\mathrm{cm}$}{$a\doteq 13{,}6\,\mathrm{cm}$, $v\doteq 108{,}6\,\mathrm{cm}$}{}{}}
\LANGes{
\Question{
¿Qué dimensiones (en centímetros) debe tener un acuario de vidrio en forma de ortoedro con fondo cuadrado para que su volumen sea de \(20\) litros y la superficie del acuario sea lo más pequeña posible? (Consideramos que el ortoedro carece de tapa).}
{$a\doteq 34.2\,\mathrm{cm}$, $v\doteq 17.1\,\mathrm{cm}$}{$a\doteq 27.1\,\mathrm{cm}$, $v\doteq 27.1\,\mathrm{cm}$}{$a\doteq 63.2\,\mathrm{cm}$, $v\doteq 5\,\mathrm{cm}$}{$a\doteq 13.6\,\mathrm{cm}$, $v\doteq 108.6\,\mathrm{cm}$}{}{}}
\End

\Data\ID{2010017806}
\Updated{2022-11-02 09:11:42}
\Level{C}
\SubArea{Applications of derivatives}
\IMG{2010017801-figure1.pdf}
\LANGen{
\Question{
We want to lift an edge of a large square plate with a side of \(4\,\mathrm{m}\) so that it creates a shelter (see the picture). To what height \(h\) do we have to lift the edge of the plate, if the created shelter shall be of the largest possible volume?}
{$h=2\sqrt2\,\mathrm{m}$}{$h=4\cdot \sqrt{\frac23}\,\mathrm{m}$}{$h=\frac43\sqrt3\,\mathrm{m}$}{$h=\left( -\frac12 + \sqrt{65}\right)\,\mathrm{m}$}{}{}}
\LANGpl{
\Question{
Chcemy podnieść krawędź dużej kwadratowej płyty o boku \(4\,\mathrm{m}\) tak, aby tworzyła schronienie (patrz rysunek). Na jaką wysokość \(h\) musimy podnieść krawędź płyty, jeśli stworzony schron ma mieć największą możliwą objętość?}
{$h=2\sqrt2\,\mathrm{m}$}{$h=4\cdot \sqrt{\frac23}\,\mathrm{m}$}{$h=\frac43\sqrt3\,\mathrm{m}$}{$h=\left( -\frac12 + \sqrt{65}\right)\,\mathrm{m}$}{}{}}
\LANGcs{
\Question{
Velkou desku čtvercového tvaru o straně \(4\,\mathrm{m}\) chceme na jedné straně zvednout tak, aby vzniknul přístřešek (viz obrázek). Do jaké výšky \(h\) musíme stranu desky zvednout, aby vzniklý přístřešek měl co největší objem?}
{$h=2\sqrt2\,\mathrm{m}$}{$h=4\cdot \sqrt{\frac23}\,\mathrm{m}$}{$h=\frac43\sqrt3\,\mathrm{m}$}{$h=\left( -\frac12 + \sqrt{65}\right)\,\mathrm{m}$}{}{}}
\LANGsk{
\Question{
Veľkú dosku štvorcového tvaru so stranou \(4\,\mathrm{m}\) chceme na jednej strane zdvihnúť tak, aby vznikol prístrešok (pozri obrázok). Do akej výšky \(h\) musíme stranu dosky zvihnúť, aby vzniknutý prístrešok mal čo najväčší objem?}
{$h=2\sqrt2\,\mathrm{m}$}{$h=4\cdot \sqrt{\frac23}\,\mathrm{m}$}{$h=\frac43\sqrt3\,\mathrm{m}$}{$h=\left( -\frac12 + \sqrt{65}\right)\,\mathrm{m}$}{}{}}
\LANGes{
\Question{
Queremos levantar una lona cuadrada cuyo lado mide \(4\,\mathrm{m}\) para crear un refugio (mira la imagen). ¿Hasta qué altura \(h\) debemos levantar la lona, si el refugio creado debe tener el mayor volumen posible?}
{$h=2\sqrt2\,\mathrm{m}$}{$h=4\cdot \sqrt{\frac23}\,\mathrm{m}$}{$h=\frac43\sqrt3\,\mathrm{m}$}{$h=\left( -\frac12 + \sqrt{65}\right)\,\mathrm{m}$}{}{}}
\End

\Data\ID{2010020005}
\Updated{2023-01-23 10:01:09}
\Level{C}
\SubArea{Applications of derivatives}
\IMG{obrz.pdf}
\LANGen{
\Question{
The picture shows the part of the graph of a function \(f\). Choose the formula of the function \(f\) whose graph corresponds to this picture.}
{\(f(x)=\frac{x^2+4}{x^2-4}\)}{\(f(x)=\frac{x}{x^2-4}+1\)}{\(f(x)=-\frac{x^2}{4-x^2}-1\)}{\(f(x)=\frac{x^2}{4-x^2}\)}{}{}}
\LANGpl{
\Question{
Rysunek przedstawia część wykresu funkcji  \(f\). Wybierz wzór funkcji \(f\), której wykres odpowiada temu rysunkowi.}
{\(f(x)=\frac{x^2+4}{x^2-4}\)}{\(f(x)=\frac{x}{x^2-4}+1\)}{\(f(x)=-\frac{x^2}{4-x^2}-1\)}{\(f(x)=\frac{x^2}{4-x^2}\)}{}{}}
\LANGcs{
\Question{
Na obrázku je část grafu funkce \(f\). Vyberte předpis funkce, jejíž graf odpovídá tomuto obrázku.}
{\(f(x)=\frac{x^2+4}{x^2-4}\)}{\(f(x)=\frac{x}{x^2-4}+1\)}{\(f(x)=-\frac{x^2}{4-x^2}-1\)}{\(f(x)=\frac{x^2}{4-x^2}\)}{}{}}
\LANGsk{
\Question{
Na obrázku je časť grafu funkcie \(f\). Vyberte predpis funkcie \(f\), ktorej graf vidíte na obrázku.}
{\(f(x)=\frac{x^2+4}{x^2-4}\)}{\(f(x)=\frac{x}{x^2-4}+1\)}{\(f(x)=-\frac{x^2}{4-x^2}-1\)}{\(f(x)=\frac{x^2}{4-x^2}\)}{}{}}
\LANGes{
\Question{
La imagen representa la parte de la gráfica de la función \(f\). Elige la fórmula de la función \(f\) cuya gráfica corresponde a esta imagen.}
{\(f(x)=\frac{x^2+4}{x^2-4}\)}{\(f(x)=\frac{x}{x^2-4}+1\)}{\(f(x)=-\frac{x^2}{4-x^2}-1\)}{\(f(x)=\frac{x^2}{4-x^2}\)}{}{}}
\End

\Data\ID{2010020006}
\Updated{2023-01-23 10:01:09}
\Level{C}
\SubArea{Applications of derivatives}
\IMG{obrz_0.pdf}
\LANGen{
\Question{
The picture shows the part of the graph of a function \(f\). Choose the formula of the function \(f\) whose graph corresponds to this picture.}
{\(f(x)=-\frac{x^2}{4-x^2}-1\)}{\(f(x)=\frac{x^2+4}{x^2-4}\)}{\(f(x)=\frac{x}{x^2-4}+1\)}{\(f(x)=\frac{x^2}{4-x^2}\)}{}{}}
\LANGpl{
\Question{
Rysunek przedstawia część wykresu funkcji \(f\). Wybierz wzór funkcji \(f\), której wykres odpowiada temu rysunkowi.}
{\(f(x)=-\frac{x^2}{4-x^2}-1\)}{\(f(x)=\frac{x^2+4}{x^2-4}\)}{\(f(x)=\frac{x}{x^2-4}+1\)}{\(f(x)=\frac{x^2}{4-x^2}\)}{}{}}
\LANGcs{
\Question{
Na obrázku je část grafu funkce \(f\). Vyberte předpis funkce, jejíž graf odpovídá tomuto obrázku.}
{\(f(x)=-\frac{x^2}{4-x^2}-1\)}{\(f(x)=\frac{x^2+4}{x^2-4}\)}{\(f(x)=\frac{x}{x^2-4}+1\)}{\(f(x)=\frac{x^2}{4-x^2}\)}{}{}}
\LANGsk{
\Question{
Na obrázku je časť grafu funkcie \(f\). Vyberte predpis funkcie \(f\), ktorej graf vidíte na obrázku.}
{\(f(x)=-\frac{x^2}{4-x^2}-1\)}{\(f(x)=\frac{x^2+4}{x^2-4}\)}{\(f(x)=\frac{x}{x^2-4}+1\)}{\(f(x)=\frac{x^2}{4-x^2}\)}{}{}}
\LANGes{
\Question{
La imagen representa la parte de la gráfica de la función \(f\). Elige la fórmula de la función \(f\) cuya gráfica corresponde a esta imagen.}
{\(f(x)=-\frac{x^2}{4-x^2}-1\)}{\(f(x)=\frac{x^2+4}{x^2-4}\)}{\(f(x)=\frac{x}{x^2-4}+1\)}{\(f(x)=\frac{x^2}{4-x^2}\)}{}{}}
\End

\Data\ID{2010020007}
\Updated{2023-01-23 10:01:09}
\Level{C}
\SubArea{Applications of derivatives}
\LANGen{
\Question{
A function \(f\) is given by the formula \(f(x)=\frac{2x^2+3}{x+1}\). For the function \(f\) identify a slant or horizontal asymptote.}
{\(y=2x-2\)}{\(y=2\)}{\(y=-2x\)}{This function has neither horizontal nor slant asymptote.}{}{}}
\LANGpl{
\Question{
Funkcja \(f\) wyrażona jest wzorem \(f(x)=\frac{2x^2+3}{x+1}\). Dla funkcji \(f\) wskaż asymptotę ukośną lub poziomą.}
{\(y=2x-2\)}{\(y=2\)}{\(y=-2x\)}{Ta funkcja nie ma asymptoty poziomej ani ukośnej.}{}{}}
\LANGcs{
\Question{
Funkce \(f\) je určena předpisem \(f(x)=\frac{2x^2+3}{x+1}\). Pro danou funkci určete asymptotu se směrnicí.}
{\(y=2x-2\)}{\(y=2\)}{\(y=-2x\)}{Tato funkce nemá asymptotu se směrnicí.}{}{}}
\LANGsk{
\Question{
Funkcia \(f\) je daná predpisom \(f(x)=\frac{2x^2+3}{x+1}\).  Pre danú funkciu \(f\) určte rovnicu asymptoty so smernicou.}
{\(y=2x-2\)}{\(y=2\)}{\(y=-2x\)}{Táto funkcia nemá asymptotu so smernicou.}{}{}}
\LANGes{
\Question{
La función \(f\) viene dada por la fórmula \(f(x)=\frac{2x^2+3}{x+1}\). Identifica una asíntota oblicua u horizontal para la función \(f\).}
{\(y=2x-2\)}{\(y=2\)}{\(y=-2x\)}{La función no tiene asíntota horizontal ni oblicua.}{}{}}
\End

\Data\ID{2010020008}
\Updated{2023-01-23 10:01:09}
\Level{C}
\SubArea{Applications of derivatives}
\LANGen{
\Question{
A function \(f\) is given by the formula \(f(x)=\frac{3x^2-1}{x-2}\). For the function \(f\) identify a slant or horizontal asymptote.}
{\(y=3x+6\)}{\(y=3\)}{\(y=3x-3\)}{This function has neither horizontal nor slant asymptote.}{}{}}
\LANGpl{
\Question{
Funkcja \(f\) wyrażona jest wzorem \(f(x)=\frac{3x^2-1}{x-2}\). Dla funkcji \(f\) określ asymptotę ukośną lub poziomą.}
{\(y=3x+6\)}{\(y=3\)}{\(y=3x-3\)}{Ta funkcja nie ma asymptoty poziomej ani ukośnej.}{}{}}
\LANGcs{
\Question{
Funkce \(f\) je určena předpisem \(f(x)=\frac{3x^2-1}{x-2}\). Pro danou funkci určete asymptotu se směrnicí.}
{\(y=3x+6\)}{\(y=3\)}{\(y=3x-3\)}{Tato funkce nemá asymptotu se směrnicí.}{}{}}
\LANGsk{
\Question{
Funkcia \(f\) je daná predpisom \(f(x)=\frac{3x^2-1}{x-2}\).  Pre danú funkciu \(f\) určte rovnicu asymptoty so smernicou.}
{\(y=3x+6\)}{\(y=3\)}{\(y=3x-3\)}{Táto funkcia nemá asymptotu so smernicou.}{}{}}
\LANGes{
\Question{
La función \(f\) viene dada por la fórmula \(f(x)=\frac{3x^2-1}{x-2}\). Identifica una asíntota oblicua u horizontal para la función \(f\).}
{\(y=3x+6\)}{\(y=3\)}{\(y=3x-3\)}{La función no tiene asíntota horizontal ni oblicua.}{}{}}
\End

\Data\ID{2010020009}
\Updated{2023-01-23 10:01:09}
\Level{C}
\SubArea{Applications of derivatives}
\LANGen{
\Question{
A function \(f\) is given by the formula \(f(x)=\frac{\ln{x}}{2-x^2}\). For the function \(f\) identify all vertical asymptotes.}
{\(x=\sqrt2,\quad x=0\)}{\(x=\sqrt2,\quad x=-\sqrt{2},\quad x=0\)}{\(x=\sqrt2,\quad x=-\sqrt{2}\)}{This function has no vertical asymptote.}{}{}}
\LANGpl{
\Question{
Funkcja \(f\) dana jest wzorem  \(f(x)=\frac{\ln{x}}{2-x^2}\). Dla funkcji  \(f\) wkaż wszystkie pionowe asymptoty.}
{\(x=\sqrt2,\quad x=0\)}{\(x=\sqrt2,\quad x=-\sqrt{2},\quad x=0\)}{\(x=\sqrt2,\quad x=-\sqrt{2}\)}{Ta funkcja nie ma pionowej asymptoty.}{}{}}
\LANGcs{
\Question{
Funkce \(f\) je určena předpisem \(f(x)=\frac{\ln{x}}{2-x^2}\). Určete všechny svislé asymptoty dané funkce.}
{\(x=\sqrt2,\quad x=0\)}{\(x=\sqrt2,\quad x=-\sqrt{2},\quad x=0\)}{\(x=\sqrt2,\quad x=-\sqrt{2}\)}{Tato funkce nemá žádnou svislou asymptotu.}{}{}}
\LANGsk{
\Question{
Funkcia \(f\) je daná predpisom \(f(x)=\frac{\ln{x}}{2-x^2}\). Pre danú funkciu \(f\) určte všetky zvislé asymptoty.}
{\(x=\sqrt2,\quad x=0\)}{\(x=\sqrt2,\quad x=-\sqrt{2},\quad x=0\)}{\(x=\sqrt2,\quad x=-\sqrt{2}\)}{Táto funkcia nemá zvislú asymptotu.}{}{}}
\LANGes{
\Question{
La función \(f\) viene dada por la fórmula \(f(x)=\frac{\ln{x}}{2-x^2}\). Identifica todas las asíntotas verticales para la función \(f\).}
{\(x=\sqrt2,\quad x=0\)}{\(x=\sqrt2,\quad x=-\sqrt{2},\quad x=0\)}{\(x=\sqrt2,\quad x=-\sqrt{2}\)}{La función no tiene ninguna asíntota vertical.}{}{}}
\End

\Data\ID{2010020010}
\Updated{2023-01-23 10:01:09}
\Level{C}
\SubArea{Applications of derivatives}
\LANGen{
\Question{
A function \(f\) is given by the formula \(f(x)=\frac{1+\ln{x}}{x^2-3}\). For the function \(f\) identify all vertical asymptotes.}
{\(x=\sqrt3,\quad x=0\)}{\(x=\sqrt3,\quad x=-\sqrt3,\quad x=0\)}{\(x=\sqrt3,\quad x=-\sqrt3\)}{This function has no vertical asymptote.}{}{}}
\LANGpl{
\Question{
Funkcja \(f\) dana jest wzorem \(f(x)=\frac{1+\ln{x}}{x^2-3}\). Dla funkcji \(f\) wskaż wszystkie pionowe asymptoty.}
{\(x=\sqrt3,\quad x=0\)}{\(x=\sqrt3,\quad x=-\sqrt3,\quad x=0\)}{\(x=\sqrt3,\quad x=-\sqrt3\)}{Ta funkcja nie ma pionowej asymptoty.}{}{}}
\LANGcs{
\Question{
Funkce \(f\) je určena předpisem \(f(x)=\frac{1+\ln{x}}{x^2-3}\). Najděte všechny svislé asymptoty dané funkce.}
{\(x=\sqrt3,\quad x=0\)}{\(x=\sqrt3,\quad x=-\sqrt3,\quad x=0\)}{\(x=\sqrt3,\quad x=-\sqrt3\)}{Tato funkce nemá žádnou svislou asymptotu.}{}{}}
\LANGsk{
\Question{
Funkcia \(f\) je daná predpisom \(f(x)=\frac{1+\ln{x}}{x^2-3}\). Pre danú funkciu \(f\) určte všetky zvislé asymptoty.}
{\(x=\sqrt3,\quad x=0\)}{\(x=\sqrt3,\quad x=-\sqrt3,\quad x=0\)}{\(x=\sqrt3,\quad x=-\sqrt3\)}{Táto funkcia nemá zvislú asymptotu.}{}{}}
\LANGes{
\Question{
La función \(f\) viene dada por la fórmula \(f(x)=\frac{1+\ln{x}}{x^2-3}\). Identifica todas las asíntotas verticales para la función \(f\).}
{\(x=\sqrt3,\quad x=0\)}{\(x=\sqrt3,\quad x=-\sqrt3,\quad x=0\)}{\(x=\sqrt3,\quad x=-\sqrt3\)}{La función no tiene ninguna asíntota vertical.}{}{}}
\End

\Data\ID{2010020011}
\Updated{2023-01-23 10:01:09}
\Level{C}
\SubArea{Applications of derivatives}
\LANGen{
\Question{
How many of the following functions have exactly two asymptotes?
\[f(x)=\left(\frac{4+x}{4-x}\right)^4,\quad g(x)=\frac{x^3}{(x-2)^2},\quad h(x)=\sqrt{6-4x},\quad i(x)=\frac{x^2}{4-x^2}\]}
{2}{1}{3}{None of these functions has exactly two asymptotes.}{}{}}
\LANGpl{
\Question{
Ile z poniższych funkcji ma dokładnie dwie asymptoty?
\[f(x)=\left(\frac{4+x}{4-x}\right)^4,\quad g(x)=\frac{x^3}{(x-2)^2},\quad h(x)=\sqrt{6-4x},\quad i(x)=\frac{x^2}{4-x^2}\]}
{2}{1}{3}{Żadna z tych funkcji nie ma dokładnie dwóch asymptot.}{}{}}
\LANGcs{
\Question{
Kolik z uvedených funkcí má právě dvě asymptoty?
\[f(x)=\left(\frac{4+x}{4-x}\right)^4,\quad g(x)=\frac{x^3}{(x-2)^2},\quad h(x)=\sqrt{6-4x},\quad i(x)=\frac{x^2}{4-x^2}\]}
{2}{1}{3}{Žádná z uvedených funkcí nemá právě dvě asymptoty.}{}{}}
\LANGsk{
\Question{
Koľko z daných funkcií má práve dve asymptoty?
\[f(x)=\left(\frac{4+x}{4-x}\right)^4,\quad g(x)=\frac{x^3}{(x-2)^2},\quad h(x)=\sqrt{6-4x},\quad i(x)=\frac{x^2}{4-x^2}\]}
{2}{1}{3}{Žiadna z uvedených funkcií nemá práve dve asymptoty.}{}{}}
\LANGes{
\Question{
¿Cuántas de las siguientes funciones tienen exactamente dos asíntotas?
\[f(x)=\left(\frac{4+x}{4-x}\right)^4,\quad g(x)=\frac{x^3}{(x-2)^2},\quad h(x)=\sqrt{6-4x},\quad i(x)=\frac{x^2}{4-x^2}\]}
{2}{1}{3}{Ninguna de estas funciones tiene exactamente dos asíntotas.}{}{}}
\End

\Data\ID{2010020012}
\Updated{2023-01-23 10:01:12}
\Level{C}
\SubArea{Applications of derivatives}
\LANGen{
\Question{
How many of the following functions have exactly two asymptotes?
\[f(x)=\frac{-x^3+1}{(x-3)^2},\quad h(x)=\left(\frac{x+2}{x-2}\right)^4,\quad i(x)=-\frac{x^2}{x^2-2},\quad g(x)=\sqrt{6x+2}\]}
{2}{1}{3}{None of these functions has exactly two asymptotes.}{}{}}
\LANGpl{
\Question{
Ile z poniższych funkcji ma dokładnie dwie asymptoty?
\[f(x)=\frac{-x^3+1}{(x-3)^2},\quad h(x)=\left(\frac{x+2}{x-2}\right)^4,\quad i(x)=-\frac{x^2}{x^2-2},\quad g(x)=\sqrt{6x+2}\]}
{2}{1}{3}{Żadna z tych funkcji nie ma dokładnie dwóch asymptot.}{}{}}
\LANGcs{
\Question{
Kolik z uvedených funkcí má právě dvě asymptoty?
\[f(x)=\frac{-x^3+1}{(x-3)^2},\quad h(x)=\left(\frac{x+2}{x-2}\right)^4,\quad i(x)=-\frac{x^2}{x^2-2},\quad g(x)=\sqrt{6x+2}\]}
{2}{1}{3}{Žádná z uvedených funkcí nemá právě dvě asymptoty.}{}{}}
\LANGsk{
\Question{
Koľko z daných funkcií má práve dve asymptoty?
\[f(x)=\frac{-x^3+1}{(x-3)^2},\quad h(x)=\left(\frac{x+2}{x-2}\right)^4,\quad i(x)=-\frac{x^2}{x^2-2},\quad g(x)=\sqrt{6x+2}\]}
{2}{1}{3}{Žiadna z uvedených funkcií nemá práve dve asymptoty.}{}{}}
\LANGes{
\Question{
¿Cuántas de las siguientes funciones tienen exactamente dos asíntotas?
\[f(x)=\frac{-x^3+1}{(x-3)^2},\quad h(x)=\left(\frac{x+2}{x-2}\right)^4,\quad i(x)=-\frac{x^2}{x^2-2},\quad g(x)=\sqrt{6x+2}\]}
{2}{1}{3}{Ninguna de estas funciones tiene exactamente dos asíntotas.}{}{}}
\End

\Data\ID{9000062406}
\Updated{2020-07-15 03:07:37}
\Level{C}
\SubArea{Applications of derivatives}
\LANGen{
\Question{
Identify an asymptote to the graph of the following function. 
\[ 
 f(x) = \frac{x^{2} + 4} 
 {x}
\]}
{\(y = x\)}{\(y = x + 1\)}{\(y = x - 1\)}{\(y = -x\)}{}{}}
\LANGpl{
\Question{
Wyznacz asymptotę  wykresu funkcji.
\[ 
 f\colon y = \frac{x^{2} + 4} 
 {x}
\]}
{\(y = x\)}{\(y = x + 1\)}{\(y = x - 1\)}{\(y = -x\)}{}{}}
\LANGcs{
\Question{
Která z uvedených přímek je asymptotou grafu funkce
\(f\colon y = \frac{x^{2}+4} 
 {x} \)?}
{\(y = x\)}{\(y = x + 1\)}{\(y = x - 1\)}{\(y = -x\)}{}{}}
\LANGsk{
\Question{
Ktorá z uvedených priamok je asymptotou grafu funkcie
\(f\colon y = \frac{x^{2}+4} 
 {x} \)?}
{\(y = x\)}{\(y = x + 1\)}{\(y = x - 1\)}{\(y = -x\)}{}{}}
\LANGes{
\Question{
Identifica una asíntota en la gráfica de la siguiente función. 
\[ 
 f(x) = \frac{x^{2} + 4} 
 {x}
\]}
{\(y = x\)}{\(y = x + 1\)}{\(y = x - 1\)}{\(y = -x\)}{}{}}
\End

\Data\ID{9000062407}
\Updated{2021-08-12 10:08:46}
\Level{C}
\SubArea{Applications of derivatives}
\LANGen{
\Question{
Find the tangent to the function \(f(x) =\ln x\) 
at the point \(T = [1;y_{0}]\).}
{\(y = x - 1\)}{\(y = x\)}{\(y = x + 1\)}{\(y = -x\)}{}{}}
\LANGpl{
\Question{
Wyznacz styczną do wykresu funkcji \(f\colon y =\ln x\) 
w punkcie  \(T = [1;y_{0}]\).}
{\(y = x - 1\)}{\(y = x\)}{\(y = x + 1\)}{\(y = -x\)}{}{}}
\LANGcs{
\Question{
Určete rovnici tečny grafu funkce \(f\colon y =\ln x\) 
v bodě \(T = [1;y_{0}]\).}
{\(y = x - 1\)}{\(y = x\)}{\(y = x + 1\)}{\(y = -x\)}{}{}}
\LANGsk{
\Question{
Určte rovnicu dotyčnice grafu funkcie \(f\colon y =\ln x\) 
v bode \(T = [1;y_{0}]\).}
{\(y = x - 1\)}{\(y = x\)}{\(y = x + 1\)}{\(y = -x\)}{}{}}
\LANGes{
\Question{
Halla la ecuación de la recta tangente a la función \(f(x) =\ln x\) 
en el punto \(T = [1;y_{0}]\).}
{\(y = x - 1\)}{\(y = x\)}{\(y = x + 1\)}{\(y = -x\)}{}{}}
\End

\Data\ID{9000062408}
\Updated{2021-08-12 10:08:15}
\Level{C}
\SubArea{Applications of derivatives}
\LANGen{
\Question{
Find the points in which the tangent to the curve
\(y = x^{3}\) has a
slope \(m = 3\)?}
{\(T_{1} = [1;1],\ T_{2} = [-1;-1]\)}{\(T_{1} = [1;-1],\ T_{2} = [-1;1]\)}{\(T_{1} = [-1;1],\ T_{2} = [-1;-1]\)}{\(T_{1} = [1;-1],\ T_{2} = [-1;-1]\)}{}{}}
\LANGpl{
\Question{
Wyznacz punkty tak, aby styczna do krzywej 
\(y = x^{3}\) miała nachylenie 
 \(m = 3\).}
{\(T_{1} = [1;1],\ T_{2} = [-1;-1]\)}{\(T_{1} = [1;-1],\ T_{2} = [-1;1]\)}{\(T_{1} = [-1;1],\ T_{2} = [-1;-1]\)}{\(T_{1} = [1;-1],\ T_{2} = [-1;-1]\)}{}{}}
\LANGcs{
\Question{
V kterých bodech má tečna křivky, která je dána předpisem
\(y = x^{3}\), směrnici
\(k = 3\)?}
{\(T_{1} = [1;1],\ T_{2} = [-1;-1]\)}{\(T_{1} = [1;-1],\ T_{2} = [-1;1]\)}{\(T_{1} = [-1;1],\ T_{2} = [-1;-1]\)}{\(T_{1} = [1;-1],\ T_{2} = [-1;-1]\)}{}{}}
\LANGsk{
\Question{
V ktorých bodoch má dotyčnica krivky, ktorá je daná predpisom
\(y = x^{3}\), smernicou
\(k = 3\)?}
{\(T_{1} = [1;1],\ T_{2} = [-1;-1]\)}{\(T_{1} = [1;-1],\ T_{2} = [-1;1]\)}{\(T_{1} = [-1;1],\ T_{2} = [-1;-1]\)}{\(T_{1} = [1;-1],\ T_{2} = [-1;-1]\)}{}{}}
\LANGes{
\Question{
Halla los puntos en los que la tangente a la curva
\(y = x^{3}\) tiene una pendiente \(m = 3\)?}
{\(T_{1} = [1;1],\ T_{2} = [-1;-1]\)}{\(T_{1} = [1;-1],\ T_{2} = [-1;1]\)}{\(T_{1} = [-1;1],\ T_{2} = [-1;-1]\)}{\(T_{1} = [1;-1],\ T_{2} = [-1;-1]\)}{}{}}
\End

\Data\ID{9000062410}
\Updated{2021-08-12 10:08:07}
\Level{C}
\SubArea{Applications of derivatives}
\LANGen{
\Question{
Find the line perpendicular to the graph of the function
\(f(x) = x^{3}\) at the
point \(T = [-1;y_{0}]\).}
{\(x + 3y + 4 = 0\)}{\(3x - y + 2 = 0\)}{\(3x + y - 4 = 0\)}{\(x - 3y - 2 = 0\)}{}{}}
\LANGpl{
\Question{
Wskaż prostą prostopadłą do wykresu funkcji 
\(f\colon y = x^{3}\)  w punkcie \(T = [-1;y_{0}]\).}
{\(x + 3y + 4 = 0\)}{\(3x - y + 2 = 0\)}{\(3x + y - 4 = 0\)}{\(x - 3y - 2 = 0\)}{}{}}
\LANGcs{
\Question{
Určete rovnici normály grafu funkce
\(f\colon y = x^{3}\) v bodě
\(T = [-1;y_{0}]\).}
{\(x + 3y + 4 = 0\)}{\(3x - y + 2 = 0\)}{\(3x + y - 4 = 0\)}{\(x - 3y - 2 = 0\)}{}{}}
\LANGsk{
\Question{
Určte rovnicu normály grafu funkcie
\(f\colon y = x^{3}\) v bode
\(T = [-1;y_{0}]\).}
{\(x + 3y + 4 = 0\)}{\(3x - y + 2 = 0\)}{\(3x + y - 4 = 0\)}{\(x - 3y - 2 = 0\)}{}{}}
\LANGes{
\Question{
Halla la recta perpendicular a la gráfica de la función
\(f(x) = x^{3}\) en el punto \(T = [-1;y_{0}]\).}
{\(x + 3y + 4 = 0\)}{\(3x - y + 2 = 0\)}{\(3x + y - 4 = 0\)}{\(x - 3y - 2 = 0\)}{}{}}
\End

\Data\ID{9000064101}
\Updated{2020-07-15 04:07:14}
\Level{C}
\SubArea{Applications of derivatives}
\LANGen{
\Question{
Find the slope of the tangent to the graph of
\(f(x) = x^{2} + 3x - 2\) at the
point \([1;2]\).}
{\(-\frac{1} 
{5}\)}{\(5\)}{\(- 5\)}{\(\frac{1}
{5}\)}{}{}}
\LANGpl{
\Question{
Oblicz  nachylenie  stycznej do wykresu funkcji
\(f\colon y = x^{2} + 3x - 2\) w punkcie
 \([1;2]\).}
{\(-\frac{1} 
{5}\)}{\(5\)}{\(- 5\)}{\(\frac{1}
{5}\)}{}{}}
\LANGcs{
\Question{
Je dána funkce \(f\colon y = x^{2} + 3x - 2\). Směrnice
normály grafu funkce \(f\) 
v bodě \(T = \left [1;2\right ]\) 
je rovna:}
{\(-\frac{1} 
{5}\)}{\(5\)}{\(- 5\)}{\(\frac{1}
{5}\)}{}{}}
\LANGsk{
\Question{
Je daná funkcia \(f\colon y = x^{2} + 3x - 2\). Smernica
normály grafu funkcie \(f\) 
v bode \(T = \left [1;2\right ]\) 
je rovná:}
{\(-\frac{1} 
{5}\)}{\(5\)}{\(- 5\)}{\(\frac{1}
{5}\)}{}{}}
\LANGes{
\Question{
Halla la pendiente de la tangente a la gráfica de la función
\(f(x) = x^{2} + 3x - 2\) en el punto \([1;2]\).}
{\(-\frac{1} 
{5}\)}{\(5\)}{\(- 5\)}{\(\frac{1}
{5}\)}{}{}}
\End

\Data\ID{9000064102}
\Updated{2021-08-12 10:08:47}
\Level{C}
\SubArea{Applications of derivatives}
\LANGen{
\Question{
Find the tangent to the graph of the function
\(f(x) = \frac{x+1} 
{x-1}\) at the
point \([2;3]\).}
{\(2x + y - 7 = 0\)}{\(2x - y - 1 = 0\)}{\(- 2x + y + 1 = 0\)}{\(x + 2y - 9 = 0\)}{}{}}
\LANGpl{
\Question{
Wyznacz styczną do wykresu funkcji 
\(f\colon y = \frac{x+1} 
{x-1}\) w punkcie  \([2;3]\).}
{\(2x + y - 7 = 0\)}{\(2x - y - 1 = 0\)}{\(- 2x + y + 1 = 0\)}{\(x + 2y - 9 = 0\)}{}{}}
\LANGcs{
\Question{
Je dána funkce \(f\colon y = \frac{x+1} 
{x-1}\).
Tečna grafu funkce \(f\) 
v bodě \(T = \left [2;3\right ]\) 
má rovnici:}
{\(2x + y - 7 = 0\)}{\(2x - y - 1 = 0\)}{\(- 2x + y + 1 = 0\)}{\(x + 2y - 9 = 0\)}{}{}}
\LANGsk{
\Question{
Je daná funkcia \(f\colon y = \frac{x+1} 
{x-1}\).
Dotyčnica grafu funkcie \(f\) 
v bode \(T = \left [2;3\right ]\) 
má rovnici:}
{\(2x + y - 7 = 0\)}{\(2x - y - 1 = 0\)}{\(- 2x + y + 1 = 0\)}{\(x + 2y - 9 = 0\)}{}{}}
\LANGes{
\Question{
Halla la tangente a la gráfica de la función
\(f(x) = \frac{x+1} 
{x-1}\) en el punto \((2;3)\).}
{\(2x + y - 7 = 0\)}{\(2x - y - 1 = 0\)}{\(- 2x + y + 1 = 0\)}{\(x + 2y - 9 = 0\)}{}{}}
\End

\Data\ID{9000064103}
\Updated{2021-08-12 10:08:55}
\Level{C}
\SubArea{Applications of derivatives}
\LANGen{
\Question{
Find the normal line to the graph of the function
\(f(x) = 2x^{2} - 2x + 1\) at the
point \([2;5]\).}
{\(x + 6y - 32 = 0\)}{\(6x - y - 7 = 0\)}{\(x + 6y - 4 = 0\)}{\(- 6x + y - 7 = 0\)}{}{}}
\LANGpl{
\Question{
Wyznacz prostą prawidłową do wykresu funkcji 
\(f\colon y = 2x^{2} - 2x + 1\) w punkcie 
 \([2;5]\).}
{\(x + 6y - 32 = 0\)}{\(6x - y - 7 = 0\)}{\(x + 6y - 4 = 0\)}{\(- 6x + y - 7 = 0\)}{}{}}
\LANGcs{
\Question{
Je dána funkce \(f\colon y = 2x^{2} - 2x + 1\).
Normála grafu funkce \(f\) 
v bodě \(T = \left [2;5\right ]\) 
má rovnici:}
{\(x + 6y - 32 = 0\)}{\(6x - y - 7 = 0\)}{\(x + 6y - 4 = 0\)}{\(- 6x + y - 7 = 0\)}{}{}}
\LANGsk{
\Question{
Je daná funkcia \(f\colon y = 2x^{2} - 2x + 1\).
Normála grafu funkcie \(f\) 
v bode \(T = \left [2;5\right ]\) 
má rovnicu:}
{\(x + 6y - 32 = 0\)}{\(6x - y - 7 = 0\)}{\(x + 6y - 4 = 0\)}{\(- 6x + y - 7 = 0\)}{}{}}
\LANGes{
\Question{
Halla la recta normal a la gráfica de la función
\(f(x) = 2x^{2} - 2x + 1\) en el punto \([2;5]\).}
{\(x + 6y - 32 = 0\)}{\(6x - y - 7 = 0\)}{\(x + 6y - 4 = 0\)}{\(- 6x + y - 7 = 0\)}{}{}}
\End

\Data\ID{9000064104}
\Updated{2021-08-12 10:08:33}
\Level{C}
\SubArea{Applications of derivatives}
\LANGen{
\Question{
Let \(p\) be the tangent to
the graph of the function \(f(x) = x^{2} - x - 6\) 
parallel to the line \(y = 3x + 1\).
Find the point \(A\) 
where \(p\) touches
the graph of \(f\).}
{\(A = \left [2;-4\right ]\)}{\(A = \left [2;4\right ]\)}{\(A = \left [1;6\right ]\)}{\(A = \left [-1;-4\right ]\)}{}{}}
\LANGpl{
\Question{
Dana jest styczna  \(p\)  do wykresu funkcji
 \(f\colon y = x^{2} - x - 6\) 
równoległa do prostej  \(y = 3x + 1\).
Wskaż punkt \(A\) 
tak, aby  \(p\) stykała się z wykresem funkcji
 \(f\).}
{\(A = \left [2;-4\right ]\)}{\(A = \left [2;4\right ]\)}{\(A = \left [1;6\right ]\)}{\(A = \left [-1;-4\right ]\)}{}{}}
\LANGcs{
\Question{
Je dána funkce \(f\colon y = x^{2} - x - 6\). Pro dotykový
bod tečny grafu funkce \(f\) 
rovnoběžné s přímkou \(p\colon y = 3x + 1\) 
platí:}
{\(A = \left [2;-4\right ]\)}{\(A = \left [2;4\right ]\)}{\(A = \left [1;6\right ]\)}{\(A = \left [-1;-4\right ]\)}{}{}}
\LANGsk{
\Question{
Je daná funkcie \(f\colon y = x^{2} - x - 6\). Pre dotykový
bod dotyčnice grafu funkcie \(f\) 
rovnobežnej s priamkou \(p\colon y = 3x + 1\) 
platí:}
{\(A = \left [2;-4\right ]\)}{\(A = \left [2;4\right ]\)}{\(A = \left [1;6\right ]\)}{\(A = \left [-1;-4\right ]\)}{}{}}
\LANGes{
\Question{
Dada la tangente \(p\) a la gráfica de la función \(f(x) = x^{2} - x - 6\) 
paralela a la recta \(y = 3x + 1\).
Halla el punto \(A\) 
donde \(p\) toca la gráfica de \(f\).}
{\(A = \left [2;-4\right ]\)}{\(A = \left [2;4\right ]\)}{\(A = \left [1;6\right ]\)}{\(A = \left [-1;-4\right ]\)}{}{}}
\End

\Data\ID{9000064105}
\Updated{2020-07-15 04:07:14}
\Level{C}
\SubArea{Applications of derivatives}
\LANGen{
\Question{
Find the tangent to the graph of the function
\(f(x) = x\sin x\) at the
point \(\left [ \frac{\pi }{2}; \frac{\pi } {2}\right ]\).}
{\(y = x\)}{\(y = x + 1\)}{\(y = 0\)}{\(y =\pi -x\)}{}{}}
\LANGpl{
\Question{
Wyznacz styczną do wykresu funkcji 
\(f\colon y = x\sin x\)  w punkcie 
 \(\left [ \frac{\pi }{2}; \frac{\pi } {2}\right ]\).}
{\(y = x\)}{\(y = x + 1\)}{\(y = 0\)}{\(y =\pi -x\)}{}{}}
\LANGcs{
\Question{
Je dána funkce \(f\colon y = x\sin x\).
Tečna grafu funkce \(f\) 
v bodě \(A = \left [ \frac{\pi }{2}; \frac{\pi } {2}\right ]\) 
má rovnici:}
{\(y = x\)}{\(y = x + 1\)}{\(y = 0\)}{\(y =\pi -x\)}{}{}}
\LANGsk{
\Question{
Je daná funkcia \(f\colon y = x\sin x\).
Dotyčnica grafu funkcie \(f\) 
v bode \(A = \left [ \frac{\pi }{2}; \frac{\pi } {2}\right ]\) 
má rovnicu:}
{\(y = x\)}{\(y = x + 1\)}{\(y = 0\)}{\(y =\pi -x\)}{}{}}
\LANGes{
\Question{
Halla la tangente a la gráfica de la función
\(f(x) = x\sin x\) en el punto \(\left [ \frac{\pi }{2}; \frac{\pi } {2}\right ]\).}
{\(y = x\)}{\(y = x + 1\)}{\(y = 0\)}{\(y =\pi -x\)}{}{}}
\End

\Data\ID{9000064106}
\Updated{2022-04-23 05:04:17}
\Level{C}
\SubArea{Applications of derivatives}
\LANGen{
\Question{
Let \(p\) be the tangent to
the graph of the function \(f(x) = x^{2} + 4x - 2\) 
perpendicular to the line \(x + 6y + 2 = 0\).
Find the point \(A\) where
\(p\) touches the graph
of the function \(f\).}
{\(A = \left [1;3\right ]\)}{\(A = \left [-5;3\right ]\)}{\(A = \left [-3;-5\right ]\)}{\(A = \left [0;-2\right ]\)}{}{}}
\LANGpl{
\Question{
Niech  \(p\) będzie styczną do wykresu funkcji
 \(f\colon y = x^{2} + 4x - 2\) 
prostopadłą do prostej \(x + 6y + 2 = 0\).
Wyznacz punkt  \(A\) tak, aby
\(p\) stykała się  z wykresem funkcji  \(f\).}
{\(A = \left [1;3\right ]\)}{\(A = \left [-5;3\right ]\)}{\(A = \left [-3;-5\right ]\)}{\(A = \left [0;-2\right ]\)}{}{}}
\LANGcs{
\Question{
Je dána funkce \(f\colon y = x^{2} + 4x - 2\).
Tečna grafu funkce \(f\) 
kolmá na přímku \(p\colon x + 6y + 2 = 0\) se
dotýká grafu funkce \(f\) 
v bodě:}
{\(\left [1;3\right ]\)}{\(\left [-5;3\right ]\)}{\(\left [-3;-5\right ]\)}{\(\left [0;-2\right ]\)}{}{}}
\LANGsk{
\Question{
Je daná funkcia \(f\colon y = x^{2} + 4x - 2\).
Dotyčnica grafu funkcie \(f\) 
kolmá na priamku \(p\colon x + 6y + 2 = 0\) sa
dotýka grafu funkcie \(f\) 
v bode:}
{\(\left [1;3\right ]\)}{\(\left [-5;3\right ]\)}{\(\left [-3;-5\right ]\)}{\(\left [0;-2\right ]\)}{}{}}
\LANGes{
\Question{
Dada la tangente \(p\) a la gráfica de la función \(f(x) = x^{2} + 4x - 2\) 
perpendicular a la recta \(x + 6y + 2 = 0\).
Halla el punto \(A\) donde
\(p\) toca la gráfica de la función \(f\).}
{\(A = \left [1;3\right ]\)}{\(A = \left [-5;3\right ]\)}{\(A = \left [-3;-5\right ]\)}{\(A = \left [0;-2\right ]\)}{}{}}
\End

\Data\ID{9000064107}
\Updated{2021-08-12 11:08:45}
\Level{C}
\SubArea{Applications of derivatives}
\LANGen{
\Question{
Find the tangent to the graph of the function
\(f(x) = x^{2} + 4x - 2\) parallel to
the line \(2x + y + 1 = 0\).}
{\(2x + y + 11 = 0\)}{\(2x - y - 1 = 0\)}{\(2x + y - 1 = 0\)}{\(2x - y - 11 = 0\)}{}{}}
\LANGpl{
\Question{
Wyznacz styczną do wykresu funkcji
\(f\colon y = x^{2} + 4x - 2\) równoległą do prostej
 \(2x + y + 1 = 0\).}
{\(2x + y + 11 = 0\)}{\(2x - y - 1 = 0\)}{\(2x + y - 1 = 0\)}{\(2x - y - 11 = 0\)}{}{}}
\LANGcs{
\Question{
Je dána funkce \(f\colon y = x^{2} + 4x - 2\).
Tečna grafu funkce \(f\) 
rovnoběžná s přímkou \(p\colon 2x + y + 1 = 0\) 
má rovnici:}
{\(2x + y + 11 = 0\)}{\(2x - y - 1 = 0\)}{\(2x + y - 1 = 0\)}{\(2x - y - 11 = 0\)}{}{}}
\LANGsk{
\Question{
Je daná funkcia \(f\colon y = x^{2} + 4x - 2\).
Dotyčnica grafu funkcie \(f\) 
rovnobežná s priamkou \(p\colon 2x + y + 1 = 0\) 
má rovnicu:}
{\(2x + y + 11 = 0\)}{\(2x - y - 1 = 0\)}{\(2x + y - 1 = 0\)}{\(2x - y - 11 = 0\)}{}{}}
\LANGes{
\Question{
Halla la tangente a la gráfica de la función
\(f(x) = x^{2} + 4x - 2\) paralela a la recta \(2x + y + 1 = 0\).}
{\(2x + y + 11 = 0\)}{\(2x - y - 1 = 0\)}{\(2x + y - 1 = 0\)}{\(2x - y - 11 = 0\)}{}{}}
\End

\Data\ID{9000064108}
\Updated{2021-08-12 11:08:23}
\Level{C}
\SubArea{Applications of derivatives}
\LANGen{
\Question{
Find the normal line to the graph of the function
\(f(x) = 2x^{2} - 7x\) parallel to
the line \(y = -x\).}
{\(x + y + 4 = 0\)}{\(- x + y + 4 = 0\)}{\(x - y - 8 = 0\)}{\(x + y - 8 = 0\)}{}{}}
\LANGpl{
\Question{
Wskaż prostą prawidłową do wykresu funkcji 
\(f\colon y = 2x^{2} - 7x\) równoległą do prostej
 \(y = -x\).}
{\(x + y + 4 = 0\)}{\(- x + y + 4 = 0\)}{\(x - y - 8 = 0\)}{\(x + y - 8 = 0\)}{}{}}
\LANGcs{
\Question{
Je dána funkce \(f\colon y = 2x^{2} - 7x\).
Normála grafu funkce \(f\) 
rovnoběžná s osou II. a IV. kvadrantu má rovnici:}
{\(x + y + 4 = 0\)}{\(- x + y + 4 = 0\)}{\(x - y - 8 = 0\)}{\(x + y - 8 = 0\)}{}{}}
\LANGsk{
\Question{
Je daná funkcia \(f\colon y = 2x^{2} - 7x\).
Normála grafu funkcie \(f\) 
rovnobežná s osou II. a IV. kvadrante má rovnicu:}
{\(x + y + 4 = 0\)}{\(- x + y + 4 = 0\)}{\(x - y - 8 = 0\)}{\(x + y - 8 = 0\)}{}{}}
\LANGes{
\Question{
Halla la recta normal a la gráfica de la función
\(f(x) = 2x^{2} - 7x\) paralela a la recta \(y = -x\).}
{\(x + y + 4 = 0\)}{\(- x + y + 4 = 0\)}{\(x - y - 8 = 0\)}{\(x + y - 8 = 0\)}{}{}}
\End

\Data\ID{9000064109}
\Updated{2022-04-23 05:04:17}
\Level{C}
\SubArea{Applications of derivatives}
\LANGen{
\Question{
Find the tangent to the graph of the function
\(f(x) = 3x^{2} - 8x + 2\) perpendicular
to the line \(x + 4y + 5 = 0\).}
{\(4x - y - 10 = 0\)}{\(-4x + y + 1 = 0\)}{\(8x - 2y + 1 = 0\)}{\(-8x + 2y - 10 = 0\)}{}{}}
\LANGpl{
\Question{
Wyznacz styczną do wykresu funkcji
\(f\colon y = 3x^{2} - 8x + 2\) prostopadłą do prostej
 \(x + 4y + 5 = 0\).}
{\(4x - y - 10 = 0\)}{\(-4x + y + 1 = 0\)}{\(8x - 2y + 1 = 0\)}{\(-8x + 2y - 10 = 0\)}{}{}}
\LANGcs{
\Question{
Je dána funkce \(f\colon y = 3x^{2} - 8x + 2\).
Tečna grafu funkce \(f\) 
kolmá na přímku \(p\colon x + 4y + 5 = 0\) 
má rovnici:}
{\(4x - y - 10 = 0\)}{\(-4x + y + 1 = 0\)}{\(8x - 2y + 1 = 0\)}{\(-8x + 2y - 10 = 0\)}{}{}}
\LANGsk{
\Question{
Je daná funkcia \(f\colon y = 3x^{2} - 8x + 2\).
Dotyčnica grafu funkcie \(f\) 
kolmá na priamku \(p\colon x + 4y + 5 = 0\) 
má rovnicu:}
{\(4x - y - 10 = 0\)}{\(-4x + y + 1 = 0\)}{\(8x - 2y + 1 = 0\)}{\(-8x + 2y - 10 = 0\)}{}{}}
\LANGes{
\Question{
Halla la tangente a la gráfica de la función
\(f(x) = 3x^{2} - 8x + 2\) perpendicular
a la recta \(x + 4y + 5 = 0\).}
{\(4x - y - 10 = 0\)}{\(4x + y - 6 = 0\)}{\(4x + y - 10 = 0\)}{\(x + 4y + 6 = 0\)}{}{}}
\End

\Data\ID{9000064110}
\Updated{2021-08-15 08:08:39}
\Level{C}
\SubArea{Applications of derivatives}
\LANGen{
\Question{
Identify a true statement related to the function
\(f(x) = \frac{x-1} 
{x+1}\).}
{The tangent at \(T = [-3;2]\) 
is parallel to \(x - 2y + 1 = 0\).}{The tangent at \(T = [-3;2]\) 
contains the point \(A = \left [1;-4\right ]\).}{The slope of the tangent at \(T = [-3;2]\) 
is \(2\).}{The tangent at \(T = [-3;2]\) is
perpendicular to \(x + 2y + 1 = 0\).}{}{}}
\LANGpl{
\Question{
Oznacz zdanie prawdziwe biorąc po uwagę  funkcję 
\(f\colon y = \frac{x-1} 
{x+1}\).}
{Styczna \(T = [-3;2]\) 
jest równoległa do  \(x - 2y + 1 = 0\).}{Styczna \(T = [-3;2]\) 
zawiera punkt  \(A = \left [1;-4\right ]\).}{Nachylenie stycznej \(T = [-3;2]\) 
jest równe  \(2\).}{Styczna \(T = [-3;2]\) jest 
prostopadła do \(x + 2y + 1 = 0\).}{}{}}
\LANGcs{
\Question{
Je dána funkce \(f\colon y = \frac{x-1} 
{x+1}\).
Z následujících tvrzení vyberte to, které je pravdivé:}
{Tečna grafu funkce \(f\) v
bodě \(T = [-3;2]\) je rovnoběžná
s přímkou \(x - 2y + 1 = 0\).}{Tečna grafu funkce \(f\) 
v bodě \(T = [-3;2]\) 
prochází bodem \(A = \left [1;-4\right ]\).}{Tečna grafu funkce \(f\) 
v bodě \(T = [-3;2]\) má
směrnici \(2\).}{Tečna grafu funkce \(f\) 
v bodě \(T = [-3;2]\) je kolmá
na přímku \(x + 2y + 1 = 0\).}{}{}}
\LANGsk{
\Question{
Je daná funkcia \(f\colon y = \frac{x-1} 
{x+1}\).
Z nasledujúcich tvrdení vyberte to, ktoré je pravdivé:}
{Dotyčnica grafu funkcie \(f\) v
bode \(T = [-3;2]\) je rovnobežná
s priamkou \(x - 2y + 1 = 0\).}{Dotyčnica grafu funkcie \(f\) 
v bode \(T = [-3;2]\) 
prechádza bodom \(A = \left [1;-4\right ]\).}{Dotyčnica grafu funkcie \(f\) 
v bode \(T = [-3;2]\) má
smernicu \(2\).}{Dotyčnica grafu funkcie \(f\) 
v bode \(T = [-3;2]\) je kolmá
na priamku \(x + 2y + 1 = 0\).}{}{}}
\LANGes{
\Question{
Identifica la proposición lógica sobre la función
\(f(x) = \frac{x-1} 
{x+1}\).}
{La tangente en \(T = (-3;2)\) 
es paralela a \(x - 2y + 1 = 0\).}{La tangente en \(T = (-3;2)\) 
contiene el punto \(A = \left [1;-4\right ]\).}{La pendiente de la tangente en \(T = (-3;2)\) 
es \(2\).}{La tangente en \(T = (-3;2)\) es
perpendicular a \(x + 2y + 1 = 0\).}{}{}}
\End
